\documentclass[a4paper]{report}

\ifdefined\XeTeXversion
  \usepackage{fontspec}
  \setmonofont{JuliaMono}
\fi


\usepackage{amsmath}
\usepackage{amssymb}
\usepackage{amsthm}
\usepackage[colorlinks=true, linkcolor=blue, citecolor=blue, urlcolor=magenta]{hyperref}
\usepackage{microtype}
\usepackage{natbib}
\usepackage[capitalise]{cleveref}

\theoremstyle{plain}
\newtheorem{theorem}{Theorem}[chapter]
\newtheorem{lemma}[theorem]{Lemma}
\newtheorem{proposition}[theorem]{Proposition}
\newtheorem{corollary}[theorem]{Corollary}
\theoremstyle{definition}
\newtheorem{definition}[theorem]{Definition}

\theoremstyle{remark}
\newtheorem{assumptions}[theorem]{Conventions}

\newcommand{\Ll}{\mathcal{L}}
\newcommand{\Kk}{\mathcal{K}}
\newcommand{\Ran}{\mathcal{R}}
\newcommand{\Ker}{\mathcal{N}}
\newcommand{\Dom}{\mathcal{D}}
\newcommand{\Id}{\mathrm{Id}}
\newcommand{\dg}{^{\dagger}}
\newcommand{\wc}{\rightharpoonup}

 


\newif\iflabels\labelstrue
\newif\ifproofs\proofsfalse

\makeatletter
\@ifundefined{ifplastex}{\newif\ifplastex}{}
\makeatother

\ifplastex\else

\ifproofs\else
\usepackage{comment}
\excludecomment{proof}
\fi

\iflabels
\usepackage{seqsplit,accsupp,xcolor}
\makeatletter
\definecolor{runkonostate}{HTML}{FFFFFF}
\newcommand*\labelstate[2]{\@namedef{labelstate@#1}{#2}}
\newcommand*\labelcolour[3]{\@namedef{labelcolour@#1}{#2}\@namedef{labelborder@#1}{#3}}
\IfFileExists{labelstate}{\input{labelstate}}{}
\providecommand*\showkeyslabelformat[1]{\parbox{107pt}{\parbox{110pt}{\raggedleft\normalfont\tiny\ttfamily
\setlength\fboxrule{0.6pt}\@ifundefined{labelstate@#1}{\fcolorbox{runkonostate}{runkonostate}{\seqsplit{#1}}}{\edef\runko@border{\@nameuse{labelborder@\@nameuse{labelstate@#1}}}\edef\runko@colour{\@nameuse{labelcolour@\@nameuse{labelstate@#1}}}\fcolorbox{\runko@border}{\runko@colour}{\seqsplit{#1}}}}\BeginAccSupp{method=escape,ActualText={#1}}\small$|$\EndAccSupp{}}}
\makeatother
\usepackage[notref, notcite]{showkeys}
\fi

\fi
 
\title{Regularization of inverse problems with preliminaries in functional analysis}
\author{Tapio Helin, Medet Nursultanov and Lauri Oksanen}
\begin{document}
\maketitle

\chapter*{Background}
\addcontentsline{toc}{chapter}{Background}

This project auto-formalizes the classical theory of regularization for ill-posed inverse problems in the Lean~4 proof assistant. It began as the authors' entry point into the fascinating world of formalized mathematics. We emphasize that this is a work in progress: our objective is to continue to extend the formalized material, and we warmly invite colleagues to contribute and discuss how to make it useful. We hope that, as it matures, the library will serve as a foundation both for formalizing modern research-level results in the field as well as for educational purposes. This blueprint records the mathematical statements and proofs, to be matched node by node against the Lean development.

The regularization theory itself (Part~\ref{part_2_reg}) largely follows the lecture notes \emph{Regularization of Inverse Problems} by \citet{Clason2021}, and we gratefully acknowledge our debt to this clear and self-contained exposition. Although the theory is well established and its structure largely settled, having such a precise presentation at hand proved invaluable for proofreading the output of the (auto)formalization. The numbering in this document does not follow that of the lecture notes, but each node indicates the result to which it corresponds. At the beginning of each chapter, we aim to cite the key earlier works that led to the results presented there; any omission of relevant work is entirely our own oversight, and we would be grateful to have it brought to our attention.

Part~\ref{part_1_funcanal}, the functional-analytic preliminaries to the regularization
theory, is classical material gathered from various (classical) sources rather than from
a single reference. Much of the needed functional analysis is already formalized in Mathlib and we've done our best to avoid overlapping formalization. We aim for a design that adapts easily to the growth of Mathlib: as Mathlib expands, we intend to shrink our own functional analysis layer and rely increasingly on Mathlib instead.

At the time of writing, the library covers only linear ill-posed problems, but we aim to extend it to nonlinear problems and beyond.

\tableofcontents
\part{Preliminaries in functional analysis}
\label{part_1_funcanal}
\chapter{Bounded operators on normed and Banach spaces}
\label{ch_01_bounded_operators}

\begin{assumptions}
  \label{asm_bounded_operators_normed}
  $X$, $Y$, $Z$ denote normed vector spaces over
$\mathbb{K}$, and $T \in \mathcal{L}(X,Y)$.
\end{assumptions}

\begin{definition}[Bounded operator, operator norm]
  \label{def_clm}
  A linear map $T \colon X \to Y$ between normed spaces is \emph{bounded} if there is
  $C \ge 0$ with $\|Tx\|_Y \le C\|x\|_X$ for all $x \in X$.  The space of such maps is
  $\mathcal{L}(X,Y)$, normed by
  \[
    \|T\|_{\mathcal{L}(X,Y)} = \inf\{\, C \ge 0 : \|Tx\|_Y \le C\|x\|_X \text{ for all } x \,\}.
  \]
  For linear maps, boundedness and continuity coincide.
\end{definition}


\begin{lemma}[The operator norm is the least bound]
  \label{lem_isleast_opnorm}
  $\|T\|$ is the least element of $\{\,C \ge 0 : \|Tx\| \le C\|x\| \ \forall x\,\}$; in
  particular $\|Tx\| \le \|T\|\,\|x\|$.
\end{lemma}


\begin{lemma}[Completeness of $\mathcal{L}(X,Y)$]
  \label{lem_clm_complete}
  If $Y$ is a Banach space then so is $\mathcal{L}(X,Y)$, for any normed $X$.  In
  particular the dual $X' := \mathcal{L}(X,\mathbb{K})$ is a Banach space even when $X$ is
  not complete.
\end{lemma}




\begin{lemma}[$\mathcal{L}(X,X)$ is a Banach algebra]
  \label{lem_banach_algebra}
  $\|ST\| \le \|S\|\,\|T\|$; in particular the powers $T^n$ are bounded, and if $X$ is
  Banach then $\mathcal{L}(X,X)$ is a Banach algebra.
\end{lemma}


\begin{lemma}[Operator norm as a supremum]
  \label{lem_opnorm_sup}
  For $X$ a normed space over $\mathbb{K} \in \{\mathbb{R},\mathbb{C}\}$,
  \[
    \|T\| = \sup_{\|x\| \le 1}\|Tx\| = \sup_{\|x\| = 1}\|Tx\| .
  \]
\end{lemma}



\begin{lemma}[Finite-dimensional subspaces are closed]
  \label{lem_findim_closed}
  A finite-dimensional subspace of a normed space over a \emph{complete} scalar field is
  closed.
\end{lemma}


\begin{lemma}[Riesz's lemma]
  \label{lem_riesz_lemma}
  Let $X$ be a normed space over $\mathbb{K} \in \{\mathbb{R},\mathbb{C}\}$, let
  $U \subsetneq X$ be a closed proper subspace, and let $\theta < 1$.  Then there is $x$
  with $\|x\| = 1$ and $\operatorname{dist}(x,U) \ge \theta$.
\end{lemma}


\begin{theorem}[Riesz: compact unit ball]
  \label{thm_riesz_ball}
  Let $X$ be a normed space over $\mathbb{K} \in \{\mathbb{R},\mathbb{C}\}$.  The closed
  unit ball of $X$ is compact if and only if $\dim X < \infty$.
\end{theorem}
\begin{proof}
  ($\Leftarrow$)  If $\dim X < \infty$ then $X$ is a proper metric space, and in a proper space every closed ball is compact.  The case $X = \{0\}$ is included: there
  the ball is a single point.

  ($\Rightarrow$)  We prove the contrapositive.  Suppose $\dim X = \infty$, and set
  \[
    U_n := \operatorname{span}\{x_1,\dots,x_{n-1}\} .
  \]
  We construct unit vectors $x_1, x_2, \dots$ with
  $\operatorname{dist}(x_n, U_n) \ge 1/2$.  Given the first $n-1$ of them, $U_n$ is
  finite dimensional, hence closed by \cref{lem_findim_closed}, and
  proper because $\dim X = \infty$; so \cref{lem_riesz_lemma} with $\theta = 1/2$ supplies
  $x_n$.  For $n = 1$ this reads $U_1 = \{0\}$, which is closed and proper for the same
  reason.

  For $m < n$ we have $x_m \in U_n$, hence
  $\|x_n - x_m\| \ge \operatorname{dist}(x_n,U_n) \ge 1/2$.  So $(x_n)$ is a sequence in
  the closed unit ball with no Cauchy subsequence, and therefore no convergent
  subsequence.  The closed unit ball is a metric space, in which compactness and sequential
  compactness agree, so it is not compact.
\end{proof}


\begin{theorem}[Open mapping theorem]
  \label{thm_open_mapping}
  If $X$ and $Y$ are Banach spaces and $T \in \mathcal{L}(X,Y)$ is surjective, then $T$ is
  an open map.
\end{theorem}


\begin{theorem}[Bounded inverse theorem]
  \label{thm_bounded_inverse}
  If $X,Y$ are Banach spaces and $T \in \mathcal{L}(X,Y)$ is bijective, then
  $T^{-1} \in \mathcal{L}(Y,X)$.
\end{theorem}


\begin{theorem}[Closed graph theorem]
  \label{thm_closed_graph}
  A linear map between Banach spaces whose graph is closed is continuous.
\end{theorem}


\begin{theorem}[Bounded extension]
  \label{thm_blt}
  A bounded linear operator defined on a dense subspace of a normed space, with values in
  a Banach space, extends uniquely to a bounded operator on the whole space, with no
  larger norm.
\end{theorem}




\begin{lemma}[Neumann series]
  \label{lem_neumann}
  Let $X$ be a Banach space and $T \in \mathcal{L}(X,X)$ with $\|T\| < 1$.  Then
  $\mathrm{Id}-T$ is invertible,
  \[
    (\mathrm{Id}-T)^{-1} = \sum_{k \ge 0} T^k, \qquad
    \|(\mathrm{Id}-T)^{-1}\| \le \frac{1}{1-\|T\|}.
  \]
\end{lemma}



\begin{theorem}[Banach fixed-point theorem]
  \label{thm_banach_fixed_point}
  A contraction of a nonempty \emph{complete metric space} has a unique fixed point, and
  the Picard iterates converge to it geometrically.
\end{theorem}


\begin{theorem}[Banach--Steinhaus]
  \label{thm_banach_steinhaus}
  Let $X$ be a Banach space and $Y$ a normed space.  If $\{T_i\}_{i \in I}$ in
  $\mathcal{L}(X,Y)$ is pointwise bounded, then $\sup_{i \in I}\|T_i\| < \infty$.
\end{theorem}


\begin{theorem}[Banach--Steinhaus, resonance form]
  \label{thm_resonance}
  Let $X$ be a Banach space, $Y$ a normed space, and $\{T_i\}_{i \in I}$ a family in
  $\mathcal{L}(X,Y)$ with $\sup_{i \in I}\|T_i\| = \infty$.  Then there is $x \in X$ with
  $\sup_{i \in I}\|T_ix\|_Y = \infty$.
\end{theorem}
\begin{proof}
  Suppose no such $x$ exists, i.e.\ $\sup_{i \in I}\|T_ix\|_Y < \infty$ for every
  $x \in X$.  That is exactly pointwise boundedness of $\{T_i\}_{i \in I}$, so
  \cref{thm_banach_steinhaus} gives $\sup_{i \in I}\|T_i\| < \infty$, contradicting the
  hypothesis.  Completeness of $X$ is consumed there and nowhere else.
\end{proof}



\begin{definition}[The limit operator]
  \label{def_limit_operator}
  Let $X$ be Banach, $Y$ normed, and let $T_n \in \mathcal{L}(X,Y)$ converge pointwise:
  $T_nx \to Tx$ for every $x \in X$.  The \emph{limit operator} is the map
  $x \mapsto \lim_n T_nx$, and it is an element of $\mathcal{L}(X,Y)$ --- linearity is
  immediate and boundedness is Banach--Steinhaus (\cref{thm_banach_steinhaus}), which is
  what makes the definition well posed.
\end{definition}


\begin{proposition}[Pointwise convergence along a filter]
  \label{prop_dense_ptwise_filter}
  Let $X,Y$ be normed spaces, let $\mathcal{F}$ be a filter on an index set $\mathcal{I}$,
  and let $\{T_i\}_{i \in \mathcal{I}}$ be a family in $\mathcal{L}(X,Y)$ with
  $\sup_{i}\|T_i\| < \infty$.  Let $T \in \mathcal{L}(X,Y)$ and let $M \subseteq X$ be
  total.  If $T_ix \to Tx$ along $\mathcal{F}$ for every $x \in M$, then $T_ix \to Tx$
  along $\mathcal{F}$ for every $x \in X$.  Neither space need be complete.
\end{proposition}
\begin{proof}
  This is the ``$\Leftarrow$'' half of the classical criterion --- Banach--Steinhaus against
  a total set, over a Banach $X$ --- verbatim, repeated rather than cited because that
  statement assumes completeness of $X$ and this one must not, with
  the index
  $n \to \infty$ replaced by an arbitrary filter $\mathcal{F}$; we repeat the
  $3\varepsilon$ argument in full to make visible that no property of the index beyond the
  filter axioms is used.  Put $D := \max\{\sup_{i}\|T_i\|,\ \|T\|\} < \infty$.

  \emph{Step 1.}  $T_iu \to Tu$ along $\mathcal{F}$ for every $u \in \operatorname{span}M$:
  writing $u = \sum_{k=1}^m\lambda_kx_k$ with $x_k \in M$ gives
  $T_iu - Tu = \sum_{k=1}^m \lambda_k(T_ix_k - Tx_k)$, and a finite sum of functions
  tending to $0$ along a filter tends to $0$, because a finite intersection of members of a
  filter is a member of it.  The empty sum $u = 0$ is covered.

  \emph{Step 2.}  Let $x \in X$ and $\varepsilon > 0$.  Choose $u \in \operatorname{span}M$
  with $\|x - u\| < \varepsilon/(3(D+1))$, which is possible since $\operatorname{span}M$
  is dense, and then $A \in \mathcal{F}$ with $\|T_iu - Tu\| < \varepsilon/3$ for $i \in A$.
  For those $i$, using \cref{lem_isleast_opnorm} on the outer two terms,
  \[
    \|T_ix - Tx\| \;\le\; \|T_i(x-u)\| + \|T_iu - Tu\| + \|T(u-x)\|
      \;<\; 2D\,\frac{\varepsilon}{3(D+1)} + \frac{\varepsilon}{3} \;<\; \varepsilon ,
  \]
  and $A \in \mathcal{F}$ was arbitrary in the sense that this exhibits, for each
  $\varepsilon$, a member of $\mathcal{F}$ on which the estimate holds.

  Two remarks on hypotheses.  The uniform bound $\sup_i\|T_i\| < \infty$ is a hypothesis
  here, not a conclusion: the Banach--Steinhaus half of the classical criterion plays no
  role, and consequently completeness of $X$ is not consumed anywhere in this proof --- so
  this node, unlike that one, is available over a merely normed $X$, which
  is the whole reason \cref{prop_ptwise_limit_exists} may cite it.  And the converse
  implication does \emph{not} transfer for free: pointwise convergence along a filter does
  not by itself yield a uniform bound, since \cref{thm_banach_steinhaus} is a statement
  about a family and gives no information along $\mathcal{F}$.
\end{proof}


\begin{proposition}[Existence of the limit operator]
  \label{prop_ptwise_limit_exists}
  Let $X$ be a normed space, $Y$ a Banach space, $T_n \in \mathcal{L}(X,Y)$, and
  $M \subseteq X$ total.  If $\|T_n\| \le C$ for every $n$ and $T_nx$ converges for every
  $x \in M$, then there is $T \in \mathcal{L}(X,Y)$ with $\|T\| \le C$ and
  $T_n \to T$ pointwise.
\end{proposition}
\begin{proof}
  Let $V := \operatorname{span}M$, a dense subspace of $X$; $C$ is the given uniform
  bound, $\|T_n\| \le C$ for all $n$.

  For $u = \sum_{k=1}^m\lambda_kx_k \in V$ with $x_k \in M$, the sequence
  $T_nu = \sum_{k=1}^m \lambda_kT_nx_k$ converges, being a finite linear combination of
  convergent sequences; and its limit depends only on $u$, so
  \[
    Su \;:=\; \lim_{n\to\infty} T_nu \qquad (u \in V)
  \]
  is well defined, independently of the chosen representation.  No completeness is needed
  for this: the limits along $M$ exist by hypothesis.  $S$ is linear, since each $T_n$ is
  and limits are linear, and bounded on $V$: $\|T_nu\| \le C\|u\|$ for every $n$ by
  \cref{lem_isleast_opnorm}, and letting $n \to \infty$ gives $\|Su\| \le C\|u\|$.  (If
  $M = \emptyset$ then $V = \{0\}$, whose density forces $X = \{0\}$ and $S = 0$; the
  argument is unaffected.)

  Since $V$ is dense in $X$ and $Y$ is complete, \cref{thm_blt} extends $S$ to
  $T \in \mathcal{L}(X,Y)$ with $\|T\| \le \|S\| \le C$ --- ``no larger'', which is all
  \cref{thm_blt} claims and all that is needed.  This is the only step that uses
  completeness of $Y$, and it is what distinguishes this proposition from the classical
  criterion, where the limit operator is given rather than produced.

  Finally $M \subseteq V$, so $T_nx \to Sx = Tx$ for every $x \in M$, and
  $\|T_n\| \le C$ for all $n$; \cref{prop_dense_ptwise_filter}, applied with
  $\mathcal{F} = \mathrm{atTop}$ on $\mathbb{N}$, therefore gives $T_nx \to Tx$ for every
  $x \in X$.  That node consumes no completeness of $X$, which is why $X$ is only assumed
  normed here; the classical criterion could \emph{not} be cited in its place, being
  stated over a Banach $X$.
\end{proof}


\begin{definition}[Quotient space]
  \label{def_quotient}
  For a closed subspace $M$ of a normed space $X$, the quotient $X/M$ carries the norm
  $\|x+M\| = \inf_{y \in M}\|x+y\|$, and is a Banach space when $X$ is.
\end{definition}


\begin{lemma}[Kernels of bounded operators are closed]
  \label{lem_ker_closed}
  For $T \in \mathcal{L}(X,Y)$ between normed spaces, $\mathcal{N}(T)$ is a closed
  subspace.
\end{lemma}


\begin{theorem}[Closed range criterion]
  \label{thm_closed_range}
  Let $X,Y$ be Banach spaces and $T \in \mathcal{L}(X,Y)$.  Then $\mathcal{R}(T)$ is
  closed if and only if there is $c > 0$ with
  \[
    \|Tx\|_Y \;\ge\; c \inf_{z \in \mathcal{N}(T)} \|x + z\|_X
    \;=\; c\,\|x + \mathcal{N}(T)\|_{X/\mathcal{N}(T)}
    \qquad \text{for all } x \in X,
  \]
  i.e.\ the induced injective operator $X/\mathcal{N}(T) \to Y$ is bounded below.
\end{theorem}
\begin{proof}
  Write $N := \mathcal{N}(T)$, a closed subspace by \cref{lem_ker_closed}.  By
  \cref{def_quotient}, $X/N$ is therefore a normed space, and a Banach space because $X$
  is.

  Let $\hat T(x+N) := Tx$.  This is well defined ($x - x' \in N$ gives $Tx = Tx'$), linear,
  injective ($\hat T(x+N) = 0$ means $x \in N$), and has the same range as $T$.  It is
  bounded: for every $z \in N$, $\|\hat T(x+N)\| = \|T(x+z)\| \le \|T\|\,\|x+z\|$ by
  \cref{lem_isleast_opnorm}, and taking the infimum over $z \in N$ gives
  $\|\hat T(x+N)\| \le \|T\|\,\|x+N\|$.  The displayed inequality of the statement says
  exactly that $\hat T$ is bounded below by $c$.

  ($\Leftarrow$)  Assume the inequality holds with some $c > 0$, and let
  $y \in \overline{\mathcal{R}(T)}$, say $y = \lim_k \hat T\xi_k$ with $\xi_k \in X/N$.
  Then $(\hat T\xi_k)_k$ is Cauchy, and $c\,\|\xi_k - \xi_j\| \le \|\hat T\xi_k -
  \hat T\xi_j\|$, so $(\xi_k)_k$ is Cauchy in $X/N$, which is complete because $X$ is
  (\cref{def_quotient}).  Hence $\xi_k \to \xi$ for some $\xi$, and continuity of $\hat T$
  gives $y = \hat T\xi \in \mathcal{R}(T)$.  \emph{This direction consumes completeness of
  $X$ only}; $Y$ is not used, since the candidate limit $y$ is handed to us.

  ($\Rightarrow$)  Assume $\mathcal{R}(T)$ is closed.  Then $\mathcal{R}(T)$ is a closed
  subspace of the Banach space $Y$, hence itself complete, and $X/N$ is complete because
  $X$ is.  So $\hat T \colon X/N \to \mathcal{R}(T)$ is a bounded bijection between Banach
  spaces, and \cref{thm_bounded_inverse} makes $\hat T^{-1}$ bounded, say
  $\|\hat T^{-1}\| \le C$.  With $c := (C+1)^{-1} > 0$ we get, for every $x$,
  \[
    \|x+N\| \;=\; \bigl\|\hat T^{-1}\hat T(x+N)\bigr\| \;\le\; C\,\|Tx\|
      \;\le\; c^{-1}\|Tx\| ,
  \]
  which is the asserted inequality; the shift from $C$ to $C+1$ only covers the degenerate
  case $X/N = \{0\}$, where $C = 0$ is admissible while $c$ must be positive.
  \emph{This direction consumes completeness of both $X$ and $Y$.}

  No completeness hypothesis can be dropped.  For ``$\Rightarrow$'' without $Y$ complete:
  take $X = \ell^1$ with $\|\cdot\|_1$, $Y = \ell^1$ with $\|\cdot\|_\infty$ (normed but
  not complete), and $T = \mathrm{id}$, which is bounded since
  $\|x\|_\infty \le \|x\|_1$; then $\mathcal{R}(T) = Y$ is closed while a bound
  $\|x\|_\infty \ge c\|x\|_1$ would make the two norms equivalent, which they are not.
  For ``$\Leftarrow$'' without $X$ complete: take
  $X = c_{00}$ with the $\ell^1$ norm, $Y = \ell^1$, and $T$ the inclusion; then
  $N = \{0\}$ and $\|Tx\| = \|x\|$, so the inequality holds with $c = 1$, yet
  $\mathcal{R}(T) = c_{00}$ is dense and proper in $\ell^1$, hence not closed.  For
  ``$\Rightarrow$'' without $X$ complete: on an infinite-dimensional Banach space $Y$ pick
  a discontinuous linear functional $\varphi$ and let $X$ be $Y$ with the strictly finer
  norm $\|y\|_X = \|y\| + |\varphi(y)|$; the identity $T \colon X \to Y$ is bounded and
  bijective, so $\mathcal{R}(T) = Y$ is closed, but a lower bound $\|y\| \ge c\|y\|_X$
  would make $\varphi$ continuous.
\end{proof}


\chapter{Duality}
\label{ch_02_duality}

\begin{assumptions}
  \label{asm_duality}
  $X$, $Y$ are normed spaces over $\mathbb{K}$;
$X' := \mathcal{L}(X,\mathbb{K})$ is the dual.
\end{assumptions}



\begin{definition}[Dual space and duality pairing]
  \label{def_dual_space}
  $X' = \mathcal{L}(X,\mathbb{K})$ with $\langle x, f\rangle := f(x)$ and
  $|\langle x,f\rangle| \le \|x\|\,\|f\|_{X'}$.
\end{definition}


\begin{theorem}[Hahn--Banach, norm-preserving extension]
  \label{thm_hahn_banach}
  Let $X_0 \subseteq X$ be a subspace of a normed space and $f_0 \in X_0'$.  Then there is
  $f \in X'$ extending $f_0$ with $\|f\|_{X'} = \|f_0\|_{X_0'}$.
\end{theorem}


\begin{corollary}[Norming functionals]
  \label{cor_norming_functional}
  For $x \ne 0$ there is $f \in X'$ with $\|f\| = 1$ and $f(x) = \|x\|$.  Consequently
  $\|x\| = \sup_{\|f\|\le 1}|\langle x,f\rangle|$, and $X'$ separates the points of $X$.
\end{corollary}


\begin{definition}[Dual operator]
  \label{def_dual_operator}
  For $T \in \mathcal{L}(X,Y)$ the \emph{dual} (or Banach-adjoint) operator
  $T' \in \mathcal{L}(Y',X')$ is $T'g = g \circ T$, characterised by
  $\langle Tx, g\rangle = \langle x, T'g\rangle$ for all $x \in X$, $g \in Y'$.
\end{definition}



\begin{theorem}[Hahn--Banach separation]
  \label{thm_hahn_banach_separation}
  Let $X$ be a normed space over $\mathbb{K} \in \{\mathbb{R},\mathbb{C}\}$ that is in
  addition a normed space over $\mathbb{R}$ with compatible scalars and locally convex, let
  $C \subseteq X$ be closed and convex, and let $x_0 \notin C$.  Then there are $f \in X'$
  and $u \in \mathbb{R}$ with
  \[
    \operatorname{Re}f(w) < u < \operatorname{Re}f(x_0) \qquad \text{for all } w \in C .
  \]
\end{theorem}



\chapter{Hilbert spaces}
\label{ch_03_hilbert_spaces}

\begin{assumptions}
  \label{asm_hilbert_spaces}
  $X$, $Y$ are Hilbert spaces over $\mathbb{K}$ with
$\mathbb{K} = \mathbb{R}$ or $\mathbb{K} = \mathbb{C}$, i.e.\ complete inner product spaces, and
$T \in \mathcal{L}(X,Y)$.  $U$ denotes a closed subspace of $X$.  Statements needing only
an inner product space, without completeness, say so.
\end{assumptions}

\begin{definition}[Inner product space, Hilbert space]
  \label{def_hilbert}
  A Hilbert space is a complete inner product space.
\end{definition}


\begin{theorem}[Cauchy--Schwarz]
  \label{thm_cauchy_schwarz}
  $|\langle x, y\rangle| \le \|x\|\,\|y\|$, in any inner product space.
\end{theorem}


\begin{lemma}[Pythagoras and the parallelogram law]
  \label{lem_pythagoras}
  In an inner product space,
  $\|x+y\|^2 = \|x\|^2 + 2\operatorname{Re}\langle x,y\rangle + \|y\|^2$; hence
  $\|x+y\|^2 = \|x\|^2+\|y\|^2$ when $x \perp y$, and
  $\|x+y\|^2 + \|x-y\|^2 = 2(\|x\|^2+\|y\|^2)$.
\end{lemma}



\begin{definition}[Orthogonal complement, orthogonal projection]
  \label{def_orth_proj}
  $U^\perp = \{x \in X : \langle x,u\rangle = 0 \ \forall u \in U\}$, always a closed
  subspace.  For $U$ closed, $P_U \in \mathcal{L}(X,X)$ is the orthogonal projection
  onto $U$.
\end{definition}


\begin{theorem}[Orthogonal decomposition]
  \label{thm_orth_decomp}
  For $U$ a closed subspace of a Hilbert space $X$: $X = U \oplus U^\perp$ and
  $(U^\perp)^\perp = U$.  For a general subspace, $(U^\perp)^\perp = \overline{U}$.
\end{theorem}


\begin{lemma}[Orthogonal complement sees only the closure]
  \label{lem_orthogonal_closure}
  For any subspace $M$ of a Hilbert space, $M^\perp = (\overline{M})^\perp$.  Consequently
  $M$ is dense if and only if $M^\perp = \{0\}$.
\end{lemma}


\begin{proposition}[Properties of the orthogonal projection]
  \label{prop_proj_props}
  For $U$ a closed subspace: $P_U$ is self-adjoint; $\|P_U\| = 1$ if $U \ne \{0\}$;
  $\mathrm{Id}-P_U = P_{U^\perp}$; and $\|x-P_Ux\| = \min_{u \in U}\|x-u\|$.
\end{proposition}


\begin{lemma}[Idempotence and nesting]
  \label{lem_proj_nested_comp}
  For closed subspaces $U \subseteq W$ of a Hilbert space, $P_U^2 = P_U$ and
  $P_U P_W = P_U$.
\end{lemma}


\begin{lemma}[Characterisation of the orthogonal projection]
  \label{lem_proj_characterisation}
  Let $U \subseteq X$ be a \emph{closed} subspace of a Hilbert space and $x \in X$.  Then
  $z = P_Ux$ if and only if $z \in U$ and $z-x \in U^\perp$.
\end{lemma}
\begin{proof}
  Closedness of $U$ is exactly what makes both sides meaningful: a closed subspace of the
  complete space $X$ is complete.

  ($\Rightarrow$)  Let $z = P_Ux$; then $z \in U$ by \cref{def_orth_proj}.  Let $u \in U$
  with $u \ne 0$ (for $u = 0$ there is nothing to check) and $t \in \mathbb{K}$.  Since
  $z+tu \in U$, \cref{prop_proj_props} (minimality of the projection) gives
  $\|x-z\| \le \|x-z-tu\|$, while \cref{lem_pythagoras} applied to $-tu$ and $x-z$
  expands the right-hand side as
  \[
    \|x-z-tu\|^2
      = \|x-z\|^2 - 2\operatorname{Re}\bigl(\bar t\,\langle u, x-z\rangle\bigr)
        + |t|^2\|u\|^2 .
  \]
  Choosing $t = \langle u, x-z\rangle/\|u\|^2$ turns this into
  $\|x-z\|^2 - |\langle u, x-z\rangle|^2/\|u\|^2$, so minimality forces
  $\langle u, x-z\rangle = 0$.  As $u \in U$ was arbitrary, $x - z \in U^\perp$, hence also
  $z - x \in U^\perp$.

  ($\Leftarrow$)  Suppose $z \in U$ and $z-x \in U^\perp$, and set $w = z - P_Ux$.  Both
  $z$ and $P_Ux$ lie in $U$, so $w \in U$; and $w = (z-x) + (x - P_Ux)$ is a sum of two
  elements of $U^\perp$, the second by the previous paragraph, so $w \in U^\perp$.  Then
  $\|w\|^2 = \langle w,w\rangle = 0$, i.e.\ $w = 0$, that is $z = P_Ux$.
\end{proof}


\begin{lemma}[Dimension counting]
  \label{lem_dim_count}
  Let $X$ be a Hilbert space and let $U, V \subseteq X$ be subspaces with $U$
  \emph{finite-dimensional} and $\dim U < \dim V$.  Then $V \cap U^\perp \ne \{0\}$.
\end{lemma}
\begin{proof}
  Since $U$ is finite-dimensional
  it is closed (\cref{lem_findim_closed}) and complete, so the orthogonal projection $P_U$
  exists and $x - P_Ux \in U^\perp$ for every $x$ (\cref{thm_orth_decomp},
  \cref{lem_proj_characterisation}).  The completeness required is that of $U$ alone --- so the ambient space need not be complete, and no
  codimension count in $X$, which may be infinite-dimensional, is involved.

  Consider the linear map $S \colon V \to U$, $Sv = P_Uv$, between finite-dimensional
  spaces.  Its range is a subspace of $U$, so
  $\dim \mathcal{R}(S) \le \dim U$, and
  rank--nullity gives
  \[
    \dim \mathcal{N}(S) = \dim V - \dim \mathcal{R}(S) \ge \dim V - \dim U \ge 1 .
  \]
  In particular $\mathcal{N}(S) \ne \{0\}$; pick $v \in V$, $v \ne 0$, with $P_Uv = 0$.
  Applying \cref{lem_proj_characterisation} with $x = v$ and $z = 0 = P_Uv$ yields
  $0 - v \in U^\perp$, and $U^\perp$ is a subspace, so $v \in U^\perp$.  Hence
  $0 \ne v \in V \cap U^\perp$.
\end{proof}







\begin{theorem}[Fr\'echet--Riesz representation]
  \label{thm_frechet_riesz}
  For $\lambda \in X'$ there is a unique $z_\lambda \in X$ with
  $\lambda(x) = \langle z_\lambda, x\rangle$ for all $x$, and
  $\|\lambda\|_{X'} = \|z_\lambda\|$.  The map $\lambda \mapsto z_\lambda$ is a
  conjugate-linear isometric bijection $X' \to X$.
\end{theorem}


\begin{definition}[Adjoint]
  \label{def_adjoint}
  For $T \in \mathcal{L}(X,Y)$ between Hilbert spaces, $T^* \in \mathcal{L}(Y,X)$ is
  determined by $\langle T^*y, x\rangle_X = \langle y, Tx\rangle_Y$.
\end{definition}


\begin{proposition}[Properties of the adjoint]
  \label{prop_adjoint_props}
  $(T^*)^* = T$; $\|T^*\| = \|T\|$; $\|T^*T\| = \|T\|^2$; $(\alpha T)^* = \bar\alpha T^*$;
  $(ST)^* = T^*S^*$.
\end{proposition}


\begin{theorem}[The four subspace relations]
  \label{thm_four_subspaces}
  For $T \in \mathcal{L}(X,Y)$ between Hilbert spaces,
  \[
    \mathcal{R}(T)^\perp = \mathcal{N}(T^*), \qquad
    \mathcal{N}(T)^\perp = \overline{\mathcal{R}(T^*)} .
  \]
\end{theorem}


\begin{lemma}[Closed range transfers to the adjoint]
  \label{lem_closed_range_adjoint}
  Let $T \in \mathcal{L}(X,Y)$ between Hilbert spaces have closed range.  Then
  \[
    \mathcal{R}(T^*) = \mathcal{N}(T)^\perp ,
  \]
  and in particular $\mathcal{R}(T^*)$ is closed.
\end{lemma}
\begin{proof}
  Write $N := \mathcal{N}(T)$, closed by \cref{lem_ker_closed}.  Since $N^\perp$ is an
  orthogonal complement it is closed, so the displayed identity gives the ``in
  particular'' at once.

  \emph{The inclusion $\subseteq$.}  \cref{thm_four_subspaces} gives
  $N^\perp = \overline{\mathcal{R}(T^*)} \supseteq \mathcal{R}(T^*)$, with no hypothesis
  on the range of $T$.  Only the reverse inclusion has content.

  \emph{$T$ is bounded below on $N^\perp$.}  Both spaces are Banach and
  $\mathcal{R}(T)$ is closed, so \cref{thm_closed_range} supplies $c > 0$ with
  $\|Tx\|_Y \ge c\inf_{z \in N}\|x+z\|_X$ for every $x \in X$.  For $u \in N^\perp$ and
  $z \in N$, \cref{lem_pythagoras} gives $\|u+z\|^2 = \|u\|^2 + \|z\|^2 \ge \|u\|^2$, so
  that infimum equals $\|u\|$, attained at $z = 0$.  Hence
  \begin{equation}
    \|Tu\|_Y \;\ge\; c\,\|u\|_X \quad \text{for all } u \in N^\perp . \tag{$*$}
  \end{equation}

  \emph{Every $w \in \mathcal{R}(T)$ has a unique preimage in $N^\perp$.}  Given
  $w = Tv$, decompose $v = u + z$ with $u \in N^\perp$, $z \in N$ by
  \cref{thm_orth_decomp}; then $Tu = Tv = w$.  If $Tu = Tu'$ with $u,u' \in N^\perp$ then
  $u - u' \in N \cap N^\perp = \{0\}$ by \cref{thm_orth_decomp}.  Write $u = Sw$ for this
  preimage; $S : \mathcal{R}(T) \to N^\perp$ is linear, being the inverse of the linear
  bijection $u \mapsto Tu$.

  \emph{The inclusion $\supseteq$.}  Let $x \in N^\perp$ and define
  $f : \mathcal{R}(T) \to \mathbb{K}$ by $f(w) := \langle x, Sw\rangle_X$, which is linear
  because $S$ is and the inner product is linear in its second argument.  By
  \cref{thm_cauchy_schwarz} and $(*)$,
  \[
    |f(w)| \;\le\; \|x\|_X\,\|Sw\|_X \;\le\; c^{-1}\|x\|_X\,\|w\|_Y ,
  \]
  so $f$ is bounded.  As $\mathcal{R}(T)$ is a closed subspace of the Hilbert space $Y$
  it is itself a Hilbert space, so \cref{thm_frechet_riesz} yields
  $y \in \mathcal{R}(T) \subseteq Y$ with $f(w) = \langle y, w\rangle_Y$ on
  $\mathcal{R}(T)$.  Taking $w = Tu$ with $u \in N^\perp$, so that $Sw = u$, and using
  \cref{def_adjoint},
  \[
    \langle x, u\rangle_X \;=\; f(Tu) \;=\; \langle y, Tu\rangle_Y
      \;=\; \langle T^*y, u\rangle_X .
  \]
  Thus $\langle x - T^*y, u\rangle_X = 0$ for every $u \in N^\perp$, i.e.\
  $x - T^*y \in (N^\perp)^\perp = N$ by \cref{thm_orth_decomp}.  But $x \in N^\perp$ by
  choice and $T^*y \in \mathcal{R}(T^*) \subseteq N^\perp$ by the first step, so
  $x - T^*y$ lies in $N \cap N^\perp = \{0\}$.  Hence $x = T^*y \in \mathcal{R}(T^*)$.
\end{proof}




\begin{lemma}[The normal operator has the same kernel and co-range]
  \label{lem_kernel_normal}
  For $T \in \mathcal{L}(X,Y)$ between Hilbert spaces,
  \[
    \mathcal{N}(T^*T) = \mathcal{N}(T),
    \qquad
    \overline{\mathcal{R}(T^*T)} = \overline{\mathcal{R}(T^*)} .
  \]
\end{lemma}
\begin{proof}
  Throughout, $T^*T \in \mathcal{L}(X,X)$ and $T^* \in \mathcal{L}(Y,X)$, so every
  orthogonal complement below is taken in $X$.

  For the first identity, $\mathcal{N}(T) \subseteq \mathcal{N}(T^*T)$ is immediate.
  Conversely, if $T^*Tx = 0$ then, by \cref{def_adjoint} with $y = Tx$,
  \[
    0 = \langle T^*Tx, x\rangle_X = \langle Tx, Tx\rangle_Y = \|Tx\|^2 ,
  \]
  so $x \in \mathcal{N}(T)$.

  For the second, set $S = T^*T \in \mathcal{L}(X,X)$.  By \cref{prop_adjoint_props},
  $S^* = (T^*T)^* = T^*(T^*)^* = T^*T = S$.  Applying \cref{thm_four_subspaces} to $S$,
  viewed as an operator $X \to X$ so that $\mathcal{N}(S)^\perp$ and
  $\overline{\mathcal{R}(S^*)}$ are both subspaces of $X$, gives
  \[
    \overline{\mathcal{R}(T^*T)} = \overline{\mathcal{R}(S^*)}
      = \mathcal{N}(S)^\perp = \mathcal{N}(T^*T)^\perp = \mathcal{N}(T)^\perp ,
  \]
  the last step by the first identity.  Applying \cref{thm_four_subspaces} to $T$ itself
  gives $\mathcal{N}(T)^\perp = \overline{\mathcal{R}(T^*)}$, again a complement inside
  $X$.  Combining the two displays yields
  $\overline{\mathcal{R}(T^*T)} = \overline{\mathcal{R}(T^*)}$.
\end{proof}



\chapter{Orthonormal series}
\label{ch_04_orthonormal_series}

\begin{assumptions}
  \label{asm_orthonormal_series}
  $X$ is a Hilbert space over $\mathbb{K}$;
$\{e_n\}$, $\{u_i\}$ denote orthonormal families in $X$, indexed by an arbitrary type
unless stated otherwise.  Completeness is used throughout.
\end{assumptions}

\begin{theorem}[Bessel's inequality]
  \label{thm_bessel}
  For an orthonormal system $\{u_i\}$ and any $x$,
  $\sum_i |\langle u_i,x\rangle|^2 \le \|x\|^2$, and the family is summable.
\end{theorem}


\begin{lemma}[Riesz--Fischer for orthonormal systems]
  \label{lem_riesz_fischer}
  For an orthonormal system $\{e_n\}$ and scalars $c_n$, the series $\sum_n c_ne_n$
  converges in $X$ \emph{if and only if} $\sum_n|c_n|^2 < \infty$; in that case
  $\|\sum_n c_ne_n\|^2 = \sum_n|c_n|^2$ and the coefficients of the sum are the $c_n$.
\end{lemma}
\begin{proof}
  The index set may be countably infinite and carries no order, so $\sum_n$ is read as an
  unordered sum: convergence means that the net $F \mapsto s_F := \sum_{n \in F}c_ne_n$,
  indexed by the finite subsets $F$ ordered by inclusion, converges.  For a family indexed by $\mathbb{N}$ this is unconditional
  convergence of the partial sums, and it is the notion under which the three assertions
  below hold without choosing an enumeration.  If the index set is finite (in particular
  empty) every sum is a finite sum, both conditions hold, and the assertions reduce to the
  finite identity of the next paragraph.

  \emph{Finite identity.}  For $F$ finite, orthonormality and \cref{lem_pythagoras}
  (induction on $|F|$, the new term being orthogonal to the previous ones) give
  $\|s_F\|^2 = \sum_{n \in F}|c_n|^2$.  The inner product is linear in its second slot,
  so also $\langle e_m, s_F\rangle = c_m$ for every $m \in F$.

  ($\Rightarrow$)  Suppose $s_F \to x$.  Fix an index $m$.  For every finite $F \ni m$,
  \[
    |\langle e_m, x\rangle - c_m| = |\langle e_m, x - s_F\rangle|
      \le \|e_m\|\,\|x-s_F\| = \|x - s_F\|
  \]
  by \cref{thm_cauchy_schwarz}, and the right-hand side tends to $0$ along the net; hence
  $\langle e_m,x\rangle = c_m$.  This proves the last assertion --- the coefficients of the
  sum are the $c_n$, so two expansions of the same vector in $\{e_n\}$ have the same
  coefficients --- and \cref{thm_bessel} then gives
  $\sum_n|c_n|^2 = \sum_n|\langle e_n,x\rangle|^2 \le \|x\|^2 < \infty$.

  ($\Leftarrow$)  Suppose $\sum_n|c_n|^2 < \infty$.  Then the net of finite partial sums of
  the nonnegative reals $|c_n|^2$ converges, hence is Cauchy, so for $\varepsilon>0$ there
  is a finite $F_0$ with $\sum_{n \in F}|c_n|^2 < \varepsilon^2$ for every finite $F$
  disjoint from $F_0$.  By the finite
  identity $\|s_F\| < \varepsilon$ for such $F$, i.e.\ the net $(s_F)$ is Cauchy.  Since
  $X$ is complete (\cref{def_hilbert}) it converges.  Completeness is used here and nowhere
  else, and cannot be dropped: in a dense proper subspace of a Hilbert space one can choose
  square-summable coefficients whose sum lies outside.

  \emph{Norm identity.}  If $s_F \to x$ then $\|s_F\|^2 \to \|x\|^2$ by continuity of the
  norm, while $\|s_F\|^2 = \sum_{n\in F}|c_n|^2 \to \sum_n|c_n|^2$ by the finite identity
  and the summability just established; hence $\|x\|^2 = \sum_n|c_n|^2$.
\end{proof}


\begin{lemma}[A bounded operator commutes with a series]
  \label{lem_operator_through_series}
  If $T \in \mathcal{L}(X,Y)$ between normed spaces and $\sum_n x_n$ converges, then
  $T(\sum_n x_n) = \sum_n Tx_n$.
\end{lemma}


\begin{lemma}[The inner product commutes with a series]
  \label{lem_inner_through_series}
  If $\sum_n x_n$ converges then
  $\langle \sum_n x_n, z\rangle = \sum_n \langle x_n, z\rangle$.  Consequently the
  coefficients of an expansion in an orthonormal system are unique.
\end{lemma}
\begin{proof}
  Our convention is that inner product is conjugate-linear in the \emph{first} argument and linear in
  the second, so the slot through which \cref{lem_operator_through_series} may be pushed
  directly is the second: for fixed $z$ the map $y \mapsto \langle z,y\rangle$ is linear,
  and $|\langle z,y\rangle| \le \|z\|\,\|y\|$ by \cref{thm_cauchy_schwarz}, so it is an
  element of $\mathcal{L}(X,\mathbb{K})$ of norm at most $\|z\|$.
  Hence, whenever $\sum_n x_n$ converges,
  \begin{equation}
    \Bigl\langle z, \sum_n x_n\Bigr\rangle = \sum_n \langle z, x_n\rangle . \tag{$*$}
  \end{equation}

  For the first slot, conjugate: complex conjugation is an isometric additive automorphism
  of $\mathbb{K}$, hence continuous, so it carries convergent unordered sums to convergent
  unordered sums; with
  $\langle a,b\rangle = \overline{\langle b,a\rangle}$ this
  gives
  \[
    \Bigl\langle \sum_n x_n, z\Bigr\rangle
      = \overline{\Bigl\langle z, \sum_n x_n\Bigr\rangle}
      = \overline{\sum_n \langle z,x_n\rangle}
      = \sum_n \overline{\langle z,x_n\rangle}
      = \sum_n \langle x_n, z\rangle .
  \]

  For the consequence, let $\{e_n\}$ be an orthonormal system and suppose that
  $\sum_n c_ne_n = \sum_n d_ne_n = x$, both series convergent.  Apply $(*)$ with
  $z = e_m$ to $x_n = c_ne_n$; the second slot is linear, so no conjugation appears on the
  coefficients and
  \[
    \langle e_m, x\rangle = \sum_n c_n\langle e_m,e_n\rangle = c_m ,
  \]
  the family on the right being supported at $n=m$; the same computation gives
  $\langle e_m,x\rangle = d_m$.  Hence $c_m = d_m$ for every $m$.
\end{proof}


\begin{definition}[Orthonormal basis]
  \label{def_hilbert_basis}
  An orthonormal basis of $X$ is a maximal orthonormal system; equivalently an orthonormal
  family whose closed span is $X$ (\cref{lem_onb_criterion}).
\end{definition}


\begin{theorem}[Orthonormal bases: existence, expansion, Parseval]
  \label{thm_onb}
  Every Hilbert space has an orthonormal basis; $x = \sum_i \langle u_i,x\rangle u_i$; and
  $\sum_i \langle x,u_i\rangle\langle u_i,y\rangle = \langle x,y\rangle$.
\end{theorem}


\begin{lemma}[Recognising an orthonormal basis]
  \label{lem_onb_criterion}
  Let $\{u_i\}_{i \in \iota}$ be an orthonormal family in a Hilbert space $X$.  The
  following are equivalent, and each makes $\{u_i\}$ an orthonormal basis of $X$:
  \begin{enumerate}
    \item $\overline{\operatorname{span}}\{u_i\} = X$;
    \item $\{u_i\}^\perp = \{0\}$;
    \item $\{u_i\}$ is maximal among orthonormal families.
  \end{enumerate}
\end{lemma}



\begin{theorem}[Projection onto a closed subspace via a basis]
  \label{thm_proj_expansion}
  Let $X$ be complete and let $U \subseteq X$ be a complete subspace.  If $\{u_j\}$ is an orthonormal basis
  of $U$, then $P_Ux = \sum_j \langle u_j,x\rangle u_j$.
\end{theorem}


\chapter{Weak convergence}
\label{ch_05_weak_convergence}

\begin{assumptions}
  \label{asm_weak_convergence}
  $X$, $Y$ are normed spaces; weak convergence is
defined through $X'$, so no inner product is needed except where a Hilbert space is named.
\cref{thm_radon_riesz} is genuinely Hilbert; \cref{thm_weak_compactness} is genuinely
reflexive.
\end{assumptions}


\begin{definition}[Weak convergence]
  \label{def_weak_conv}
  A sequence $\{x_n\}$ in a normed space $X$ converges weakly to $x$, written
  $x_n \rightharpoonup x$, if $\langle x_n, f\rangle \to \langle x, f\rangle$ for every
  $f \in X'$.  In a Hilbert space this reads
  $\langle x_n, z\rangle \to \langle x, z\rangle$ for all $z \in X$.
\end{definition}



\begin{lemma}[Strong convergence implies weak convergence]
  \label{lem_strong_to_weak}
  If $x_n \to x$ in norm then $x_n \rightharpoonup x$.
\end{lemma}
\begin{proof}
  Let $f \in X'$.  By \cref{lem_isleast_opnorm} applied to $f$,
  \[
    |\langle x_n,f\rangle - \langle x,f\rangle| = |f(x_n-x)| \le \|f\|_{X'}\,\|x_n-x\|
    \longrightarrow 0 .
  \]
  Since $f$ was arbitrary, $x_n \rightharpoonup x$ by \cref{def_weak_conv}.  No
  completeness and no inner product are used.
\end{proof}


\begin{lemma}[Weakly convergent sequences are bounded]
  \label{lem_weak_bounded}
  If $x_n \rightharpoonup x$ in a normed space then $\sup_n\|x_n\| < \infty$.
\end{lemma}
\begin{proof}
  For $x \in X$ let $Jx \colon X' \to \mathbb{K}$ be evaluation, $(Jx)(f) = f(x)$.  It is
  linear in $f$ and $|(Jx)(f)| \le \|x\|\,\|f\|_{X'}$, so
  $Jx \in \mathcal{L}(X',\mathbb{K})$ with $\|Jx\| \le \|x\|$; conversely
  $\|Jx\| \ge \|x\|$ because for $x \ne 0$ \cref{cor_norming_functional} supplies $f$ with
  $\|f\| = 1$ and $f(x) = \|x\|$, while for $x = 0$ both sides vanish.  Hence
  $\|Jx\| = \|x\|$.

  Now let $x_n \rightharpoonup x$.  For each fixed $f \in X'$ the scalar sequence
  $(Jx_n)(f) = \langle x_n,f\rangle$ converges by \cref{def_weak_conv}, hence is bounded;
  so the family $\{Jx_n\}$ is pointwise bounded on $X'$.  The space $X'$ is a Banach space
  by \cref{lem_clm_complete} \emph{even though $X$ need not be} --- this is exactly why the
  argument is run on the dual side --- so \cref{thm_banach_steinhaus}, whose hypothesis is
  completeness of the domain, gives $\sup_n\|Jx_n\| < \infty$, i.e.\
  $\sup_n\|x_n\| < \infty$.
\end{proof}


\begin{lemma}[Bounded operators are weakly sequentially continuous]
  \label{lem_weak_seq_cont}
  If $T \in \mathcal{L}(X,Y)$ and $x_n \rightharpoonup x$, then
  $Tx_n \rightharpoonup Tx$.
\end{lemma}
\begin{proof}
  Let $g \in Y'$.  By \cref{def_dual_operator}, $T'g \in X'$ and
  $\langle Tx_n, g\rangle = \langle x_n, T'g\rangle$ for every $n$; the right-hand side
  converges to $\langle x, T'g\rangle = \langle Tx, g\rangle$ because
  $x_n \rightharpoonup x$ and $T'g$ is one admissible test functional in
  \cref{def_weak_conv}.  As $g \in Y'$ was arbitrary, $Tx_n \rightharpoonup Tx$.

  In the Hilbert case one may argue with \cref{def_adjoint} instead: for $z \in Y$,
  \[
    \langle Tx_n, z\rangle_Y = \overline{\langle z, Tx_n\rangle_Y}
      = \overline{\langle T^*z, x_n\rangle_X} = \langle x_n, T^*z\rangle_X
      \longrightarrow \langle x, T^*z\rangle_X = \langle Tx, z\rangle_Y .
  \]
\end{proof}


\begin{lemma}[Uniqueness of weak limits]
  \label{lem_weak_limit_unique}
  A sequence in a normed space has at most one weak limit.
\end{lemma}
\begin{proof}
  Suppose $x_n \rightharpoonup x$ and $x_n \rightharpoonup y$.  For every $f \in X'$ the
  scalar sequence $\langle x_n,f\rangle$ has both limits $\langle x,f\rangle$ and
  $\langle y,f\rangle$, so $f(x-y) = 0$.  If $x \ne y$, \cref{cor_norming_functional}
  applied to $x - y \ne 0$ gives $f \in X'$ with $f(x-y) = \|x-y\| > 0$, a contradiction.
  Hence $x = y$.
\end{proof}


\begin{theorem}[Weak sequential compactness of bounded sets]
  \label{thm_weak_compactness}
  Every bounded sequence in a Hilbert space has a weakly convergent subsequence.
\end{theorem}
\begin{proof}
  Let $\{x_n\}_{n \in \mathbb{N}}$ satisfy $M := \sup_n\|x_n\| < \infty$.

  \emph{Step 1: a separable closed subspace.}  Put
  $X_0 := \overline{\operatorname{span}}\{x_n : n \in \mathbb{N}\}$.  Let
  $Q \subseteq \mathbb{K}$ be a countable dense subset.  The finite $Q$-linear combinations
  of the $x_n$ form a countable set, dense in $\operatorname{span}\{x_n\}$ and hence in
  $X_0$, so $X_0$ is separable.  Being closed in the complete space $X$ it is itself
  complete, hence a Hilbert space.  If $X_0 = \{0\}$ then every $x_n = 0$ and the full
  sequence converges weakly to $0$; assume from now on $X_0 \ne \{0\}$.

  \emph{Step 2: a countable orthonormal basis.}  Being separable, the Hilbert space $X_0$
  has a countable orthonormal basis $\{u_j\}_{j \in J}$, with $J$
  finite or countably infinite.

  \emph{Step 3: diagonal extraction.}  By \cref{thm_cauchy_schwarz},
  $|\langle x_n, u_j\rangle| \le M$ for all $n,j$, so each scalar sequence
  $n \mapsto \langle x_n,u_j\rangle$ lies in the closed ball of radius $M$ in $\mathbb{K}$,
  which is compact because $\mathbb{K}$ is finite-dimensional over $\mathbb{R}$ and hence
  proper; every such sequence therefore has
  a convergent subsequence.  Enumerate
  $J = \{j_1,j_2,\dots\}$ and choose recursively infinite sets
  $\mathbb{N} \supseteq N_1 \supseteq N_2 \supseteq \cdots$ such that
  $\langle x_n, u_{j_m}\rangle$ converges as $n \to \infty$ along $N_m$.  Let $n_k$ be the
  least element of $N_k$ larger than $n_{k-1}$.  Then $(n_k)$ is strictly increasing and,
  for each $m$, $(n_k)_{k \ge m}$ runs inside $N_m$, so $\langle x_{n_k},u_{j_m}\rangle$
  converges as $k \to \infty$ for every $m$.  (If $J$ is finite the construction stops
  after finitely many stages and no diagonal is needed.)  Write $y_k := x_{n_k}$.

  \emph{Step 4: the limit functional.}  Since the inner product is linear in its second
  argument, $f_k := \langle y_k, \cdot\rangle$ restricted to $X_0$ is a linear functional,
  and $|f_k(z)| \le \|y_k\|\,\|z\| \le M\|z\|$ by \cref{thm_cauchy_schwarz}, so
  $f_k \in X_0'$ with $\sup_k\|f_k\| \le M$.  The family $\{u_j\}$ is total in $X_0$ by
  \cref{def_hilbert_basis}, and
  $f_k(u_j) = \langle y_k, u_j\rangle = \overline{\langle u_j, y_k\rangle}$ converges for
  every $j$ by Step 3, conjugation being continuous.  Both $X_0$ and $\mathbb{K}$ are
  Banach spaces, so \cref{prop_ptwise_limit_exists} yields $f \in X_0'$ with
  $f_k(z) \to f(z)$ for every $z \in X_0$.

  \emph{Step 5: representing the limit.}  Apply \cref{thm_frechet_riesz} on the Hilbert
  space $X_0$: there is $\bar x \in X_0$ with $f(z) = \langle \bar x, z\rangle$ for all
  $z \in X_0$.  Thus $\langle y_k, z\rangle \to \langle \bar x, z\rangle$ for $z \in X_0$.

  \emph{Step 6: from $X_0$ to $X$.}  By \cref{thm_orth_decomp},
  $X = X_0 \oplus X_0^\perp$; write $z = z_0 + z_1$ accordingly.  By
  \cref{def_orth_proj}, $z_1 \in X_0^\perp$ means $\langle z_1,u\rangle = 0$ for all
  $u \in X_0$, hence $\langle y_k,z_1\rangle = \overline{\langle z_1,y_k\rangle} = 0$ and
  likewise $\langle \bar x, z_1\rangle = 0$, since $y_k,\bar x \in X_0$.  Therefore
  \[
    \langle y_k, z\rangle = \langle y_k, z_0\rangle
      \longrightarrow \langle \bar x, z_0\rangle = \langle \bar x, z\rangle
      \qquad \text{for every } z \in X,
  \]
  which is the Hilbert reading of \cref{def_weak_conv}, equivalent to the dual reading by
  \cref{thm_frechet_riesz}.  Hence $y_k \rightharpoonup \bar x$, and $\{y_k\}$ is a
  subsequence of $\{x_n\}$.
\end{proof}



\begin{lemma}[Weak lower semicontinuity of the norm]
  \label{lem_weak_lsc}
  If $x_n \rightharpoonup x$ in a normed space then
  $\|x\| \le \liminf_{n\to\infty}\|x_n\|$.
\end{lemma}
\begin{proof}
  If $x = 0$ the claim is $0 \le \liminf_n\|x_n\|$, which holds since norms are
  nonnegative.

  If $x \ne 0$, take $f \in X'$ with $\|f\| = 1$ and $f(x) = \|x\|$ from
  \cref{cor_norming_functional}.  By \cref{def_weak_conv},
  $\langle x_n,f\rangle \to \langle x,f\rangle$, hence
  $|\langle x_n,f\rangle| \to |\langle x,f\rangle| = \|x\|$, and for a convergent sequence
  the $\liminf$ is the limit.  By \cref{lem_isleast_opnorm},
  $|\langle x_n,f\rangle| \le \|f\|\,\|x_n\| = \|x_n\|$ for every $n$, and $\liminf$ is
  monotone, so
  \[
    \|x\| = \liminf_{n\to\infty}|\langle x_n,f\rangle| \le \liminf_{n\to\infty}\|x_n\| .
  \]
  Only the normed structure is used.
\end{proof}


\begin{theorem}[Radon--Riesz property]
   \label{thm_radon_riesz}
  In a \emph{Hilbert} space, if $x_n \rightharpoonup x$ and $\|x_n\| \to \|x\|$ then
  $x_n \to x$ strongly.
\end{theorem}
\begin{proof}
  By \cref{lem_pythagoras} with the pair $(x_n, -x)$,
  \[
    \|x_n - x\|^2 = \|x_n\|^2 - 2\operatorname{Re}\langle x_n, x\rangle + \|x\|^2 .
  \]
  Testing the Hilbert reading of \cref{def_weak_conv} --- the dual formulation transported
  by \cref{thm_frechet_riesz} --- at the particular vector $z = x$ gives
  $\langle x_n, x\rangle \to \langle x,x\rangle = \|x\|^2$, hence
  $\operatorname{Re}\langle x_n,x\rangle \to \|x\|^2$.  Together with the hypothesis
  $\|x_n\| \to \|x\|$, and so $\|x_n\|^2 \to \|x\|^2$,
  \[
    \|x_n-x\|^2 \longrightarrow \|x\|^2 - 2\|x\|^2 + \|x\|^2 = 0 .
  \]
  The inner product is what makes this work: it is the cross term
  $\operatorname{Re}\langle x_n,x\rangle$, controlled by weak convergence alone, that
  cancels the two norms.
\end{proof}



\chapter{Compact operators}
\label{ch_06_compact_operators}

\begin{assumptions}
  \label{asm_compact_operators}
  $X$, $Y$, $Z$ are normed spaces over $\mathbb{K}$
unless completeness is stated, and $K \in \mathcal{K}(X,Y)$ denotes a compact operator.

\end{assumptions}

\begin{definition}[Compact operator]
  \label{def_compact_op}
  $T \colon X \to Y$ is compact if the image of some neighbourhood of $0$ has compact
  closure; equivalently, if the image of the unit ball is relatively compact.
\end{definition}


\begin{lemma}[Sequential characterisation]
  \label{lem_compact_seq}
  $T$ is compact if and only if every bounded sequence $\{x_n\} \subseteq X$ has a
  subsequence with $\{Tx_{n_k}\}$ convergent in $Y$.
\end{lemma}
\begin{proof}
  Write $B$ for the closed unit ball of $X$ and $W := \overline{T(B)}$.  By
  \cref{def_compact_op}, $T$ is compact exactly when $W$ is compact.  Throughout, $X$ and
  $Y$ are only normed; no completeness enters.

  ($\Rightarrow$)  Let $\{x_n\}$ be bounded, say $\|x_n\| \le R$.  If $R = 0$ then
  $Tx_n = 0$ for all $n$ and the full sequence converges.  Otherwise $x_n/R \in B$, so
  $Tx_n \in R\,W$, and $R\,W$ is compact, being the image of the compact $W$ under the
  homeomorphism $y \mapsto Ry$.  A compact subset of a metric space is sequentially
  compact, so $\{Tx_n\}$ has a convergent subsequence.

  ($\Leftarrow$)  We show $W$ is sequentially compact.  Let $\{w_m\} \subseteq W$.  Each
  $w_m$ is a limit of points of $T(B)$, so there is $x_m \in B$ with
  $\|Tx_m - w_m\| < 1/(m+1)$.  The sequence $\{x_m\}$ is bounded, so by hypothesis some
  subsequence has $Tx_{m_i} \to v$ in $Y$; then $w_{m_i} \to v$ as well, and $v \in W$
  because $W$ is closed.  Thus $W$ is sequentially compact, hence compact --- the one place
  where the metric structure of $Y$ is used --- and $T$ is compact by \cref{def_compact_op}.
\end{proof}


\begin{lemma}[Finite rank implies compact]
  \label{lem_finrank_compact}
  A bounded linear operator between normed spaces with finite-dimensional range is
  compact.
\end{lemma}
\begin{proof}
  Let $F := \mathcal{R}(T)$, a subspace of $Y$ with $\dim F < \infty$, and let $B$ be the
  closed unit ball of $X$.  By \cref{lem_isleast_opnorm}, $\|Tx\| \le \|T\|$ for $x \in B$,
  so $T(B)$ is contained in the closed ball $B_F$ of radius $\|T\|$ of the normed space
  $F$.  By \cref{lem_findim_closed}, $F$ is closed in $Y$, so the closure of $T(B)$ taken
  in $Y$ is contained in $F$ and agrees with its closure taken in $F$.

  By \cref{thm_riesz_ball} the closed unit ball of $F$ is compact, $\dim F$ being finite;
  scaling by $\|T\|$ (a homeomorphism of $F$, and the singleton $\{0\}$ if $T = 0$) shows
  $B_F$ is compact.  Hence $\overline{T(B)}$ is a closed subset of a compact set, therefore
  compact, and $T$ is compact by \cref{def_compact_op}.  The degenerate cases $T = 0$,
  $F = \{0\}$ and $X = \{0\}$ are covered.  No completeness is used.
\end{proof}


\begin{lemma}[$\mathcal{K}(X,Y)$ is a closed subspace of $\mathcal{L}(X,Y)$]
  \label{lem_compact_closed}
  The compact operators form a submodule of $\mathcal{L}(X,Y)$, closed when $Y$ is
  complete.
\end{lemma}


\begin{lemma}[Composition]
  \label{lem_compact_comp}
  If $T \in \mathcal{K}(X,Y)$ then $AT$ is compact for every $A \in \mathcal{L}(Y,Z)$ and
  $TB$ is compact for every $B \in \mathcal{L}(Z,X)$.  No completeness is needed, and
  $A$, $B$ are not assumed compact.
\end{lemma}


\begin{theorem}[The identity is compact only in finite dimensions]
  \label{thm_id_compact}
  Let $\mathbb{K}$ be locally compact.  Then $\mathrm{Id}_X$ is compact if and only
  if $\dim X < \infty$.  No completeness of $X$ is needed, and $X$ need not even be
  normed: a Hausdorff topological $\mathbb{K}$-module suffices.
\end{theorem}


\begin{theorem}[Compact operators have non-closed range]
  \label{thm_compact_range_not_closed}
  Let $X,Y$ be Banach spaces over $\mathbb{K} \in \{\mathbb{R},\mathbb{C}\}$ and
  $K \in \mathcal{K}(X,Y)$ with $\dim\mathcal{R}(K) = \infty$.  Then $\mathcal{R}(K)$ is
  not closed in $Y$.
\end{theorem}
\begin{proof}
  Suppose, for contradiction, that $\mathcal{R}(K)$ is closed.  Since $X$ and $Y$ are
  Banach, the ``$\Rightarrow$'' direction of \cref{thm_closed_range} applies --- the only
  place completeness of both spaces is used --- and gives $c > 0$ with
  \[
    \|Kx\| \;\ge\; c \inf_{z \in \mathcal{N}(K)}\|x+z\| \qquad \text{for all } x \in X .
  \]

  Let $B_0$ be the closed unit ball of the normed space $\mathcal{R}(K)$ and let
  $y \in B_0$, say $y = Kx$.  If $y = 0$ take the representative $0$.  Otherwise
  $\inf_{z \in \mathcal{N}(K)}\|x+z\| \le \|y\|/c \le 1/c < 2/c$, so there is
  $z \in \mathcal{N}(K)$ with $\|x+z\| \le 2/c$, and $K(x+z) = Kx = y$.  Hence
  $B_0 \subseteq (2/c)\,K(B_X)$ for $B_X$ the closed unit ball of $X$, and therefore
  $B_0 \subseteq (2/c)\,\overline{K(B_X)}$, a compact set by \cref{def_compact_op}
  (rescaling is a homeomorphism).  Moreover $B_0$ is closed in $\mathcal{R}(K)$, which is
  closed in $Y$ by assumption, so $B_0$ is closed in $Y$ and, as a closed subset of a
  compact set, compact.

  Thus the closed unit ball of the normed space $\mathcal{R}(K)$ is compact, and
  \cref{thm_riesz_ball} forces $\dim\mathcal{R}(K) < \infty$, contradicting the
  hypothesis.

  The argument is routed through \cref{thm_riesz_ball} rather than \cref{thm_id_compact},
  and that choice is what keeps it available over an abstract
  scalar field.  Completeness of $\mathcal{R}(K)$, which the identity-operator
  route would additionally need, is therefore not used, though it does hold: under the
  contradiction hypothesis $\mathcal{R}(K)$ is closed in the Banach space $Y$.
\end{proof}


\begin{theorem}[Complete continuity]
  \label{thm_complete_continuity}
  Let $X,Y$ be normed spaces and $T \in \mathcal{L}(X,Y)$.  If $T$ is compact then $T$
  maps weakly convergent sequences to norm convergent ones.  The converse holds when $X$
  is a Hilbert space, so that there the two are equivalent.
\end{theorem}
\begin{proof}
  \emph{Compact $\Rightarrow$ completely continuous.}  Let $x_n \rightharpoonup x$ and
  suppose $Tx_n \not\to Tx$ in norm.  Then there are $\varepsilon > 0$ and a subsequence
  with $\|Tx_{n_k} - Tx\| \ge \varepsilon$ for all $k$.  By \cref{lem_weak_bounded} the
  sequence $\{x_n\}$ is bounded, so \cref{lem_compact_seq} yields a further subsequence
  with $Tx_{n_{k_l}} \to w$ in norm, and hence $Tx_{n_{k_l}} \rightharpoonup w$ by
  \cref{lem_strong_to_weak}.  On the other hand \cref{lem_weak_seq_cont} gives
  $Tx_n \rightharpoonup Tx$, and every subsequence of a weakly convergent sequence
  converges weakly to the same limit --- immediate from \cref{def_weak_conv}, since each
  scalar sequence $\langle Tx_n,g\rangle$ converges.  So
  $Tx_{n_{k_l}} \rightharpoonup Tx$, and \cref{lem_weak_limit_unique} forces $w = Tx$,
  contradicting $\|Tx_{n_{k_l}} - Tx\| \ge \varepsilon$.  This direction is general: it
  uses no completeness of $X$, because the Banach--Steinhaus argument behind
  \cref{lem_weak_bounded} is run on $X'$, which is complete for any normed $X$.

  \emph{Completely continuous $\Rightarrow$ compact, for $X$ a Hilbert space.}  Let
  $\{x_n\}$ be a bounded sequence in $X$.  By \cref{thm_weak_compactness} there is a
  subsequence with $x_{n_k} \rightharpoonup \bar x$, and the hypothesis turns this into
  $Tx_{n_k} \to T\bar x$ in norm.  Thus every bounded sequence has a subsequence with
  norm-convergent image, so $T$ is compact by \cref{lem_compact_seq}.  The Hilbert
  hypothesis enters in exactly one place, the production of the weakly convergent
  subsequence; nothing else in this direction needs it.  The same argument works verbatim
  in a reflexive Banach space once a weak compactness theorem is available there, and it
  is false without reflexivity.
\end{proof}


\begin{theorem}[Schauder]
  \label{thm_schauder}
  Let $X$ be a normed space, $Y$ a Banach space, and $T \in \mathcal{L}(X,Y)$.  Then $T$
  is compact if and only if the dual operator $T' \in \mathcal{L}(Y',X')$ is compact.  In
  a Hilbert space the same holds for the adjoint $T^*$.
\end{theorem}
\begin{proof}
  Write $B_X$, $B_{Y'}$, $B_{X''}$ for the closed unit balls of the spaces indicated.  The
  proof avoids Arzel\`a--Ascoli by running its two ingredients --- a diagonal extraction and
  a finite net --- directly.

  \emph{Step A: $T$ compact $\Rightarrow$ $T'$ compact.}  No completeness is needed here.
  Let $S := \overline{T(B_X)}$, compact by \cref{def_compact_op}.  Let
  $\{g_n\} \subseteq Y'$ be bounded, $\|g_n\| \le R$; if $R = 0$ all $T'g_n$ vanish, so
  assume $R > 0$ and, after dividing by $R$, that $\|g_n\| \le 1$.

  Being compact, $S$ is totally bounded: for each
  $m \ge 1$ choose a finite $F_m \subseteq S$ with
  $S \subseteq \bigcup_{y \in F_m}\{\,y' : \|y'-y\| < 1/m\,\}$, and set
  $D := \bigcup_m F_m$, a countable subset of $S$.  For each $y \in D$ the scalar sequence
  $g_n(y)$ is bounded by $\|y\|$, hence has a convergent subsequence, the closed balls of
  $\mathbb{K}$ being compact.  Enumerating $D$, extracting nested subsequences
  and taking the diagonal produces a subsequence $(h_k)$ of $(g_n)$ such that $h_k(y)$
  converges for every $y \in D$.

  We claim $(h_k)$ is uniformly Cauchy on $S$.  Given $\varepsilon > 0$ pick $m$ with
  $1/m < \varepsilon/3$ and then $K$ such that $|h_k(y')-h_l(y')| < \varepsilon/3$ for all
  $k,l \ge K$ and all $y'$ in the finite set $F_m$.  For $y \in S$ choose $y' \in F_m$ with
  $\|y-y'\| < 1/m$; then, using $\|h_k\|,\|h_l\| \le 1$ and \cref{lem_isleast_opnorm},
  \[
    |h_k(y)-h_l(y)| \le |h_k(y-y')| + |h_k(y')-h_l(y')| + |h_l(y'-y)|
      < \tfrac{\varepsilon}{3}+\tfrac{\varepsilon}{3}+\tfrac{\varepsilon}{3} .
  \]
  Since $T(B_X) \subseteq S$, \cref{def_dual_operator} and \cref{lem_opnorm_sup} give
  \[
    \|T'h_k - T'h_l\|_{X'} = \sup_{\|x\|\le 1}|h_k(Tx)-h_l(Tx)|
      \le \sup_{y \in S}|h_k(y)-h_l(y)| \le \varepsilon
  \]
  for all $k,l \ge K$.  So $(T'h_k)$ is Cauchy in $X'$, which is complete by
  \cref{lem_clm_complete}, hence convergent.  By \cref{lem_compact_seq}, $T'$ is compact.

  \emph{Step B: $T'$ compact $\Rightarrow$ $T$ compact.}  Apply Step A to
  $T' \in \mathcal{L}(Y',X')$: the bidual operator
  $T'' := (T')' \in \mathcal{L}(X'',Y'')$ is compact.  Let $J_X \colon X \to X''$,
  $(J_Xx)(f) := f(x)$, and likewise $J_Y$; these are linear, and isometric by
  \cref{cor_norming_functional}.  For $x \in X$ and $\varphi \in Y'$, two applications of
  \cref{def_dual_operator} give
  \[
    (T''J_Xx)(\varphi) = (J_Xx)(T'\varphi) = (T'\varphi)(x) = \varphi(Tx)
      = (J_Y Tx)(\varphi),
  \]
  so $J_Y \circ T = T'' \circ J_X$.  Hence
  $J_Y(T(B_X)) = T''(J_X(B_X)) \subseteq T''(B_{X''})$, which is relatively compact and
  therefore totally bounded; so $J_Y(T(B_X))$ is totally bounded, and $T(B_X)$ is totally
  bounded in $Y$ because $J_Y$ is an isometry.  Then $\overline{T(B_X)}$ is totally bounded and closed in the \emph{complete} space $Y$, hence
  complete, hence compact.  By
  \cref{def_compact_op}, $T$ is compact.  Completeness of $Y$ is used exactly here.

  \emph{Step C: the Hilbert adjoint.}  Let $\Phi_X \colon X' \to X$ and
  $\Phi_Y \colon Y' \to Y$ be the conjugate-linear isometric bijections of
  \cref{thm_frechet_riesz}, so $\Phi_Y^{-1}y = \langle y,\cdot\rangle$.  For $y \in Y$ and
  all $x \in X$,
  \[
    \bigl(T'(\Phi_Y^{-1}y)\bigr)(x) = \langle y, Tx\rangle_Y = \langle T^*y, x\rangle_X
  \]
  by \cref{def_adjoint}, whence $T^* = \Phi_X \circ T' \circ \Phi_Y^{-1}$.  Both $\Phi$'s
  are isometric bijections, hence homeomorphisms carrying unit ball onto unit ball, so
  $\overline{T^*(B_Y)} = \Phi_X\bigl(\overline{T'(B_{Y'})}\bigr)$ and one is compact if and
  only if the other is.  Combining with Steps A and B: $T$ compact $\iff$ $T'$ compact
  $\iff$ $T^*$ compact.
\end{proof}



\chapter{Spectral theory}
\label{ch_07_spectral_theory}

\begin{assumptions}
  \label{asm_spectral_theory}
  $X$ is a complex Banach space and
$T \in \mathcal{L}(X,X)$, except where a Hilbert space is named.
Complexity is essential
for \cref{thm_spectrum_compact}: the spectrum of a rotation of $\mathbb{R}^2$ is empty.
\end{assumptions}

\begin{definition}[Resolvent set, spectrum, point spectrum, spectral radius]
  \label{def_spectrum}
  $\rho(T) = \{\lambda : (\lambda\,\mathrm{Id}-T)^{-1} \in \mathcal{L}(X,X)\}$,
  $\sigma(T) = \mathbb{C}\setminus\rho(T)$,
  $\sigma_p(T) = \{\lambda : \mathcal{N}(\lambda\,\mathrm{Id}-T) \ne \{0\}\}$, and
  $r(T) = \sup\{|\lambda| : \lambda \in \sigma(T)\}$.
\end{definition}


\begin{theorem}[The spectrum is compact and bounded by the norm]
  \label{thm_spectrum_compact}
  $\rho(T)$ is open.  If the scalar field is proper then
  $\sigma(T)$ is compact; if $X \ne \{0\}$ then
  $\sigma(T) \subseteq \{|\lambda| \le \|T\|\}$, so $r(T) \le \|T\|$.  If moreover the
  scalars are $\mathbb{C}$ and $X \ne \{0\}$, then $\sigma(T) \ne \emptyset$.
\end{theorem}


\begin{theorem}[Gelfand's spectral radius formula]
  \label{thm_gelfand}
  Let $X$ be a \emph{complex} Banach space.  Then
  $r(T) = \lim_{n\to\infty}\|T^n\|^{1/n}$.
\end{theorem}


\begin{definition}[Self-adjoint operator]
  \label{def_selfadjoint}
  $T \in \mathcal{L}(X,X)$ on a Hilbert space is self-adjoint if $T^* = T$.
\end{definition}


\begin{definition}[Positive operator]
  \label{def_positive}
  $T \in \mathcal{L}(X,X)$ is \emph{positive} if it is symmetric and
  $\operatorname{Re}\langle Tx,x\rangle \ge 0$ for every $x$.
\end{definition}


\begin{lemma}[Eigenvalues of a self-adjoint operator are real]
  \label{lem_eigenvalue_real}
  Let $X$ be a complex Hilbert space and $T \in \mathcal{L}(X,X)$ self-adjoint.  Then
  every eigenvalue of $T$ is real.
\end{lemma}
\begin{proof}
  For every $x \in X$, self-adjointness (\cref{def_selfadjoint}) and conjugate symmetry
  give
  \[
    \langle x,Tx\rangle = \langle T^*x,x\rangle = \langle Tx,x\rangle
      = \overline{\langle x,Tx\rangle},
  \]
  so $\langle x,Tx\rangle \in \mathbb{R}$.  If $Tx = \mu x$ with $x \ne 0$, then
  $\langle x,Tx\rangle = \mu\|x\|^2$ is real and $\|x\|^2 > 0$, hence $\mu \in \mathbb{R}$.
\end{proof}




\begin{theorem}[Norm, numerical range and spectral radius]
  \label{thm_spectral_radius}
  For $T$ self-adjoint on a Hilbert space $X$,
  \[
    \|T\| = \sup_{\|x\|=1}\bigl|\langle Tx,x\rangle\bigr| = r(T) .
  \]
\end{theorem}



\begin{lemma}[The spectral radius is attained]
  \label{lem_spectral_radius_attained}
  Let $X$ be a complex Banach space, $X \ne \{0\}$, and $T \in \mathcal{L}(X,X)$.  Then
  there is $\lambda \in \sigma(T)$ with $|\lambda| = r(T)$.
\end{lemma}


\begin{lemma}[Spectrum of a compact operator]
  \label{lem_eigenvalue_null}
  Let $X$ be a Hilbert space over $\mathbb{K}$ and $K \in \mathcal{K}(X,X)$.  Then
  \begin{enumerate}
    \item for every $\varepsilon > 0$ the set
      $\{\lambda \in \sigma(K) : |\lambda| \ge \varepsilon\}$ is \emph{finite};
    \item consequently $\sigma(K)$ is at most countable and its only possible accumulation
      point is $0$;
    \item every nonzero $\lambda \in \sigma(K)$ is an eigenvalue, with
      $\dim\mathcal{N}(\lambda\,\mathrm{Id}-K) < \infty$.
  \end{enumerate}
\end{lemma}
\begin{proof}
  \emph{Finite multiplicity.}  Let $\lambda \ne 0$ and
  $E := \mathcal{N}(\lambda\,\mathrm{Id}-K)$, a closed subspace by \cref{lem_ker_closed}.
  Every $x$ in the closed unit ball $B_E$ of $E$ satisfies $x = \lambda^{-1}Kx$ with
  $\|x\| \le 1$, so $B_E \subseteq \lambda^{-1}\overline{K(B_X)} \cap E$, where $B_X$ is
  the closed unit ball of $X$.  The set $\overline{K(B_X)}$ is compact by
  \cref{def_compact_op} and $E$ is closed, so the right-hand side is compact and its closed
  subset $B_E$ is compact.  By \cref{thm_riesz_ball}, $\dim E < \infty$.

  \emph{Nonzero spectral points are eigenvalues.}  This is the Riesz--Schauder
  alternative: for compact $K$ and $\lambda \ne 0$, injectivity of
  $\lambda\,\mathrm{Id}-K$ forces surjectivity, so a non-eigenvalue $\lambda \ne 0$ lies in
  the resolvent set of \cref{def_spectrum}.  This clause is \emph{not} proved here.

  \emph{Only $0$ can accumulate, and countability.}
  Fix $\varepsilon > 0$ and suppose
  there were infinitely many distinct $\lambda_1,\lambda_2,\dots \in \sigma(K)$ with
  $|\lambda_j| \ge \varepsilon$.  Each is an eigenvalue by the previous paragraph; pick
  $x_j \ne 0$ with $Kx_j = \lambda_jx_j$.  Eigenvectors for distinct eigenvalues are
  linearly independent: if $x_1,\dots,x_m$ are independent and
  $x_{m+1} = \sum_{j\le m}c_jx_j$, then applying $K-\lambda_{m+1}\mathrm{Id}$ gives
  $\sum_{j\le m}c_j(\lambda_j-\lambda_{m+1})x_j = 0$, so every $c_j = 0$ and $x_{m+1} = 0$,
  a contradiction; induction gives the claim.  Hence
  $U_m := \operatorname{span}\{x_1,\dots,x_m\}$ is a strictly increasing chain of
  finite-dimensional, therefore closed (\cref{lem_findim_closed}), subspaces.  By
  \cref{lem_riesz_lemma} choose $y_m \in U_m$ with $\|y_m\| = 1$ and
  $\operatorname{dist}(y_m,U_{m-1}) \ge \tfrac12$.  For $y = \sum_{j\le m}c_jx_j \in U_m$,
  \[
    (\lambda_m\,\mathrm{Id}-K)y = \sum_{j\le m}c_j(\lambda_m-\lambda_j)x_j \in U_{m-1},
  \]
  the $j = m$ term vanishing.  So for $l < m$ the vector
  $z := (\lambda_m\,\mathrm{Id}-K)y_m + Ky_l$ lies in $U_{m-1}$ --- note
  $Ky_l \in U_l \subseteq U_{m-1}$ --- and
  \[
    \|Ky_m - Ky_l\| = \|\lambda_my_m - z\|
      = |\lambda_m|\,\bigl\|y_m - \lambda_m^{-1}z\bigr\|
      \ge |\lambda_m|\operatorname{dist}(y_m,U_{m-1}) \ge \tfrac{\varepsilon}{2} .
  \]
  Thus $\{Ky_m\}$ has no convergent subsequence although $\{y_m\}$ is bounded,
  contradicting \cref{lem_compact_seq}.  On a Hilbert space the orthogonal projection
  replaces \cref{lem_riesz_lemma} and does better: taking $y_m$ to be the normalised
  component of $x_m$ off $U_{m-1}$ makes it a unit vector \emph{orthogonal} to
  $U_{m-1}$, so $\operatorname{dist}(y_m,U_{m-1}) = 1$ and the separation constant is
  $\varepsilon$ rather than $\varepsilon/2$.  That is the route the formalisation takes;
  every other step is the one above.  Therefore
  $\{\lambda \in \sigma(K) : |\lambda| \ge \varepsilon\}$ is finite for every
  $\varepsilon > 0$, which is clause (1).  Hence $\sigma(K)\setminus\{0\}$ is the countable union over
  $k \in \mathbb{N}$ of these finite sets for $\varepsilon = 1/k$, so $\sigma(K)$ is at
  most countable, and no point of modulus $> 0$ can be an accumulation point of
  $\sigma(K)$.

  No step above uses complex scalars: each clause holds verbatim over $\mathbb{R}$, and the
  only completeness used is that of $X$, in the Riesz--Schauder step.  What the complex
  hypothesis buys is that $\sigma(K)$ is nonempty (\cref{thm_spectrum_compact}), which is
  used in \cref{thm_max_eigenvalue}, not here.
\end{proof}


\begin{theorem}[Existence of a largest eigenvalue]
  \label{thm_max_eigenvalue}
  A nonzero compact self-adjoint operator $K$ on a Hilbert space has an eigenvalue
  $\lambda_1$ with $|\lambda_1| = \|K\|$.
\end{theorem}
\begin{proof}
  Since $K \ne 0$ we have $X \ne \{0\}$ and $\|K\| > 0$.  By \cref{thm_spectral_radius},
  $r(K) = \|K\| > 0$, and by \cref{lem_spectral_radius_attained} there is
  $\lambda_1 \in \sigma(K)$ with $|\lambda_1| = r(K) = \|K\|$.

  In particular $\lambda_1 \ne 0$, and this is exactly what $K \ne 0$ buys: a nonzero
  self-adjoint operator has positive norm, so the spectral point of maximal modulus cannot
  be $0$.  Thus $\lambda_1$ is a nonzero point of the spectrum of the compact operator $K$
  (\cref{def_compact_op}), and \cref{lem_eigenvalue_null} makes every such point an
  eigenvalue: there is $u \ne 0$ with $Ku = \lambda_1u$ and $|\lambda_1| = \|K\|$.  By
  \cref{lem_eigenvalue_real}, $\lambda_1$ is real, so it is $\pm\|K\|$.

  Both cited spectral facts need $X$ to be a \emph{complex} Hilbert space --- nonemptiness
  of the spectrum in \cref{lem_spectral_radius_attained}, and the Fredholm alternative in
  \cref{lem_eigenvalue_null} --- which is the standing convention of this chapter.
\end{proof}


\begin{theorem}[Spectral theorem, qualitative form]
  \label{thm_spectral_qualitative}
  For compact self-adjoint $K$ on a Hilbert space the eigenspaces span a dense
  subspace.  Every eigenspace for a nonzero eigenvalue is finite dimensional by
  \cref{lem_eigenvalue_null}, which needs no self-adjointness.
\end{theorem}


\begin{theorem}[Spectral theorem, diagonalised form]
  \label{thm_spectral_diagonal}
  Let $K$ be compact and self-adjoint on a Hilbert space $X$.  Then there are an
  \emph{initial segment} $I \subseteq \mathbb{N}$ (finite, or all of $\mathbb{N}$), an
  orthonormal system $\{u_n\}_{n \in I}$ and nonzero reals
  $\{\lambda_n\}_{n \in I}$ with $|\lambda_1| \ge |\lambda_2| \ge \cdots$ such that
  \[
    Kx = \sum_{n \in I} \lambda_n \langle u_n,x\rangle\,u_n \qquad \text{for all } x \in X,
  \]
  and $\{u_n\}_{n\in I}$ is an orthonormal basis of $\overline{\mathcal{R}(K)}$.  If $I$
  is infinite then $\lambda_n \to 0$.  The list $\{\lambda_n\}$ repeats each nonzero
  eigenvalue according to the (finite) dimension of its eigenspace.
\end{theorem}
\begin{proof}
  \emph{Step 1: the nonzero eigenvalues form a countable set clustering only at $0$.}  Let
  $\Lambda$ be the set of nonzero eigenvalues of $K$ and put
  $E_\lambda := \mathcal{N}(K-\lambda\,\mathrm{Id})$.  Every $\lambda \in \Lambda$ is real
  by \cref{lem_eigenvalue_real}, and $d_\lambda := \dim E_\lambda < \infty$ by
  \cref{thm_spectral_qualitative}.  Fix $\varepsilon > 0$.  Since
  $\Lambda \subseteq \sigma(K)$, the set
  $\Lambda_\varepsilon := \{\lambda \in \Lambda : |\lambda| \ge \varepsilon\}$ is contained
  in $\{\lambda \in \sigma(K) : |\lambda| \ge \varepsilon\}$, which is \emph{finite} by
  \cref{lem_eigenvalue_null}(1).  So $\Lambda_\varepsilon$ is finite, and
  $\Lambda = \bigcup_{k\ge1}\Lambda_{1/k}$ is at most countable.

  \emph{Step 2: the ordered list with multiplicity.}  Each $E_\lambda$ is a
  finite-dimensional, hence complete, inner product space; by \cref{thm_onb} choose an
  orthonormal basis $u_{\lambda,1},\dots,u_{\lambda,d_\lambda}$ of it.  Put
  $S := \{(\lambda,j) : \lambda \in \Lambda,\ 1 \le j \le d_\lambda\}$, at most countable by
  Step 1, and give $(\lambda,j)$ the weight $|\lambda|$.  For every $\varepsilon > 0$ only
  finitely many elements of $S$ have weight $\ge \varepsilon$: there are finitely many such
  $\lambda$, and each contributes $d_\lambda < \infty$ elements.  Consequently
  \[
    s \longmapsto \#\{s' \in S : \operatorname{weight}(s') > \operatorname{weight}(s)\}
      + \bigl(\text{position of } s \text{ in its finite tie class}\bigr)
  \]
  is a well-defined bijection of $S$ onto an initial segment $I \subseteq \mathbb{N}$ ($I$
  finite if $S$ is finite, $I = \mathbb{N}$ otherwise) along which the weights are
  nonincreasing; both counts are finite by the preceding sentence.  Writing
  $(\lambda_n,u_n)$ for the eigenvalue and vector attached to the $n$-th element, we obtain
  nonzero reals with $|\lambda_1| \ge |\lambda_2| \ge \cdots$ in which each
  $\lambda \in \Lambda$ occurs exactly $d_\lambda = \dim E_\lambda$ times.  If $I$ is
  infinite then for each $\varepsilon > 0$ only finitely many $n$ have
  $|\lambda_n| \ge \varepsilon$, which is to say $\lambda_n \to 0$.

  \emph{Step 3: orthonormality.}  Vectors attached to one $\lambda$ are orthonormal by
  construction.  If $\lambda \ne \mu$ are eigenvalues (possibly $\mu = 0$),
  $u \in E_\lambda$ and $v \in E_\mu$, then both eigenvalues are real by
  \cref{lem_eigenvalue_real}, so no conjugate appears in
  \[
    \lambda\langle u,v\rangle = \langle Ku,v\rangle = \langle K^*u,v\rangle
      = \langle u,Kv\rangle = \mu\langle u,v\rangle ,
  \]
  and therefore $\langle u,v\rangle = 0$.  Hence $\{u_n\}_{n \in I}$ is orthonormal, and
  every $u_n$ is orthogonal to $\mathcal{N}(K)$.

  \emph{Step 4: the closed span is $\overline{\mathcal{R}(K)}$.}  Put
  $W := \overline{\operatorname{span}}\{u_n\}_{n\in I}$; by Step 2,
  $\operatorname{span}\{u_n\}_{n\in I} = \sum_{\lambda\in\Lambda}E_\lambda$.  Every
  eigenspace of $K$ is either $\mathcal{N}(K)$ or one of the $E_\lambda$, so
  \cref{thm_spectral_qualitative} says that $\mathcal{N}(K)+\operatorname{span}\{u_n\}$ is
  dense, that is,
  \[
    \mathcal{N}(K)^\perp \cap W^\perp
      = \bigl(\mathcal{N}(K) + \operatorname{span}\{u_n\}\bigr)^\perp = \{0\} .
  \]
  By Step 3, $W \subseteq \mathcal{N}(K)^\perp$.  Conversely let
  $x \in \mathcal{N}(K)^\perp$ and split $x = P_Wx + (x-P_Wx)$ using \cref{def_orth_proj}
  and \cref{thm_orth_decomp}.  Then $x-P_Wx \in W^\perp$, and also
  $x-P_Wx \in \mathcal{N}(K)^\perp$, since both $x$ and $P_Wx \in W$ lie in that subspace;
  the display forces $x = P_Wx \in W$.  Hence $W = \mathcal{N}(K)^\perp$, and since
  $K^* = K$, \cref{thm_four_subspaces} gives
  \[
    W = \mathcal{N}(K)^\perp = \overline{\mathcal{R}(K^*)} = \overline{\mathcal{R}(K)} .
  \]
  So $\{u_n\}$ is an orthonormal family in the Hilbert space $\overline{\mathcal{R}(K)}$
  whose closed span is all of it, and \cref{lem_onb_criterion}(1) makes it an orthonormal
  basis of $\overline{\mathcal{R}(K)}$.

  \emph{Step 5: the series.}  $\mathcal{N}(K)$ is closed by \cref{lem_ker_closed}, so
  $W^\perp = \mathcal{N}(K)^{\perp\perp} = \mathcal{N}(K)$ by \cref{thm_orth_decomp}.  Let
  $x \in X$.  By \cref{thm_proj_expansion} the series
  $\sum_{n\in I}\langle u_n,x\rangle u_n$ converges, with sum $P_Wx$; and
  $x - P_Wx \in W^\perp = \mathcal{N}(K)$, so $Kx = KP_Wx$.  Applying $K$ term by term is
  licensed by \cref{lem_operator_through_series}:
  \[
    Kx = KP_Wx = \sum_{n\in I}\langle u_n,x\rangle\,Ku_n
       = \sum_{n\in I}\lambda_n\langle u_n,x\rangle\,u_n .
  \]

  \emph{Step 6: the degenerate case.}  If $K = 0$ then $\Lambda = \emptyset$ and
  $I = \emptyset$: the series is the empty sum, $\overline{\mathcal{R}(K)} = \{0\}$, and
  the empty family is an orthonormal basis of it, so every clause holds.  If $K \ne 0$ then
  $W = \overline{\mathcal{R}(K)} \ne \{0\}$ forces $I \ne \emptyset$.  In both cases the
  index set is the initial segment produced in Step 2.

  \emph{A remark, not a clause.}  Every eigenvalue satisfies
  $|\lambda|\,\|u\| = \|Ku\| \le \|K\|\,\|u\|$, hence $|\lambda| \le \|K\|$, and over a
  \emph{complex} $X$ one may add that \cref{thm_max_eigenvalue} produces an eigenvalue of
  modulus exactly $\|K\|$, so that $|\lambda_1| = \|K\|$ by the ordering of Step 2.  This is
  deliberately not part of the statement: \cref{thm_max_eigenvalue} needs $\sigma(K)$
  nonempty and so is available only over $\mathbb{C}$, whereas Steps 1--5 are valid over
  $\mathbb{K} \in \{\mathbb{R},\mathbb{C}\}$ and \cref{thm_svd} applies this node in a
  chapter stated over $\mathbb{K}$.  The corresponding fact for singular values ---
  the bound $\sigma_n \le \|K\|$, which \emph{is} used --- is \cref{lem_sigma_bound}, and
  it is proved from the singular system directly, with no appeal to nonemptiness of any
  spectrum.  The sharper $\sigma_1 = \|K\|$ is nowhere needed.
\end{proof}



\chapter{Singular value decomposition}
\label{ch_08_svd}

\begin{assumptions}
  \label{asm_singular_value_decomposition}
  $X$, $Y$ are Hilbert spaces over $\mathbb{K}$ and
$K \in \mathcal{K}(X,Y)$.
\end{assumptions}

\begin{lemma}[$T^*T$ is positive]
  \label{lem_adjoint_self_positive}
  For $T \in \mathcal{L}(X,Y)$ between Hilbert spaces, $T^*T$ is positive.
\end{lemma}



\begin{definition}[Singular system]
  \label{def_singular_system}
  A \emph{singular system} for $K \in \mathcal{K}(X,Y)$ is a family
  $\{(\sigma_n,v_n,u_n)\}_{n \in I}$ indexed by an initial segment
  $I \subseteq \mathbb{N}$ (finite, or all of $\mathbb{N}$) such that
  \begin{enumerate}
    \item $\sigma_1 \ge \sigma_2 \ge \cdots > 0$, and $\sigma_n \to 0$ if $I$ is infinite;
    \item $\{v_n\}_{n \in I}$ is an orthonormal basis of
      $\overline{\mathcal{R}(K^*)} = \mathcal{N}(K)^\perp$;
    \item $\{u_n\}_{n \in I}$ is an orthonormal basis of $\overline{\mathcal{R}(K)}$;
    \item $Kv_n = \sigma_n u_n$ and $K^*u_n = \sigma_n v_n$ for every $n \in I$.
  \end{enumerate}
\end{definition}



\begin{theorem}[Singular value decomposition of a compact operator]
  \label{thm_svd}
  Every $K \in \mathcal{K}(X,Y)$ admits a singular system, indexed by an initial segment
  $I \subseteq \mathbb{N}$ and \emph{ordered} in the sense of
  \cref{def_singular_system}(1), and $Kx = \sum_n \sigma_n\langle v_n,x\rangle u_n$ for
  all $x \in X$.
\end{theorem}
\begin{proof}
  \emph{The operator $T := K^*K$.}  It is compact by \cref{lem_compact_comp} ($K$ compact,
  $K^*$ bounded), self-adjoint since $(K^*K)^* = K^*(K^*)^* = K^*K$ by
  \cref{prop_adjoint_props}, and positive by \cref{lem_adjoint_self_positive}.  Apply
  \cref{thm_spectral_diagonal} to $T$: it yields an initial segment
  $I \subseteq \mathbb{N}$, an orthonormal family $\{v_n\}_{n\in I}$ and nonzero reals
  $\mu_n$ with $|\mu_1| \ge |\mu_2| \ge \cdots$ such that
  $Tx = \sum_{n\in I}\mu_n\langle v_n,x\rangle v_n$, such that $\{v_n\}$ is an orthonormal
  basis of $\overline{\mathcal{R}(T)}$, and with $\mu_n \to 0$ when $I$ is infinite; the
  list repeats each nonzero eigenvalue of $T$ according to the dimension of its
  eigenspace.

  \emph{Positivity of the $\mu_n$, without circularity.}  Using $\|v_n\| = 1$, that $\mu_n$
  is real, and \cref{def_adjoint},
  \[
    \mu_n = \langle \mu_nv_n,v_n\rangle = \langle K^*Kv_n,v_n\rangle
          = \langle Kv_n,Kv_n\rangle = \|Kv_n\|^2 \ge 0 .
  \]
  Since $\mu_n \ne 0$, this gives $\mu_n > 0$ and $Kv_n \ne 0$.  The only inputs are the
  defining identity of the adjoint and the nonvanishing of the eigenvalues supplied by
  \cref{thm_spectral_diagonal}: nothing about the singular system under construction is
  presupposed.  Set $\sigma_n := \mu_n^{1/2} > 0$, so that $\sigma_n = \|Kv_n\|$.  As the
  $\mu_n$ are positive and ordered by decreasing modulus,
  $\mu_1 \ge \mu_2 \ge \cdots > 0$, and the square root is increasing, so
  $\sigma_1 \ge \sigma_2 \ge \cdots > 0$; and if $I$ is infinite then $\mu_n \to 0$ gives
  $\sigma_n \to 0$.  This is \cref{def_singular_system}(1).

  \emph{The vectors $u_n$.}  Put $u_n := \sigma_n^{-1}Kv_n$.  By \cref{def_adjoint},
  \[
    \langle u_m,u_n\rangle
      = (\sigma_m\sigma_n)^{-1}\langle K^*Kv_m,v_n\rangle
      = \frac{\mu_m}{\sigma_m\sigma_n}\,\langle v_m,v_n\rangle = \delta_{mn},
  \]
  so $\{u_n\}_{n\in I}$ is orthonormal.  By construction $Kv_n = \sigma_nu_n$, and
  \[
    K^*u_n = \sigma_n^{-1}K^*Kv_n = \sigma_n^{-1}\mu_nv_n = \sigma_nv_n :
  \]
  this is \cref{def_singular_system}(4).

  \emph{Condition (2).}  By \cref{lem_kernel_normal} and \cref{thm_four_subspaces},
  $\overline{\mathcal{R}(T)} = \overline{\mathcal{R}(K^*K)} = \overline{\mathcal{R}(K^*)}
  = \mathcal{N}(K)^\perp$, so $\{v_n\}$ is an orthonormal basis of $\mathcal{N}(K)^\perp$.

  \emph{The series.}  $\mathcal{N}(K)$ is closed (\cref{lem_ker_closed}), so
  $(\mathcal{N}(K)^\perp)^\perp = \mathcal{N}(K)$ by \cref{thm_orth_decomp}.  Let
  $x \in X$ and let $P$ be the orthogonal projection onto $\mathcal{N}(K)^\perp$.  By
  \cref{thm_proj_expansion} the series $\sum_{n\in I}\langle v_n,x\rangle v_n$ converges,
  with sum $Px$, and $x - Px \in \mathcal{N}(K)$, so $Kx = KPx$; passing $K$ through the
  series by \cref{lem_operator_through_series},
  \[
    Kx = \sum_{n\in I}\langle v_n,x\rangle\,Kv_n
       = \sum_{n\in I}\sigma_n\langle v_n,x\rangle\,u_n .
  \]

  \emph{Condition (3).}  Each $u_n = \sigma_n^{-1}Kv_n$ lies in $\mathcal{R}(K)$, so
  $\overline{\operatorname{span}}\{u_n\} \subseteq \overline{\mathcal{R}(K)}$; conversely
  the series just displayed exhibits every $Kx$ as a limit of finite linear combinations of
  the $u_n$, so $\overline{\mathcal{R}(K)} \subseteq
  \overline{\operatorname{span}}\{u_n\}$.  The two closed subspaces coincide, and
  \cref{lem_onb_criterion}(1), applied inside the Hilbert space
  $\overline{\mathcal{R}(K)}$, makes $\{u_n\}$ an orthonormal basis of it.

  \emph{Degenerate and finite cases.}  If $K = 0$ then $\mathcal{N}(K)^\perp = \{0\}$, so
  the orthonormal family produced by \cref{thm_spectral_diagonal} lives in $\{0\}$ and must
  be empty: $I = \emptyset$, the singular system is empty, both closed ranges are $\{0\}$,
  of which the empty family is an orthonormal basis, the four conditions hold vacuously,
  and the series is the empty sum $0 = Kx$.  If $I$ is finite --- equivalently $K$ has
  finite rank, since condition (3) gives $\dim\overline{\mathcal{R}(K)} = \#I$ --- all sums
  above are finite and the null-sequence clause of \cref{def_singular_system}(1) is not
  imposed.  In every case $I$ is an initial segment of $\mathbb{N}$.
\end{proof}


\begin{lemma}[Companion systems]
  \label{lem_svd_companion}
  Let $K \in \mathcal{K}(X,Y)$ have singular system $\{(\sigma_n,v_n,u_n)\}_{n \in I}$.
  Then for every $n \in I$
  \[
    K^*Kv_n = \sigma_n^2 v_n, \qquad KK^*u_n = \sigma_n^2 u_n ,
  \]
  and for every $x \in X$ and $y \in Y$
  \[
    Kx = \sum_{n \in I} \sigma_n \langle v_n, x\rangle_X\, u_n, \qquad
    K^*y = \sum_{n \in I} \sigma_n \langle u_n, y\rangle_Y\, v_n,
    \qquad
    K^*Kx = \sum_{n \in I} \sigma_n^2 \langle v_n, x\rangle_X\, v_n .
  \]
  Moreover a real $\lambda \ne 0$ is an eigenvalue of $K^*K$ if and only if
  $\lambda = \sigma_n^2$ for some $n \in I$.
  Equivalently: $\{(\sigma_n,u_n,v_n)\}_{n\in I}$ is a singular system of $K^*$,
  $\{(\sigma_n^2,v_n,v_n)\}$ one of $K^*K$, and $\{(\sigma_n^2,u_n,u_n)\}$ one of $KK^*$.
\end{lemma}
\begin{proof}
  Throughout, (1)--(4) refer to the conditions of \cref{def_singular_system}.

  \emph{The two eigenrelations.}  By (4), used twice,
  \[
    K^*Kv_n = K^*(\sigma_nu_n) = \sigma_nK^*u_n = \sigma_n^2v_n,
  \]
  \[
    KK^*u_n = K(\sigma_nv_n) = \sigma_nKv_n = \sigma_n^2u_n .
  \]

  \emph{The expansion of $K$.}  By (2), $\{v_n\}$ is an orthonormal basis of
  $\mathcal{N}(K)^\perp$, which is closed, so by \cref{thm_proj_expansion} the series
  $\sum_{n\in I}\langle v_n,x\rangle v_n$ converges to $Px$, $P$ the orthogonal projection
  onto $\mathcal{N}(K)^\perp$.  Since $\mathcal{N}(K)$ is closed (\cref{lem_ker_closed}),
  $(\mathcal{N}(K)^\perp)^\perp = \mathcal{N}(K)$ by \cref{thm_orth_decomp}, so
  $x - Px \in \mathcal{N}(K)$ and $Kx = KPx$.  Passing $K$ through the series with
  \cref{lem_operator_through_series} and using (4),
  \[
    Kx = \sum_{n\in I}\langle v_n,x\rangle\,Kv_n
       = \sum_{n\in I}\sigma_n\langle v_n,x\rangle\,u_n .
  \]
  This is the same computation as in \cref{thm_svd}, but run against the given system
  rather than the constructed one, which is the point of stating it here.

  \emph{The expansion of $K^*$.}  By (3), $\{u_n\}$ is an orthonormal basis of the closed
  subspace $\overline{\mathcal{R}(K)}$, so by \cref{thm_proj_expansion} the series
  $\sum_{n\in I}\langle u_n,y\rangle u_n$ converges to $Py$, where $P$ is the orthogonal
  projection onto $\overline{\mathcal{R}(K)}$.  By \cref{lem_orthogonal_closure} and
  \cref{thm_four_subspaces},
  $\overline{\mathcal{R}(K)}^\perp = \mathcal{R}(K)^\perp = \mathcal{N}(K^*)$, so
  $y - Py \in \mathcal{N}(K^*)$ and $K^*y = K^*Py$.  Passing $K^*$ through the series with
  \cref{lem_operator_through_series} and using (4),
  \[
    K^*y = \sum_{n\in I}\langle u_n,y\rangle\,K^*u_n
         = \sum_{n\in I}\sigma_n\langle u_n,y\rangle\,v_n .
  \]

  \emph{The expansion of $K^*K$.}  Apply the previous display to $y = Kx$, and use
  \cref{def_adjoint} with (4): $\langle u_n,Kx\rangle = \langle K^*u_n,x\rangle =
  \sigma_n\langle v_n,x\rangle$.  Hence
  \[
    K^*Kx = \sum_{n\in I}\sigma_n\langle u_n,Kx\rangle\,v_n
          = \sum_{n\in I}\sigma_n^2\langle v_n,x\rangle\,v_n .
  \]

  \emph{The eigenvalues of $K^*K$.}  If $\lambda = \sigma_n^2$ then $v_n \ne 0$ and
  $K^*Kv_n = \lambda v_n$ by the first eigenrelation, so $\lambda$ is an eigenvalue.
  Conversely let $\lambda \ne 0$ and $K^*Kw = \lambda w$ with $w \ne 0$.  Then
  $w = \lambda^{-1}K^*(Kw) \in \mathcal{R}(K^*) \subseteq \mathcal{N}(K)^\perp$
  (\cref{thm_four_subspaces}), so by (2) and \cref{thm_proj_expansion}
  $w = \sum_{n}\langle v_n,w\rangle v_n$ with not all coefficients zero.  Applying the
  expansion of $K^*K$ just proved and comparing coefficients with
  \cref{lem_riesz_fischer}, $\sigma_n^2\langle v_n,w\rangle = \lambda\langle v_n,w\rangle$
  for every $n$; picking an $n$ with $\langle v_n,w\rangle \ne 0$ gives
  $\lambda = \sigma_n^2$.

  \emph{The companion systems.}  For $K^*$: it is compact by \cref{thm_schauder}; (1) is
  unchanged; condition (2) for $K^*$ asks that $\{u_n\}$ be an orthonormal basis of
  $\overline{\mathcal{R}(K^{**})} = \overline{\mathcal{R}(K)}$, which is (3) for $K$, and
  \cref{thm_four_subspaces} applied to $K^*$ together with $K^{**} = K$
  (\cref{prop_adjoint_props}) identifies that space with $\mathcal{N}(K^*)^\perp$;
  condition (3) for $K^*$ asks that $\{v_n\}$ be an orthonormal basis of
  $\overline{\mathcal{R}(K^*)} = \mathcal{N}(K)^\perp$, which is (2) for $K$; and (4) for
  $K^*$ reads $K^*u_n = \sigma_nv_n$ and $K^{**}v_n = Kv_n = \sigma_nu_n$, again (4) for
  $K$.

  For $K^*K$: it is compact by \cref{lem_compact_comp} and self-adjoint by
  \cref{prop_adjoint_props}.  The numbers $\sigma_n^2$ are positive, nonincreasing, and
  tend to $0$ when $I$ is infinite, so (1) holds.  Conditions (2) and (3) both ask that
  $\{v_n\}$ be an orthonormal basis of $\overline{\mathcal{R}(K^*K)}$, which equals
  $\overline{\mathcal{R}(K^*)} = \mathcal{N}(K)^\perp$ by \cref{lem_kernel_normal} and
  \cref{thm_four_subspaces}; that is (2) for $K$.  And (4) is the first eigenrelation
  above, in both of the two required forms because $K^*K$ is self-adjoint.  Applying this
  last statement to $K^*$ in place of $K$ --- legitimate, since $\{(\sigma_n,u_n,v_n)\}$
  has just been shown to be a singular system of $K^*$ --- gives that
  $\{(\sigma_n^2,u_n,u_n)\}$ is a singular system of $K^{**}K^* = KK^*$.
\end{proof}



\begin{corollary}[Singular values vanish in the limit]
  \label{cor_sigma_null}
  If $\dim\mathcal{R}(K) = \infty$ then the index set $I$ is infinite and
  $\sigma_n \to 0$.
\end{corollary}
\begin{proof}
  By \cref{thm_svd}, $K$ has a singular system $\{(\sigma_n,v_n,u_n)\}_{n\in I}$ in which
  $\{u_n\}_{n\in I}$ is an orthonormal basis of $\overline{\mathcal{R}(K)}$.  Were $I$
  finite, $\overline{\mathcal{R}(K)}$ would equal $\operatorname{span}\{u_n\}_{n\in I}$,
  which is finite-dimensional and already closed (\cref{lem_findim_closed}), forcing
  $\dim\mathcal{R}(K) \le \#I < \infty$.  So $I$ is infinite, i.e.\ $I = \mathbb{N}$, and
  \cref{def_singular_system}(1) gives $\sigma_n \to 0$.  That this clause can be met is not
  a matter of convention: in the system produced by \cref{thm_svd} the $\sigma_n^2$ are the
  nonzero eigenvalues of the compact operator $K^*K$ repeated by multiplicity, and
  \cref{lem_eigenvalue_null}(1) says that for each $\varepsilon > 0$ only finitely many
  points of $\sigma(K^*K)$ have modulus $\ge \varepsilon$, and clause (3) that each has
  finite multiplicity --- so only finitely many indices $n$ have $\sigma_n^2 \ge \varepsilon$.
  Clause (1) is the one cited: from countability and ``$0$ is the only accumulation point''
  alone the finiteness would not follow.

  Since $\sigma_n > 0$, it follows that $\sigma_n^{-1} \to \infty$, and two things follow
  for the formal inverse $y \mapsto \sum_n\sigma_n^{-1}\langle u_n,y\rangle v_n$.  First,
  it is not everywhere defined: by \cref{lem_riesz_fischer} the series converges exactly
  when $\sum_n\sigma_n^{-2}|\langle u_n,y\rangle|^2 < \infty$, and choosing indices $n_k$
  with $\sigma_{n_k} \le 2^{-k}$ and setting $y := \sum_k k^{-1}u_{n_k}$ --- a convergent
  series by \cref{lem_riesz_fischer}, since $\sum_kk^{-2} < \infty$ --- gives
  $\sum_k\sigma_{n_k}^{-2}k^{-2} \ge \sum_k 4^k/k^2 = \infty$.  Second, where it is defined
  it is unbounded: $y = u_n$ has $\|u_n\| = 1$ while its image $\sigma_n^{-1}v_n$ has norm
  $\sigma_n^{-1} \to \infty$, so the map admits no bounded extension to
  $\overline{\mathcal{R}(K)}$.
\end{proof}




\chapter{Functions of a compact operator}
\label{ch_09_functions_of_compact_operator}

\begin{assumptions}
  \label{asm_functions_compact_operator}
  $X$, $Y$ are Hilbert spaces,
$K \in \mathcal{K}(X,Y)$ with a fixed singular system $\{(\sigma_n,v_n,u_n)\}$, and
$\kappa = \|K\|^2$.  $\varphi$ denotes a \emph{real}-valued function on $[0,\infty)$,
bounded on $[0,\kappa]$.
\end{assumptions}



\begin{lemma}[Singular values are bounded by the norm]
  \label{lem_sigma_bound}
  Let $K \in \mathcal{L}(X,Y)$ have singular system $\{(\sigma_n,v_n,u_n)\}_{n \in I}$,
  ordered or not.  Then $\sigma_n \le \|K\|$, and hence $\sigma_n^2 \in [0,\|K\|^2]$ for
  every $n \in I$.
\end{lemma}
\begin{proof}
  Each $v_n$ is a unit vector with $Kv_n = \sigma_n u_n$ and $\|u_n\| = 1$
  (\cref{def_singular_system}), so $\sigma_n = \|Kv_n\| \le \|K\|\,\|v_n\| = \|K\|$.
  Squaring gives $\sigma_n^2 \le \|K\|^2$, and $\sigma_n^2 \ge 0$ is immediate.
\end{proof}


\begin{definition}[Functions of $K^*K$]
  \label{def_fun_calc}
  For real-valued $\varphi$ bounded on $[0,\kappa]$, define
  $\varphi(K^*K) \in \mathcal{L}(X,X)$ by
  \[
    \varphi(K^*K)x = \sum_n \varphi(\sigma_n^2)\,\langle v_n,x\rangle\,v_n
      \;+\; \varphi(0)\,P_{\mathcal{N}(K)}x .
  \]
\end{definition}
\begin{proof}
  Write $N = \mathcal{N}(K)$.  A singular system exists by \cref{thm_svd}; $N$ is closed,
  so $P_N$ is defined, and $\{v_n\}_{n \in I}$ is an orthonormal basis of $N^\perp$ by
  \cref{def_singular_system}(2).  Three things must be checked: that the series converges
  for every $x$, that the resulting map is linear, and that it does not depend on which
  singular system of $K$ is used.

  \emph{Convergence.}  If $K = 0$ then $N^\perp = \{0\}$ contains no unit vector, so
  $I = \emptyset$, the series is empty, and the formula reads
  $\varphi(K^*K)x = \varphi(0)x$.  Assume $K \ne 0$.  By \cref{def_singular_system} and
  \cref{lem_sigma_bound} we have $\sigma_n^2 \in (0,\kappa]$ for every $n \in I$, and $\varphi$, being locally bounded on
  $[0,\infty)$, is bounded on the compact interval $[0,\kappa]$; put
  $M = \sup_{\lambda \in [0,\kappa]}|\varphi(\lambda)| < \infty$, which also bounds
  $|\varphi(0)|$.  For $x \in X$ the family $(\langle v_n,x\rangle)_{n \in I}$ is
  square-summable with $\sum_n |\langle v_n,x\rangle|^2 \le \|x\|^2$ by \cref{thm_bessel},
  whence
  \[
    \sum_{n \in I} \bigl|\varphi(\sigma_n^2)\langle v_n,x\rangle\bigr|^2
      \;\le\; M^2\|x\|^2 \;<\; \infty ,
  \]
  so the series converges in $X$ by \cref{lem_riesz_fischer}, and its coefficients against
  $\{v_n\}$ are the $\varphi(\sigma_n^2)\langle v_n,x\rangle$.  Only boundedness of
  $\varphi$ on $[0,\kappa]$ was used, so the argument applies to any bounded
  $\mathbb{K}$-valued function there.

  \emph{Linearity.}  Each $x \mapsto \langle v_n,x\rangle$ is linear and $P_N$ is linear.
  Given $x,y \in X$ the two series converge, so their termwise sum converges to the sum of
  the limits, which is the series for $x+y$; and for
  $c \in \mathbb{K}$, applying the bounded operator $z \mapsto cz$ through the series
  (\cref{lem_operator_through_series}) gives the series for $cx$.  Hence
  $\varphi(K^*K)$ is a linear map $X \to X$; boundedness is established separately in
  \cref{lem_fun_calc_bounded}.

  \emph{Independence of the singular system.}  Singular systems are not unique: a repeated
  singular value permits any orthonormal basis of the corresponding eigenspace, and each
  pair $(v_n,u_n)$ may be multiplied by a unimodular scalar.  Let
  $\{(\tau_m,w_m,z_m)\}_{m \in J}$ be a second singular system for the same $K$ and set
  \[
    Tx = \sum_{n \in I}\varphi(\sigma_n^2)\langle v_n,x\rangle\,v_n , \qquad
    Sx = \sum_{m \in J}\varphi(\tau_m^2)\langle w_m,x\rangle\,w_m ,
  \]
  both convergent by the first paragraph; note that $\kappa$ and $N$ depend on $K$ alone,
  so $\tau_m^2 \in (0,\kappa]$ as well and the kernel term $\varphi(0)P_Nx$ involves no
  choice.  It therefore suffices to prove $T = S$.

  Fix $m \in J$ and $n \in I$.  By \cref{lem_svd_companion}, $K^*Kv_n = \sigma_n^2v_n$ and
  $K^*Kw_m = \tau_m^2w_m$, and $(K^*K)^* = K^*K$ by \cref{prop_adjoint_props}; so by
  \cref{def_adjoint}
  \[
    \sigma_n^2\langle w_m,v_n\rangle
      = \langle w_m, K^*Kv_n\rangle
      = \langle K^*Kw_m, v_n\rangle
      = \tau_m^2\langle w_m,v_n\rangle ,
  \]
  the last step because $\tau_m^2$ is real.  Hence $\langle w_m,v_n\rangle \ne 0$ forces
  $\sigma_n^2 = \tau_m^2$, and consequently
  $\varphi(\sigma_n^2)\langle w_m,v_n\rangle = \varphi(\tau_m^2)\langle w_m,v_n\rangle$
  for \emph{every} $n \in I$, both sides being $0$ when $\langle w_m,v_n\rangle = 0$.
  This is the point at which the freedom in the choice of eigenvectors is absorbed:
  $\varphi$ is evaluated only at eigenvalues, which are determined by $K$.  Taking the
  inner product with $w_m$ through the series (\cref{lem_inner_through_series}) and using
  \cref{thm_proj_expansion} for the basis $\{v_n\}$ of $N^\perp$,
  \[
    \langle w_m, Tx\rangle
      = \sum_{n \in I}\varphi(\sigma_n^2)\langle v_n,x\rangle\langle w_m,v_n\rangle
      = \varphi(\tau_m^2)\,\langle w_m, P_{N^\perp}x\rangle .
  \]
  Now $P_{N^\perp}w_m = w_m$ by \cref{lem_proj_characterisation}, and $P_{N^\perp}$ is
  self-adjoint (\cref{prop_proj_props}), so
  $\langle w_m, P_{N^\perp}x\rangle = \langle w_m,x\rangle$ and
  $\langle w_m, Tx\rangle = \varphi(\tau_m^2)\langle w_m,x\rangle$.  The same computation
  for $S$, using orthonormality of $\{w_m\}$, gives
  $\langle w_m, Sx\rangle = \varphi(\tau_m^2)\langle w_m,x\rangle$.  Finally $Tx$ and $Sx$
  lie in $N^\perp$, since every partial sum does and $N^\perp$ is closed; applying
  \cref{thm_proj_expansion} to the basis $\{w_m\}_{m \in J}$ of $N^\perp$,
  \[
    Tx - Sx = P_{N^\perp}(Tx-Sx)
      = \sum_{m \in J}\langle w_m, Tx-Sx\rangle\,w_m = 0 .
  \]
  As $x$ was arbitrary, $T = S$, and $\varphi(K^*K)$ is well defined.
\end{proof}


\begin{lemma}[The functional calculus is bounded]
  \label{lem_fun_calc_bounded}
  $\varphi(K^*K) \in \mathcal{L}(X,X)$ with
  $\|\varphi(K^*K)\| \le \sup_n|\varphi(\sigma_n^2)| + |\varphi(0)|
   \le 2\sup_{\lambda\in[0,\kappa]}|\varphi(\lambda)|$.
\end{lemma}
\begin{proof}
  Linearity is part of \cref{def_fun_calc}.  Write $N = \mathcal{N}(K)$ and
  \[
    Sx = \sum_{n\in I}\varphi(\sigma_n^2)\langle v_n,x\rangle\,v_n ,
  \]
  so that $\varphi(K^*K)x = Sx + \varphi(0)P_Nx$, and put
  \[
    a = \sup_{n \in I}|\varphi(\sigma_n^2)| , \quad b = |\varphi(0)| ,
  \]
  with $a = 0$ when $I = \emptyset$.  We prove the sharp form
  $\|\varphi(K^*K)\| \le \max\{a,b\}$; the two printed bounds follow from it.

  Every $v_n$ lies in $N^\perp$ and $N^\perp$ is closed, so $Sx \in N^\perp$, while
  $P_Nx \in N$; the two summands are therefore orthogonal and \cref{lem_pythagoras} gives
  $\|\varphi(K^*K)x\|^2 = \|Sx\|^2 + b^2\|P_Nx\|^2$.  The coefficient family of $Sx$ is
  square-summable (\cref{thm_bessel} and $a < \infty$), so by the norm identity of
  \cref{lem_riesz_fischer} and \cref{thm_proj_expansion},
  \[
    \|Sx\|^2 = \sum_{n \in I}|\varphi(\sigma_n^2)|^2\,|\langle v_n,x\rangle|^2
      \le a^2\sum_{n \in I}|\langle v_n,x\rangle|^2
      = a^2\,\|P_{N^\perp}x\|^2 .
  \]
  Since $x = P_Nx + P_{N^\perp}x$ with orthogonal summands (\cref{prop_proj_props},
  \cref{lem_pythagoras}), we have $\|P_Nx\|^2 + \|P_{N^\perp}x\|^2 = \|x\|^2$, and so
  $\|\varphi(K^*K)x\|^2 \le \max\{a,b\}^2\,\|x\|^2$.  By \cref{lem_isleast_opnorm} the
  operator is bounded, with $\|\varphi(K^*K)\| \le \max\{a,b\}$.

  The maximum cannot be improved.  For $n \in I$ we have $P_Nv_n = 0$ by
  \cref{lem_proj_characterisation}, so $\varphi(K^*K)v_n = \varphi(\sigma_n^2)v_n$ and
  $\|\varphi(K^*K)\| \ge a$; and if $N \ne \{0\}$, a unit $x \in N$ has $Sx = 0$ by
  \cref{thm_proj_expansion}, so $\|\varphi(K^*K)x\| = b$ and
  $\|\varphi(K^*K)\| \ge b$.  If $N = \{0\}$ the second term is absent and the norm is $a$.

  Finally $\max\{a,b\} \le a + b$, which is the first printed bound; and by
  \cref{def_singular_system} and \cref{lem_sigma_bound} every $\sigma_n^2$ lies
  in $(0,\kappa]$, while
  $0 \in [0,\kappa]$, so $a$ and $b$ are both at most
  $\sup_{\lambda\in[0,\kappa]}|\varphi(\lambda)|$ and
  $a + b \le 2\sup_{\lambda\in[0,\kappa]}|\varphi(\lambda)|$.  (For $K = 0$ one has
  $I = \emptyset$, $\kappa = 0$ and $\varphi(K^*K) = \varphi(0)\mathrm{Id}$, and the chain
  still reads $b \le b \le 2b$.)  The absolute value on $\varphi(0)$ is not cosmetic: for
  $\varphi \equiv -1$ the operator has norm $1$, so a bound written with
  $+\,\varphi(0)$ in place of $+\,|\varphi(0)|$ would assert $1 \le 0$.
\end{proof}


\begin{lemma}[Basic instances of the calculus]
  \label{lem_fun_calc_basic}
  With $\mathbf{1}(t) = 1$, $\mathrm{id}(t) = t$ and $\sqrt{\phantom{t}}(t) = \sqrt{t}$ on
  $[0,\kappa]$:
  \[
    \mathbf{1}(K^*K) = \mathrm{Id}_X, \qquad
    \mathrm{id}(K^*K) = K^*K, \qquad
    \bigl(\sqrt{\phantom{t}}\,\bigr)(K^*K) \circ
      \bigl(\sqrt{\phantom{t}}\,\bigr)(K^*K) = K^*K .
  \]
  The operator $(\sqrt{\phantom{t}}\,)(K^*K)$ is what \cref{def_frac_power} will name
  $|K|$, so the last identity is $|K|^2 = K^*K$.
\end{lemma}
\begin{proof}
  Write $N = \mathcal{N}(K)$; by \cref{def_singular_system}(2) and
  \cref{thm_four_subspaces}, $\{v_n\}_{n\in I}$ is an orthonormal basis of
  $N^\perp = \overline{\mathcal{R}(K^*)}$.

  For $\varphi = \mathbf{1}$ the coefficients are $\langle v_n,x\rangle$ and the kernel
  term is $P_Nx$, so by \cref{thm_proj_expansion}
  \[
    \mathbf{1}(K^*K)x = \sum_{n\in I}\langle v_n,x\rangle v_n + P_Nx
      = P_{N^\perp}x + P_Nx = x ,
  \]
  the last equality by $\mathrm{Id} - P_N = P_{N^\perp}$ (\cref{prop_proj_props}), which is
  the decomposition $X = N \oplus N^\perp$ of \cref{thm_orth_decomp}.  This is where the
  kernel term is indispensable: without it one would get $P_{N^\perp}$, not $\mathrm{Id}$.

  For $\varphi = \mathrm{id}$ one has $\varphi(0) = 0$, so the kernel term drops and
  $\mathrm{id}(K^*K)x = \sum_{n\in I}\sigma_n^2\langle v_n,x\rangle v_n = K^*Kx$ by the
  last display of \cref{lem_svd_companion}.

  The third identity is the definition of the symbol $|K|$: it \emph{is} $\varphi(K^*K)$
  for $\varphi = \sqrt{\phantom{t}}$, which is bounded on $[0,\kappa]$, so
  \cref{def_fun_calc} applies and, again because $\sqrt{0} = 0$,
  $|K|x = \sum_{n \in I}\sigma_n\langle v_n,x\rangle\,v_n$.  For $|K|^2 = K^*K$ we compute
  directly rather than invoking multiplicativity.  The series converges by
  \cref{thm_bessel} and \cref{lem_riesz_fischer}, and taking the inner product with $v_m$
  through it (\cref{lem_inner_through_series}) gives
  $\langle v_m, |K|x\rangle = \sigma_m\langle v_m,x\rangle$; moreover $|K|x \in N^\perp$,
  so the kernel term vanishes when $|K|$ is applied a second time.  Hence
  \[
    |K|\bigl(|K|x\bigr) = \sum_{n\in I}\sigma_n\langle v_n,|K|x\rangle v_n
      = \sum_{n\in I}\sigma_n^2\langle v_n,x\rangle v_n = K^*Kx ,
  \]
  by \cref{lem_svd_companion} once more.
\end{proof}


\begin{lemma}[Termwise homomorphism property]
  \label{lem_fun_calc_hom}
  Let $g,h$ be real-valued and bounded on $[0,\kappa]$, as in \cref{def_fun_calc}, and let
  $c \in \mathbb{R}$ --- so that $cg$ is again in that class, which is what makes the
  second identity a statement about a defined object.  Then for \emph{every} $x \in X$,
  \begin{align*}
    (g+h)(K^*K)x &= g(K^*K)x + h(K^*K)x, \\
    (c\,g)(K^*K)x &= c\,g(K^*K)x, \\
    (gh)(K^*K)x &= g(K^*K)\bigl(h(K^*K)x\bigr) .
  \end{align*}
  In particular $\varphi \mapsto \varphi(K^*K)$ is an $\mathbb{R}$-algebra homomorphism
  from the bounded real functions on $[0,\kappa]$ into $\mathcal{L}(X,X)$, and
  $\varphi_\nu(K^*K)v_n = \sigma_n^\nu v_n$ for every $n \in I$ and $\nu > 0$, where
  $\varphi_\nu(t) = t^{\nu/2}$.
\end{lemma}
\begin{proof}
  Write $N = \mathcal{N}(K)$ and, for a bounded function $f$ on $[0,\kappa]$,
  $S_fx = \sum_{n\in I}f(\sigma_n^2)\langle v_n,x\rangle v_n$, so that
  $f(K^*K)x = S_fx + f(0)P_Nx$.  The convergence argument of \cref{def_fun_calc} used only
  boundedness of the function on $[0,\kappa]$, and $g+h$, $cg$, $gh$ are again bounded
  there.  No restriction of $x$ to $N^\perp$ is made anywhere below.

  \emph{Additivity and homogeneity.}  Both series converge, so their termwise sum converges
  to the sum of the limits; since
  $(g+h)(\sigma_n^2) = g(\sigma_n^2) + h(\sigma_n^2)$ and $(g+h)(0) = g(0)+h(0)$, adding
  the kernel terms gives $(g+h)(K^*K)x = g(K^*K)x + h(K^*K)x$.  Applying the bounded map
  $z \mapsto cz$ through the series (\cref{lem_operator_through_series}) gives
  $(cg)(K^*K)x = c\,g(K^*K)x$.

  \emph{Multiplicativity.}  Put $y = h(K^*K)x = S_hx + h(0)P_Nx$.  The two summands are
  separated by the decomposition of $X$: every $v_n$ lies in $N^\perp$, which is closed
  (\cref{def_orth_proj}), so $S_hx \in N^\perp$, while $P_Nx \in N$.  Two consequences, and
  they are exactly the points where multiplicativity could fail.

  First, the coefficients of $y$.  For $m \in I$ we have $P_Nv_m = 0$
  (\cref{lem_proj_characterisation}) and $P_N$ is self-adjoint (\cref{prop_proj_props}), so
  $\langle v_m, P_Nx\rangle = \langle P_Nv_m, x\rangle = 0$: the kernel term contributes
  nothing to the coefficients.  With \cref{lem_inner_through_series} and orthonormality,
  \[
    \langle v_m, y\rangle = \langle v_m, S_hx\rangle + h(0)\langle v_m,P_Nx\rangle
      = h(\sigma_m^2)\langle v_m,x\rangle .
  \]
  Second, the kernel component of $y$.  Since $S_hx \in N^\perp$ we get $P_NS_hx = 0$
  (\cref{lem_proj_characterisation}), and $P_N^2 = P_N$ (\cref{lem_proj_nested_comp}), so
  $P_Ny = h(0)P_Nx$: the two kernel terms compose by idempotence, and the cross terms
  vanish.  Therefore
  \[
    g(K^*K)y
      = \sum_{n \in I}g(\sigma_n^2)h(\sigma_n^2)\langle v_n,x\rangle v_n
        + g(0)h(0)\,P_Nx
      = (gh)(K^*K)x ,
  \]
  because $(gh)(\sigma_n^2) = g(\sigma_n^2)h(\sigma_n^2)$ and $(gh)(0) = g(0)h(0)$; the
  series converges by \cref{thm_bessel} and \cref{lem_riesz_fischer}, $gh$ being bounded.

  \emph{The homomorphism.}  The bounded real-valued functions on $[0,\kappa]$ form a
  unital $\mathbb{R}$-algebra under pointwise operations; the map $\varphi \mapsto \varphi(K^*K)$
  lands in $\mathcal{L}(X,X)$ by \cref{lem_fun_calc_bounded}, is linear and multiplicative
  by the above, and sends $\mathbf{1}$ to $\mathrm{Id}_X$ by \cref{lem_fun_calc_basic}.

  \emph{Eigenvectors.}  For $\nu > 0$ let $\varphi_\nu(t) = t^{\nu/2}$.  Fix $n \in I$.
  Then $\varphi_\nu(0) = 0$ and $\langle v_m,v_n\rangle = \delta_{mn}$, so exactly one term
  of the series survives:
  $\varphi_\nu(K^*K)v_n = \varphi_\nu(\sigma_n^2)v_n = (\sigma_n^2)^{\nu/2}v_n
    = \sigma_n^\nu v_n$,
  using $\sigma_n > 0$ (\cref{def_singular_system}(1)).  The operator $\varphi_\nu(K^*K)$
  is the one \cref{def_frac_power} will name $|K|^\nu$; writing it that way here would be a
  forward use of a symbol this chapter has not introduced, which is why both the statement
  and this proof spell it $\varphi_\nu$.
\end{proof}




\begin{definition}[Functions of $KK^*$]
  \label{def_fun_calc_y}
  For real-valued $\psi$ bounded on $[0,\kappa]$, as in \cref{def_fun_calc}, define
  $\psi(KK^*) \in \mathcal{L}(Y,Y)$ by
  \[
    \psi(KK^*)y \;=\; \sum_{n \in I} \psi(\sigma_n^2)\,\langle u_n,y\rangle_Y\,u_n
      \;+\; \psi(0)\,P_{\mathcal{N}(K^*)}y .
  \]
\end{definition}
\begin{proof}
  There is nothing new to prove: the displayed formula is \cref{def_fun_calc} applied to
  the operator $K^*$ in place of $K$.  Indeed $K^* \in \mathcal{K}(Y,X)$ by
  \cref{thm_schauder}, and by \cref{lem_svd_companion} the family
  $\{(\sigma_n,u_n,v_n)\}_{n\in I}$ is a singular system of $K^*$: the roles of the two
  orthonormal families are exchanged, the singular values unchanged.  Since
  $(K^*)^*K^* = KK^*$ and $\|K^*\| = \|K\|$ (\cref{prop_adjoint_props}), the operator whose
  functions are being defined is $KK^*$ and the relevant interval is the same $[0,\kappa]$;
  and the kernel of $K^*$ is $\mathcal{N}(K^*)$, which is why the projection term is
  $P_{\mathcal{N}(K^*)}$ and not $P_{\mathcal{N}(K)}$.  Instantiating \cref{def_fun_calc}
  at $K^*$ therefore yields exactly
  \[
    \psi(KK^*)y = \sum_{n\in I}\psi(\sigma_n^2)\langle u_n,y\rangle_Y u_n
      + \psi(0)P_{\mathcal{N}(K^*)}y ,
  \]
  and with it the three assertions proved there: the series converges for every $y \in Y$
  (\cref{thm_bessel} and \cref{lem_riesz_fischer}, with $\sigma_n^2 \in (0,\kappa]$ by
  \cref{def_singular_system} and \cref{lem_sigma_bound}, or $I = \emptyset$ when
  $K = 0$), the map is linear, and its value
  is independent of the singular system chosen --- the independence statement of
  \cref{def_fun_calc} quantifies over all singular systems of the operator it is applied
  to, here $K^*$, and every singular system of $K$ produces one of $K^*$ by
  \cref{lem_svd_companion}.
\end{proof}


\begin{lemma}[Transfer to $KK^*$]
  \label{lem_fun_calc_y}
  Everything asserted of $\varphi(K^*K)$ holds of $\psi(KK^*)$ with $u_n$ in place of
  $v_n$ and $\mathcal{N}(K^*)$ in place of $\mathcal{N}(K)$ --- in particular the bound of
  \cref{lem_fun_calc_bounded}, the basic instances of \cref{lem_fun_calc_basic} and the
  homomorphism identities of \cref{lem_fun_calc_hom}, all read on the same $[0,\kappa]$
  because $\|K^*\| = \|K\|$.
\end{lemma}
\begin{proof}
  The transfer is the instantiation of \cref{def_fun_calc} at $K^*$ carried out in
  \cref{def_fun_calc_y}, applied to each of \cref{lem_fun_calc_bounded},
  \cref{lem_fun_calc_basic} and \cref{lem_fun_calc_hom}: boundedness with the same
  estimate, $\mathbf{1}(KK^*) = \mathrm{Id}_Y$ and $\mathrm{id}(KK^*) = KK^*$, and the
  three algebra identities.  Note that $\mathrm{id}(KK^*) = (K^*)^*K^* = KK^*$, so the
  second basic instance reads on this side as it does on the other.
\end{proof}


\begin{lemma}[Intertwining]
  \label{lem_fun_calc_intertwine}
  For every bounded $\varphi$ on $[0,\kappa]$,
  \[
    K\,\varphi(K^*K) = \varphi(KK^*)\,K
    \qquad\text{and}\qquad
    \varphi(K^*K)\,K^* = K^*\,\varphi(KK^*) .
  \]
\end{lemma}
\begin{proof}
  Both identities are computed termwise; the kernel terms disappear on both sides, but for
  different reasons.

  \emph{First identity.}  Let $x \in X$.  Apply $K$ through the convergent series, as
  \cref{lem_operator_through_series} permits, using $Kv_n = \sigma_nu_n$ from
  \cref{def_singular_system}(4) together with $KP_{\mathcal{N}(K)}x = 0$:
  \[
    K\varphi(K^*K)x
      = \sum_{n \in I}\sigma_n\varphi(\sigma_n^2)\langle v_n,x\rangle\,u_n .
  \]
  On the other side, $\langle u_n, Kx\rangle_Y = \langle K^*u_n, x\rangle_X
  = \sigma_n\langle v_n,x\rangle$ by \cref{def_adjoint} and
  \cref{def_singular_system}(4), while $P_{\mathcal{N}(K^*)}Kx = 0$: by
  \cref{thm_four_subspaces}, $\mathcal{N}(K^*) = \mathcal{R}(K)^\perp$, so $Kx$ lies in
  $\mathcal{N}(K^*)^\perp$ and \cref{lem_proj_characterisation} applies.  Hence
  \cref{def_fun_calc_y} gives
  $\varphi(KK^*)Kx = \sum_{n}\sigma_n\varphi(\sigma_n^2)\langle v_n,x\rangle u_n$, the same
  vector.  The series converges by \cref{thm_bessel} and \cref{lem_riesz_fischer}, the
  coefficients being bounded multiples of $\langle v_n,x\rangle$.

  \emph{Second identity.}  Let $y \in Y$.  By \cref{def_adjoint} and
  \cref{def_singular_system}(4),
  $\langle v_n, K^*y\rangle_X = \overline{\langle K^*y,v_n\rangle_X}
   = \overline{\langle y, Kv_n\rangle_Y} = \sigma_n\langle u_n,y\rangle_Y$; and
  $P_{\mathcal{N}(K)}K^*y = 0$, since
  $K^*y \in \mathcal{R}(K^*) \subseteq \overline{\mathcal{R}(K^*)} =
  \mathcal{N}(K)^\perp$ (\cref{thm_four_subspaces}) and
  \cref{lem_proj_characterisation} applies again.  So
  \[
    \varphi(K^*K)K^*y
      = \sum_{n\in I}\sigma_n\varphi(\sigma_n^2)\langle u_n,y\rangle\,v_n .
  \]
  Conversely, applying $K^*$ through the series of \cref{def_fun_calc_y}
  (\cref{lem_operator_through_series}) and using $K^*u_n = \sigma_nv_n$ gives the same
  vector, provided $K^*P_{\mathcal{N}(K^*)}y = 0$ --- which holds because
  $P_{\mathcal{N}(K^*)}y \in \mathcal{N}(K^*)$.  This is the reason no hypothesis on
  $\varphi(0)$ is needed: the kernel term is annihilated by $K^*$ itself.
\end{proof}


\begin{definition}[Fractional powers]
  \label{def_frac_power}
  $|K|^\nu := \varphi(K^*K)$ with $\varphi(t) = t^{\nu/2}$, for $\nu > 0$.
\end{definition}


\begin{lemma}[$\||K|x\| = \|Kx\|$ and $\mathcal{R}(|K|) = \mathcal{R}(K^*)$]
  \label{lem_abs_norm}
  For all $x \in X$ one has $\||K|x\|_X = \|Kx\|_Y$.  Moreover
  $\mathcal{R}(|K|) = \mathcal{R}(K^*)$, and $|K|^r$ is self-adjoint for every $r > 0$.
\end{lemma}
\begin{proof}
  By \cref{def_frac_power} with $\varphi(t) = t^{r/2}$, and since $\varphi(0) = 0$ for
  $r > 0$ kills the kernel term,
  \[
    |K|^rx = \sum_{n\in I}\sigma_n^r\langle v_n,x\rangle\,v_n \qquad (r>0),
  \]
  in particular $|K|x = \sum_n\sigma_n\langle v_n,x\rangle v_n$; all these series converge
  by \cref{thm_bessel} and \cref{lem_riesz_fischer}, the $\sigma_n^r$ being bounded by
  $\sigma_1^r$.  If $K = 0$ then $I = \emptyset$ and every assertion below is trivial.

  \emph{(i) The norm identity.}  By the expansion of $K$ in \cref{lem_svd_companion} ---
  which holds for the fixed but arbitrary singular system of this chapter, whereas
  \cref{thm_svd} asserts it only of the system it constructs, and singular systems are not
  unique --- $Kx = \sum_n\sigma_n\langle v_n,x\rangle u_n$ with $\{u_n\}$ orthonormal in
  $Y$ by \cref{def_singular_system}(3), so the
  norm identity of \cref{lem_riesz_fischer}, applied on each side to the orthonormal
  families $\{v_n\}$ and $\{u_n\}$ with the \emph{same} coefficients
  $\sigma_n\langle v_n,x\rangle$, gives
  \[
    \bigl\||K|x\bigr\|_X^2 = \sum_{n\in I}\sigma_n^2|\langle v_n,x\rangle|^2
      = \|Kx\|_Y^2 .
  \]

  \emph{(ii) $\mathcal{R}(|K|) = \mathcal{R}(K^*)$.}  Both inclusions are proved by
  transporting coefficients between the two orthonormal families; throughout we use that,
  by \cref{lem_riesz_fischer}, a convergent series $\sum_n c_ne_n$ over an orthonormal
  system has coefficients $\langle e_m, \sum_n c_ne_n\rangle = c_m$.

  Let $w \in X$.  The family $(\langle v_n,w\rangle)_{n\in I}$ is square-summable
  (\cref{thm_bessel}), so $y := \sum_{n\in I}\langle v_n,w\rangle u_n$ converges in $Y$
  (\cref{lem_riesz_fischer}) and $\langle u_n,y\rangle = \langle v_n,w\rangle$ for every
  $n$.  By the expansion of $K^*$ in \cref{lem_svd_companion},
  \[
    K^*y = \sum_{n\in I}\sigma_n\langle u_n,y\rangle\,v_n
      = \sum_{n\in I}\sigma_n\langle v_n,w\rangle\,v_n = |K|w ,
  \]
  so $\mathcal{R}(|K|) \subseteq \mathcal{R}(K^*)$.  Conversely let $y \in Y$.  The family
  $(\langle u_n,y\rangle)_{n\in I}$ is square-summable (\cref{thm_bessel} in $Y$), so
  $w := \sum_{n\in I}\langle u_n,y\rangle v_n$ converges in $X$ with
  $\langle v_n,w\rangle = \langle u_n,y\rangle$, and then
  \[
    |K|w = \sum_{n\in I}\sigma_n\langle v_n,w\rangle\,v_n
      = \sum_{n\in I}\sigma_n\langle u_n,y\rangle\,v_n = K^*y ,
  \]
  so $\mathcal{R}(K^*) \subseteq \mathcal{R}(|K|)$.  Equivalently, $|K| = K^*V$ and
  $K^* = |K|V^*$ for the partial isometry $V \colon w \mapsto \sum_n\langle
  v_n,w\rangle u_n$.  No intermediate characterisation of $\mathcal{R}(K^*)$ is needed;
  as a by-product both ranges equal
  $\{x \in \mathcal{N}(K)^\perp : \sum_n\sigma_n^{-2}|\langle v_n,x\rangle|^2 < \infty\}$.

  \emph{(iii) Self-adjointness of $|K|^r$.}  $|K|^r \in \mathcal{L}(X,X)$ by
  \cref{lem_fun_calc_bounded}.  For $x,z \in X$, \cref{lem_inner_through_series} applied to
  the convergent series above gives
  \[
    \langle |K|^rx, z\rangle
      = \sum_{n\in I}\sigma_n^r\,\overline{\langle v_n,x\rangle}\,\langle v_n,z\rangle ,
  \]
  the $\sigma_n^r$ being real and positive, and the same computation with the roles of $x$
  and $z$ exchanged gives $\langle x, |K|^rz\rangle$ equal to the same sum, since
  $\langle x,v_n\rangle = \overline{\langle v_n,x\rangle}$.  Thus $|K|^r$ is symmetric, and
  being bounded it satisfies $(|K|^r)^* = |K|^r$ by \cref{def_adjoint}, i.e.\ it is
  self-adjoint in the sense of \cref{def_selfadjoint}.
\end{proof}


\begin{theorem}[Interpolation inequality]
  \label{thm_interpolation}
  For $r > s \ge 0$,
  \[
    \||K|^sx\| \le \||K|^rx\|^{s/r}\,\|x\|^{1-s/r}.
  \]
\end{theorem}
\begin{proof}
  Write $c_n = \langle v_n,x\rangle$ and $N = \mathcal{N}(K)$.  As in
  \cref{def_frac_power}, for $t > 0$ the function $\lambda \mapsto \lambda^{t/2}$ vanishes
  at $0$, so the kernel term of \cref{def_fun_calc} drops and
  $|K|^tx = \sum_{n\in I}\sigma_n^tc_nv_n$; by the norm identity of
  \cref{lem_riesz_fischer} and \cref{thm_bessel},
  \[
    \bigl\||K|^tx\bigr\|^2 = \sum_{n\in I}\sigma_n^{2t}|c_n|^2 , \qquad
    \sum_{n\in I}|c_n|^2 = \|P_{N^\perp}x\|^2 \le \|x\|^2 ,
  \]
  the middle equality by \cref{thm_proj_expansion}.

  \emph{The case $0 < s < r$.}  Use H\"older's inequality for families of nonnegative reals
  with the conjugate exponents $p = r/s$ and $q = r/(r-s)$, for which
  $1/p + 1/q = 1$, applied to $a_n = \sigma_n^{2s}|c_n|^{2s/r}$ and
  $b_n = |c_n|^{2(r-s)/r}$, so that $a_n^p = \sigma_n^{2r}|c_n|^2$ and $b_n^q = |c_n|^2$:
  \[
    \bigl\||K|^sx\bigr\|^2 = \sum_{n\in I}a_nb_n
      \le \Bigl(\sum_{n\in I}\sigma_n^{2r}|c_n|^2\Bigr)^{s/r}
          \Bigl(\sum_{n\in I}|c_n|^2\Bigr)^{(r-s)/r} ,
  \]
  and the right-hand side is at most
  \[
    \bigl\||K|^rx\bigr\|^{2s/r}\,\|x\|^{2(1-s/r)} .
  \]
  Taking square roots is the assertion.

  \emph{The endpoint $s = 0$.}  Here $|K|^0$ is not covered by \cref{def_frac_power} and a
  convention must be fixed.  We read $\lambda^0 \equiv 1$ on $[0,\kappa]$, i.e.\
  $0^0 = 1$; then $|K|^0 = \mathbf{1}(K^*K) = \mathrm{Id}_X$ by
  \cref{lem_fun_calc_basic}, and the same convention read on the right-hand side makes
  $\||K|^rx\|^0 = 1$ even when $\||K|^rx\| = 0$.  The asserted inequality is then
  $\|x\| \le \|x\|$.

  The statement is also true under the other convention, $0^0 = 0$, provided it is used
  consistently.  Then $|K|^0x = \sum_n\langle v_n,x\rangle v_n = P_{N^\perp}x$ by
  \cref{thm_proj_expansion}.  If $\||K|^rx\| \ne 0$ the right-hand side is $\|x\|$ and
  $\|P_{N^\perp}x\| \le \|x\|$ by \cref{thm_bessel}; if $\||K|^rx\| = 0$ the right-hand
  side is $0$, but then $\sum_n\sigma_n^{2r}|c_n|^2 = 0$ with every $\sigma_n > 0$, so all
  $c_n = 0$ and $P_{N^\perp}x = 0$, and the left-hand side is $0$ as well.  What is not
  admissible is mixing the two: with $|K|^0 = \mathrm{Id}$ and $0^0 = 0$ on the right, the
  inequality fails for any $x \in \mathcal{N}(K)\setminus\{0\}$.
\end{proof}



\chapter{Best approximation, widths, and approximation numbers}
\label{ch_10_best_approximation}

\begin{assumptions}
  \label{asm_best_approximation_widths}
  $X$, $Y$ are normed spaces over $\mathbb{K}$, Banach
or Hilbert where a statement says so; $W$ denotes a subspace of $X$ and $K$ a bounded
subset.  $\mathcal{F}_n(X)$ is the set of subspaces of $X$ of dimension at most $n$.

\end{assumptions}

\section{The approximation error}

\begin{definition}[Approximation error]
  \label{def_approx_error}
  For $x \in X$ and $W \subseteq X$ a subspace, the \emph{approximation error} is
  \[
    E(x,W) \;:=\; \inf_{w \in W}\|x-w\| .
  \]
\end{definition}



\section{Nested families}

\emph{Conventions for this section, and for this section only.}  $X$ is a Hilbert space,
$X_1 \subseteq X_2 \subseteq \cdots$ are closed subspaces with
$\overline{\bigcup_n X_n} = X$, and $P_n$ is the orthogonal projection onto $X_n$.
Nestedness is a hypothesis \emph{here} and is never a standing assumption elsewhere in this
chapter or the next: the families a downstream approximation scheme produces need not be
nested, and every statement outside this section is written against $E(x,W_n) \to 0$
directly.

\begin{theorem}[Nested projections converge pointwise]
  \label{thm_nested_proj}
  $\|(\mathrm{Id}-P_n)x\| \to 0$ for every $x \in X$.
\end{theorem}


\begin{theorem}[Nested projections onto a proper limit subspace]
  \label{thm_nested_proj_proper}
  Let $\{V_k\}$ be an increasing family of closed subspaces of a Hilbert space $X$ and let
  $V := \overline{\bigcup_k V_k}$, which need not be $X$.  Then
  $P_{V_k}x \to P_Vx$ for every $x \in X$.
\end{theorem}



 \part{Regularization theory}
\label{part_2_reg}

\chapter{Ill-posedness and the Moore--Penrose inverse}
\label{ch_11_moore_penrose}

\noindent\emph{Background.} Well-posedness in the sense of Hadamard --- existence,
uniqueness, and continuous dependence of the solution on the data --- was introduced by
\citet{Hadamard1902} as the defining property of a properly posed mathematical model. The
Moore--Penrose inverse, recovering the least-squares, minimum-norm solution whenever one
exists, goes back to \citet{Moore1920} and was rediscovered, together with its
now-standard characterisation via the four defining equations, by \citet{Penrose1955}.
The operator-theoretic notion of ill-posedness used throughout the sequel --- non-closed
range, rather than mere discontinuity of $T\dg$ on its domain --- is due to
\citet{Nashed1987}; it is this failure of closedness, rather than unboundedness as such,
that makes $T\dg$ genuinely pathological and regularization necessary.

\begin{assumptions}
  \label{asm_ill_posedness_moore}
  $X$ and $Y$ are Hilbert spaces over $\mathbb{K}$ and
$T \in \Ll(X,Y)$, except in \cref{def_well_posed}, where they are merely normed and $T$
is not assumed bounded.  $P_N$ denotes the orthogonal projection
onto $\Ker(T)$ and $P_R$ the orthogonal projection onto $\overline{\Ran(T)}$; both
subspaces are closed, the first because $T$ is bounded.
\end{assumptions}

\begin{definition}[Well-posedness in the sense of Hadamard]
  \label{def_well_posed}
  Let $X$, $Y$ be normed spaces and $T \colon X \to Y$ linear.  The equation $Tx = y$ is
  \emph{well-posed} if
  \begin{enumerate}
    \item for every $y \in Y$ there is $x \in X$ with $Tx = y$;
    \item that $x$ is unique;
    \item the solution depends continuously on the datum: for all $x \in X$ and all
      sequences $\{x_n\} \subset X$ with $Tx_n \to Tx$, one has $x_n \to x$.
  \end{enumerate}
  Otherwise the equation is \emph{ill-posed}.
\end{definition}


\begin{proposition}[Two conditions imply the third on Banach spaces]
  \label{prop_well_posed_banach}
  Let $X$, $Y$ be Banach spaces and $T \in \Ll(X,Y)$ satisfy (i) and (ii) of
  \cref{def_well_posed}.  Then $T$ is bijective, $T^{-1} \in \Ll(Y,X)$, and (iii) holds.
\end{proposition}

\begin{proof}
  Condition (i) says $T$ is surjective.  For injectivity, let $Tx = Tz$ and put
  $y := Tx$; then $x$ and $z$ both solve $Tw = y$, so (ii) forces $z = x$.  Hence $T$ is
  bijective, and since $X$ and $Y$ are Banach spaces the bounded inverse theorem
  \cref{thm_bounded_inverse} gives $T^{-1} \in \Ll(Y,X)$.

  For (iii), let $x \in X$ and $\{x_n\} \subset X$ with $Tx_n \to Tx$.  Continuity of
  $T^{-1}$ gives $x_n = T^{-1}(Tx_n) \to T^{-1}(Tx) = x$.
\end{proof}


\begin{definition}[Least squares and minimum norm solutions]
  \label{def_lsq_minnorm}
  Let $y \in Y$.  An element $\hat{x} \in X$ is a \emph{least squares solution} of
  $Tx = y$ if $\|T\hat{x} - y\|_Y \le \|Tz - y\|_Y$ for all $z \in X$.  It is a
  \emph{minimum norm solution} if in addition $\|\hat{x}\|_X \le \|z\|_X$ for every least
  squares solution $z$.
\end{definition}


\begin{lemma}[The minimum norm solution is unique]
  \label{lem_minnorm_unique}
  Let $y \in Y$.  Any two minimum norm solutions of $Tx = y$ in the sense of
  \cref{def_lsq_minnorm} are equal.
\end{lemma}

\begin{proof}
  Let $\hat{x}$ and $\hat{x}'$ be minimum norm solutions.  Both are least squares
  solutions, so $T\hat{x} = T\hat{x}'$: their difference lies in $\Ran(T)$ and, by
  \cref{lem_proj_characterisation}, the two residuals $y - T\hat{x}$ and $y - T\hat{x}'$
  lie in $\Ran(T)^{\perp}$, whence so does $T\hat{x} - T\hat{x}'$, and
  $\Ran(T) \cap \Ran(T)^{\perp} = \{0\}$.  Hence $\hat{x} - \hat{x}' \in \Ker(T)$.  A
  minimum norm solution lies in $\Ker(T)^{\perp}$, so $\hat{x} - \hat{x}'$ lies there too,
  and $\Ker(T) \cap \Ker(T)^{\perp} = \{0\}$ gives $\hat{x} = \hat{x}'$.
\end{proof}


\begin{definition}[The Moore--Penrose inverse]
  \label{def_moore_penrose}
  Let $T \in \Ll(X,Y)$.  The \emph{Moore--Penrose domain} of $T$ is
  \[
    \Dom(T\dg) := \Ran(T) \oplus \Ran(T)^{\perp} \subseteq Y ,
  \]
  and the \emph{Moore--Penrose inverse} $T\dg$ of $T$ is the linear map with domain
  $\Dom(T\dg)$ sending $y$ to the minimum norm solution of $Tx = y$ in the sense of
  \cref{def_lsq_minnorm}.  Outside that domain $T\dg$ takes no value.  It is in general
  unbounded.
\end{definition}

\begin{proof}
  \emph{The sum is direct.}  If $y \in \Ran(T) \cap \Ran(T)^{\perp}$ then
  $\langle y,y\rangle = 0$, so $y = 0$.  Every $y \in \Dom(T\dg)$ therefore has a
  \emph{unique} decomposition $y = y_1 + y_2$ with $y_1 \in \Ran(T)$ and
  $y_2 \in \Ran(T)^{\perp}$.

  \emph{The value exists.}  Write $y = y_1 + y_2$ as above and pick $x_0$ with
  $Tx_0 = y_1$.  Then $\hat{x} := (\Id - P_N)x_0$ satisfies $T\hat{x} = y_1$, so
  $y - T\hat{x} = y_2 \in \Ran(T)^{\perp}$, and $\hat{x} \in \Ker(T)^{\perp}$; by
  \cref{lem_proj_characterisation} these two conditions say exactly that $\hat{x}$ is a
  minimum norm solution of $Tx = y$.  It is the only one by
  \cref{lem_minnorm_unique}, so the assignment $y \mapsto T\dg y$ is well defined on
  $\Dom(T\dg)$.

  \emph{The assignment is linear.}  Let $y, z \in \Dom(T\dg)$ and $c$ a scalar.  The sum
  $T\dg y + T\dg z$ has residual $(y + z) - T(T\dg y + T\dg z)$, the sum of two elements of
  $\Ran(T)^{\perp}$, and lies in $\Ker(T)^{\perp}$, so it is a minimum norm solution of
  $Tx = y + z$ and \cref{lem_minnorm_unique} identifies it with $T\dg(y+z)$.  The same
  argument with $c\,T\dg y$ gives $T\dg(cy) = c\,T\dg y$.  Hence $T\dg$ is linear on
  $\Dom(T\dg)$, which is a subspace since $\Ran(T)$ and $\Ran(T)^{\perp}$ are.
\end{proof}


\begin{theorem}[Existence and uniqueness of the Moore--Penrose inverse]
  \label{thm_mp_inverse}
  Let $T \in \Ll(X,Y)$ and let $\tilde{T} := T|_{\Ker(T)^{\perp}} \colon \Ker(T)^{\perp}
  \to \Ran(T)$.  Then $\tilde{T}$ is bijective, and the Moore--Penrose inverse $T\dg$ of
  \cref{def_moore_penrose} is the unique linear extension of $\tilde{T}^{-1}$ with
  \[
    \Dom(T\dg) = \Ran(T) \oplus \Ran(T)^{\perp},
    \qquad
    \Ker(T\dg) = \Ran(T)^{\perp}.
  \]
\end{theorem}

\begin{proof}
  \emph{$\tilde{T}$ is bijective.}  $\Ker(T)$ is closed (\cref{lem_ker_closed}), so the
  orthogonal projection $P_N$ onto it exists and $X = \Ker(T) \oplus \Ker(T)^{\perp}$
  (\cref{thm_orth_decomp}).  If $x \in \Ker(T)^{\perp}$ and $\tilde{T}x = 0$ then
  $x \in \Ker(T) \cap \Ker(T)^{\perp}$, so $\langle x, x\rangle = 0$ and hence
  $x = 0$; thus $\tilde{T}$ is injective.  For surjectivity onto
  $\Ran(T)$, let $y = Tx$ and split $x = P_Nx + (\Id - P_N)x$.  Then
  $(\Id - P_N)x \in \Ker(T)^{\perp}$ and $T(\Id-P_N)x = Tx - TP_Nx = Tx = y$, since
  $P_Nx \in \Ker(T)$.  Hence $\tilde{T}$ is a linear bijection and $\tilde{T}^{-1} \colon
  \Ran(T) \to \Ker(T)^{\perp}$ is a well-defined linear map.

  \emph{$T\dg$ extends $\tilde{T}^{-1}$.}  Let $y \in \Ran(T)$, so that the decomposition
  of \cref{def_moore_penrose} has $y_1 = y$ and $y_2 = 0$.  The value $T\dg y$ is the
  minimum norm solution, which lies in $\Ker(T)^{\perp}$ and satisfies $T(T\dg y) = y$;
  these are exactly the two conditions defining $\tilde{T}^{-1}y$, and $\tilde{T}$ is
  injective, so $T\dg y = \tilde{T}^{-1}y$.

  \emph{The kernel.}  For $y = y_1 + y_2$ as in \cref{def_moore_penrose}, $T\dg y$ is the
  minimum norm solution, so $T(T\dg y) = y_1$ and $T\dg y \in \Ker(T)^{\perp}$.  Hence
  $T\dg y = 0$ forces $y_1 = 0$, that is $y \in \Ran(T)^{\perp}$; conversely if
  $y \in \Ran(T)^{\perp}$ then $y_1 = 0$ and $0$ is a minimum norm solution of $Tx = y$,
  so $T\dg y = 0$ by \cref{lem_minnorm_unique}.  So
  $\Ker(T\dg) = \Ran(T)^{\perp}$ as claimed.

  \emph{Uniqueness of the extension.}  Let $S$ be linear on $\Dom(T\dg)$, extend
  $\tilde{T}^{-1}$, and satisfy $\Ker(S) = \Ran(T)^{\perp}$.  For $y = y_1 + y_2$ as
  above, $Sy = Sy_1 + Sy_2 = \tilde{T}^{-1}y_1 + 0 = T\dg y$.

  Boundedness is not claimed, and in general fails; that is the subject of the
  continuity theorem below.
\end{proof}


\begin{lemma}[Characterisation of the domain]
  \label{lem_mp_domain}
  For $y \in Y$ one has $y \in \Dom(T\dg)$ if and only if $P_R y \in \Ran(T)$.
\end{lemma}

\begin{proof}
  We first record that $\Ran(T)^{\perp} = \overline{\Ran(T)}^{\perp}$.  Indeed
  $\supseteq$ is clear, and if $z \perp \Ran(T)$ and $w = \lim w_n$ with
  $w_n \in \Ran(T)$, then $\langle z, w\rangle = \lim \langle z, w_n\rangle = 0$ because
  $|\langle z, w - w_n\rangle| \le \|z\|\,\|w - w_n\|$ (\cref{thm_cauchy_schwarz}).

  ($\Rightarrow$)  Let $y = y_1 + y_2$ with $y_1 \in \Ran(T)$ and
  $y_2 \in \Ran(T)^{\perp} = \overline{\Ran(T)}^{\perp}$.  Since
  $y_1 \in \overline{\Ran(T)}$, this exhibits $y$ as the sum of an element of the closed
  subspace $\overline{\Ran(T)}$ and an element of its orthogonal complement; that
  decomposition is unique (\cref{thm_orth_decomp}), so $P_Ry = y_1 \in \Ran(T)$.

  ($\Leftarrow$)  If $P_Ry \in \Ran(T)$, write $y = P_Ry + (y - P_Ry)$.  The second summand
  lies in $\overline{\Ran(T)}^{\perp} = \Ran(T)^{\perp}$ (\cref{lem_proj_characterisation}),
  so $y \in \Ran(T) \oplus \Ran(T)^{\perp} = \Dom(T\dg)$.
\end{proof}


\begin{lemma}[Moore--Penrose equations]
  \label{lem_mp_equations}
  $\Ran(T\dg) = \Ker(T)^{\perp}$, and
  \begin{enumerate}
    \item $T T\dg T = T$ on $X$;
    \item $T\dg T T\dg = T\dg$ on $\Dom(T\dg)$;
    \item $T\dg T = \Id - P_N$ on $X$;
    \item $T T\dg = P_R|_{\Dom(T\dg)}$.
  \end{enumerate}
\end{lemma}

\begin{proof}
  Throughout, for $y \in \Dom(T\dg)$ write $y = y_1 + y_2$ as in \cref{thm_mp_inverse};
  by \cref{lem_mp_domain}, $y_1 = P_Ry$, so the definition of $T\dg$ reads
  \begin{equation}\label{eq_mp_via_proj}
    T\dg y = \tilde{T}^{-1}P_Ry, \qquad y \in \Dom(T\dg),
  \end{equation}
  and $P_Ry \in \Ran(T)$ --- not merely in $\overline{\Ran(T)}$.  This is the property
  that the rest of the chapter uses.

  \emph{$\Ran(T\dg) = \Ker(T)^{\perp}$.}  By \eqref{eq_mp_via_proj},
  $T\dg y \in \Ran(\tilde{T}^{-1}) = \Ker(T)^{\perp}$, giving $\subseteq$.  Conversely, for
  $x \in \Ker(T)^{\perp}$ we have $Tx \in \Ran(T) \subseteq \Dom(T\dg)$ and
  $T\dg Tx = \tilde{T}^{-1}\tilde{T}x = x$, so $x \in \Ran(T\dg)$.

  \emph{(iv).}  For $y \in \Dom(T\dg)$, \eqref{eq_mp_via_proj} gives
  $TT\dg y = T\tilde{T}^{-1}P_Ry$.  Now $\tilde{T}^{-1}P_Ry \in \Ker(T)^{\perp}$ and $T$
  agrees with $\tilde{T}$ there, so
  $TT\dg y = \tilde{T}\tilde{T}^{-1}P_Ry = P_Ry$.

  \emph{(iii).}  For $x \in X$ we have $Tx \in \Ran(T) \subseteq \Dom(T\dg)$ and
  $P_RTx = Tx$, so \eqref{eq_mp_via_proj} gives $T\dg Tx = \tilde{T}^{-1}Tx$.  Splitting
  $x = P_Nx + (\Id - P_N)x$ and using $TP_Nx = 0$,
  \[
    T\dg Tx = \tilde{T}^{-1}T(\Id - P_N)x = \tilde{T}^{-1}\tilde{T}(\Id - P_N)x
      = (\Id - P_N)x .
  \]

  \emph{(ii).}  For $y \in \Dom(T\dg)$, item (iv) and \eqref{eq_mp_via_proj} give
  \[
    T\dg TT\dg y = T\dg P_Ry = \tilde{T}^{-1}P_R P_R y = \tilde{T}^{-1}P_Ry = T\dg y,
  \]
  using $P_R^2 = P_R$ (\cref{prop_proj_props}) and
  $P_Ry \in \Ran(T) \subseteq \Dom(T\dg)$.

  \emph{(i).}  By (iii), $TT\dg Tx = T(\Id - P_N)x = Tx - TP_Nx = Tx$ for all $x \in X$.
\end{proof}


\begin{lemma}[The domain is dense]
  \label{lem_mp_domain_dense}
  $\Dom(T\dg)$ is dense in $Y$, and $\Dom(T\dg) = Y$ if and only if $\Ran(T)$ is closed.
  For $y \in \Ran(T)^{\perp}$ one has $T\dg y = 0$.
\end{lemma}

\begin{proof}
  \emph{Density.}  Let $z \perp \Dom(T\dg)$.  Since $\Ran(T) \subseteq \Dom(T\dg)$ we get
  $z \in \Ran(T)^{\perp}$, and $\Ran(T)^{\perp} \subseteq \Dom(T\dg)$, so $z \perp z$ and
  therefore $z = 0$.  Thus
  $\Dom(T\dg)^{\perp} = \{0\}$.  As in \cref{lem_mp_domain},
  $\Dom(T\dg)^{\perp} = \overline{\Dom(T\dg)}^{\perp}$, and $\overline{\Dom(T\dg)}$ is a
  closed subspace, so $Y = \overline{\Dom(T\dg)} \oplus \overline{\Dom(T\dg)}^{\perp}
  = \overline{\Dom(T\dg)}$ by \cref{thm_orth_decomp}.

  \emph{$\Dom(T\dg) = Y$ iff $\Ran(T)$ is closed.}  If $\Ran(T)$ is closed then
  $\Ran(T) \oplus \Ran(T)^{\perp} = Y$ by \cref{thm_orth_decomp}.  Conversely, if
  $\Dom(T\dg) = Y$ then every $y \in \overline{\Ran(T)}$ lies in $\Dom(T\dg)$, so
  $y = P_Ry \in \Ran(T)$ by \cref{lem_mp_domain}; hence
  $\overline{\Ran(T)} \subseteq \Ran(T)$.

  \emph{The last clause} is $\Ker(T\dg) = \Ran(T)^{\perp}$ from \cref{thm_mp_inverse}.
\end{proof}


\begin{theorem}[The Moore--Penrose inverse computes the minimum norm solution]
  \label{thm_mp_minnorm}
  Let $y \in \Dom(T\dg)$.  Then
  \begin{enumerate}
    \item least squares solutions of $Tx = y$ exist, and are exactly the solutions of
      $Tx = P_R y$;
    \item the set of all least squares solutions is $x\dg + \Ker(T)$, where
      $x\dg := T\dg y$.
  \end{enumerate}
\end{theorem}

\begin{proof}
  Write $U := \overline{\Ran(T)}$, a closed subspace, and recall from
  \cref{lem_mp_domain} that $P_Ry \in \Ran(T)$.

  \emph{(i).}  Since $P_Ry \in \Ran(T)$, the equation $Tx = P_Ry$ has a solution.  For any
  $x \in X$ we have $Tx \in U$, so by the minimising property of the orthogonal projection
  onto the closed \emph{subspace} $U$ (\cref{prop_proj_props}, last clause)
  \[
    \|P_Ry - y\| = \min_{w \in U}\|w - y\| \le \|Tx - y\| .
  \]
  Hence every solution $z$ of $Tz = P_Ry$ satisfies $\|Tz - y\| \le \|Tx - y\|$ for all
  $x$, i.e.\ is a least squares solution.  Conversely, let $z$ be a least squares solution.
  Then $\|Tz - y\| \le \|Tx - y\|$ for every $x$, and choosing $x$ with $Tx = P_Ry$ gives
  $\|Tz - y\| \le \|P_Ry - y\| = \min_{w \in U}\|w - y\|$.  As $Tz \in U$, the reverse
  inequality also holds, so $Tz$ attains the minimum.  That the minimiser is \emph{unique}
  needs no convexity theorem here: $Tz - P_Ry \in U$ while $y - P_Ry \in U^{\perp}$
  (\cref{lem_proj_characterisation}), so by Pythagoras (\cref{lem_pythagoras})
  \[
    \|Tz - y\|^2 = \|Tz - P_Ry\|^2 + \|P_Ry - y\|^2 ,
  \]
  and equality of the two sides forces $Tz = P_Ry$.



  \emph{(ii) and (iii).}  By (i) the set of least squares solutions is the affine set
  $\{x : Tx = P_Ry\}$, which is nonempty; if $Tx = Tx' = P_Ry$ then $x - x' \in \Ker(T)$,
  so that set is $x_0 + \Ker(T)$ for any one solution $x_0$.  Decompose a least squares
  solution as $x = \bar{x} + x_0$ with $\bar{x} = (\Id - P_N)x \in \Ker(T)^{\perp}$ and
  $x_0 = P_Nx \in \Ker(T)$ (\cref{thm_orth_decomp}).  The component $\bar{x}$ does not depend
  on the choice of $x$: if $\bar{x}, \bar{x}'$ both arise this way then
  $\tilde{T}\bar{x} = P_Ry = \tilde{T}\bar{x}'$, and $\tilde{T}$ is injective
  (\cref{thm_mp_inverse}).  By Pythagoras (\cref{lem_pythagoras}),
  \[
    \|x\|^2 = \|\bar{x}\|^2 + \|x_0\|^2 \ge \|\bar{x}\|^2 ,
  \]
  with equality only for $x_0 = 0$.  So $\bar{x}$ is the unique minimum norm solution, and
  the set of all least squares solutions is $\bar{x} + \Ker(T)$.

  Finally $\bar{x} \in \Ker(T)^{\perp}$ and $T\bar{x} = P_Ry$ give, using
  \cref{lem_mp_equations}(iii) and then (ii),
  \[
    \bar{x} = (\Id - P_N)\bar{x} = T\dg T\bar{x} = T\dg P_Ry = T\dg TT\dg y = T\dg y,
  \]
  the last step by \cref{lem_mp_equations}(iv) and (ii).
\end{proof}


\begin{corollary}[Normal equation]
  \label{cor_normal_equation}
  For $y \in Y$, an element $x \in X$ is a least squares solution of $Tx = y$ if and only
  if $T^{*}Tx = T^{*}y$.  If $x$ is a least squares solution with $x \in \Ker(T)^{\perp}$,
  then $y \in \Dom(T\dg)$ and $x = T\dg y$.
\end{corollary}

\begin{proof}
  Let $y \in Y$ be arbitrary and put $U := \overline{\Ran(T)}$.

  \emph{First statement.}  As in the proof of \cref{thm_mp_minnorm}, for every $x \in X$
  one has $Tx \in U$ and
  $\|P_Ry - y\| = \min_{w \in U}\|w - y\| \le \|Tx - y\|$ (\cref{prop_proj_props}, last
  clause), and
  $\inf_{x \in X}\|Tx - y\| = \inf_{w \in \Ran(T)}\|w - y\| = \min_{w \in U}\|w - y\|$,
  the last equality because $\Ran(T)$ is dense in $U$.  Hence $x$ is a least squares
  solution if and only if $\|Tx - y\|$ attains $\min_{w \in U}\|w-y\|$, which by the
  Pythagoras argument of \cref{thm_mp_minnorm}(i) --- and not by any convexity theorem,
  for the reason given there --- happens if and only if $Tx = P_Ry$.  By the
  characterisation of the projection (\cref{lem_proj_characterisation}), $Tx = P_Ry$ holds if
  and only if $Tx \in U$ --- automatic --- and $y - Tx \perp U$.  Finally
  $U^{\perp} = \Ran(T)^{\perp} = \Ker(T^{*})$ (\cref{thm_four_subspaces}), so this says
  $T^{*}(y - Tx) = 0$, i.e.\ $T^{*}Tx = T^{*}y$.

  Note that no hypothesis on $y$ was used: when $y \notin \Dom(T\dg)$ the projection $P_Ry$
  simply fails to lie in $\Ran(T)$ (\cref{lem_mp_domain}) and both sides are unsatisfiable.

  \emph{Second statement.}  Let $x$ be a least squares solution with
  $x \in \Ker(T)^{\perp}$.  By the first part $Tx = P_Ry$, so $P_Ry \in \Ran(T)$ and
  therefore $y \in \Dom(T\dg)$ by \cref{lem_mp_domain}.  Since $x \in \Ker(T)^{\perp}$ and
  $\tilde{T}x = Tx = P_Ry$, the definition of $T\dg$ gives
  $T\dg y = \tilde{T}^{-1}P_Ry = x$.
\end{proof}


\begin{lemma}[Minimum norm solution of the normal equation]
  \label{lem_normal_mp}
  For $y \in \Dom(T\dg)$ one has $T\dg y = (T^{*}T)\dg T^{*} y$.
\end{lemma}

\begin{proof}
  Write $S := T^{*}T \in \Ll(X,X)$ and $x\dg := T\dg y$.

  By \cref{thm_mp_minnorm}, $x\dg$ is a least squares solution of $Tx = y$, and by
  \cref{lem_mp_equations} it lies in $\Ran(T\dg) = \Ker(T)^{\perp}$.  Hence
  \cref{cor_normal_equation} gives
  \[
    Sx\dg = T^{*}Tx\dg = T^{*}y .
  \]
  In particular $T^{*}y \in \Ran(S)$, so $T^{*}y \in \Dom(S\dg)$ and
  $P_{\overline{\Ran(S)}}\,T^{*}y = T^{*}y$.

  Moreover the normal operator has the same kernel as $T$ itself,
  $\Ker(S) = \Ker(T^{*}T) = \Ker(T)$, by \cref{lem_kernel_normal};
  so $x\dg \in \Ker(T)^{\perp} = \Ker(S)^{\perp}$.  Thus $x\dg$ is the unique element of
  $\Ker(S)^{\perp}$ mapped by $S$ to $T^{*}y$, which is precisely
  $\tilde{S}^{-1}T^{*}y$.  By the definition of the Moore--Penrose
  inverse (\cref{thm_mp_inverse}) applied to $S$,
  \[
    S\dg T^{*}y = \tilde{S}^{-1}P_{\overline{\Ran(S)}}T^{*}y
      = \tilde{S}^{-1}T^{*}y = x\dg = T\dg y . \qedhere
  \]
\end{proof}


\begin{theorem}[Continuity of $T\dg$ characterises closed range]
  \label{thm_mp_continuous_iff_closed_range}
  $T\dg \colon \Dom(T\dg) \to X$ is bounded if and only if $\Ran(T)$ is closed.
\end{theorem}

\begin{proof}
  \emph{$T\dg$ is a closed operator.}  Let $\{y_n\} \subset \Dom(T\dg)$ with $y_n \to y$
  in $Y$ and $T\dg y_n \to x$ in $X$.  The projection $P_R$ is bounded
  (\cref{prop_proj_props}), so \cref{lem_mp_equations}(iv) gives
  \[
    P_Ry = \lim_{n} P_Ry_n = \lim_n TT\dg y_n = Tx,
  \]
  the last equality by continuity of $T$.  In particular $P_Ry \in \Ran(T)$, so
  $y \in \Dom(T\dg)$ by \cref{lem_mp_domain} --- this step has to be taken before
  \cref{thm_mp_minnorm} may be invoked.  Next, $T\dg y_n \in \Ran(T\dg) =
  \Ker(T)^{\perp}$ (\cref{lem_mp_equations}), a closed subspace, so $x \in
  \Ker(T)^{\perp}$.  Since $Tx = P_Ry$, $x$ is a least squares solution of $Tx = y$ by
  \cref{thm_mp_minnorm}(i), and being in $\Ker(T)^{\perp}$ it equals $T\dg y$ by
  \cref{cor_normal_equation}.  Hence the graph of $T\dg$ is closed.

  \emph{($\Leftarrow$)}  If $\Ran(T)$ is closed then $\Dom(T\dg) = Y$
  (\cref{lem_mp_domain_dense}), so $T\dg$ is an everywhere-defined closed linear map
  between the Hilbert --- hence Banach --- spaces $Y$ and $X$, and the closed graph theorem
  \cref{thm_closed_graph} makes it bounded.

  \emph{($\Rightarrow$)}  Suppose $T\dg$ is bounded on $\Dom(T\dg)$.  That domain is dense
  (\cref{lem_mp_domain_dense}) and $X$ is complete, so $T\dg$ extends to a bounded
  $S \in \Ll(Y,X)$ with $S|_{\Dom(T\dg)} = T\dg$ (\cref{thm_blt}).  Let
  $y \in \overline{\Ran(T)}$ and choose $y_n \in \Ran(T)$ with $y_n \to y$.  For each $n$,
  $y_n \in \Dom(T\dg)$ and $P_Ry_n = y_n$, so \cref{lem_mp_equations}(iv) gives
  $TT\dg y_n = y_n$.  Passing to the limit, using continuity of $S$ and of $T$,
  \[
    y = \lim_n y_n = \lim_n T S y_n = T S y \in \Ran(T).
  \]
  Hence $\overline{\Ran(T)} \subseteq \Ran(T)$, i.e.\ $\Ran(T)$ is closed.
\end{proof}


\begin{corollary}[Compact operators with infinite-dimensional range]
  \label{cor_mp_compact_discontinuous}
  If $K \in \Kk(X,Y)$ has infinite-dimensional range, then $K\dg$ is unbounded.
  Equivalently: for compact $K$, $\Ran(K)$ is closed if and only if it is
  finite-dimensional.
\end{corollary}

\begin{proof}
  Suppose $K\dg$ were bounded.  By \cref{thm_mp_continuous_iff_closed_range}, $\Ran(K)$ is
  then closed.  But a compact operator with closed range has finite-dimensional range
  (\cref{thm_compact_range_not_closed}), contradicting the hypothesis.

  The equivalence is the same statement read twice: if $\Ran(K)$ is closed then it is
  finite-dimensional by \cref{thm_compact_range_not_closed}, and if it is finite-dimensional
  then it is closed (\cref{lem_findim_closed}).
\end{proof}


\begin{definition}[Ill-posedness in the sense of Nashed]
  \label{def_nashed}
  The equation $Tx = y$ is \emph{ill-posed in the sense of Nashed} if $\Ran(T)$ is not
  closed.
\end{definition}


\begin{lemma}[The four equations characterise the Moore--Penrose inverse]
  \label{lem_mp_characterisation}
  Let $\Ran(T)$ be closed and let $S \in \Ll(Y,X)$ satisfy $TST = T$, $STS = S$, and
  $ST$, $TS$ self-adjoint.  Then $S = T\dg$.
\end{lemma}

\begin{proof}
  Since $\Ran(T)$ is closed, $\Dom(T\dg) = Y$ and $T\dg \in \Ll(Y,X)$
  (\cref{lem_mp_domain_dense}, \cref{thm_mp_continuous_iff_closed_range}), so all products
  below are everywhere defined.

  \emph{A self-adjoint idempotent is an orthogonal projection.}  Let $P \in \Ll(H,H)$
  satisfy $P^2 = P$ and $P^{*} = P$.  Then $\Ran(P) = \Ker(\Id - P)$ is closed
  (\cref{lem_ker_closed}), and for $w \in H$ and $v \in H$,
  $\langle Pv, w - Pw\rangle = \langle v, P(w - Pw)\rangle = \langle v, Pw - P^2w\rangle
  = 0$ (\cref{def_adjoint}).  So $Pw \in \Ran(P)$ and $w - Pw \perp \Ran(P)$, whence
  $Pw = P_{\Ran(P)}w$ by \cref{lem_proj_characterisation}.

  \emph{$ST = T\dg T$.}  From $TST = T$ we get $(ST)^2 = S(TST) = ST$, so $ST$ is
  idempotent, and it is self-adjoint by hypothesis; by the previous paragraph it is the
  orthogonal projection onto $\Ran(ST)$.  Now $\Ran(ST) \subseteq \Ran(S)$, while
  $STS = S$ gives $\Ran(S) = \Ran(STS) \subseteq \Ran(ST)$, so $\Ran(ST) = \Ran(S)$.
  Also $\Ker(T) \subseteq \Ker(ST)$, while $TST = T$ gives
  $\Ker(ST) \subseteq \Ker(TST) = \Ker(T)$, so $\Ker(ST) = \Ker(T)$.  For an orthogonal
  projection $\Ran(ST) = \Ker(ST)^{\perp}$ (\cref{prop_proj_props}), hence
  $ST = P_{\Ker(T)^{\perp}} = \Id - P_N$, which is $T\dg T$ by
  \cref{lem_mp_equations}(iii).

  \emph{$TS = TT\dg$.}  Symmetrically, $(TS)^2 = T(STS) = TS$ and $TS$ is self-adjoint, so
  $TS$ is the orthogonal projection onto $\Ran(TS)$.  From $TST = T$ we get
  $\Ran(T) \subseteq \Ran(TS)$, and $\Ran(TS) \subseteq \Ran(T)$ because
  $TSy = T(Sy)$ for every $y$, so
  $\Ran(TS) = \Ran(T) = \overline{\Ran(T)}$ by closedness.  Hence $TS = P_R = TT\dg$ by
  \cref{lem_mp_equations}(iv).

  \emph{Conclusion.}  Using $STS = S$, then $TS = TT\dg$, then
  \cref{lem_mp_equations}(ii),
  \[
    S = S(TS) = S\,TT\dg = (ST)T\dg = T\dg T T\dg = T\dg . \qedhere
  \]
\end{proof}


\begin{definition}[The Moore--Penrose inverse as a bounded operator]
  \label{def_moore_penrose_bounded}
  Suppose $\Ran(T)$ is closed.  Then $\Dom(T\dg) = Y$ by \cref{lem_mp_domain_dense} and
  $T\dg$ is bounded by \cref{thm_mp_continuous_iff_closed_range}, so $T\dg$ is an element
  of $\Ll(Y,X)$.  We write $T\dg$ for that element as well.
\end{definition}


\begin{lemma}[Adjoint of the Moore--Penrose inverse]
  \label{lem_mp_adjoint}
  If $\Ran(T)$ is closed, then $T\dg \in \Ll(Y,X)$ and $(T\dg)^{*} = (T^{*})\dg$.
\end{lemma}

\begin{proof}
  Since $\Ran(T)$ is closed, $T\dg \in \Ll(Y,X)$ by
  \cref{thm_mp_continuous_iff_closed_range} and \cref{lem_mp_domain_dense}, and
  $\Ran(T^{*}) = \Ker(T)^{\perp}$ is closed as well (\cref{lem_closed_range_adjoint}, which
  is a statement about the adjoint; \cref{thm_closed_range} is a criterion on $T$ alone and
  says nothing here), so $(T^{*})\dg \in \Ll(X,Y)$ for the same reason.  Put
  $S := (T\dg)^{*} \in \Ll(X,Y)$.  We verify the four hypotheses of
  \cref{lem_mp_characterisation} for $S$ with respect to $T^{*}$, using throughout that
  adjunction reverses products, $(AB)^{*} = B^{*}A^{*}$, and is involutive
  (\cref{prop_adjoint_props}).
  \begin{itemize}
    \item $T^{*}ST^{*} = T^{*}(T\dg)^{*}T^{*} = (TT\dg T)^{*} = T^{*}$, by
      \cref{lem_mp_equations}(i).
    \item $ST^{*}S = (T\dg)^{*}T^{*}(T\dg)^{*} = (T\dg TT\dg)^{*} = (T\dg)^{*} = S$, by
      \cref{lem_mp_equations}(ii).
    \item $ST^{*} = (T\dg)^{*}T^{*} = (TT\dg)^{*} = P_R^{*} = P_R$,
      by \cref{lem_mp_equations}(iv) and self-adjointness of
      orthogonal projections (\cref{prop_proj_props}).  So $ST^{*}$ is
      self-adjoint.
    \item $T^{*}S = T^{*}(T\dg)^{*} = (T\dg T)^{*} = (\Id - P_N)^{*} = \Id - P_N$
      by \cref{lem_mp_equations}(iii) and \cref{prop_proj_props}; in particular
      $T^{*}S$ is self-adjoint.
  \end{itemize}
  As $\Ran(T^{*})$ is closed, \cref{lem_mp_characterisation} applies to $T^{*}$ and gives
  $S = (T^{*})\dg$, i.e.\ $(T\dg)^{*} = (T^{*})\dg$.
\end{proof}


\chapter{Singular values, the Picard criterion, and source conditions}
\label{ch_12_picard}

\noindent\emph{Background.} The Picard criterion --- solvability decided by the decay of
the data's coefficients against the singular values --- goes back to \citet{Picard1910},
who established it for Fredholm integral equations of the first kind; it is the precise
account of why the classical inverse $K^{-1}$, defined term-by-term on the singular
system, fails to be everywhere defined or continuous. The resulting split of ill-posed
problems into moderately (polynomial decay of $\sigma_n$) and severely, in particular
exponentially, ill-posed ones, together with the power-scale source conditions
$x\dg \in X_{\nu,\rho}$ that quantify how much smoother than merely being in
$\Dom(K\dg)$ the sought solution is assumed to be, are standard in the regularization
literature and are treated systematically in the monograph of
\citet{EnglHankeNeubauer1996}, the main source \citet{Clason2021} follows.

\begin{assumptions}
  \label{asm_singular_values_picard}
  $X$, $Y$ are Hilbert spaces and $K \in \Kk(X,Y)$ is
compact.  $\{(\sigma_n, u_n, v_n)\}_{n \in I}$ denotes a singular system of $K$ in the
sense \cref{def_singular_system}, with the second and
third components written in the opposite order to that document --- the roles are spelled
out at once and never change, $v_n \in X$ and $u_n \in Y$ in both --- so: $I \subseteq
\mathbb{N}$ is finite or all of $\mathbb{N}$, $\sigma_1 \ge \sigma_2 \ge \cdots > 0$ is
non-increasing and is a null sequence when $I$ is infinite, $\{v_n\}$ is an orthonormal
basis of $\overline{\Ran(K^{*})} = \Ker(K)^{\perp}$, $\{u_n\}$ is an orthonormal basis of
$\overline{\Ran(K)}$, and $Kv_n = \sigma_n u_n$, $K^{*}u_n = \sigma_n v_n$.  Existence is
\cref{thm_svd}; the functional calculus $\varphi(K^{*}K)$ is \cref{def_fun_calc} and the
absolute value $|K| = (K^{*}K)^{1/2}$ is \cref{def_frac_power}.  Note that the whole of this
chapter and the four that follow depends on the singular system being available, hence on
$K$ being \emph{compact}; $\kappa := \|K\|^{2} = \|K^{*}K\|$ throughout, the second
equality being \cref{prop_adjoint_props}.
\end{assumptions}

\begin{theorem}[Picard criterion]
  \label{thm_picard}
  Let $y \in \overline{\Ran(K)}$.  Then $y \in \Ran(K)$ if and only if
  \[
    \sum_{n \in I} \sigma_n^{-2}\,|\langle u_n, y\rangle_Y|^{2} < \infty ,
  \]
  and in that case $K\dg y = \sum_{n \in I} \sigma_n^{-1}\langle u_n, y\rangle_Y\, v_n$.
\end{theorem}

\begin{proof}
  Throughout, $\{(\sigma_n, u_n, v_n)\}_{n \in I}$ is the singular system fixed in the
  conventions of this chapter, and $P_R$ denotes the orthogonal projection onto
  $\overline{\Ran(K)}$.  If $I$ is finite, every series below is a finite sum and every
  convergence assertion is vacuous; the argument is written for $I = \mathbb{N}$ and reads
  verbatim in the finite case.

  \emph{A coefficient identity.}  For $x \in X$ and $n \in I$ the adjoint relation
  (\cref{def_adjoint}), the identity $K^{*}u_n = \sigma_n v_n$, and the fact that
  $\sigma_n > 0$ is real give
  \begin{equation}\label{eq_picard_coeff}
    \langle u_n, Kx\rangle_Y = \langle K^{*}u_n, x\rangle_X
      = \langle \sigma_n v_n, x\rangle_X = \sigma_n\langle v_n, x\rangle_X .
  \end{equation}

  \emph{($\Rightarrow$).}  Let $y = Kx$ for some $x \in X$.  Then
  \eqref{eq_picard_coeff} reads $\sigma_n^{-1}\langle u_n, y\rangle_Y = \langle v_n,
  x\rangle_X$ for every $n \in I$, whence
  \[
    \sum_{n \in I}\sigma_n^{-2}\,|\langle u_n, y\rangle_Y|^{2}
      = \sum_{n \in I}|\langle v_n, x\rangle_X|^{2} \le \|x\|_X^{2} < \infty
  \]
  by Bessel's inequality (\cref{thm_bessel}) for the orthonormal system $\{v_n\}$.

  \emph{($\Leftarrow$).}  Assume the sum is finite, and for $n \in I$ set
  $c_n := \sigma_n^{-1}\langle u_n, y\rangle_Y$; then $\sum_{n \in I}|c_n|^{2} <
  \infty$.  Since $\{v_n\}$ is orthonormal, \cref{lem_riesz_fischer} shows that
  \[
    x := \sum_{n \in I} c_n v_n
  \]
  converges in $X$.  Pushing the bounded operator $K$ through the series
  (\cref{lem_operator_through_series}) and using $Kv_n = \sigma_n u_n$ together with
  $c_n\sigma_n = \langle u_n, y\rangle_Y$,
  \[
    Kx = \sum_{n \in I} c_n\,Kv_n = \sum_{n \in I}\langle u_n, y\rangle_Y\, u_n .
  \]
  As $\{u_n\}$ is an orthonormal basis of $\overline{\Ran(K)}$, the right-hand side is
  $P_R y$ (\cref{thm_proj_expansion}).  Since $y \in \overline{\Ran(K)}$ and
  $y - y = 0 \perp \overline{\Ran(K)}$, the characterisation of the orthogonal projection
  (\cref{lem_proj_characterisation}) gives $P_R y = y$.  Hence $Kx = y$, so
  $y \in \Ran(K)$.

  \emph{The formula.}  Suppose the two equivalent conditions hold; by the paragraph above
  there is $x = \sum_{n \in I}c_n v_n$ with $Kx = y$.  This $x$ lies in
  $\Ker(K)^{\perp}$: each $v_n$ belongs to $\Ker(K)^{\perp}$, because $\{v_n\}$ is an
  orthonormal basis of that subspace, so for $z \in \Ker(K)$ term-by-term evaluation of
  the inner product (\cref{lem_inner_through_series}) gives
  \[
    \langle z, x\rangle_X = \sum_{n \in I} c_n\,\langle z, v_n\rangle_X = 0 .
  \]
  Moreover $\|Kx - y\|_Y = 0 \le \|Kz - y\|_Y$ for every $z \in X$, so $x$ is a least
  squares solution of $Kx = y$ in the sense of \cref{def_lsq_minnorm}.  Being a least
  squares solution lying in $\Ker(K)^{\perp}$, it satisfies $y \in \Dom(K\dg)$ and
  \[
    K\dg y = x = \sum_{n \in I}\sigma_n^{-1}\langle u_n, y\rangle_Y\, v_n
  \]
  by \cref{cor_normal_equation}.
\end{proof}


\begin{corollary}[Picard representation on the whole domain]
  \label{cor_picard_domain}
  For every $y \in \Dom(K\dg)$,
  \[
    K\dg y = \sum_{n \in I} \sigma_n^{-1}\langle u_n, y\rangle_Y\, v_n .
  \]
\end{corollary}

\begin{proof}
  Let $y \in \Dom(K\dg)$ and let $P_R$ be the orthogonal projection onto
  $\overline{\Ran(K)}$.  By \cref{lem_mp_domain} we have $P_R y \in \Ran(K)$; in
  particular $P_R y \in \overline{\Ran(K)}$, and
  $P_R y \in \Ran(K) \subseteq \Ran(K) \oplus \Ran(K)^{\perp} = \Dom(K\dg)$
  (\cref{thm_mp_inverse}).

  \emph{The coefficients are unchanged by $P_R$.}  Each $u_n$ lies in
  $\overline{\Ran(K)}$ and $u_n - u_n = 0 \perp \overline{\Ran(K)}$, so $P_R u_n = u_n$
  (\cref{lem_proj_characterisation}); and $P_R$ is self-adjoint (\cref{prop_proj_props}).
  Hence, for every $n \in I$,
  \[
    \langle u_n, y\rangle_Y = \langle P_R u_n, y\rangle_Y
      = \langle u_n, P_R y\rangle_Y .
  \]

  \emph{Reduction to \cref{thm_picard}.}  The datum $P_R y$ lies in
  $\overline{\Ran(K)}$ and in $\Ran(K)$, so \cref{thm_picard} applies to it and gives
  \[
    K\dg P_R y = \sum_{n \in I}\sigma_n^{-1}\langle u_n, P_R y\rangle_Y\, v_n
      = \sum_{n \in I}\sigma_n^{-1}\langle u_n, y\rangle_Y\, v_n ,
  \]
  the series converging in $X$.  Finally \cref{lem_mp_equations}(iv), applied to
  $y \in \Dom(K\dg)$, gives $P_R y = KK\dg y$, and therefore
  \[
    \sum_{n \in I}\sigma_n^{-1}\langle u_n, y\rangle_Y\, v_n
      = K\dg P_R y = K\dg K K\dg y = K\dg y ,
  \]
  the last equality by \cref{lem_mp_equations}(ii).
\end{proof}


\begin{lemma}[The Moore--Penrose inverse on singular vectors]
  \label{lem_mp_singular_vector}
  For every $n \in I$ one has $u_n \in \Dom(K\dg)$ and $K\dg u_n = \sigma_n^{-1} v_n$.
\end{lemma}

\begin{proof}
  Fix $n \in I$.  Since $\sigma_n > 0$ and $Kv_n = \sigma_n u_n$, we have
  $u_n = K(\sigma_n^{-1}v_n) \in \Ran(K)$, and
  $\Ran(K) \subseteq \Ran(K) \oplus \Ran(K)^{\perp} = \Dom(K\dg)$
  (\cref{thm_mp_inverse}); hence $u_n \in \Dom(K\dg)$.

  \cref{cor_picard_domain} applied to the datum $u_n$ now gives
  \[
    K\dg u_n = \sum_{m \in I}\sigma_m^{-1}\langle u_m, u_n\rangle_Y\, v_m
      = \sigma_n^{-1} v_n ,
  \]
  because $\{u_m\}$ is orthonormal, so $\langle u_m, u_n\rangle_Y = 0$ for $m \ne n$ and
  $\langle u_n, u_n\rangle_Y = 1$: every term of the series except the $n$-th vanishes.
  Since $\sigma_n^{-1} > 0$ is real and $\|v_n\|_X = 1$, this yields
  $\|K\dg u_n\|_X = \sigma_n^{-1}\|v_n\|_X = \sigma_n^{-1}$.
\end{proof}




\begin{definition}[Degree of ill-posedness]
  \label{def_degree_illposed}
  $K$ is \emph{moderately ill-posed} if there are $c, r > 0$ with $\sigma_n \ge c\,n^{-r}$
  for all $n \in I$, and \emph{severely ill-posed} otherwise.  It is \emph{exponentially
  ill-posed} if there are $c, r > 0$ with $\sigma_n \le c\,e^{-n^{r}}$ for all $n \in I$.
\end{definition}


\begin{definition}[Source sets]
  \label{def_source_set}
  For $\nu, \rho > 0$ put
  \[
    X_{\nu,\rho} := \{\, |K|^{\nu} w : w \in X,\ \|w\|_X \le \rho \,\},
    \qquad
    X_{\nu} := \bigcup_{\rho > 0} X_{\nu,\rho} = \Ran(|K|^{\nu}) .
  \]
  The assumption $x\dg \in X_{\nu,\rho}$ is a \emph{source condition}, and any $w$ with
  $|K|^{\nu}w = x\dg$ is a \emph{source representer}.
\end{definition}


\begin{lemma}[Source sets lie in the orthogonal complement of the kernel]
  \label{lem_source_in_kerperp}
  $X_{\nu,\rho} \subseteq \overline{\Ran(K^{*})} = \Ker(K)^{\perp}$ for all
  $\nu, \rho > 0$.
\end{lemma}

\begin{proof}
  Fix $\nu, \rho > 0$ and let $x \in X_{\nu,\rho}$, say $x = |K|^{\nu}w$ with $w \in X$
  and $\|w\|_X \le \rho$ (\cref{def_source_set}).  By \cref{def_frac_power} the fractional
  power $|K|^{\nu}$ is the functional calculus (\cref{def_fun_calc}) of $K^{*}K$ for the
  function $\varphi(t) = t^{\nu/2}$; as $\nu > 0$ we have $\varphi(0) = 0$, so the summand
  supported on $\Ker(K)$ is absent and
  \begin{equation}\label{eq_frac_power_series}
    |K|^{\nu}w = \sum_{n \in I}\sigma_n^{\nu}\langle v_n, w\rangle_X\, v_n ,
  \end{equation}
  the series converging in $X$.

  Let $z \in \Ker(K)$.  Each $v_n$ lies in $\Ker(K)^{\perp}$, since $\{v_n\}$ is an
  orthonormal basis of $\Ker(K)^{\perp}$, so $\langle z, v_n\rangle_X = 0$ for every
  $n \in I$.  Evaluating the inner product of $z$ against
  \eqref{eq_frac_power_series} term by term (\cref{lem_inner_through_series}),
  \[
    \langle z, x\rangle_X
      = \sum_{n \in I}\sigma_n^{\nu}\langle v_n, w\rangle_X\,\langle z, v_n\rangle_X = 0 .
  \]
  As $z \in \Ker(K)$ was arbitrary, $x \in \Ker(K)^{\perp}$.  Since $x \in X_{\nu,\rho}$
  was arbitrary, $X_{\nu,\rho} \subseteq \Ker(K)^{\perp}$, and
  $\Ker(K)^{\perp} = \overline{\Ran(K^{*})}$ by \cref{thm_four_subspaces}.
\end{proof}


\begin{lemma}[Source condition as a strengthened Picard condition]
  \label{lem_source_picard}
  Let $\nu > 0$ and $y \in \Dom(K\dg)$.  Then $K\dg y \in X_{\nu}$ if and only if
  \[
    \sum_{n \in I} \sigma_n^{-2-2\nu}\,|\langle u_n, y\rangle_Y|^{2} < \infty .
  \]
\end{lemma}

\begin{proof}
  Write $x\dg := K\dg y$ and $b_n := \langle u_n, y\rangle_Y$ for $n \in I$.  If $I$ is
  finite, all series below are finite sums and every convergence assertion is vacuous.

  \emph{Coefficients of $x\dg$.}  By \cref{cor_picard_domain},
  $x\dg = \sum_{n \in I}\sigma_n^{-1}b_n\,v_n$, the series converging in $X$.  Evaluating
  the inner product against $v_m$ term by term (\cref{lem_inner_through_series}) and using
  orthonormality of $\{v_n\}$,
  \begin{equation}\label{eq_source_coeff}
    \langle v_m, x\dg\rangle_X = \sigma_m^{-1}b_m \qquad (m \in I).
  \end{equation}

  \emph{Coefficients of $|K|^{\nu}w$.}  By \cref{def_frac_power} the fractional power
  $|K|^{\nu}$ is the functional calculus (\cref{def_fun_calc}) of $K^{*}K$ for the function
  $\varphi(t) = t^{\nu/2}$; as $\nu > 0$ we have $\varphi(0) = 0$, so the summand
  supported on $\Ker(K)$ is absent and
  \begin{equation}\label{eq_source_frac}
    |K|^{\nu}w = \sum_{n \in I}\sigma_n^{\nu}\langle v_n, w\rangle_X\, v_n ,
      \qquad w \in X ,
  \end{equation}
  the series converging in $X$.  Again by \cref{lem_inner_through_series} and
  orthonormality of $\{v_n\}$, this gives
  $\langle v_m, |K|^{\nu}w\rangle_X = \sigma_m^{\nu}\langle v_m, w\rangle_X$ for every
  $m \in I$.

  \emph{($\Rightarrow$).}  Let $x\dg \in X_{\nu} = \Ran(|K|^{\nu})$
  (\cref{def_source_set}) and pick $w \in X$ with $|K|^{\nu}w = x\dg$.  Comparing
  coefficients, the last display and \eqref{eq_source_coeff} give
  $\sigma_m^{\nu}\langle v_m, w\rangle_X = \sigma_m^{-1}b_m$, that is
  $\langle v_m, w\rangle_X = \sigma_m^{-1-\nu}b_m$, for every $m \in I$.  Bessel's
  inequality (\cref{thm_bessel}) for the orthonormal system $\{v_n\}$ now yields
  \[
    \sum_{n \in I}\sigma_n^{-2-2\nu}\,|b_n|^{2}
      = \sum_{n \in I}|\langle v_n, w\rangle_X|^{2} \le \|w\|_X^{2} < \infty .
  \]

  \emph{($\Leftarrow$).}  Assume $\sum_{n \in I}\sigma_n^{-2-2\nu}|b_n|^{2} < \infty$ and
  put $c_n := \sigma_n^{-1-\nu}b_n$, so that $\sum_{n \in I}|c_n|^{2} < \infty$.  Since
  $\{v_n\}$ is orthonormal, \cref{lem_riesz_fischer} shows that $w := \sum_{n \in I}c_n v_n$
  converges in $X$, and $\langle v_m, w\rangle_X = c_m$ for every $m \in I$ by
  \cref{lem_inner_through_series} and orthonormality.  Substituting this into
  \eqref{eq_source_frac} and using $\sigma_n^{\nu}c_n = \sigma_n^{-1}b_n$,
  \[
    |K|^{\nu}w = \sum_{n \in I}\sigma_n^{\nu}c_n\,v_n
      = \sum_{n \in I}\sigma_n^{-1}b_n\,v_n = x\dg ,
  \]
  so $x\dg \in \Ran(|K|^{\nu}) = X_{\nu}$ (\cref{def_source_set}).
\end{proof}


\chapter{Regularization and parameter choice}
\label{ch_13_regularization}

\noindent\emph{Background.} The abstract framework of this chapter --- approximating the
ill-posed inverse $T\dg$ by a family $\{R_{\alpha}\}$ of continuous operators, and
choosing $\alpha$ from the noise level and, if needed, the data --- is the basic
machinery of regularization theory as laid out by \citet{TikhonovArsenin1977}, the book
that gave the field its name. The discrepancy principle, choosing $\alpha$ so that the
residual matches the noise level, is due to \citet{Morozov1966}. That a heuristic rule,
one using the data alone and never the noise level $\delta$, cannot yield a convergent
method unless $T\dg$ is already continuous --- the Bakushinski\u{\i} veto, motivating why
a parameter choice rule is allowed to see $\delta$ at all --- is due to
\citet{Bakushinskii1985}.

\begin{assumptions}
  \label{asm_regularization_parameter_choice}
  $X$, $Y$ are Hilbert spaces and $T \in \Ll(X,Y)$;
compactness is \emph{not} assumed anywhere in this chapter.  The exact datum is
$y \in \Dom(T\dg)$ with minimum norm solution $x\dg := T\dg y$, and the noisy datum is any
$y^{\delta} \in B_{\delta}(y)$, i.e.\ $\|y - y^{\delta}\|_Y \le \delta$ with
$\delta > 0$; note $y^{\delta}$ is \emph{not} assumed to lie in $\Dom(T\dg)$.  We write
$x^{\delta}_{\alpha} := R_{\alpha} y^{\delta}$.
\end{assumptions}

\begin{definition}[Regularization]
  \label{def_regularization}
  A family $\{R_{\alpha}\}_{\alpha > 0} \subseteq \Ll(Y,X)$ is a \emph{regularization} of
  $T\dg$ if $R_{\alpha} y \to T\dg y$ as $\alpha \to 0$, for every $y \in \Dom(T\dg)$.
\end{definition}


\begin{lemma}[Error decomposition]
  \label{lem_error_decomposition}
  For $y \in \Dom(T\dg)$, $y^{\delta} \in B_{\delta}(y)$ and $\alpha > 0$,
  \[
    \|R_{\alpha} y^{\delta} - T\dg y\|_X
      \;\le\; \delta\,\|R_{\alpha}\|_{\Ll(Y,X)} \;+\; \|R_{\alpha} y - T\dg y\|_X .
  \]
  The first term is the \emph{propagated data error}, the second the
  \emph{approximation error}.
\end{lemma}

\begin{proof}
  Since $y \in \Dom(T\dg)$, the element $T\dg y$ is defined (\cref{thm_mp_inverse}), and
  $R_{\alpha} \in \Ll(Y,X)$ by \cref{def_regularization}, so
  $\|R_{\alpha}\|_{\Ll(Y,X)}$ is finite.  Inserting $R_{\alpha}y$ and using the triangle
  inequality,
  \[
    \|R_{\alpha}y^{\delta} - T\dg y\|_X
      \le \|R_{\alpha}y^{\delta} - R_{\alpha}y\|_X + \|R_{\alpha}y - T\dg y\|_X .
  \]
  By linearity $R_{\alpha}y^{\delta} - R_{\alpha}y = R_{\alpha}(y^{\delta} - y)$, and the
  operator norm bounds $\|R_{\alpha}z\|_X \le \|R_{\alpha}\|_{\Ll(Y,X)}\|z\|_Y$
  (\cref{lem_isleast_opnorm}).  With $z := y^{\delta} - y$, whose norm is at most $\delta$
  because $y^{\delta} \in B_{\delta}(y)$,
  \[
    \|R_{\alpha}y^{\delta} - R_{\alpha}y\|_X \le \delta\,\|R_{\alpha}\|_{\Ll(Y,X)} .
      \qedhere
  \]
\end{proof}


\begin{theorem}[A regularization is never uniformly bounded]
  \label{thm_reg_not_uniformly_bounded}
  Let $\{R_{\alpha}\}_{\alpha>0}$ be a regularization of $T\dg$, and suppose that
  $T\dg$ is unbounded.  Then
  \[
    \limsup_{\alpha \to 0}\ \|R_{\alpha}\|_{\Ll(Y,X)} = \infty ,
  \]
  and there exist $y \in Y$ and a null sequence $\{\alpha_n\}$ along which
  $\|R_{\alpha_n} y\|_X \to \infty$.
\end{theorem}

\begin{proof}
  Throughout, the limes superior along $\alpha \to 0$ is understood as
  \[
    \limsup_{\alpha \to 0}\|R_{\alpha}\|_{\Ll(Y,X)}
      := \inf_{\beta > 0}\ \sup_{0 < \alpha \le \beta}\|R_{\alpha}\|_{\Ll(Y,X)}
      \in [0,\infty] .
  \]

  \emph{The limes superior is infinite.}  Suppose not.  Then there are $\beta > 0$ and
  $M < \infty$ with $\|R_{\alpha}\|_{\Ll(Y,X)} \le M$ for all $\alpha \in (0,\beta]$.  Let
  $y \in \Dom(T\dg)$ and $\varepsilon > 0$.  By \cref{def_regularization} there is
  $\alpha \in (0,\beta]$ with $\|R_{\alpha}y - T\dg y\|_X \le \varepsilon$, whence
  \[
    \|T\dg y\|_X \le \|R_{\alpha}y\|_X + \varepsilon
      \le \|R_{\alpha}\|_{\Ll(Y,X)}\|y\|_Y + \varepsilon \le M\|y\|_Y + \varepsilon
  \]
  by \cref{lem_isleast_opnorm}.  As $\varepsilon > 0$ was arbitrary,
  $\|T\dg y\|_X \le M\|y\|_Y$ for every $y$ in the subspace $\Dom(T\dg)$
  (\cref{thm_mp_inverse}), i.e.\ $T\dg$ is bounded, contrary to hypothesis.

  \emph{Choice of a null sequence of parameters.}  By the previous paragraph, for
  \emph{every} $\beta > 0$ the supremum of $\|R_{\alpha}\|_{\Ll(Y,X)}$ over
  $\alpha \in (0,\beta]$ is infinite.
  Hence we may choose recursively $\alpha_1 > \alpha_2 > \cdots > 0$ such that
  $\alpha_n \le 1/n$ and $\|R_{\alpha_n}\|_{\Ll(Y,X)} \ge 4^{n}$ for every $n$: apply
  the previous sentence with $\beta := 1$ when $n = 1$, and with
  $\beta := \min\{1/n,\ \alpha_{n-1}/2\}$ when $n > 1$.  Write
  $A_n := R_{\alpha_n} \in \Ll(Y,X)$; then $\{\alpha_n\}$ is a null sequence and
  $\|A_n\| \ge 4^{n} > 0$.

  \emph{The resonance point.}  The family $\{A_n\}_{n \in \mathbb{N}}$ satisfies
  $\sup_n\|A_n\|_{\Ll(Y,X)} = \infty$, and $Y$ is complete (\cref{def_hilbert}).  By the
  resonance form of the uniform boundedness principle, \cref{thm_resonance}, there is
  therefore a single $y \in Y$ with $\sup_n\|A_ny\|_X = \infty$; passing to a subsequence
  of $\{\alpha_n\}$ if necessary, $\|A_ny\|_X \to \infty$.

  Hence $\|R_{\alpha_n}y\|_X \to \infty$ along the null sequence $\{\alpha_n\}$.
\end{proof}


\begin{theorem}[Divergence outside the domain]
  \label{thm_reg_diverges_outside_domain}
  Let $\{R_{\alpha}\}_{\alpha>0}$ be a regularization of $T\dg$ with
  $\sup_{\alpha>0}\|TR_{\alpha}\|_{\Ll(Y,Y)} < \infty$.  Then
  $\|R_{\alpha}y\|_X \to \infty$ as $\alpha \to 0$, for every $y \notin \Dom(T\dg)$.
\end{theorem}

\begin{proof}
  Write $C := \sup_{\alpha > 0}\|TR_{\alpha}\|_{\Ll(Y,Y)} < \infty$.

  \emph{Step 1: $TR_{\alpha}z \to P_Rz$ as $\alpha \to 0$, for every $z \in Y$.}  Let
  first $\tilde{y} \in \Dom(T\dg)$.  Then $R_{\alpha}\tilde{y} \to T\dg\tilde{y}$ by
  \cref{def_regularization}, and $T$ is continuous, so
  $TR_{\alpha}\tilde{y} \to TT\dg\tilde{y} = P_R\tilde{y}$ by
  \cref{lem_mp_equations}(iv).  Thus $\{TR_{\alpha}\}$ converges to $P_R$ pointwise on
  $M := \Dom(T\dg)$, a linear subspace of $Y$ (\cref{thm_mp_inverse}) that is dense
  (\cref{lem_mp_domain_dense}) and therefore total.  Now $Y$ is a Hilbert space and so
  Banach, $\{TR_{\alpha}\}_{\alpha>0} \subseteq \Ll(Y,Y)$ satisfies
  $\sup_{\alpha>0}\|TR_{\alpha}\| = C < \infty$, and $P_R \in \Ll(Y,Y)$; so
  \cref{prop_dense_ptwise_filter}, applied at the filter $\mathcal{N}(0^{+})$ used here,
  gives $TR_{\alpha}z \to P_Rz$ for every $z \in Y$.

  \emph{Step 2.}  Let $y \notin \Dom(T\dg)$ and suppose, contrary to the claim, that
  $\|R_{\alpha}y\|_X \not\to \infty$ as $\alpha \to 0$.  Unfolding that negation: there is
  $M < \infty$ such that for every $n \in \mathbb{N}$ there is $\alpha_n \in (0,1/n)$ with
  $\|R_{\alpha_n}y\|_X \le M$.  Thus $\{\alpha_n\}$ is a null sequence and the sequence
  $x_n := R_{\alpha_n}y$ is bounded in $X$, so it has a weakly convergent subsequence
  $x_{n_k} \wc x \in X$ (\cref{thm_weak_compactness}).  Bounded operators are weakly
  sequentially continuous (\cref{lem_weak_seq_cont}), whence $Tx_{n_k} \wc Tx$.

  On the other hand $\alpha_{n_k} \to 0$, so Step 1 gives
  $Tx_{n_k} = TR_{\alpha_{n_k}}y \to P_Ry$ in norm, hence $Tx_{n_k} \wc P_Ry$
  (\cref{lem_strong_to_weak}).  Weak limits are unique (\cref{lem_weak_limit_unique}), so
  $Tx = P_Ry$.  In particular $P_Ry \in \Ran(T)$, and therefore $y \in \Dom(T\dg)$ by
  \cref{lem_mp_domain} --- contradicting the choice of $y$.
\end{proof}


\begin{definition}[Parameter choice rule; convergent regularization method]
  \label{def_parameter_choice}
  A \emph{parameter choice rule} is a map $\alpha \colon (0,\infty) \times Y \to
  (0,\infty)$.  It is \emph{a priori} if it depends only on $\delta$, \emph{a posteriori}
  if it depends on both arguments, and \emph{heuristic} if it depends only on
  $y^{\delta}$.  Given a regularization $\{R_{\alpha}\}$, the pair $(R_{\alpha}, \alpha)$
  is a \emph{convergent regularization method} if for every $y \in \Dom(T\dg)$
  \[
    \lim_{\delta \to 0}\ \sup_{y^{\delta} \in B_{\delta}(y)}
      \bigl\| R_{\alpha(\delta, y^{\delta})} y^{\delta} - T\dg y \bigr\|_X = 0 .
  \]
\end{definition}


\begin{theorem}[Every regularization admits an a priori rule]
  \label{thm_apriori_exists}
  Let $\{R_{\alpha}\}_{\alpha>0}$ be a regularization of $T\dg$.  Then there is an a
  priori choice rule $\alpha$ such that $(R_{\alpha}, \alpha)$ is a convergent
  regularization method.
\end{theorem}

\begin{proof}
  \emph{Construction of the rule.}  For $n \in \mathbb{N}$, $n \ge 1$, put
  $c_n := \|R_{1/n}\|_{\Ll(Y,X)}$, which is finite because $R_{1/n} \in \Ll(Y,X)$
  (\cref{def_regularization}), and
  \[
    d_n := \frac{1}{n\,(1 + c_n)} \in (0, 1/n], \qquad
    e_n := \min\{d_1, \dots, d_n\} > 0 .
  \]
  Then $\{e_n\}$ is non-increasing, positive, and $e_n \le d_n \le 1/n$, so $e_n \to 0$.
  For $\delta > 0$ the set $\{n \ge 1 : \delta < e_n\}$ is therefore finite; if it is
  empty put $n(\delta) := 1$, and otherwise let $n(\delta)$ be its largest element.
  Define the a priori rule
  \[
    \alpha(\delta) := \frac{1}{n(\delta)} \in (0,\infty) .
  \]

  \emph{$\alpha(\delta) \to 0$ as $\delta \to 0$.}  Let $N \ge 1$.  If $\delta < e_N$ then
  $N$ belongs to $\{n : \delta < e_n\}$, so $n(\delta) \ge N$ and
  $\alpha(\delta) \le 1/N$.  Since $e_N > 0$, this proves the claim.

  \emph{$\delta\,\|R_{\alpha(\delta)}\| \to 0$ as $\delta \to 0$.}  Let $\delta < e_1$.
  Then $\{n : \delta < e_n\}$ is nonempty and $n := n(\delta)$ satisfies $\delta < e_n$,
  so, using $e_n \le d_n$ and $\alpha(\delta) = 1/n$,
  \[
    \delta\,\|R_{\alpha(\delta)}\|_{\Ll(Y,X)} \le d_n\,c_n
      = \frac{c_n}{n\,(1 + c_n)} \le \frac{1}{n(\delta)} = \alpha(\delta) ,
  \]
  which tends to $0$ by the previous paragraph.

  \emph{$(R_{\alpha},\alpha)$ is a convergent regularization method.}  Let
  $y \in \Dom(T\dg)$, let $\varepsilon > 0$, and let $\delta > 0$ with
  $y^{\delta} \in B_{\delta}(y)$.  By \cref{lem_error_decomposition},
  \[
    \|R_{\alpha(\delta)}y^{\delta} - T\dg y\|_X
      \le \delta\,\|R_{\alpha(\delta)}\|_{\Ll(Y,X)}
        + \|R_{\alpha(\delta)}y - T\dg y\|_X ,
  \]
  and the right-hand side does not depend on $y^{\delta}$, so it bounds the supremum in
  \cref{def_parameter_choice}.  By \cref{def_regularization} there is $\beta > 0$ with
  $\|R_{\alpha}y - T\dg y\|_X \le \varepsilon/2$ for all $\alpha \in (0,\beta]$; by the
  second paragraph there is $\delta_1 > 0$ with $\alpha(\delta) \le \beta$ for
  $0 < \delta \le \delta_1$, and by the third there is $\delta_2 > 0$ with
  $\delta\|R_{\alpha(\delta)}\| \le \varepsilon/2$ for $0 < \delta \le \delta_2$.  For
  $0 < \delta \le \min\{\delta_1,\delta_2\}$ both terms are at most $\varepsilon/2$, so
  \[
    \sup_{y^{\delta} \in B_{\delta}(y)}
      \|R_{\alpha(\delta)}y^{\delta} - T\dg y\|_X \le \varepsilon . \qedhere
  \]
\end{proof}


\begin{theorem}[Characterisation of convergent a priori rules]
  \label{thm_apriori_characterisation}
  Let $\{R_{\alpha}\}_{\alpha>0}$ be a regularization of $T\dg$ and let $\alpha$ be an a
  priori rule.
  \begin{enumerate}
    \item If, as $\delta \to 0$,
      \[
        \alpha(\delta) \to 0
        \qquad\text{and}\qquad
        \delta\,\|R_{\alpha(\delta)}\|_{\Ll(Y,X)} \to 0 ,
      \]
      then $(R_{\alpha}, \alpha)$ is a convergent regularization method.
    \item Conversely, if $(R_{\alpha}, \alpha)$ is a convergent regularization method,
      then, as $\delta \to 0$,
      \[
        \delta\,\|R_{\alpha(\delta)}\|_{\Ll(Y,X)} \to 0
        \qquad\text{and}\qquad
        \|R_{\alpha(\delta)} y - T\dg y\|_X \to 0
      \]
      for every $y \in \Dom(T\dg)$.
  \end{enumerate}
\end{theorem}

\begin{proof}
  \emph{(i).}  Let $y \in \Dom(T\dg)$ and $\varepsilon > 0$.  For every $\delta > 0$ and
  every $y^{\delta} \in B_{\delta}(y)$, \cref{lem_error_decomposition} gives
  \[
    \|R_{\alpha(\delta)}y^{\delta} - T\dg y\|_X
      \le \delta\,\|R_{\alpha(\delta)}\|_{\Ll(Y,X)}
        + \|R_{\alpha(\delta)}y - T\dg y\|_X ,
  \]
  whose right-hand side is independent of $y^{\delta}$ and hence bounds the supremum
  over $y^{\delta} \in B_{\delta}(y)$ occurring in \cref{def_parameter_choice}.  By
  \cref{def_regularization} there is $\beta > 0$ with
  $\|R_{\alpha}y - T\dg y\|_X \le \varepsilon/2$ for all $\alpha \in (0,\beta]$; since
  $\alpha(\delta) \to 0$ and $\alpha$ takes values in $(0,\infty)$, there is
  $\delta_1 > 0$ such that $\alpha(\delta) \le \beta$ whenever
  $0 < \delta \le \delta_1$; the second term is then at most $\varepsilon/2$.  By
  hypothesis there is $\delta_2 > 0$
  with $\delta\|R_{\alpha(\delta)}\|_{\Ll(Y,X)} \le \varepsilon/2$ for
  $0 < \delta \le \delta_2$.  For $0 < \delta \le \min\{\delta_1,\delta_2\}$ the supremum
  is therefore at most $\varepsilon$, which is the limit demanded by
  \cref{def_parameter_choice}.

  \emph{(ii), second assertion.}  Let $y \in \Dom(T\dg)$.  Since $y \in B_{\delta}(y)$ for
  every $\delta > 0$,
  \[
    \|R_{\alpha(\delta)}y - T\dg y\|_X
      \le \sup_{y^{\delta} \in B_{\delta}(y)}\|R_{\alpha(\delta)}y^{\delta} - T\dg y\|_X ,
  \]
  and the right-hand side tends to $0$ as $\delta \to 0$ by \cref{def_parameter_choice}.

  \emph{(ii), first assertion.}  We test the definition at the exact datum $y := 0$, which
  lies in $\Dom(T\dg)$ because $0 = T0 \in \Ran(T) \subseteq \Dom(T\dg)$, and for which
  $T\dg 0 = 0$ by linearity of $T\dg$.  Suppose
  $\delta\,\|R_{\alpha(\delta)}\|_{\Ll(Y,X)} \not\to 0$.  Then there are $\varepsilon > 0$
  and a null sequence $\{\delta_n\} \subset (0,\infty)$ with
  $\delta_n\|R_{\alpha(\delta_n)}\|_{\Ll(Y,X)} \ge \varepsilon$ for all $n$.  In
  particular $\|R_{\alpha(\delta_n)}\| > 0$, so by \cref{lem_opnorm_sup} there is
  $z_n \in Y$ with $\|z_n\|_Y \le 1$ and
  $\|R_{\alpha(\delta_n)}z_n\|_X \ge \tfrac12\|R_{\alpha(\delta_n)}\|_{\Ll(Y,X)}$.  Put
  $y_n := \delta_n z_n$; then $\|y_n - 0\|_Y \le \delta_n$, i.e.\
  $y_n \in B_{\delta_n}(0)$, and by linearity
  \[
    \|R_{\alpha(\delta_n)}y_n - T\dg 0\|_X = \delta_n\|R_{\alpha(\delta_n)}z_n\|_X
      \ge \tfrac12\,\delta_n\|R_{\alpha(\delta_n)}\|_{\Ll(Y,X)}
      \ge \tfrac{\varepsilon}{2} .
  \]
  Hence $\sup_{y^{\delta} \in B_{\delta_n}(0)}\|R_{\alpha(\delta_n)}y^{\delta}
  - T\dg 0\|_X \ge \varepsilon/2$ for every $n$, while $\delta_n \to 0$; this contradicts
  the limit of \cref{def_parameter_choice} at $y = 0$.  Therefore
  $\delta\|R_{\alpha(\delta)}\|_{\Ll(Y,X)} \to 0$.
\end{proof}


\begin{definition}[Morozov discrepancy principle]
  \label{def_discrepancy_principle}
  Fix $\tau > 1$.  An a posteriori rule $\alpha$ satisfies the \emph{discrepancy
  principle} if for all $\delta > 0$ and all $y^{\delta} \in Y$ with
  $\|y^{\delta}\|_Y > \tau\delta$,
  \[
    \bigl\| T R_{\alpha(\delta,y^{\delta})} y^{\delta} - y^{\delta} \bigr\|_Y \le \tau\delta
    \;<\; \bigl\| T R_{\alpha} y^{\delta} - y^{\delta} \bigr\|_Y
    \quad\text{for all } \alpha > \alpha(\delta,y^{\delta}) ,
  \]
  i.e.\ $\alpha$ is chosen as large as possible subject to the residual falling below
  $\tau\delta$.
\end{definition}


\begin{theorem}[The discrepancy principle is realisable]
  \label{thm_discrepancy_well_defined}
  Let $\{R_{\alpha}\}_{\alpha>0}$ be a regularization of $T\dg$ with $\Ran(T)$ dense in
  $Y$, let $\{\alpha_n\}$ be a strictly decreasing null sequence, let $\tau > 0$, and
  suppose $\sup_{\alpha>0}\|TR_{\alpha}\|_{\Ll(Y,Y)} < \infty$.  Then for all
  $\delta > 0$ and $y^{\delta} \in Y$ there is a smallest $n^{*}$ with
  \[
    \|TR_{\alpha_{n^{*}}} y^{\delta} - y^{\delta}\|_Y \le \tau\delta
    < \|TR_{\alpha_n} y^{\delta} - y^{\delta}\|_Y \quad\text{for all } n < n^{*} .
  \]
\end{theorem}

\begin{proof}
  Write $C := \sup_{\alpha > 0}\|TR_{\alpha}\|_{\Ll(Y,Y)} < \infty$.

  \emph{Step 1: $TR_{\alpha}z \to P_Rz$ as $\alpha \to 0$, for every $z \in Y$.}  For
  $\tilde{y} \in \Dom(T\dg)$ we have $R_{\alpha}\tilde{y} \to T\dg\tilde{y}$
  (\cref{def_regularization}), so, $T$ being continuous,
  $TR_{\alpha}\tilde{y} \to TT\dg\tilde{y} = P_R\tilde{y}$ by
  \cref{lem_mp_equations}(iv).  That is pointwise convergence of $\{TR_{\alpha}\}$ to
  $P_R$ on $M := \Dom(T\dg)$, which is a linear subspace of $Y$ (\cref{thm_mp_inverse})
  and dense (\cref{lem_mp_domain_dense}), hence total.  Here $Y$ is a Hilbert space and
  so Banach, $\{TR_{\alpha}\}_{\alpha>0} \subseteq \Ll(Y,Y)$ has
  $\sup_{\alpha>0}\|TR_{\alpha}\| = C < \infty$, and $P_R \in \Ll(Y,Y)$.  So the
  extension principle for a continuous parameter, \cref{prop_dense_ptwise_filter}, applies
  verbatim to the limit $\alpha \to 0^{+}$ taken here and gives
  $TR_{\alpha}z \to P_Rz$ for every $z \in Y$.  This is Step 1.

  \emph{Step 2: $P_R = \Id$.}  By hypothesis $\Ran(T)$ is dense, so
  $\overline{\Ran(T)} = Y$.  For $z \in Y$ the element $z$ itself lies in
  $\overline{\Ran(T)}$ and satisfies $z - z = 0 \perp \overline{\Ran(T)}$, so
  $P_Rz = z$ by \cref{lem_proj_characterisation}.

  \emph{Step 3: existence of the index.}  Fix $\delta > 0$ and $y^{\delta} \in Y$.  Since
  $\{\alpha_n\}$ is a null sequence, Steps 1 and 2 give
  \[
    \|TR_{\alpha_n}y^{\delta} - y^{\delta}\|_Y
      = \|TR_{\alpha_n}y^{\delta} - P_Ry^{\delta}\|_Y \longrightarrow 0
      \qquad (n \to \infty) .
  \]
  As $\tau\delta > 0$, the set
  $S := \{n : \|TR_{\alpha_n}y^{\delta} - y^{\delta}\|_Y \le \tau\delta\}$ is therefore
  nonempty, and being a nonempty set of natural numbers it has a least element $n^{*}$
  (well-ordering).  For $n < n^{*}$ we have $n \notin S$, i.e.\
  $\|TR_{\alpha_n}y^{\delta} - y^{\delta}\|_Y > \tau\delta$, while $n^{*} \in S$ gives
  $\|TR_{\alpha_{n^{*}}}y^{\delta} - y^{\delta}\|_Y \le \tau\delta$.  This is the asserted
  pair of inequalities, and $n^{*}$ is by construction the smallest index with that
  property.
\end{proof}


\begin{theorem}[Bakushinski\u{\i} veto]
  \label{thm_bakushinskii}
  Let $\{R_{\alpha}\}_{\alpha>0}$ be any family of maps $Y \to X$.  If there is a
  \emph{heuristic} rule $\alpha$ such that $(R_{\alpha}, \alpha)$ satisfies the limit of
  \cref{def_parameter_choice}, then $T\dg$ is bounded.
\end{theorem}

\begin{proof}
  Since the rule is heuristic it is a map $\alpha \colon Y \to (0,\infty)$ depending on
  $y^{\delta}$ alone, so
  \[
    R \colon Y \to X, \qquad Ry := R_{\alpha(y)}y
  \]
  is well defined (possibly nonlinear and unbounded).  The hypothesis is that for every
  $y \in \Dom(T\dg)$,
  \begin{equation}\label{eq_bakush_limit}
    \lim_{\delta \to 0}\ \sup_{y^{\delta} \in B_{\delta}(y)}
      \|R_{\alpha(y^{\delta})}y^{\delta} - T\dg y\|_X = 0 .
  \end{equation}

  \emph{Step 1: $R$ agrees with $T\dg$ on $\Dom(T\dg)$.}  Let
  $\tilde{y} \in \Dom(T\dg)$.  For every $\delta > 0$ we have
  $\tilde{y} \in B_{\delta}(\tilde{y})$, so the quantity
  $\|R\tilde{y} - T\dg\tilde{y}\|_X = \|R_{\alpha(\tilde{y})}\tilde{y}
  - T\dg\tilde{y}\|_X$ --- which does not depend on $\delta$ --- is bounded by the
  supremum in \eqref{eq_bakush_limit} written at $y := \tilde{y}$.  Letting
  $\delta \to 0$ gives $\|R\tilde{y} - T\dg\tilde{y}\|_X = 0$, i.e.\
  $R\tilde{y} = T\dg\tilde{y}$.

  \emph{Step 2: $T\dg$ is continuous at every point of its domain.}  Let
  $y \in \Dom(T\dg)$ and $\varepsilon > 0$.  By \eqref{eq_bakush_limit} there is
  $\delta_0 > 0$ with
  \[
    \sup_{y^{\delta} \in B_{\delta_0}(y)}
      \|R_{\alpha(y^{\delta})}y^{\delta} - T\dg y\|_X \le \varepsilon .
  \]
  Let $z \in \Dom(T\dg)$ with $\|z - y\|_Y \le \delta_0$.  Then $z \in B_{\delta_0}(y)$,
  so the displayed bound applies to $y^{\delta} := z$ and, by Step 1,
  \[
    \|T\dg z - T\dg y\|_X = \|Rz - T\dg y\|_X
      = \|R_{\alpha(z)}z - T\dg y\|_X \le \varepsilon .
  \]
  Note that no case distinction is needed: the parameter $\delta_0$ is used in place of
  $\|z - y\|_Y$, which may vanish while \cref{def_parameter_choice} excludes
  $\delta = 0$.

  \emph{Step 3: continuity implies boundedness.}  $\Dom(T\dg)$ is a subspace of $Y$ and
  $T\dg$ is linear on it, with $T\dg 0 = 0$ (\cref{thm_mp_inverse}).  Apply Step 2 at
  $y := 0$ with $\varepsilon := 1$: there is $\delta_0 > 0$ such that
  $\|T\dg z\|_X \le 1$ for all $z \in \Dom(T\dg)$ with $\|z\|_Y \le \delta_0$.  Given
  $0 \ne w \in \Dom(T\dg)$, the element $z := \delta_0 w/\|w\|_Y$ lies in $\Dom(T\dg)$ and
  has norm $\delta_0$, so $\|T\dg z\|_X \le 1$; by linearity
  \[
    \|T\dg w\|_X = \frac{\|w\|_Y}{\delta_0}\,\|T\dg z\|_X
      \le \frac{1}{\delta_0}\,\|w\|_Y ,
  \]
  and this holds trivially for $w = 0$ as well.  Hence $T\dg$ is bounded.
\end{proof}


\begin{definition}[Heuristic rules]
  \label{def_heuristic_rules}
  Three heuristic rules, each defined for $y^{\delta} \in Y$:
  \begin{enumerate}
    \item \emph{quasi-optimality}: fix a finite strictly decreasing
      $\{\alpha_n\}_{n \le N}$ and take $\alpha(y^{\delta}) := \alpha_{n^{*}}$ with
      $n^{*} \in \arg\min_{1 \le n < N} \|R_{\alpha_{n+1}}y^{\delta} -
      R_{\alpha_n}y^{\delta}\|_X$;
    \item \emph{Hanke--Raus}: $\alpha(y^{\delta}) \in \arg\min_{\alpha > 0}
      \alpha^{-1/2}\|TR_{\alpha}y^{\delta} - y^{\delta}\|_Y$;
    \item \emph{L-curve}: $\alpha(y^{\delta}) \in \arg\min_{\alpha > 0}
      \|R_{\alpha}y^{\delta}\|_X \cdot \|TR_{\alpha}y^{\delta} - y^{\delta}\|_Y$.
  \end{enumerate}
\end{definition}


\chapter{Convergence rates and order optimality}
\label{ch_14_order_optimality}

\noindent\emph{Background.} Convergence of a regularization method is, by itself, pointwise
and can be arbitrarily slow, so a genuine rate needs an extra smoothness assumption on
$x\dg$; this chapter makes that precise via the worst-case error over noise balls and the
power-type source sets $X_{\nu,\rho}$ of the previous chapter. The sharp two-sided estimate
$\varepsilon(X_{\nu,\rho},\delta) \asymp \delta^{\nu/(\nu+1)}\rho^{1/(\nu+1)}$, and the
resulting notions of optimality and order optimality built on it, are due to
\citet{Natterer1984}, who first established this convergence-rate exponent, together with
its minimax optimality, for linear ill-posed problems under Hölder-type source
conditions.

\begin{assumptions}
  \label{asm_convergence_rates_order}
  As in the previous chapter, with $(R_{\alpha},
\alpha)$ a fixed regularization method.  $B_Y$ denotes the closed unit ball of $Y$.  In
\cref{thm_modulus_source_upper} onwards, $T = K$ is compact and $X_{\nu,\rho}$ is as in
\cref{def_source_set}.
\end{assumptions}

\begin{definition}[Worst-case error and modulus of conditional continuity]
  \label{def_worst_case_error}
  For $y \in \Dom(T\dg)$ and $\delta > 0$,
  \[
    E(y,\delta) := \sup_{y^{\delta} \in B_{\delta}(y)}
      \bigl\| R_{\alpha(\delta,y^{\delta})} y^{\delta} - T\dg y \bigr\|_X .
  \]
  For $M \subseteq X$ and $\delta > 0$,
  \[
    \varepsilon(M,\delta) := \sup\{\, \|x\|_X : x \in M,\ \|Tx\|_Y \le \delta \,\} .
  \]
\end{definition}


\begin{theorem}[No uniform convergence rates]
  \label{thm_no_uniform_rates}
  Let $\{R_{\alpha}\}_{\alpha>0}$ be a family of linear maps $Y \to X$ and $\alpha$ a
  parameter choice rule.  If there is
  $\psi \colon (0,\infty) \to (0,\infty)$ with $\psi(t) \to 0$ as $t \to 0$ and
  \[
    \sup_{y \in \Dom(T\dg) \cap B_Y} E(y,\delta) \le \psi(\delta)
    \qquad \text{for all small } \delta > 0,
  \]
  then $T\dg$ is bounded.
\end{theorem}

\begin{proof}
  Let $\delta_1 > 0$ be such that the hypothesis
  $\sup_{y \in \Dom(T\dg) \cap B_Y}E(y,\delta) \le \psi(\delta)$ holds for all
  $\delta \in (0,\delta_1]$.  Recall that $\alpha$ maps $(0,\infty) \times Y$ into
  $(0,\infty)$ (\cref{def_parameter_choice}), so $R_{\alpha(\delta,z)}$ is defined for
  every $\delta > 0$ and $z \in Y$.

  \emph{A modulus of continuity on the unit ball.}  Let $y, z \in \Dom(T\dg) \cap B_Y$
  with $0 < \|y - z\|_Y \le \delta_1$, and put $\delta := \|y - z\|_Y$.  Then
  $z \in B_{\delta}(z)$ and $z \in B_{\delta}(y)$, so by \cref{def_worst_case_error},
  applied once at the exact datum $z$ and once at the exact datum $y$ --- in both cases
  with the noisy datum $z$ --- we get
  \[
    \|T\dg z - R_{\alpha(\delta,z)}z\|_X \le E(z,\delta), \qquad
    \|R_{\alpha(\delta,z)}z - T\dg y\|_X \le E(y,\delta) .
  \]
  Both $y$ and $z$ lie in $\Dom(T\dg) \cap B_Y$, so both right-hand sides are at most
  $\psi(\delta)$, and the triangle inequality gives
  \begin{equation}\label{eq_no_uniform_modulus}
    \|T\dg z - T\dg y\|_X \le 2\psi(\|y - z\|_Y) .
  \end{equation}

  \emph{Boundedness.}  Since $\psi(t) \to 0$ as $t \to 0$, there is
  $t_0 \in (0, \min\{1,\delta_1\}]$ with $\psi(t_0) \le 1$.  Note
  $0 \in \Dom(T\dg) \cap B_Y$ and $T\dg 0 = 0$, the domain being a subspace on which
  $T\dg$ is linear (\cref{thm_mp_inverse}).  Let $0 \ne w \in \Dom(T\dg)$ and put
  $z := t_0 w/\|w\|_Y$.  Then $z \in \Dom(T\dg)$, $\|z\|_Y = t_0 \le 1$, so
  $z \in \Dom(T\dg) \cap B_Y$, and $0 < \|z - 0\|_Y = t_0 \le \delta_1$.  Applying
  \eqref{eq_no_uniform_modulus} with $y := 0$,
  \[
    \|T\dg z\|_X \le 2\psi(t_0) \le 2 ,
  \]
  and by linearity $T\dg z = (t_0/\|w\|_Y)\,T\dg w$, whence
  \[
    \|T\dg w\|_X \le \frac{2}{t_0}\,\|w\|_Y .
  \]
  This inequality also holds for $w = 0$, so $T\dg$ is bounded.
\end{proof}


\begin{definition}[Order optimality]
  \label{def_order_optimal}
  Let $K \in \Kk(X,Y)$, let $\{R_{\alpha}\}_{\alpha>0}$ be a family of maps $Y \to X$, let
  $\beta$ be a parameter choice rule (\cref{def_parameter_choice}), and let
  $\nu, \rho > 0$ and $c \ge 1$.  Write
  \[
    E_{\beta}(y,\delta) := \sup_{y^{\delta} \in B_{\delta}(y)}
      \bigl\| R_{\beta(\delta,y^{\delta})} y^{\delta} - K\dg y \bigr\|_X
  \]
  for the worst-case error of \cref{def_worst_case_error} formed with $\{R_{\alpha}\}$ and
  $\beta$.  The pair $(\{R_{\alpha}\}, \beta)$ is \emph{order optimal with constant $c$
  for $\nu$ and $\rho$} if
  \[
    E_{\beta}(Kx\dg, \delta) \;\le\; c\,\delta^{\nu/(\nu+1)}\rho^{1/(\nu+1)}
    \qquad \text{for all } \delta > 0 \text{ and all } x\dg \in X_{\nu,\rho} .
  \]
  It is \emph{order optimal for $\nu$} if for every $x\dg \in X_{\nu}$ there is
  $c = c(x\dg) \ge 1$ with $E_{\beta}(Kx\dg,\delta) \le c\,\delta^{\nu/(\nu+1)}$ for all
  $\delta > 0$.  No convergence of $\{R_{\alpha}\}$ or of $(\{R_{\alpha}\},\beta)$ is
  presupposed.
\end{definition}


\begin{theorem}[The modulus is a lower bound for the worst-case error]
  \label{thm_worstcase_lower}
  Let $\{R_{\alpha}\}_{\alpha>0}$ be any family of \emph{linear} operators $Y \to X$ and
  $\alpha$ any parameter choice rule.  Then for all $\delta > 0$ and
  $M \subseteq \Ker(T)^{\perp}$,
  \[
    \sup\{\, E(y,\delta) : y \in \Dom(T\dg),\ T\dg y \in M \,\}
      \;\ge\; \varepsilon(M,\delta) .
  \]
\end{theorem}

\begin{proof}
  Write $A := \{x \in M : \|Tx\|_Y \le \delta\}$, so that
  $\varepsilon(M,\delta) = \sup\{\|x\|_X : x \in A\}$, and let
  $B := \{E(y,\delta) : y \in \Dom(T\dg),\ T\dg y \in M\}$, so that the assertion is
  $\sup B \ge \sup\{\|x\|_X : x \in A\}$.  If $A$ is empty then
  $\varepsilon(M,\delta) = 0$ with the convention $\sup\emptyset = 0$ and the inequality
  holds.  It therefore suffices to show that $\|x\|_X \le \sup B$ for every $x \in A$.

  So let $x \in A$ and put $y := Tx$.  Then $y \in \Ran(T) \subseteq \Dom(T\dg)$
  (\cref{thm_mp_inverse}).  Since $x \in M \subseteq \Ker(T)^{\perp}$, the projection
  $P_N$ onto $\Ker(T)$ annihilates $x$: indeed $0 \in \Ker(T)$ and $x - 0 = x \perp
  \Ker(T)$, so $P_Nx = 0$ by \cref{lem_proj_characterisation}.  Hence
  \cref{lem_mp_equations}(iii) gives
  \[
    T\dg y = T\dg Tx = (\Id - P_N)x = x \in M ,
  \]
  so $E(y,\delta)$ is one of the numbers in $B$.  Moreover $\|y - 0\|_Y = \|Tx\|_Y \le
  \delta$, i.e.\ the noisy datum $y^{\delta} := 0$ lies in $B_{\delta}(y)$, and
  $R_{\alpha(\delta,0)}0 = 0$ because $R_{\alpha(\delta,0)}$ is linear.  By
  \cref{def_worst_case_error},
  \[
    E(y,\delta) \ge \|R_{\alpha(\delta,0)}0 - T\dg y\|_X = \|x\|_X .
  \]
  Thus $\|x\|_X \le \sup B$ for every $x \in A$, and taking the supremum over $x \in A$
  yields $\varepsilon(M,\delta) \le \sup B$, which is the assertion.
\end{proof}


\begin{theorem}[Upper bound for the modulus on source sets]
  \label{thm_modulus_source_upper}
  Let $K \in \Kk(X,Y)$ and $\nu, \rho > 0$.  Then
  $\varepsilon(X_{\nu,\rho},\delta) \le \delta^{\nu/(\nu+1)}\rho^{1/(\nu+1)}$ for all
  $\delta > 0$.
\end{theorem}

\begin{proof}
  The modulus $\varepsilon(X_{\nu,\rho},\delta)$ is the supremum of $\|x\|_X$ over those
  $x \in X_{\nu,\rho}$ with $\|Kx\|_Y \le \delta$ (\cref{def_worst_case_error}).  If there
  are none, then $\varepsilon(X_{\nu,\rho},\delta) = 0$ with the convention
  $\sup\emptyset = 0$ and there is nothing to prove.  Otherwise, let
  $x \in X_{\nu,\rho}$ with $\|Kx\|_Y \le \delta$.  By
  \cref{def_source_set} there is $w \in X$ with $x = |K|^{\nu}w$ and $\|w\|_X \le \rho$.

  The interpolation inequality \cref{thm_interpolation}, applied with $s := \nu$ and
  $r := \nu + 1$, gives
  \[
    \|x\|_X = \bigl\||K|^{\nu}w\bigr\|_X
      \le \bigl\||K|^{\nu+1}w\bigr\|_X^{\frac{\nu}{\nu+1}}\,
          \|w\|_X^{\frac{1}{\nu+1}} .
  \]
  By the semigroup law for fractional powers (\cref{lem_fun_calc_hom}),
  $|K|^{\nu+1}w = |K|\,|K|^{\nu}w = |K|x$, and $\||K|x\|_X = \|Kx\|_Y$
  (\cref{lem_abs_norm}).  Since $t \mapsto t^{\theta}$ is non-decreasing on $[0,\infty)$
  for $\theta > 0$, inserting $\||K|^{\nu+1}w\|_X = \|Kx\|_Y \le \delta$ and
  $\|w\|_X \le \rho$ yields
  \[
    \|x\|_X \le \|Kx\|_Y^{\frac{\nu}{\nu+1}}\,\|w\|_X^{\frac{1}{\nu+1}}
      \le \delta^{\frac{\nu}{\nu+1}}\rho^{\frac{1}{\nu+1}} .
  \]
  Taking the supremum over all such $x$ gives
  $\varepsilon(X_{\nu,\rho},\delta) \le \delta^{\nu/(\nu+1)}\rho^{1/(\nu+1)}$.
\end{proof}


\begin{theorem}[The bound is attained along a sequence]
  \label{thm_modulus_attained}
  Let $K \in \Kk(X,Y)$ have infinite-dimensional range and $\nu, \rho > 0$.  Then there is
  a null sequence $\{\delta_n\}$ with
  $\varepsilon(X_{\nu,\rho},\delta_n) = \delta_n^{\nu/(\nu+1)}\rho^{1/(\nu+1)}$.
\end{theorem}
\begin{proof}
  \emph{The singular system is infinite.}  Let $\{(\sigma_n,u_n,v_n)\}_{n \in I}$ be a
  singular system of $K$ (\cref{thm_svd}).  By \cref{def_singular_system}, $\{u_n\}_{n \in I}$
  is an orthonormal basis of $\overline{\Ran(K)}$; if $I$ were finite then
  $\overline{\Ran(K)}$, and hence its subspace $\Ran(K)$, would be finite-dimensional,
  contrary to hypothesis.  So $I = \mathbb{N}$ and $\{\sigma_n\}$ is a null sequence of
  strictly positive numbers.

  \emph{The sequence.}  Put
  \[
    \delta_n := \rho\,\sigma_n^{\nu+1} > 0, \qquad
    w_n := \rho v_n, \qquad x_n := |K|^{\nu}w_n .
  \]
  Since $t \mapsto \rho t^{\nu+1}$ is continuous at $0$ and $\sigma_n \to 0$, the sequence
  $\{\delta_n\}$ is a null sequence.  From $Kv_n = \sigma_nu_n$ and $K^{*}u_n =
  \sigma_nv_n$ we get $K^{*}Kv_n = \sigma_n^{2}v_n$, so the functional calculus
  (\cref{def_fun_calc}) applied to $t \mapsto t^{\nu/2}$ gives
  $|K|^{\nu}v_n = (K^{*}K)^{\nu/2}v_n = \sigma_n^{\nu}v_n$ (\cref{def_frac_power}).  Hence
  \[
    x_n = \rho\,\sigma_n^{\nu}v_n, \qquad
    \|x_n\|_X = \rho\,\sigma_n^{\nu}, \qquad
    \|Kx_n\|_Y = \rho\,\sigma_n^{\nu}\|Kv_n\|_Y = \rho\,\sigma_n^{\nu+1} = \delta_n ,
  \]
  using $\|v_n\|_X = \|u_n\|_Y = 1$.  Also $\|w_n\|_X = \rho$, so $x_n \in X_{\nu,\rho}$
  by \cref{def_source_set}.

  \emph{Conclusion.}  Because $x_n \in X_{\nu,\rho}$ and $\|Kx_n\|_Y \le \delta_n$, the
  number $\|x_n\|_X$ is one of those over which the supremum defining
  $\varepsilon(X_{\nu,\rho},\delta_n)$ is taken, and
  \[
    \|x_n\|_X = \rho\,\sigma_n^{\nu}
      = \bigl(\rho\,\sigma_n^{\nu+1}\bigr)^{\frac{\nu}{\nu+1}}\rho^{\frac{1}{\nu+1}}
      = \delta_n^{\frac{\nu}{\nu+1}}\rho^{\frac{1}{\nu+1}} ,
  \]
  since the exponents of $\rho$ add up to $\nu/(\nu+1) + 1/(\nu+1) = 1$.  Therefore
  $\varepsilon(X_{\nu,\rho},\delta_n) \ge \delta_n^{\nu/(\nu+1)}\rho^{1/(\nu+1)}$, and the
  reverse inequality is \cref{thm_modulus_source_upper}.
\end{proof}


\begin{theorem}[Order optimality implies convergence]
  \label{thm_order_optimal_implies_convergence}
  Let $K \in \Kk(X,Y)$ have infinite-dimensional range and dense range in $Y$, let
  $\{R_{\alpha}\}_{\alpha>0}$ be a family in $\Ll(Y,X)$, and let $\alpha$ be a parameter
  choice rule.  Suppose there is $\tau_0 \ge 1$ such that for every $\tau > \tau_0$ the
  pair $(R_{\alpha}, \alpha_{\tau})$ with $\alpha_{\tau}(\delta,y^{\delta}) :=
  \alpha(\tau\delta, y^{\delta})$ is order optimal for some fixed $\nu > 0$ and
  \emph{all} $\rho > 0$, with one constant $c = c(\nu)$.  Then $(R_{\alpha},
  \alpha_{\tau})$ is a convergent regularization method for every $\tau > \tau_0$.
\end{theorem}

\begin{proof}
  Fix $\tau > \tau_0$ and write, for a rule $\beta$ and $y \in \Dom(K\dg)$,
  \[
    E_{\beta}(y,\delta) := \sup_{y^{\delta} \in B_{\delta}(y)}
      \|R_{\beta(\delta,y^{\delta})}y^{\delta} - K\dg y\|_X ,
  \]
  as in \cref{def_worst_case_error} and \cref{def_order_optimal}.  We must show
  $E_{\alpha_{\tau}}(y,\delta) \to 0$ as $\delta \to 0$ for every $y \in \Dom(K\dg)$
  (\cref{def_parameter_choice}).

  \emph{Notation and standing facts.}  Let $\{(\sigma_n,u_n,v_n)\}_{n \in I}$ be a
  singular system of $K$ (\cref{thm_svd}).  By \cref{def_singular_system}, $\{v_n\}$ is an
  orthonormal basis of $\Ker(K)^{\perp}$ and $\{u_n\}$ one of $\overline{\Ran(K)}$; were
  $I$ finite, $\overline{\Ran(K)}$ and hence its subspace $\Ran(K)$ would be
  finite-dimensional, so $I = \mathbb{N}$ and $\{\sigma_n\}$ is a non-increasing null
  sequence of strictly positive numbers.  As above, $K^{*}Kv_n = \sigma_n^{2}v_n$ and hence
  $|K|^{\nu}v_n = \sigma_n^{\nu}v_n$ (\cref{def_fun_calc}, \cref{def_frac_power}).

  Since $\Ran(K)$ is dense, $\Ran(K)^{\perp} = \overline{\Ran(K)}^{\perp} = \{0\}$ ---
  both equalities being \cref{lem_orthogonal_closure}, which supplies $M^{\perp} =
  (\overline{M})^{\perp}$ and $M$ dense $\iff M^{\perp} = \{0\}$ --- so
  $\Dom(K\dg) = \Ran(K) \oplus \Ran(K)^{\perp} = \Ran(K)$ (\cref{thm_mp_inverse}) and
  $P_R = \Id$.  Fix $y \in \Dom(K\dg)$ and put $x\dg := K\dg y$.  By
  \cref{lem_mp_equations} we have $x\dg \in \Ran(K\dg) = \Ker(K)^{\perp}$ and, by item
  (iv), $Kx\dg = KK\dg y = P_Ry = y$.  Expanding in the orthonormal basis $\{v_n\}$ of
  $\Ker(K)^{\perp}$ (\cref{thm_onb}) we write
  \[
    x\dg = \sum_{n \in \mathbb{N}} c_nv_n, \qquad c_n := \langle v_n, x\dg\rangle_X,
    \qquad \sum_{n} |c_n|^{2} = \|x\dg\|_X^{2} < \infty .
  \]

  \emph{Truncation.}  For $\varepsilon > 0$ let $N_{\varepsilon} := \#\{n : \sigma_n \ge
  \varepsilon\}$, a finite number because $\sigma_n \to 0$.  Since $\{\sigma_n\}$ is
  non-increasing \emph{along an index set that is an initial segment of $\mathbb{N}$}
  (\cref{def_singular_system}), $\sigma_n \ge \varepsilon$ holds exactly for
  $n \le N_{\varepsilon}$; over an arbitrary index set the two conditions would not
  coincide, which is why that clause is a field of the definition and not an informal
  remark.  Put
  \[
    x_{\varepsilon} := \sum_{n \le N_{\varepsilon}} c_nv_n, \qquad
    y_{\varepsilon} := Kx_{\varepsilon} = \sum_{n \le N_{\varepsilon}}
      \sigma_nc_nu_n .
  \]
  Both are finite sums.  Since $x_{\varepsilon} \in \Ker(K)^{\perp}$ we have
  $K\dg y_{\varepsilon} = K\dg Kx_{\varepsilon} = (\Id - P_N)x_{\varepsilon} =
  x_{\varepsilon}$ by \cref{lem_mp_equations}(iii), and $y_{\varepsilon} \in \Ran(K) =
  \Dom(K\dg)$.  Moreover $x\dg - x_{\varepsilon} = \sum_{n > N_{\varepsilon}}c_nv_n$ and
  $y - y_{\varepsilon} = K(x\dg - x_{\varepsilon}) = \sum_{n > N_{\varepsilon}}
  \sigma_nc_nu_n$ (\cref{lem_operator_through_series}), so by orthonormality
  (\cref{lem_riesz_fischer})
  \begin{equation}\label{eq_oo_tails}
    \|x\dg - x_{\varepsilon}\|_X^{2} = \sum_{n > N_{\varepsilon}}|c_n|^{2},
    \qquad
    \|y - y_{\varepsilon}\|_Y^{2} = \sum_{n > N_{\varepsilon}}\sigma_n^{2}|c_n|^{2}
      \le \varepsilon^{2}\,\|x\dg - x_{\varepsilon}\|_X^{2} ,
  \end{equation}
  the last step because $\sigma_n < \varepsilon$ for $n > N_{\varepsilon}$.  Both
  quantities are non-decreasing functions of $\varepsilon$, since $N_{\varepsilon}$ is
  non-increasing in $\varepsilon$.  Finally, $N_{\varepsilon} \to \infty$ as
  $\varepsilon \to 0$ --- given $N$, every $\varepsilon \le \sigma_N$ has
  $N_{\varepsilon} \ge N$ --- so $\|x\dg - x_{\varepsilon}\|_X \to 0$ and
  $\|y - y_{\varepsilon}\|_Y \to 0$ as $\varepsilon \to 0$, both being tails of
  convergent series.

  \emph{The truncations satisfy a source condition.}  Fix $\varepsilon > 0$ and set
  $w_{\varepsilon} := \sum_{n \le N_{\varepsilon}}\sigma_n^{-\nu}c_nv_n$, again a finite
  sum.  Then $|K|^{\nu}w_{\varepsilon} = \sum_{n \le N_{\varepsilon}}\sigma_n^{-\nu}c_n
  \sigma_n^{\nu}v_n = x_{\varepsilon}$, and since $\sigma_n \ge \varepsilon$ for
  $n \le N_{\varepsilon}$,
  \begin{equation}\label{eq_oo_representer}
    \|w_{\varepsilon}\|_X^{2} = \sum_{n \le N_{\varepsilon}}\sigma_n^{-2\nu}|c_n|^{2}
      \le \varepsilon^{-2\nu}\sum_{n}|c_n|^{2}
      = \bigl(\varepsilon^{-\nu}\|x\dg\|_X\bigr)^{2} .
  \end{equation}
  Thus $x_{\varepsilon} \in X_{\nu,\rho}$ (\cref{def_source_set}) for every
  $\rho \ge \varepsilon^{-\nu}\|x\dg\|_X$ with $\rho > 0$.

  \emph{Case A: only finitely many $c_n$ are nonzero.}  Let
  $F := \{n : c_n \ne 0\}$ and suppose $F$ is finite.  Put
  $w := \sum_{n \in F}\sigma_n^{-\nu}c_nv_n$, a finite sum; then $|K|^{\nu}w = x\dg$, so
  $x\dg \in X_{\nu,\rho}$ with $\rho := \|w\|_X + 1 > 0$.  Order optimality of
  $(R_{\alpha},\alpha_{\tau})$ for $\nu$ and this $\rho$ (\cref{def_order_optimal}),
  applied to the source element $x\dg$ with $Kx\dg = y$, gives
  \[
    E_{\alpha_{\tau}}(y,\delta) \le c\,\delta^{\frac{\nu}{\nu+1}}
      \rho^{\frac{1}{\nu+1}} \qquad \text{for all } \delta > 0,
  \]
  and the right-hand side tends to $0$ as $\delta \to 0$.  This settles Case A; note it
  contains the degenerate case $x\dg = 0$, i.e.\ $y = 0$, where the truncation
  construction below would run over an empty index set.

  \emph{Case B: infinitely many $c_n$ are nonzero.}  Then $x\dg \ne x_{\varepsilon}$ for
  every $\varepsilon > 0$, and since $K$ is injective on $\Ker(K)^{\perp}$
  (\cref{thm_mp_inverse}) and $x\dg - x_{\varepsilon} \in \Ker(K)^{\perp}$ is nonzero,
  \begin{equation}\label{eq_oo_proper}
    \|y - y_{\varepsilon}\|_Y = \|K(x\dg - x_{\varepsilon})\|_Y > 0
    \qquad \text{for every } \varepsilon > 0 .
  \end{equation}
  In particular $\|y\|_Y > 0$, because $y_{\varepsilon} = 0$ for
  $\varepsilon > \sigma_1$.  Put
  \[
    q := \frac{\tau - \tau_0}{\tau + \tau_0} \in (0,1),
    \qquad \tilde{\tau} := \frac{\tau + \tau_0}{2} > \tau_0,
    \qquad \delta_1 := \frac{\|y\|_Y}{q} > 0 ,
  \]
  and for $0 < \delta < \delta_1$ define
  \[
    \varepsilon(\delta) := \inf S_{\delta}, \qquad
    S_{\delta} := \{\varepsilon > 0 : \|y - y_{\varepsilon}\|_Y \ge q\delta\} .
  \]
  The set $S_{\delta}$ is nonempty: any $\varepsilon > \sigma_1$ has $y_{\varepsilon} = 0$
  and $\|y - y_{\varepsilon}\|_Y = \|y\|_Y > q\delta$.  Also $\varepsilon(\delta) > 0$:
  choose $\varepsilon_0 > 0$ with $\|y - y_{\varepsilon_0}\|_Y < q\delta$, possible
  because $\|y - y_{\varepsilon}\|_Y \to 0$ as $\varepsilon \to 0$; by monotonicity every
  $\varepsilon \le \varepsilon_0$ satisfies $\|y - y_{\varepsilon}\|_Y < q\delta$ and
  hence lies outside $S_{\delta}$, so $\varepsilon(\delta) \ge \varepsilon_0$.  Finally
  $\delta \mapsto \varepsilon(\delta)$ is non-decreasing, since $S_{\delta}$ shrinks as
  $\delta$ grows.

  \emph{The two-sided estimate.}  We claim that for $0 < \delta < \delta_1$
  \begin{equation}\label{eq_oo_sandwich}
    \|y - y_{\varepsilon(\delta)}\|_Y \le q\delta \le \|y - y_{2\varepsilon(\delta)}\|_Y .
  \end{equation}
  For the right inequality, $2\varepsilon(\delta) > \varepsilon(\delta) = \inf S_{\delta}$
  gives some $\varepsilon \in S_{\delta}$ with $\varepsilon < 2\varepsilon(\delta)$, and
  monotonicity yields $\|y - y_{2\varepsilon(\delta)}\|_Y \ge \|y - y_{\varepsilon}\|_Y
  \ge q\delta$.  For the left inequality we use that $\varepsilon \mapsto N_{\varepsilon}$
  is constant just to the left of any point: writing $N := N_{\varepsilon(\delta)}$, the
  value $\sigma_{N+1}$ is the largest $\sigma_n$ that is $< \varepsilon(\delta)$, so for
  any $\varepsilon' \in (\sigma_{N+1},\varepsilon(\delta))$ no $\sigma_n$ lies
  in $[\varepsilon',\varepsilon(\delta))$ and therefore $N_{\varepsilon'} = N$, i.e.\
  $y_{\varepsilon'} = y_{\varepsilon(\delta)}$.  Such an $\varepsilon'$ is smaller than
  $\inf S_{\delta}$, hence not in $S_{\delta}$, so
  $\|y - y_{\varepsilon(\delta)}\|_Y = \|y - y_{\varepsilon'}\|_Y < q\delta$.

  \emph{Transfer of the noise level.}  Let $0 < \delta < \delta_1$ and
  $y^{\delta} \in B_{\delta}(y)$.  By \eqref{eq_oo_sandwich},
  \[
    \|y^{\delta} - y_{\varepsilon(\delta)}\|_Y \le \|y^{\delta} - y\|_Y
      + \|y - y_{\varepsilon(\delta)}\|_Y \le (1 + q)\delta =: \tilde{\delta} ,
  \]
  so $y^{\delta} \in B_{\tilde{\delta}}(y_{\varepsilon(\delta)})$: the same noisy datum is
  admissible for the exact datum $y_{\varepsilon(\delta)}$ at the level $\tilde{\delta}$.
  Moreover
  \[
    \tilde{\tau}\tilde{\delta} = \frac{\tau + \tau_0}{2}\Bigl(1
      + \frac{\tau - \tau_0}{\tau + \tau_0}\Bigr)\delta
      = \frac{\tau + \tau_0}{2}\cdot\frac{2\tau}{\tau + \tau_0}\,\delta = \tau\delta ,
  \]
  whence $\alpha_{\tilde{\tau}}(\tilde{\delta},y^{\delta}) =
  \alpha(\tilde{\tau}\tilde{\delta},y^{\delta}) = \alpha(\tau\delta,y^{\delta}) =
  \alpha_{\tau}(\delta,y^{\delta})$: the two rules select the same parameter.

  \emph{Applying order optimality.}  By \eqref{eq_oo_representer} we have
  $x_{\varepsilon(\delta)} \in X_{\nu,\rho(\delta)}$, where
  $\rho(\delta) := \varepsilon(\delta)^{-\nu}\|x\dg\|_X > 0$ (positive because
  $x\dg \ne 0$ in Case B), and $Kx_{\varepsilon(\delta)} =
  y_{\varepsilon(\delta)}$ with $K\dg y_{\varepsilon(\delta)} =
  x_{\varepsilon(\delta)}$.  Since
  $\tilde{\tau} > \tau_0$, the pair $(R_{\alpha},\alpha_{\tilde{\tau}})$ is order
  optimal for $\nu$ and $\rho(\delta)$ with the constant $c$, so by
  \cref{def_order_optimal} and \cref{def_worst_case_error},
  \[
    \|R_{\alpha_{\tau}(\delta,y^{\delta})}y^{\delta} - x_{\varepsilon(\delta)}\|_X
      = \|R_{\alpha_{\tilde{\tau}}(\tilde{\delta},y^{\delta})}y^{\delta}
        - K\dg y_{\varepsilon(\delta)}\|_X
      \le E_{\alpha_{\tilde{\tau}}}(y_{\varepsilon(\delta)},\tilde{\delta})
      \le c\,\tilde{\delta}^{\frac{\nu}{\nu+1}}\rho(\delta)^{\frac{1}{\nu+1}} .
  \]
  With $\tilde{\delta} = (1+q)\delta \le 2\delta$ and
  $\rho(\delta)^{1/(\nu+1)} = \varepsilon(\delta)^{-\nu/(\nu+1)}
  \|x\dg\|_X^{1/(\nu+1)}$ this reads
  \[
    \|R_{\alpha_{\tau}(\delta,y^{\delta})}y^{\delta} - x_{\varepsilon(\delta)}\|_X
      \le 2c\,\Bigl(\frac{\delta}{\varepsilon(\delta)}\Bigr)^{\frac{\nu}{\nu+1}}
        \|x\dg\|_X^{\frac{1}{\nu+1}} .
  \]
  Adding $\|x_{\varepsilon(\delta)} - x\dg\|_X$ and taking the supremum over
  $y^{\delta} \in B_{\delta}(y)$ --- the right-hand side no longer depends on
  $y^{\delta}$ --- gives
  \begin{equation}\label{eq_oo_final}
    E_{\alpha_{\tau}}(y,\delta) \le 2c\,
      \Bigl(\frac{\delta}{\varepsilon(\delta)}\Bigr)^{\frac{\nu}{\nu+1}}
      \|x\dg\|_X^{\frac{1}{\nu+1}} + \|x_{\varepsilon(\delta)} - x\dg\|_X .
  \end{equation}

  \emph{$\varepsilon(\delta) \to 0$ as $\delta \to 0$.}  The function
  $\varepsilon(\cdot)$ is non-decreasing and positive, so
  $\varepsilon_{*} := \inf_{0 < \delta < \delta_1}\varepsilon(\delta) \ge 0$ exists and
  $\varepsilon(\delta) \to \varepsilon_{*}$ as $\delta \to 0$.  Suppose
  $\varepsilon_{*} > 0$.  The set $J := \{n : \sigma_n > \varepsilon_{*}\}$ is finite, so
  $\eta := \min\bigl(\{\sigma_n - \varepsilon_{*} : n \in J\} \cup \{1\}\bigr) > 0$, and
  for every $\varepsilon \in (\varepsilon_{*},\varepsilon_{*} + \eta]$ one has
  $N_{\varepsilon} = \#J$: indeed $\sigma_n \ge \varepsilon_{*} + \eta \ge \varepsilon$
  for $n \in J$, while $\sigma_n \le \varepsilon_{*} < \varepsilon$ for $n \notin J$.  Now
  either $\varepsilon(\delta_2) = \varepsilon_{*}$ for some $\delta_2$, in which case
  $\varepsilon(\delta) = \varepsilon_{*}$ for all $\delta \le \delta_2$ by monotonicity;
  or $\varepsilon(\delta) > \varepsilon_{*}$ for all $\delta$, in which case
  $\varepsilon(\delta) \in (\varepsilon_{*},\varepsilon_{*} + \eta]$ for all small
  $\delta$ and $N_{\varepsilon(\delta)} = \#J$ there.  In both cases
  $N_{\varepsilon(\delta)}$, and hence $y_{\varepsilon(\delta)}$, is one and the same for
  all sufficiently small $\delta$, while $\|y - y_{\varepsilon(\delta)}\|_Y \le q\delta
  \to 0$ by \eqref{eq_oo_sandwich}.  So that fixed vector equals $y$, contradicting
  \eqref{eq_oo_proper}.  Hence $\varepsilon_{*} = 0$.

  \emph{Conclusion.}  Since $\varepsilon(\delta) \to 0$, the second term of
  \eqref{eq_oo_final} tends to $0$, because $\|x\dg - x_{\varepsilon}\|_X \to 0$ as
  $\varepsilon \to 0$.  For the first term, combine the right half of
  \eqref{eq_oo_sandwich} with \eqref{eq_oo_tails} at $2\varepsilon(\delta)$:
  \[
    q\delta \le \|y - y_{2\varepsilon(\delta)}\|_Y
      \le 2\varepsilon(\delta)\,\|x\dg - x_{2\varepsilon(\delta)}\|_X ,
    \qquad\text{i.e.}\qquad
    \frac{\delta}{\varepsilon(\delta)}
      \le \frac{2}{q}\,\|x\dg - x_{2\varepsilon(\delta)}\|_X ,
  \]
  and the right-hand side tends to $0$ as $\delta \to 0$ because
  $2\varepsilon(\delta) \to 0$.  Both terms of \eqref{eq_oo_final} therefore vanish in the
  limit, so $E_{\alpha_{\tau}}(y,\delta) \to 0$ as $\delta \to 0$.  As
  $y \in \Dom(K\dg)$ was arbitrary, $(R_{\alpha},\alpha_{\tau})$ is a convergent
  regularization method for every $\tau > \tau_0$.
\end{proof}


\chapter{Spectral regularization}
\label{ch_15_spectral_regularization}

\noindent\emph{Background.} Truncated singular value decomposition, Tikhonov
regularization, and Landweber iteration are all instances of a single scheme,
$R_{\alpha} = \varphi_{\alpha}(K^{*}K)K^{*}$ for a filter function $\varphi_{\alpha}$
approximating $\lambda^{-1}$, and the qualification of a filter --- how much of a
Hölder-type source condition it can turn into a rate --- bounds the best order any such
method can achieve. This unifying framework is systematized in
\citet{EnglHankeNeubauer1996}. The Hanke--Raus rule, choosing $\alpha$ heuristically from
the residual alone rather than from a known noise level, is due to \citet{HankeRaus1996}.

\begin{assumptions}
  \label{asm_spectral_regularization}
  $X$, $Y$ are Hilbert spaces, $K \in \Kk(X,Y)$ is
compact and nonzero, with singular system $\{(\sigma_n, u_n, v_n)\}_{n \in I}$ and
$\kappa := \|K\|^{2}_{\Ll(X,Y)} = \|K^{*}K\|_{\Ll(X,X)}$, the equality being
\cref{prop_adjoint_props}.  The bound $\sigma_n^{2} \le \kappa$ is \cref{lem_sigma_bound}
and the strict positivity is carried by \cref{def_singular_system}, so
$\sigma_n^{2} \in (0,\kappa]$ for every $n \in I$, which turns every supremum over $n$
below into a supremum over $\lambda \in (0,\kappa]$.  That $\kappa$ is also $\sigma_1^{2}$
--- the norm is attained at the largest singular value --- is true but nowhere needed.  We
write
$x_{\alpha} := R_{\alpha} y$ and $x^{\delta}_{\alpha} := R_{\alpha} y^{\delta}$.  Every
result about the discrepancy principle in this chapter carries the standing hypothesis
$\overline{\Ran(K)} = Y$, equivalently $\Ker(K^{*}) = 0$.
\end{assumptions}

\begin{definition}[Regularizing filter]
  \label{def_filter}
  A family $\{\varphi_{\alpha}\}_{\alpha>0}$ of piecewise continuous bounded functions
  $\varphi_{\alpha} \colon [0,\kappa] \to \mathbb{R}$ is a \emph{regularizing filter} if
  \begin{enumerate}
    \item $\varphi_{\alpha}(\lambda) \to \lambda^{-1}$ as $\alpha \to 0$, for every
      $\lambda \in (0,\kappa]$;
    \item there is $C_{\varphi} > 0$ with $\lambda\,|\varphi_{\alpha}(\lambda)| \le
      C_{\varphi}$ for all $\lambda \in (0,\kappa]$ and all $\alpha > 0$.
  \end{enumerate}
\end{definition}


\begin{definition}[Spectral regularization operator]
  \label{def_spectral_reg}
  For a regularizing filter $\{\varphi_{\alpha}\}$ put $R_{\alpha} :=
  \varphi_{\alpha}(K^{*}K)K^{*}$, so that for $y \in Y$
  \[
    R_{\alpha} y \;=\; \sum_{n \in I} \varphi_{\alpha}(\sigma_n^{2})\,\sigma_n\,
      \langle u_n, y\rangle_Y\, v_n .
  \]
\end{definition}


\begin{definition}[Three standard filters]
  \label{def_standard_filters}
  \begin{enumerate}
    \item \emph{Truncated singular value decomposition}:
      $\varphi_{\alpha}(\lambda) := \lambda^{-1}$ if $\lambda \ge \alpha$ and $0$
      otherwise, so that $R_{\alpha}y$ is the truncated sum
      $\sum_{\sigma_n \ge \sqrt{\alpha}} \sigma_n^{-1}\langle u_n, y\rangle_Y v_n$.
    \item \emph{Tikhonov}: $\varphi_{\alpha}(\lambda) := (\lambda + \alpha)^{-1}$, so
      $R_{\alpha}y = \sum_{n} \sigma_n(\sigma_n^{2}+\alpha)^{-1}\langle u_n,y\rangle_Y
      v_n$.
    \item \emph{Landweber}, for $\omega \in (0,\kappa^{-1}]$:
      $\varphi_{\alpha}(\lambda) := \lambda^{-1}\bigl(1 - (1-\omega\lambda)^{1/\alpha}
      \bigr)$ for $\lambda > 0$, and $\varphi_{\alpha}(0) := \omega/\alpha$.
  \end{enumerate}
  Each is a regularizing filter, with $C_{\varphi} = 1$.
\end{definition}


\begin{lemma}[Parseval identity for orthonormal series]
  \label{lem_parseval_series}
  Let $H$ be a Hilbert space, $\{e_n\}_{n \in I}$ an orthonormal system in $H$ and
  $\{c_n\}_{n \in I}$ scalars with $\sum_{n \in I}|c_n|^{2} < \infty$.  Then
  $z := \sum_{n \in I} c_ne_n$ converges in $H$ and
  \[
    \|z\|_H^{2} \;=\; \sum_{n \in I}|c_n|^{2},
    \qquad
    \langle e_m, z\rangle_H \;=\; c_m \quad (m \in I) .
  \]
\end{lemma}

\begin{proof}
  Convergence of the series is \cref{lem_riesz_fischer}.  For the norm identity, let
  $z_N := \sum_{n \le N}c_ne_n$ be the $N$th partial sum; orthonormality and Pythagoras
  (\cref{lem_pythagoras}) give $\|z_N\|_H^{2} = \sum_{n \le N}|c_n|^{2}$ by induction on
  $N$, and $\|z_N\|_H \to \|z\|_H$ because the norm is continuous, while
  $\sum_{n \le N}|c_n|^{2} \to \sum_{n \in I}|c_n|^{2}$ by definition of the sum.  The
  coefficient identity is \cref{lem_inner_through_series} applied to the continuous
  linear functional $\langle e_m, \cdot\rangle_H$, together with
  $\langle e_m, e_n\rangle_H = \delta_{mn}$.
\end{proof}


\begin{lemma}[Spectral form of $R_{\alpha}$ and of the residual]
  \label{lem_spectral_residual_form}
  Let $\{\varphi_{\alpha}\}$ be a regularizing filter, let $\alpha > 0$ and $y \in Y$,
  and put $c_n := \langle u_n, y\rangle_Y$ and $r_{\alpha}(\lambda) := 1 -
  \lambda\varphi_{\alpha}(\lambda)$.  Let $P_R$ denote the orthogonal projection onto
  $\overline{\Ran(K)}$.  Then all series below converge and
  \begin{enumerate}
    \item $R_{\alpha}y = \sum_{n \in I}\varphi_{\alpha}(\sigma_n^{2})\sigma_nc_nv_n$ and
      $\|R_{\alpha}y\|_X^{2} = \sum_{n \in I}|\varphi_{\alpha}(\sigma_n^{2})|^{2}
      \sigma_n^{2}|c_n|^{2}$;
    \item $KR_{\alpha}y = \sum_{n \in I}\varphi_{\alpha}(\sigma_n^{2})\sigma_n^{2}
      c_nu_n$ and $\|KR_{\alpha}y\|_Y^{2} = \sum_{n \in I}
      |\varphi_{\alpha}(\sigma_n^{2})|^{2}\sigma_n^{4}|c_n|^{2}$;
    \item $y - KR_{\alpha}y = (\Id - P_R)y + \sum_{n \in I}r_{\alpha}(\sigma_n^{2})
      c_nu_n$, and
      \[
        \|y - KR_{\alpha}y\|_Y^{2} \;=\; \|(\Id - P_R)y\|_Y^{2}
          \;+\; \sum_{n \in I}|r_{\alpha}(\sigma_n^{2})|^{2}|c_n|^{2} ;
      \]
    \item $\|y\|_Y^{2} = \|(\Id - P_R)y\|_Y^{2} + \sum_{n \in I}|c_n|^{2}$.
  \end{enumerate}
\end{lemma}

\begin{proof}
  Write $\Phi_{\alpha}$ for the supremum of $|\varphi_{\alpha}|$ over $(0,\kappa]$,
  which is finite because $\varphi_{\alpha}$ is bounded (\cref{def_filter}).  Bessel's
  inequality
  (\cref{thm_bessel}) gives $\sum_n |c_n|^{2} \le \|y\|_Y^{2} < \infty$, and
  $0 < \sigma_n^{2} \le \kappa$ for every $n \in I$ (\cref{def_singular_system} and
  \cref{lem_sigma_bound}), so the three
  coefficient families $\varphi_{\alpha}(\sigma_n^{2})\sigma_nc_n$,
  $\varphi_{\alpha}(\sigma_n^{2})\sigma_n^{2}c_n$ and $r_{\alpha}(\sigma_n^{2})c_n$ are
  square summable: they are bounded multiples of $c_n$, by $\Phi_{\alpha}\sqrt{\kappa}$,
  by $\Phi_{\alpha}\kappa$ and by $1 + C_{\varphi}$ respectively, the last because
  $|r_{\alpha}(\lambda)| \le 1 + \lambda|\varphi_{\alpha}(\lambda)| \le 1 + C_{\varphi}$
  on $(0,\kappa]$ by \cref{def_filter}(ii).  Every series below therefore converges by
  \cref{lem_parseval_series}, which also supplies each norm identity, the systems
  $\{v_n\}$ and $\{u_n\}$ being orthonormal (\cref{def_singular_system}).

  \emph{(i).}  The series is \cref{def_spectral_reg}, and the norm identity is
  \cref{lem_parseval_series} for $\{v_n\}$.

  \emph{(ii).}  Applying the bounded operator $K$ term by term
  (\cref{lem_operator_through_series}) to the series of (i) and using $Kv_n = \sigma_nu_n$
  (\cref{def_singular_system}) gives the series; the norm identity is
  \cref{lem_parseval_series} for $\{u_n\}$.

  \emph{(iii) and (iv).}  Since $\{u_n\}_{n \in I}$ is an orthonormal basis of
  $\overline{\Ran(K)}$ (\cref{def_singular_system}), \cref{thm_proj_expansion} gives
  $P_Ry = \sum_{n \in I}c_nu_n$, hence $y = (\Id - P_R)y + \sum_n c_nu_n$.  The first
  summand lies in $\overline{\Ran(K)}^{\perp}$ and the second in $\overline{\Ran(K)}$
  (\cref{thm_orth_decomp}), so Pythagoras (\cref{lem_pythagoras}) together with
  \cref{lem_parseval_series} gives (iv).  Subtracting the series of (ii) term by term
  from the displayed decomposition of $y$ yields
  \[
    \begin{aligned}
      y - KR_{\alpha}y &\;=\; (\Id - P_R)y \;+\; \sum_{n \in I}
        \bigl(1 - \varphi_{\alpha}(\sigma_n^{2})\sigma_n^{2}\bigr)c_nu_n \\
        &\;=\; (\Id - P_R)y \;+\; \sum_{n \in I}r_{\alpha}(\sigma_n^{2})c_nu_n ,
    \end{aligned}
  \]
  again a sum of two orthogonal vectors, so Pythagoras and \cref{lem_parseval_series}
  give the norm identity of (iii).
\end{proof}


\begin{lemma}[Residual bound]
  \label{lem_filter_residual_bound}
  For a regularizing filter and all $\alpha > 0$,
  \[
    \|K R_{\alpha}\|_{\Ll(Y,Y)} \;\le\; \sup_{n \in I}
      |\varphi_{\alpha}(\sigma_n^{2})|\,\sigma_n^{2} \;\le\; C_{\varphi} .
  \]
\end{lemma}

\begin{proof}
  Fix $\alpha > 0$ and put
  \[
    M_{\alpha} := \sup_{n \in I}|\varphi_{\alpha}(\sigma_n^{2})|\,\sigma_n^{2} .
  \]

  \emph{The first inequality.}  Let $y \in Y$.  By \cref{lem_spectral_residual_form}(ii),
  \[
    \|KR_{\alpha}y\|_Y^{2} \;=\; \sum_{n \in I} |\varphi_{\alpha}(\sigma_n^{2})
      \sigma_n^{2}|^{2}\,|\langle u_n,y\rangle_Y|^{2}
      \;\le\; M_{\alpha}^{2} \sum_{n \in I} |\langle u_n,y\rangle_Y|^{2}
      \;\le\; M_{\alpha}^{2}\,\|y\|_Y^{2},
  \]
  the last step again by \cref{thm_bessel}.  As $y$ was arbitrary,
  $\|KR_{\alpha}\|_{\Ll(Y,Y)} \le M_{\alpha}$.

  \emph{The second inequality.}  For every $n \in I$ one has $0 < \sigma_n^{2} \le
  \kappa$ by \cref{def_singular_system} and \cref{lem_sigma_bound}, so
  $\sigma_n^{2} \in (0,\kappa]$ and \cref{def_filter}(ii)
  gives $\sigma_n^{2}|\varphi_{\alpha}(\sigma_n^{2})| \le C_{\varphi}$.  Taking the
  supremum over $n$ yields $M_{\alpha} \le C_{\varphi}$.
\end{proof}


\begin{lemma}[Operator norm bound]
  \label{lem_filter_norm_bound}
  For a regularizing filter and all $\alpha > 0$,
  \[
    \|R_{\alpha}\|_{\Ll(Y,X)} \;\le\; \sqrt{C_{\varphi}}\,
      \Bigl(\sup_{\lambda \in (0,\kappa]} |\varphi_{\alpha}(\lambda)|\Bigr)^{1/2} .
  \]
  In particular each $R_{\alpha}$ is bounded.
\end{lemma}

\begin{proof}
  Fix $\alpha > 0$, put $\Phi_{\alpha} := \sup_{\lambda \in (0,\kappa]}|\varphi_{\alpha}
  (\lambda)|$ --- finite because $\varphi_{\alpha}$ is bounded (\cref{def_filter}) ---
  and let $y \in Y$.  Write $c_n := \langle u_n, y\rangle_Y$.  By
  \cref{lem_spectral_residual_form}(i),
  \[
    \|R_{\alpha}y\|_X^{2} \;=\; \sum_{n \in I} |\varphi_{\alpha}(\sigma_n^{2})|^{2}
      \sigma_n^{2}\,|c_n|^{2} .
  \]
  By the series for $KR_{\alpha}y$ of \cref{lem_spectral_residual_form}(ii) and the
  coefficient identity of \cref{lem_parseval_series} for $\{u_n\}$,
  \[
    d_n \;:=\; \langle u_n, KR_{\alpha}y\rangle_Y \;=\;
      \varphi_{\alpha}(\sigma_n^{2})\,\sigma_n^{2}\,c_n .
  \]
  As $\varphi_{\alpha}$ is real valued, the $n$th summand above equals
  $\varphi_{\alpha}(\sigma_n^{2})\,\overline{c_n}\,d_n$, whence
  \[
    \|R_{\alpha}y\|_X^{2} \;\le\; \Phi_{\alpha} \sum_{n \in I} |c_n|\,|d_n|
      \;\le\; \Phi_{\alpha}\Bigl(\sum_{n \in I}|c_n|^{2}\Bigr)^{1/2}
        \Bigl(\sum_{n \in I}|d_n|^{2}\Bigr)^{1/2},
  \]
  where the first step needs $|\varphi_{\alpha}(\sigma_n^{2})| \le \Phi_{\alpha}$, i.e.\
  $\sigma_n^{2} \in (0,\kappa]$, which is \cref{def_singular_system} with
  \cref{lem_sigma_bound}, and the second step is the Cauchy--Schwarz inequality
  (\cref{thm_cauchy_schwarz})
  applied, for each $N$, to the vectors $(|c_n|)_{n \le N}$ and $(|d_n|)_{n \le N}$ in
  $\mathbb{K}^{N}$ with its standard inner product, followed by $N \to \infty$.  Bessel's
  inequality (\cref{thm_bessel}) bounds the two factors by $\|y\|_Y$ and
  $\|KR_{\alpha}y\|_Y$ respectively, and $\|KR_{\alpha}y\|_Y \le C_{\varphi}\|y\|_Y$ by
  \cref{lem_filter_residual_bound}.  Therefore
  \[
    \|R_{\alpha}y\|_X^{2} \;\le\; \Phi_{\alpha}\,C_{\varphi}\,\|y\|_Y^{2} .
  \]
  Taking square roots and the supremum over $\|y\|_Y \le 1$ gives
  $\|R_{\alpha}\|_{\Ll(Y,X)} \le \sqrt{C_{\varphi}}\,\Phi_{\alpha}^{1/2} < \infty$;
  in particular $R_{\alpha}$ is bounded.
\end{proof}




\begin{lemma}[Spectral form of the approximation error]
  \label{lem_approx_error_spectral}
  Let $\{\varphi_{\alpha}\}$ be a regularizing filter, $y \in \Dom(K\dg)$ and
  $x\dg := K\dg y$.  With $r_{\alpha}(\lambda) := 1 - \lambda\varphi_{\alpha}(\lambda)$,
  \[
    K\dg y - R_{\alpha}y \;=\; r_{\alpha}(K^{*}K)\,x\dg
      \;=\; \sum_{n \in I} r_{\alpha}(\sigma_n^{2})\,\langle v_n, x\dg\rangle_X\, v_n ,
  \]
  and $r_{\alpha}(\lambda) \to 0$ as $\alpha \to 0$ for each $\lambda \in (0,\kappa]$,
  with $|r_{\alpha}(\lambda)| \le 1 + C_{\varphi}$ throughout.
\end{lemma}

\begin{proof}
  \emph{The normal equation.}  By \cref{thm_mp_minnorm}(ii), $x\dg = K\dg y$ is the
  minimum norm solution of $Kx = y$, hence in particular a least squares solution
  (\cref{def_lsq_minnorm}), so $K^{*}Kx\dg = K^{*}y$ by \cref{cor_normal_equation}.
  Consequently, by \cref{def_spectral_reg},
  \[
    R_{\alpha}y \;=\; \varphi_{\alpha}(K^{*}K)K^{*}y
      \;=\; \varphi_{\alpha}(K^{*}K)\,K^{*}K\,x\dg .
  \]

  \emph{Rewriting through the functional calculus.}  Let $\iota(\lambda) := \lambda$ and
  $\mathbf{1}(\lambda) := 1$ on $[0,\kappa]$.  By \cref{lem_fun_calc_basic},
  $\iota(K^{*}K) = K^{*}K$ and $\mathbf{1}(K^{*}K) = \Id$, and by multiplicativity of the
  calculus (\cref{lem_fun_calc_hom}) $\varphi_{\alpha}(K^{*}K)\,K^{*}K =
  (\varphi_{\alpha}\iota)(K^{*}K)$.  The function $r_{\alpha} = \mathbf{1} -
  \varphi_{\alpha}\iota$ is again piecewise continuous and bounded on $[0,\kappa]$, so
  linearity of the calculus (\cref{lem_fun_calc_hom}) gives, as operators on $X$,
  \[
    \Id - \varphi_{\alpha}(K^{*}K)K^{*}K \;=\; r_{\alpha}(K^{*}K) .
  \]
  Therefore $K\dg y - R_{\alpha}y = x\dg - \varphi_{\alpha}(K^{*}K)K^{*}Kx\dg =
  r_{\alpha}(K^{*}K)x\dg$.

  \emph{The series.}  By \cref{def_fun_calc},
  \[
    r_{\alpha}(K^{*}K)x\dg \;=\; \sum_{n \in I} r_{\alpha}(\sigma_n^{2})
      \langle v_n, x\dg\rangle_X\, v_n \;+\; r_{\alpha}(0)\,P_{\Ker(K)}x\dg ,
  \]
  and $P_{\Ker(K)}x\dg = 0$ because $x\dg \in \Ran(K\dg) = \Ker(K)^{\perp}$
  (\cref{lem_mp_equations}).  This is the asserted identity.

  \emph{Properties of $r_{\alpha}$.}  For $\lambda \in (0,\kappa]$, \cref{def_filter}(i)
  gives $\lambda\varphi_{\alpha}(\lambda) \to \lambda\cdot\lambda^{-1} = 1$ as
  $\alpha \to 0$, whence $r_{\alpha}(\lambda) = 1 - \lambda\varphi_{\alpha}(\lambda) \to
  0$; and \cref{def_filter}(ii) gives $|r_{\alpha}(\lambda)| \le 1 +
  \lambda|\varphi_{\alpha}(\lambda)| \le 1 + C_{\varphi}$ for all $\alpha > 0$.
\end{proof}


\begin{lemma}[The residual constant]
  \label{lem_residual_constant}
  Let $\{\varphi_{\alpha}\}$ be a regularizing filter and, with $r_{\alpha}$ as in
  \cref{lem_approx_error_spectral}, put
  \[
    C_r \;:=\; \sup_{\alpha > 0,\ \lambda \in (0,\kappa]} |r_{\alpha}(\lambda)| .
  \]
  Then $1 \le C_r \le 1 + C_{\varphi}$, and
  \[
    \|y - KR_{\alpha}y\|_Y \;\le\; C_r\,\|y\|_Y
    \qquad \text{for all } y \in Y \text{ and all } \alpha > 0 .
  \]
\end{lemma}

\begin{proof}
  \emph{The upper bound.}  By \cref{lem_approx_error_spectral},
  $|r_{\alpha}(\lambda)| \le 1 + C_{\varphi}$ for all $\alpha > 0$ and
  $\lambda \in (0,\kappa]$, so the supremum is at most $1 + C_{\varphi}$, and in
  particular finite.

  \emph{The lower bound.}  Fix $\alpha > 0$ and $\varepsilon \in (0,1)$, and let
  $\Phi_{\alpha}$ be the supremum of $|\varphi_{\alpha}|$ over $(0,\kappa]$, finite by
  \cref{def_filter}.  Since $K \ne 0$ we have $\kappa > 0$, so
  $\lambda_{\varepsilon} := \min\{\kappa,\ \varepsilon/(\Phi_{\alpha}+1)\}$ lies in
  $(0,\kappa]$, and
  \[
    \lambda_{\varepsilon}\,|\varphi_{\alpha}(\lambda_{\varepsilon})|
      \;\le\; \frac{\varepsilon\,\Phi_{\alpha}}{\Phi_{\alpha}+1} \;<\; \varepsilon .
  \]
  Hence $|r_{\alpha}(\lambda_{\varepsilon})| \ge 1 - \lambda_{\varepsilon}
  |\varphi_{\alpha}(\lambda_{\varepsilon})| > 1 - \varepsilon$, so $C_r > 1 -
  \varepsilon$.  As $\varepsilon \in (0,1)$ was arbitrary, $C_r \ge 1$.

  \emph{The residual bound.}  Let $y \in Y$, $\alpha > 0$ and $c_n := \langle u_n,
  y\rangle_Y$.  By \cref{lem_spectral_residual_form}(iii) and
  $|r_{\alpha}(\sigma_n^{2})| \le C_r$ (legitimate because $\sigma_n^{2} \in
  (0,\kappa]$),
  \[
    \|y - KR_{\alpha}y\|_Y^{2} \;\le\; \|(\Id - P_R)y\|_Y^{2}
      + C_r^{2}\sum_{n \in I}|c_n|^{2}
    \;\le\; C_r^{2}\Bigl(\|(\Id - P_R)y\|_Y^{2} + \sum_{n \in I}|c_n|^{2}\Bigr),
  \]
  the last step because $C_r \ge 1$.  By \cref{lem_spectral_residual_form}(iv) the
  bracket equals $\|y\|_Y^{2}$, and taking square roots gives the claim.
\end{proof}


\begin{theorem}[Spectral regularizations are regularizations]
  \label{thm_spectral_is_regularization}
  Let $\{\varphi_{\alpha}\}$ be a regularizing filter.  Then $R_{\alpha}y \to K\dg y$ for
  every $y \in \Dom(K\dg)$, i.e.\ $\{R_{\alpha}\}_{\alpha>0}$ is a regularization of
  $K\dg$.  If moreover $K\dg$ is unbounded, then $\|R_{\alpha}y\|_X \to \infty$ as
  $\alpha \to 0$ for every $y \notin \Dom(K\dg)$.
\end{theorem}

\begin{proof}
  \emph{Each $R_{\alpha}$ is bounded.}  This is \cref{lem_filter_norm_bound}, so
  $\{R_{\alpha}\}_{\alpha>0} \subseteq \Ll(Y,X)$ as \cref{def_regularization} requires.

  \emph{Convergence.}  Let $y \in \Dom(K\dg)$ and $x\dg := K\dg y$.  By
  \cref{lem_approx_error_spectral} and \cref{lem_parseval_series} applied to the
  orthonormal system $\{v_n\}$,
  \[
    \|K\dg y - R_{\alpha}y\|_X^{2} \;=\; \sum_{n \in I} |r_{\alpha}(\sigma_n^{2})|^{2}\,
      |\langle v_n, x\dg\rangle_X|^{2} .
  \]
  If $x\dg = 0$ every summand vanishes and there is nothing to prove, so assume
  $x\dg \ne 0$ and let $\varepsilon > 0$.  By Bessel's inequality (\cref{thm_bessel}) the
  series $\sum_{n} |\langle v_n,x\dg\rangle_X|^{2}$ converges, so there is a finite
  $J \subseteq I$ with
  \[
    \sum_{n \in I \setminus J} |\langle v_n, x\dg\rangle_X|^{2}
      \;<\; \frac{\varepsilon^{2}}{2(1+C_{\varphi})^{2}} .
  \]
  By \cref{lem_approx_error_spectral}, $r_{\alpha}(\sigma_n^{2}) \to 0$ as $\alpha \to 0$
  for each of the finitely many $n \in J$ (note $\sigma_n^{2} \in (0,\kappa]$ by
  \cref{def_singular_system} and \cref{lem_sigma_bound}), so there
  is $\alpha_0 > 0$ with $|r_{\alpha}(\sigma_n^{2})|^{2} < \varepsilon^{2}/(2\|x\dg\|
  _X^{2})$ for all $n \in J$ and all $\alpha \in (0,\alpha_0)$; take the smallest of the
  finitely many thresholds.  Using $|r_{\alpha}| \le 1 + C_{\varphi}$ on the tail and
  \cref{thm_bessel} on the head, for such $\alpha$
  \[
    \|K\dg y - R_{\alpha}y\|_X^{2} \;\le\; \frac{\varepsilon^{2}}{2\|x\dg\|_X^{2}}
      \sum_{n \in J}|\langle v_n,x\dg\rangle_X|^{2}
      + (1+C_{\varphi})^{2}\frac{\varepsilon^{2}}{2(1+C_{\varphi})^{2}}
      \;\le\; \frac{\varepsilon^{2}}{2} + \frac{\varepsilon^{2}}{2} .
  \]
  Hence $\|K\dg y - R_{\alpha}y\|_X \le \varepsilon$ for $0 < \alpha < \alpha_0$, i.e.\
  $R_{\alpha}y \to K\dg y$.  By \cref{def_regularization}, $\{R_{\alpha}\}_{\alpha>0}$ is
  a regularization of $K\dg$.

  \emph{Divergence off the domain.}  By \cref{lem_filter_residual_bound},
  $\sup_{\alpha>0}\|KR_{\alpha}\|_{\Ll(Y,Y)} \le C_{\varphi} < \infty$, so
  \cref{thm_reg_diverges_outside_domain} applies with $T = K$ and gives
  $\|R_{\alpha}y\|_X \to \infty$ as $\alpha \to 0$ for every $y \notin \Dom(K\dg)$.
\end{proof}


\begin{lemma}[Approximation error under a source condition]
  \label{lem_omega_nu}
  Put $\omega_{\nu}(\alpha) := \sup_{\lambda \in (0,\kappa]}
  \lambda^{\nu/2}|r_{\alpha}(\lambda)|$.  Let $y \in \Dom(K\dg)$ and $x\dg = K\dg y \in
  X_{\nu,\rho}$ for some $\nu, \rho > 0$.  Then for all $\alpha > 0$
  \[
    \|x_{\alpha} - x\dg\|_X \le \omega_{\nu}(\alpha)\,\rho,
    \qquad
    \|Kx_{\alpha} - Kx\dg\|_Y \le \omega_{\nu+1}(\alpha)\,\rho .
  \]
\end{lemma}

\begin{proof}
  \emph{$\omega_{\nu}$ is finite.}  By \cref{lem_approx_error_spectral},
  $|r_{\alpha}(\lambda)| \le 1 + C_{\varphi}$ on $(0,\kappa]$, so $\omega_{\nu}(\alpha)
  \le \kappa^{\nu/2}(1+C_{\varphi}) < \infty$ for every $\nu > 0$; the same holds for
  $\omega_{\nu+1}$.

  \emph{Setup.}  By \cref{def_source_set} there is $w \in X$ with $\|w\|_X \le \rho$ and
  $x\dg = |K|^{\nu}w$.  Write $\psi_s(\lambda) := \lambda^{s/2}$ for $s > 0$, so that
  $|K|^{s} = (K^{*}K)^{s/2} = \psi_s(K^{*}K)$ by \cref{def_frac_power}.  By
  \cref{lem_approx_error_spectral} and multiplicativity of the functional calculus
  (\cref{lem_fun_calc_hom}),
  \[
    x\dg - x_{\alpha} \;=\; r_{\alpha}(K^{*}K)\,\psi_{\nu}(K^{*}K)\,w
      \;=\; (r_{\alpha}\psi_{\nu})(K^{*}K)\,w ,
  \]
  and since $(r_{\alpha}\psi_{\nu})(0) = r_{\alpha}(0)\cdot 0 = 0$, the projection term
  in \cref{def_fun_calc} drops out:
  \[
    x\dg - x_{\alpha} \;=\; \sum_{n \in I} r_{\alpha}(\sigma_n^{2})\,\sigma_n^{\nu}\,
      \langle v_n, w\rangle_X\, v_n .
  \]

  \emph{The first estimate.}  Since $\sigma_n^{2} \in (0,\kappa]$
  (\cref{def_singular_system} and \cref{lem_sigma_bound}), the definition of
  $\omega_{\nu}$ gives $\sigma_n^{\nu}|r_{\alpha}(\sigma_n^{2})| \le
  \omega_{\nu}(\alpha)$.  Hence, by \cref{lem_parseval_series} applied to $\{v_n\}$,
  \[
    \begin{aligned}
      \|x_{\alpha} - x\dg\|_X^{2}
        &\;=\; \sum_{n \in I} |r_{\alpha}(\sigma_n^{2})|^{2}\,\sigma_n^{2\nu}\,
          |\langle v_n,w\rangle_X|^{2} \\
        &\;\le\; \omega_{\nu}(\alpha)^{2}\sum_{n \in I}|\langle v_n,w\rangle_X|^{2}
          \;\le\; \omega_{\nu}(\alpha)^{2}\rho^{2},
    \end{aligned}
  \]
  the last step by Bessel's inequality (\cref{thm_bessel}) and $\|w\|_X \le \rho$.

  \emph{The second estimate.}  By \cref{lem_abs_norm}, $\|Kz\|_Y = \||K|z\|_X$ for every
  $z \in X$; applied to $z := x_{\alpha} - x\dg$ and using $|K| = \psi_1(K^{*}K)$ and
  \cref{lem_fun_calc_hom} once more,
  \[
    |K|(x\dg - x_{\alpha}) \;=\; (\psi_1 r_{\alpha}\psi_{\nu})(K^{*}K)w
      \;=\; \sum_{n \in I} r_{\alpha}(\sigma_n^{2})\,\sigma_n^{\nu+1}\,
        \langle v_n,w\rangle_X\, v_n ,
  \]
  again with no projection term, since $(\psi_1r_{\alpha}\psi_{\nu})(0) = 0$.  As above,
  $\sigma_n^{\nu+1}|r_{\alpha}(\sigma_n^{2})|$ is bounded by $\omega_{\nu+1}(\alpha)$,
  so
  \begin{align*}
      \|Kx_{\alpha} - Kx\dg\|_Y^{2} &= \||K|(x\dg - x_{\alpha})\|_X^{2} \\
        &\le \omega_{\nu+1}(\alpha)^{2}\sum_{n \in I}|\langle v_n,w\rangle_X|^{2}
        \le \omega_{\nu+1}(\alpha)^{2}\rho^{2} . \qedhere
  \end{align*}
\end{proof}


\begin{definition}[Qualification]
  \label{def_qualification}
  Say the filter satisfies condition (Q$_{\nu}$) if there is $C_{\nu} > 0$ with
  $\omega_{\nu}(\alpha) \le C_{\nu}\alpha^{\nu/2}$ for all $\alpha \in (0,\kappa]$; and condition (F)
  if there is $C > 0$ with $\sup_{\lambda \in (0,\kappa]}|\varphi_{\alpha}(\lambda)| \le
  C\alpha^{-1}$ for all $\alpha > 0$.  The \emph{qualification} of the filter is
  \[
    \nu_0 := \sup\{\, \nu > 0 : (\text{Q}_{\nu}) \text{ holds} \,\} \in (0,\infty] .
  \]
\end{definition}


\begin{theorem}[Order optimality for an a priori rule]
  \label{thm_apriori_order_optimal}
  Let $y \in \Dom(K\dg)$ and $x\dg = K\dg y \in X_{\nu,\rho}$ for $\nu,\rho > 0$.  Suppose
  the filter satisfies (F) with constant $C_F$, and (Q$_{\nu}$) with constant $C_{\nu}$
  \emph{for every $\alpha > 0$} and not only on $(0,\kappa]$ as
  \cref{def_qualification} asks, and that $\alpha$ is an a priori rule with
  \[
    c\,(\delta/\rho)^{2/(\nu+1)} \;\le\; \alpha(\delta) \;\le\;
      C\,(\delta/\rho)^{2/(\nu+1)}, \qquad c, C > 0 .
  \]
  Then
  \[
    \|x^{\delta}_{\alpha(\delta)} - x\dg\|_X \;\le\;
      \bigl(\sqrt{C_{\varphi}C_F}\,c^{-1/2} + C_{\nu}C^{\nu/2}\bigr)\,
      \delta^{\nu/(\nu+1)}\rho^{1/(\nu+1)} ;
  \]
  in particular $(R_{\alpha},\alpha)$ is order optimal for this $\nu$ and this $\rho$, and
  it is a convergent regularization method by \cref{thm_apriori_characterisation}.
\end{theorem}

\begin{proof}
  Throughout, $C_F$ is the constant of condition (F) in \cref{def_qualification} and
  $C$, $c$ are the constants of the a priori rule; these are three unrelated numbers and
  no relation between them is assumed.

  \emph{The data error.}  Let $\delta > 0$.  Condition (F) gives $\sup_{\lambda \in
  (0,\kappa]}|\varphi_{\alpha(\delta)}(\lambda)| \le C_F\alpha(\delta)^{-1}$, so
  \cref{lem_filter_norm_bound} yields
  \[
    \|R_{\alpha(\delta)}\|_{\Ll(Y,X)} \;\le\; \sqrt{C_{\varphi}}\,
      \bigl(C_F\alpha(\delta)^{-1}\bigr)^{1/2}
      \;=\; \sqrt{C_{\varphi}C_F}\,\alpha(\delta)^{-1/2} .
  \]
  The lower bound in the rule gives $\alpha(\delta)^{-1/2} \le c^{-1/2}
  (\delta/\rho)^{-1/(\nu+1)}$, whence
  \[
    \begin{aligned}
      \delta\,\|R_{\alpha(\delta)}\|_{\Ll(Y,X)}
        &\;\le\; \sqrt{C_{\varphi}C_F}\,c^{-1/2}\,\delta^{1-1/(\nu+1)}\rho^{1/(\nu+1)} \\
        &\;=\; \sqrt{C_{\varphi}C_F}\,c^{-1/2}\,\delta^{\nu/(\nu+1)}\rho^{1/(\nu+1)} .
    \end{aligned}
  \]

  \emph{The approximation error.}  By \cref{lem_omega_nu} and condition (Q$_{\nu}$) of
  \cref{def_qualification}, and then the upper bound in the rule,
  \[
    \begin{aligned}
      \|x_{\alpha(\delta)} - x\dg\|_X
        &\;\le\; \omega_{\nu}(\alpha(\delta))\,\rho
        \;\le\; C_{\nu}\,\alpha(\delta)^{\nu/2}\rho \\
        &\;\le\; C_{\nu}C^{\nu/2}(\delta/\rho)^{\nu/(\nu+1)}\rho
        \;=\; C_{\nu}C^{\nu/2}\delta^{\nu/(\nu+1)}\rho^{1/(\nu+1)} .
    \end{aligned}
  \]

  \emph{The estimate.}  Adding the two bounds and applying
  \cref{lem_error_decomposition} with $T = K$, for every $y^{\delta} \in
  B_{\delta}(y)$,
  \[
    \begin{aligned}
      \|x^{\delta}_{\alpha(\delta)} - x\dg\|_X
        &\;\le\; \delta\,\|R_{\alpha(\delta)}\|_{\Ll(Y,X)}
          + \|x_{\alpha(\delta)} - x\dg\|_X \\
        &\;\le\; \bigl(\sqrt{C_{\varphi}C_F}\,c^{-1/2} + C_{\nu}C^{\nu/2}\bigr)
          \delta^{\nu/(\nu+1)}\rho^{1/(\nu+1)},
    \end{aligned}
  \]
  which is the asserted inequality.

  \emph{Order optimality.}  Let $x\dg \in X_{\nu,\rho}$ (\cref{def_source_set}) be
  arbitrary.  Then $x\dg \in
  \Ker(K)^{\perp}$ by \cref{lem_source_in_kerperp}, so $y := Kx\dg \in \Ran(K) \subseteq
  \Dom(K\dg)$ and $K\dg y = K\dg Kx\dg = (\Id - P_N)x\dg = x\dg$ by
  \cref{lem_mp_equations}(iii); thus the hypotheses of the theorem hold for this $y$, and
  the displayed bound is uniform over $y^{\delta} \in B_{\delta}(Kx\dg)$.  Taking the
  supremum, \cref{def_worst_case_error} gives $E(Kx\dg,\delta) \le
  c_0\,\delta^{\nu/(\nu+1)}\rho^{1/(\nu+1)}$ for all $\delta > 0$ and all $x\dg \in
  X_{\nu,\rho}$, with $c_0 := \max\{1, \sqrt{C_{\varphi}C_F}c^{-1/2}
  + C_{\nu}C^{\nu/2}\}$.
  By \cref{def_order_optimal} this is order optimality with constant $c_0$ for this $\nu$
  and this $\rho$.

  \emph{Convergence.}  $\{R_{\alpha}\}_{\alpha>0}$ is a regularization of $K\dg$ by
  \cref{thm_spectral_is_regularization}.  The rule satisfies $\alpha(\delta) \le
  C(\delta/\rho)^{2/(\nu+1)} \to 0$ and, by the data error bound above,
  $\delta\|R_{\alpha(\delta)}\|_{\Ll(Y,X)} \le \sqrt{C_{\varphi}C_F}c^{-1/2}
  \delta^{\nu/(\nu+1)}\rho^{1/(\nu+1)} \to 0$ as $\delta \to 0$, because $\nu > 0$.  By
  \cref{thm_apriori_characterisation}(i), $(R_{\alpha},\alpha)$ is a convergent
  regularization method.
\end{proof}


\begin{theorem}[Rate and convergence for the discrepancy principle]
  \label{thm_discrepancy_order_optimal}
  Assume $\Ran(K)$ is dense in $Y$ and $\dim\Ran(K) = \infty$.  Let $\{\varphi_{\alpha}\}$
  be a filter with qualification $\nu_0 > 1$ satisfying (F) and (Q$_{\nu}$) for all
  $\nu \in (0,\nu_0]$, assume $\alpha \mapsto \varphi_{\alpha}(\lambda)$ is continuous on
  $(0,\infty)$ for each $\lambda \in (0,\kappa]$, and let $\tau > C_r$, with $C_r$ as in
  \cref{lem_residual_constant}.  Then the discrepancy principle of
  \cref{def_discrepancy_principle} defines an a posteriori rule $\alpha_{\tau}$, and:
  \begin{enumerate}
    \item \emph{(rate)} for every $\nu \in (0,\nu_0-1]$ there is $c = c(\nu,\tau) \ge 1$,
      depending only on $\nu$, on $\tau$ and on the filter --- in particular not on
      $\rho$ --- such that
      \[
        \bigl\|R_{\alpha_{\tau}(\delta,y^{\delta})}y^{\delta} - x\dg\bigr\|_X
          \;\le\; c\,\delta^{\nu/(\nu+1)}\rho^{1/(\nu+1)}
      \]
      for every $\rho > 0$, every $x\dg \in X_{\nu,\rho}$, every $\delta > 0$ and every
      $y^{\delta} \in B_{\delta}(Kx\dg)$ \emph{with} $\|y^{\delta}\|_Y > \tau\delta$;
    \item \emph{(convergence)} $(R_{\alpha},\alpha_{\tau})$ is a convergent
      regularization method.
  \end{enumerate}
\end{theorem}

\begin{proof}
  Let $C_F$ be the constant of condition (F) in \cref{def_qualification} and $C_{\nu}$
  those of (Q$_{\nu}$).

  \emph{Step 0: the residual identity.}  By \cref{lem_spectral_residual_form}(iii),
  for all $\alpha > 0$ and $z \in Y$,
  \begin{equation}\label{eq_resid_identity}
    \|z - KR_{\alpha}z\|_Y^{2} \;=\; \|(\Id-P_R)z\|_Y^{2}
      \;+\; \sum_{n \in I} |r_{\alpha}(\sigma_n^{2})|^{2}|\langle u_n,z\rangle_Y|^{2} .
  \end{equation}

  \emph{Step 1: the residual constant.}  By \cref{lem_residual_constant} we have
  $C_r \ge 1$, so $\tau > C_r \ge 1$ and $\tau$ is admissible in
  \cref{def_discrepancy_principle}; the same lemma gives
  \begin{equation}\label{eq_resid_uniform}
    \|z - KR_{\alpha}z\|_Y \;\le\; C_r\,\|z\|_Y \qquad
      \text{for all } z \in Y,\ \alpha > 0 .
  \end{equation}

  \emph{Step 2: the rule exists.}  Fix $\delta > 0$ and $y^{\delta} \in Y$ with
  $\|y^{\delta}\|_Y > \tau\delta$, and put $\theta(\alpha) := \|KR_{\alpha}y^{\delta} -
  y^{\delta}\|_Y$ and $S := \{\alpha > 0 : \theta(\alpha) \le \tau\delta\}$.  Since
  $\Ran(K)$ is dense in $Y$ we have $P_R = \Id$, so \eqref{eq_resid_identity} and
  \cref{lem_spectral_residual_form}(iv) read
  \begin{equation}\label{eq_theta_series}
    \theta(\alpha)^{2} \;=\; \sum_{n \in I} |r_{\alpha}(\sigma_n^{2})|^{2}
      |\langle u_n, y^{\delta}\rangle_Y|^{2},
    \qquad \sum_{n \in I}|\langle u_n,y^{\delta}\rangle_Y|^{2} = \|y^{\delta}\|_Y^{2} .
  \end{equation}

  \emph{$\theta$ is continuous.}  Each $\alpha \mapsto |r_{\alpha}(\sigma_n^{2})|^{2}
  |\langle u_n,y^{\delta}\rangle_Y|^{2}$ is continuous, by the assumed continuity of
  $\alpha \mapsto \varphi_{\alpha}(\lambda)$.  Given $\varepsilon > 0$, choose a finite
  $J \subseteq I$ with $\sum_{n \in I\setminus J}|\langle u_n,y^{\delta}\rangle_Y|^{2} <
  \varepsilon$; then for all $\alpha,\beta > 0$, using $|r| \le C_r$ on the tail,
  \[
    |\theta(\alpha)^{2} - \theta(\beta)^{2}| \;\le\;
      \sum_{n \in J}\bigl||r_{\alpha}(\sigma_n^{2})|^{2}
        - |r_{\beta}(\sigma_n^{2})|^{2}\bigr|\,|\langle u_n,y^{\delta}\rangle_Y|^{2}
      \;+\; 2C_r^{2}\varepsilon .
  \]
  The first term is a finite sum of continuous functions, so $\theta^{2}$, and hence
  $\theta$, is continuous on $(0,\infty)$.

  \emph{$S \ne \emptyset$.}  The family $\{R_{\alpha}\}$ is a regularization of $K\dg$
  by \cref{thm_spectral_is_regularization}, and
  $\sup_{\alpha>0}\|KR_{\alpha}\|_{\Ll(Y,Y)} \le C_{\varphi}$
  by \cref{lem_filter_residual_bound}; with $\Ran(K)$ dense,
  \cref{thm_discrepancy_well_defined} applied to any strictly decreasing null sequence
  $\{\alpha_n\}$ produces an index with $\theta(\alpha_{n^{*}}) \le \tau\delta$.

  \emph{$S$ is bounded above.}  Put $A := \kappa C_F\|y^{\delta}\|_Y/(\|y^{\delta}\|_Y -
  \tau\delta) > 0$ and let $\alpha > A$.  Condition (F) gives $\lambda\varphi_{\alpha}
  (\lambda) \le \kappa C_F\alpha^{-1}$ for $\lambda \in (0,\kappa]$, hence
  \[
    r_{\alpha}(\lambda) \;\ge\; 1 - \frac{\kappa C_F}{\alpha} \;>\;
      1 - \frac{\kappa C_F}{A} \;=\; \frac{\tau\delta}{\|y^{\delta}\|_Y} \;>\; 0 .
  \]
  Inserting this into \eqref{eq_theta_series} gives $\theta(\alpha) \ge
  (\tau\delta/\|y^{\delta}\|_Y)\|y^{\delta}\|_Y = \tau\delta$, and in fact
  $\theta(\alpha) > \tau\delta$ because the first inequality above is strict and
  $\|y^{\delta}\|_Y > 0$.  Thus $S \subseteq (0,A]$.  \emph{This is exactly where the
  hypothesis $\|y^{\delta}\|_Y > \tau\delta$ of \cref{def_discrepancy_principle} is
  used}; without it $S$ may be unbounded and no admissible parameter exists.

  \emph{The supremum is attained.}  Let $\alpha^{*} := \sup S \in (0,A]$.  There are
  $\alpha_k \in S$ with $\alpha_k \to \alpha^{*}$, and $\theta$ is continuous at
  $\alpha^{*} > 0$, so $\theta(\alpha^{*}) = \lim_k \theta(\alpha_k) \le \tau\delta$,
  i.e.\ $\alpha^{*} \in S$; and $\theta(\alpha) > \tau\delta$ for every $\alpha >
  \alpha^{*}$ by definition of the supremum.  Setting $\alpha_{\tau}(\delta,y^{\delta})
  := \alpha^{*}$ for such data, and $\alpha_{\tau}(\delta,y^{\delta}) := 1$ when
  $\|y^{\delta}\|_Y \le \tau\delta$ (where \cref{def_discrepancy_principle} imposes no
  condition), defines an a posteriori rule in the sense of
  \cref{def_parameter_choice} satisfying the discrepancy principle.

  \emph{Step 3: the rate.}  Fix $\nu \in (0,\nu_0-1]$, so that both (Q$_{\nu}$) and
  (Q$_{\nu+1}$) are available, let $\rho > 0$, $x\dg \in X_{\nu,\rho}$, $\delta > 0$ and
  $y^{\delta} \in B_{\delta}(y)$ with $\|y^{\delta}\|_Y > \tau\delta$, where $y :=
  Kx\dg$.  By \cref{lem_source_in_kerperp}, $x\dg \in \Ker(K)^{\perp}$, so $y \in
  \Ran(K) \subseteq \Dom(K\dg)$ and $K\dg y = (\Id - P_N)x\dg = x\dg$ by
  \cref{lem_mp_equations}(iii).  Write $\alpha := \alpha_{\tau}(\delta,y^{\delta})$ and
  split
  \[
    \|x^{\delta}_{\alpha} - x\dg\|_X \;\le\;
      \|x_{\alpha} - x\dg\|_X + \|x^{\delta}_{\alpha} - x_{\alpha}\|_X .
  \]

  \emph{The approximation error.}  By \cref{def_source_set} pick $w$ with $x\dg =
  |K|^{\nu}w$ and $\|w\|_X \le \rho$; replacing $w$ by $(\Id - P_N)w$ changes neither
  $|K|^{\nu}w$ nor the bound on the norm (\cref{prop_proj_props}), since by
  \cref{def_fun_calc} the value of $\psi(K^{*}K)w$ with $\psi(\lambda) =
  \lambda^{\nu/2}$ depends on $w$ only through the coefficients $\langle v_n,w\rangle_X$
  and $\psi(0) = 0$; so assume $w \in \Ker(K)^{\perp}$.  Put $\xi := r_{\alpha}(K^{*}K)w$.
  By \cref{lem_approx_error_spectral} and \cref{lem_fun_calc_hom},
  \[
    x\dg - x_{\alpha} \;=\; r_{\alpha}(K^{*}K)|K|^{\nu}w
      \;=\; |K|^{\nu}\,r_{\alpha}(K^{*}K)w \;=\; |K|^{\nu}\xi .
  \]
  The interpolation inequality \cref{thm_interpolation} with exponents $\nu$ and $\nu+1$
  gives
  \[
    \|x_{\alpha} - x\dg\|_X \;=\; \||K|^{\nu}\xi\|_X \;\le\;
      \||K|^{\nu+1}\xi\|_X^{\nu/(\nu+1)}\,\|\xi\|_X^{1/(\nu+1)} .
  \]
  Since $w \in \Ker(K)^{\perp}$, \cref{def_fun_calc} gives $\xi = \sum_n
  r_{\alpha}(\sigma_n^{2})\langle v_n,w\rangle_X v_n$, so $\|\xi\|_X \le C_r\|w\|_X \le
  C_r\rho$ by \cref{lem_parseval_series} and \cref{thm_bessel}.  For the other factor,
  \cref{lem_fun_calc_hom} and \cref{lem_abs_norm} give
  \[
    \||K|^{\nu+1}\xi\|_X = \||K|(x\dg - x_{\alpha})\|_X = \|K(x\dg - x_{\alpha})\|_Y
      = \|y - Kx_{\alpha}\|_Y ,
  \]
  and, adding and subtracting $y^{\delta} - Kx^{\delta}_{\alpha}$ and using linearity of
  $R_{\alpha}$,
  \[
    \|y - Kx_{\alpha}\|_Y \;\le\; \|y^{\delta} - Kx^{\delta}_{\alpha}\|_Y
      + \|(\Id - KR_{\alpha})(y - y^{\delta})\|_Y
      \;\le\; \tau\delta + C_r\delta ,
  \]
  by the discrepancy principle and \eqref{eq_resid_uniform}.  Altogether
  \[
    \|x_{\alpha} - x\dg\|_X \;\le\; C_1\,\delta^{\nu/(\nu+1)}\rho^{1/(\nu+1)},
    \qquad C_1 := (\tau + C_r)^{\nu/(\nu+1)}C_r^{1/(\nu+1)} .
  \]

  \emph{A lower bound for $\alpha$.}  Since $2\alpha > \alpha$, the discrepancy
  principle gives $\|Kx^{\delta}_{2\alpha} - y^{\delta}\|_Y > \tau\delta$, while the
  triangle inequality and \eqref{eq_resid_uniform} give
  $\|Kx^{\delta}_{2\alpha} - y^{\delta}\|_Y \le \|Kx_{2\alpha} - y\|_Y +
  \|(\Id-KR_{2\alpha})(y-y^{\delta})\|_Y \le \|Kx_{2\alpha} - y\|_Y + C_r\delta$.  Hence
  $\|Kx_{2\alpha} - Kx\dg\|_Y > (\tau - C_r)\delta$.  On the other hand
  \cref{lem_omega_nu} and (Q$_{\nu+1}$) give $\|Kx_{2\alpha} - Kx\dg\|_Y \le
  \omega_{\nu+1}(2\alpha)\rho \le C_{\nu+1}2^{(\nu+1)/2}\alpha^{(\nu+1)/2}\rho$.
  As $\tau > C_r$, combining the two,
  \[
    \delta \;<\; C_{\tau}\,\alpha^{(\nu+1)/2}\rho, \qquad\text{where}\qquad
      C_{\tau} := \frac{2^{(\nu+1)/2}C_{\nu+1}}{\tau - C_r} ,
  \]
  and therefore
  \[
    \alpha^{-1/2} \;<\; C_{\tau}^{1/(\nu+1)}\,\delta^{-1/(\nu+1)}\rho^{1/(\nu+1)} .
  \]

  \emph{The data error.}  By \cref{lem_filter_norm_bound} and (F),
  \begin{equation}\label{eq_disc_dataerr}
    \|R_{\alpha}\|_{\Ll(Y,X)} \;\le\; \sqrt{C_{\varphi}C_F}\,\alpha^{-1/2}
    \qquad\text{for every } \alpha > 0 ,
  \end{equation}
  so with the previous display
  \[
    \|x^{\delta}_{\alpha} - x_{\alpha}\|_X \;\le\; \|R_{\alpha}\|_{\Ll(Y,X)}\,\delta
      \;\le\; C_2\,\delta^{\nu/(\nu+1)}\rho^{1/(\nu+1)},
    \qquad C_2 := \sqrt{C_{\varphi}C_F}\,C_{\tau}^{1/(\nu+1)} .
  \]

  \emph{Conclusion.}  With $c := \max\{1, C_1 + C_2\}$, which depends only on the filter,
  on $\tau$ and on $\nu$,
  \begin{equation}\label{eq_disc_rate}
    \|x^{\delta}_{\alpha_{\tau}(\delta,y^{\delta})} - x\dg\|_X \;\le\;
      c\,\delta^{\nu/(\nu+1)}\rho^{1/(\nu+1)}
  \end{equation}
  for every $\rho > 0$, every $x\dg \in X_{\nu,\rho}$, every $\delta > 0$ and every
  $y^{\delta} \in B_{\delta}(Kx\dg)$ with $\|y^{\delta}\|_Y > \tau\delta$.  This is
  assertion (i).

  \emph{A remark on the excluded data.}  The restriction $\|y^{\delta}\|_Y > \tau\delta$
  in \eqref{eq_disc_rate} cannot be removed, and not merely for want of a better
  argument: it is exactly the set on which \cref{def_discrepancy_principle} constrains
  the rule.  Indeed $x\dg = 0$ lies in $X_{\nu,\rho}$ for \emph{every} $\rho > 0$, so a
  bound of the form \eqref{eq_disc_rate} valid for all $y^{\delta} \in B_{\delta}(0)$ and
  all $\rho > 0$ would force $R_{\alpha_{\tau}(\delta,y^{\delta})}y^{\delta} = 0$ for all
  $\|y^{\delta}\|_Y \le \delta$, which no spectral regularization with
  $\varphi_{\alpha} \not\equiv 0$ satisfies.  In particular the supremum over
  \emph{all} $y^{\delta} \in B_{\delta}(Kx\dg)$ that \cref{def_worst_case_error} forms
  is \emph{not} bounded by the right-hand side of \eqref{eq_disc_rate} uniformly in
  $\rho$, so \eqref{eq_disc_rate} does not deliver the hypothesis of
  \cref{thm_order_optimal_implies_convergence}.  Assertion (ii) is therefore proved
  directly, in Step 4.

  \emph{Step 4: convergence.}  Let $y \in \Dom(K\dg)$ and $x\dg := K\dg y$.  Since
  $\Ran(K)$ is dense, $\Ran(K)^{\perp} = \overline{\Ran(K)}^{\perp} = \{0\}$
  (\cref{lem_orthogonal_closure}, for both equalities), so $\Dom(K\dg) = \Ran(K)$ and $P_R =
  \Id$ (\cref{thm_mp_inverse}); by \cref{lem_mp_equations} we have $x\dg \in
  \Ran(K\dg) = \Ker(K)^{\perp}$ and, by item (iv), $Kx\dg = KK\dg y = P_Ry = y$.  Write
  $c_n := \langle v_n, x\dg\rangle_X$; since $\{v_n\}_{n \in I}$ is an orthonormal basis
  of $\Ker(K)^{\perp}$ (\cref{def_singular_system}, \cref{thm_onb}),
  $x\dg = \sum_{n}c_nv_n$ and $\sum_n|c_n|^{2} = \|x\dg\|_X^{2}$.  By
  \cref{def_parameter_choice} we must show
  \[
    D(\delta) \;:=\; \sup_{y^{\delta} \in B_{\delta}(y)}
      \bigl\|R_{\alpha_{\tau}(\delta,y^{\delta})}y^{\delta} - x\dg\bigr\|_X
    \;\longrightarrow\; 0 \qquad (\delta \to 0) .
  \]

  \emph{The degenerate case $y = 0$.}  Then $x\dg = K\dg 0 = 0$, and every
  $y^{\delta} \in B_{\delta}(0)$ has $\|y^{\delta}\|_Y \le \delta < \tau\delta$ because
  $\tau > 1$; so $\alpha_{\tau}(\delta,y^{\delta}) = 1$ by Step 2 and
  $\|R_1y^{\delta} - 0\|_X \le \|R_1\|_{\Ll(Y,X)}\delta$, which is finite by
  \cref{lem_filter_norm_bound}.  Hence $D(\delta) \le \|R_1\|_{\Ll(Y,X)}\delta \to 0$.

  \emph{The residual profile.}  Assume from now on $y \ne 0$, and let $0 < \delta <
  \|y\|_Y/(1+\tau)$; then every $y^{\delta} \in B_{\delta}(y)$ satisfies
  $\|y^{\delta}\|_Y \ge \|y\|_Y - \delta > \tau\delta$, so the discrepancy principle
  governs \emph{every} admissible datum.  For $\beta > 0$ put $\Theta(\beta) :=
  \|y - KR_{\beta}y\|_Y$.  Since $P_R = \Id$ and $\langle u_n,y\rangle_Y =
  \langle u_n, Kx\dg\rangle_Y = \langle K^{*}u_n, x\dg\rangle_X = \sigma_nc_n$
  (\cref{def_adjoint}, \cref{def_singular_system}), \cref{lem_spectral_residual_form}(iii)
  gives
  \begin{equation}\label{eq_disc_theta}
    \Theta(\beta)^{2} \;=\; \sum_{n \in I}\sigma_n^{2}\,
      |r_{\beta}(\sigma_n^{2})|^{2}\,|c_n|^{2} .
  \end{equation}
  Because $\nu_0 > 1$, condition (Q$_1$) is available; with $\sigma_n^{2} \in (0,\kappa]$
  (\cref{def_singular_system} and \cref{lem_sigma_bound}) it gives
  $\sigma_n|r_{\beta}(\sigma_n^{2})| \le \omega_1(\beta) \le C_1\beta^{1/2}$
  (\cref{lem_omega_nu}, \cref{def_qualification}), so
  \begin{equation}\label{eq_disc_theta_bdd}
    \frac{\sigma_n^{2}|r_{\beta}(\sigma_n^{2})|^{2}}{\beta} \;\le\; C_1^{2}
    \qquad\text{for all } n \in I,\ \beta > 0 .
  \end{equation}

  \emph{Claim: $\Theta(\beta)/\sqrt{\beta} \to 0$ as $\beta \to 0$.}  Put $\nu_1 :=
  \min\{2,\nu_0\} \in (1,2]$, so that (Q$_{\nu_1}$) holds, and note that it gives
  $|r_{\beta}(\lambda)| \le C_{\nu_1}\beta^{\nu_1/2}\lambda^{-\nu_1/2}$ on $(0,\kappa]$;
  hence for each fixed $n$
  \[
    \frac{\sigma_n^{2}|r_{\beta}(\sigma_n^{2})|^{2}}{\beta} \;\le\;
      C_{\nu_1}^{2}\,\sigma_n^{2-2\nu_1}\,\beta^{\nu_1-1}
    \;\longrightarrow\; 0 \qquad (\beta \to 0),
  \]
  because $\nu_1 > 1$.  If $x\dg = 0$ then $\Theta \equiv 0$ and there is nothing to
  prove; otherwise let $\varepsilon > 0$, choose a finite $J \subseteq I$ with
  $\sum_{n \in I \setminus J}|c_n|^{2} < \varepsilon^{2}/(2C_1^{2})$, and then
  $\beta_0 > 0$ so small that $\sigma_n^{2}|r_{\beta}(\sigma_n^{2})|^{2}/\beta <
  \varepsilon^{2}/(2\|x\dg\|_X^{2})$ for all $n \in J$ and all $\beta \in (0,\beta_0)$
  --- the minimum of finitely many thresholds.  Splitting \eqref{eq_disc_theta} at $J$
  and using \eqref{eq_disc_theta_bdd} on the tail gives $\Theta(\beta)^{2}/\beta <
  \varepsilon^{2}$ for $\beta \in (0,\beta_0)$, which is the claim.

  \emph{The data error tends to $0$ uniformly.}  Let $y^{\delta} \in B_{\delta}(y)$ and
  $\alpha := \alpha_{\tau}(\delta,y^{\delta})$.  Since $2\alpha > \alpha$,
  \cref{def_discrepancy_principle} gives $\|KR_{2\alpha}y^{\delta} - y^{\delta}\|_Y >
  \tau\delta$, while writing $y^{\delta} - KR_{2\alpha}y^{\delta} = (y - KR_{2\alpha}y)
  + (\Id - KR_{2\alpha})(y^{\delta}-y)$ and using \eqref{eq_resid_uniform} on the second
  summand gives $\|KR_{2\alpha}y^{\delta} - y^{\delta}\|_Y \le \Theta(2\alpha) +
  C_r\delta$.  Hence $\Theta(2\alpha) > (\tau - C_r)\delta$, whence
  \[
    \frac{\delta}{\sqrt{\alpha}} \;<\; \frac{\sqrt{2}}{\tau - C_r}\,
      \frac{\Theta(2\alpha)}{\sqrt{2\alpha}} .
  \]
  Let $\eta > 0$ and take $\beta_0 > 0$ with $\Theta(\beta)/\sqrt{\beta} \le \eta$ for
  $0 < \beta \le \beta_0$.  If $2\alpha \le \beta_0$ the display gives
  $\delta\alpha^{-1/2} \le \sqrt{2}\,\eta/(\tau - C_r)$; if $2\alpha > \beta_0$ then
  trivially $\delta\alpha^{-1/2} < \delta\sqrt{2/\beta_0}$.  Combining with
  \eqref{eq_disc_dataerr},
  \[
    \|R_{\alpha}y^{\delta} - R_{\alpha}y\|_X \;\le\;
      \sqrt{C_{\varphi}C_F}\,\frac{\delta}{\sqrt{\alpha}} \;\le\;
      \sqrt{C_{\varphi}C_F}\,\max\Bigl\{\frac{\sqrt{2}\,\eta}{\tau - C_r},\;
        \delta\sqrt{2/\beta_0}\Bigr\} ,
  \]
  a bound independent of $y^{\delta}$.  Letting $\delta \to 0$ shows that the
  $\limsup$ over $\delta \to 0$ of $\sup_{y^{\delta} \in B_{\delta}(y)}
  \|R_{\alpha}y^{\delta} - R_{\alpha}y\|_X$ is at most
  $\sqrt{2C_{\varphi}C_F}\,\eta/(\tau - C_r)$; as $\eta > 0$ was arbitrary, that
  $\limsup$ is $0$.

  \emph{The approximation error tends to $0$ uniformly.}  With $\alpha$ as above, the
  discrepancy principle gives $\|Kx^{\delta}_{\alpha} - y^{\delta}\|_Y \le \tau\delta$,
  so by the triangle inequality and \eqref{eq_resid_uniform} applied to $z := y -
  y^{\delta}$,
  \[
    \Theta(\alpha) \;=\; \|y - Kx_{\alpha}\|_Y \;\le\;
      \|y^{\delta} - Kx^{\delta}_{\alpha}\|_Y
      + \|(\Id - KR_{\alpha})(y - y^{\delta})\|_Y \;\le\; (\tau + C_r)\delta .
  \]
  Fix $N \in I$.  The sequence $\{\sigma_n\}$ is non-increasing
  (\cref{def_singular_system}), so $\sigma_n \ge \sigma_N > 0$ for $n \le N$; by
  \cref{lem_approx_error_spectral}, \cref{lem_parseval_series},
  $|r_{\alpha}| \le C_r$ and \eqref{eq_disc_theta},
  \[
    \begin{aligned}
      \|x_{\alpha} - x\dg\|_X^{2} &\;=\; \sum_{n \in I}|r_{\alpha}(\sigma_n^{2})|^{2}
        |c_n|^{2} \\
        &\;\le\; \sigma_N^{-2}\sum_{n \le N}\sigma_n^{2}|r_{\alpha}(\sigma_n^{2})|^{2}
          |c_n|^{2} \;+\; C_r^{2}\sum_{n > N}|c_n|^{2} \\
        &\;\le\; \sigma_N^{-2}(\tau + C_r)^{2}\delta^{2}
          \;+\; C_r^{2}\sum_{n > N}|c_n|^{2} .
    \end{aligned}
  \]
  Given $\varepsilon > 0$, choose $N$ with $C_r^{2}\sum_{n>N}|c_n|^{2} <
  \varepsilon^{2}/2$ --- possible because $\sum_n|c_n|^{2}$ converges --- and then
  $\delta$ so small that $\sigma_N^{-2}(\tau+C_r)^{2}\delta^{2} < \varepsilon^{2}/2$.
  The resulting bound $\|x_{\alpha} - x\dg\|_X \le \varepsilon$ does not depend on
  $y^{\delta}$, so $\sup_{y^{\delta} \in B_{\delta}(y)}\|x_{\alpha} - x\dg\|_X \to 0$.

  \emph{Conclusion of Step 4.}  For every $y^{\delta} \in B_{\delta}(y)$,
  \[
    \bigl\|R_{\alpha_{\tau}(\delta,y^{\delta})}y^{\delta} - x\dg\bigr\|_X
      \;\le\; \|R_{\alpha}y^{\delta} - R_{\alpha}y\|_X + \|x_{\alpha} - x\dg\|_X ,
  \]
  and both suprema tend to $0$; hence $D(\delta) \to 0$.  As $y \in \Dom(K\dg)$ was
  arbitrary, $(R_{\alpha},\alpha_{\tau})$ is a convergent regularization method
  (\cref{def_parameter_choice}), which is assertion (ii).
\end{proof}


\begin{definition}[Hanke--Raus rule]
  \label{def_hanke_raus}
  For $y^{\delta} \in Y$ put $\Psi(\alpha) := \alpha^{-1/2}\|Kx^{\delta}_{\alpha} -
  y^{\delta}\|_Y$ on $(0,\kappa]$ and choose $\alpha(y^{\delta}) \in
  \arg\min_{\alpha \in (0,\kappa]} \Psi(\alpha)$.
\end{definition}


\begin{lemma}[The discrepancy parameter is bounded below]
  \label{lem_discrepancy_parameter_bound}
  Under the hypotheses of \cref{thm_discrepancy_order_optimal}, for every
  $\nu \in (0,\nu_0-1]$ there is $c_{\tau,\nu} > 0$ such that for all $\rho > 0$,
  $x\dg \in X_{\nu,\rho}$, $\delta > 0$ and $y^{\delta} \in B_{\delta}(Kx\dg)$
  \emph{with} $\|y^{\delta}\|_Y > \tau\delta$,
  \[
    \alpha_{\tau}(\delta, y^{\delta}) \;\ge\; c_{\tau,\nu}\,(\delta/\rho)^{2/(\nu+1)} .
  \]
\end{lemma}

\begin{proof}
  Fix $\nu \in (0,\nu_0-1]$ and put
  \[
    c_{\tau,\nu} \;:=\; \Bigl(\frac{\tau - C_r}{2^{(\nu+1)/2}\,C_{\nu+1}}
      \Bigr)^{2/(\nu+1)} \;>\; 0 ,
  \]
  which is positive because $\tau > C_r$; here $C_{\nu+1}$ is the constant of
  (Q$_{\nu+1}$), available since $\nu + 1 \le \nu_0$ (\cref{def_qualification}).

  Let $\rho > 0$, $x\dg \in X_{\nu,\rho}$ (\cref{def_source_set}), $\delta > 0$ and
  $y^{\delta} \in B_{\delta}(Kx\dg)$ with $\|y^{\delta}\|_Y > \tau\delta$; write $\alpha :=
  \alpha_{\tau}(\delta,y^{\delta})$ and $y := Kx\dg$.  As in
  \cref{thm_discrepancy_order_optimal}, $x\dg \in \Ker(K)^{\perp}$ by
  \cref{lem_source_in_kerperp}, so $y \in \Ran(K)$ and $K\dg y = x\dg$ by
  \cref{lem_mp_equations}(iii).

  \emph{The residual bound.}  By \cref{lem_residual_constant},
  \[
    \|z - KR_{\beta}z\|_Y \;\le\; C_r\,\|z\|_Y
    \qquad\text{for every } z \in Y \text{ and } \beta > 0 .
  \]

  \emph{The comparison at $2\alpha$.}  Since $2\alpha > \alpha$,
  \cref{def_discrepancy_principle} gives $\|Kx^{\delta}_{2\alpha} - y^{\delta}\|_Y >
  \tau\delta$.  Writing $y^{\delta} - Kx^{\delta}_{2\alpha} = (y - Kx_{2\alpha}) +
  (\Id - KR_{2\alpha})(y^{\delta} - y)$ and applying the residual bound to $z :=
  y^{\delta}-y$,
  \[
    \tau\delta \;<\; \|Kx_{2\alpha} - y\|_Y + C_r\delta ,
    \qquad\text{so}\qquad
    \|Kx_{2\alpha} - Kx\dg\|_Y \;>\; (\tau - C_r)\delta .
  \]
  On the other hand \cref{lem_omega_nu} at exponent $\nu+1$ and (Q$_{\nu+1}$) give
  $\|Kx_{2\alpha} - Kx\dg\|_Y \le \omega_{\nu+1}(2\alpha)\rho \le
  C_{\nu+1}2^{(\nu+1)/2}\alpha^{(\nu+1)/2}\rho$.  Hence
  \[
    (\tau - C_r)\,\delta \;<\; C_{\nu+1}2^{(\nu+1)/2}\,\alpha^{(\nu+1)/2}\rho ,
  \]
  and raising to the power $2/(\nu+1)$, which is order preserving on $(0,\infty)$,
  \[
    \alpha \;>\; \Bigl(\frac{(\tau-C_r)\delta}{2^{(\nu+1)/2}C_{\nu+1}\rho}
      \Bigr)^{2/(\nu+1)} \;=\; c_{\tau,\nu}\,(\delta/\rho)^{2/(\nu+1)} . \qedhere
  \]
\end{proof}


\begin{theorem}[Conditional error estimate for Hanke--Raus]
  \label{thm_hanke_raus_conditional}
  Let the filter have qualification $\nu_0 > 1$ and satisfy (F) and (Q$_{\nu}$) for all
  $\nu \in (0,\nu_0]$, and fix $\nu \in (0,\nu_0-1]$.  Then there is $c = c(\nu) > 0$ such
  that the following holds for every $\rho > 0$, every $\delta > 0$, every
  $y \in \Ran(K)$ with $\|y\|_Y > \delta$ such that $x\dg := K\dg y \in X_{\nu,\rho}$, and
  every $y^{\delta} \in B_{\delta}(y)$ for which a minimiser
  $\alpha^{*} := \alpha(y^{\delta}) \in (0,\kappa]$ of $\Psi$ exists with
  $\delta^{*} := \|Kx^{\delta}_{\alpha^{*}} - y^{\delta}\|_Y > 0$:
  \[
    \|x^{\delta}_{\alpha^{*}} - x\dg\|_X \;\le\;
      c\Bigl(1 + \frac{\delta}{\delta^{*}}\Bigr)
      \max\{\delta,\delta^{*}\}^{\nu/(\nu+1)}\,\rho^{1/(\nu+1)} .
  \]
\end{theorem}

\begin{proof}
  Let $C_F$ be the constant of (F) in \cref{def_qualification}, let $C_{\nu}$,
  $C_{\nu+1}$ be those of (Q$_{\nu}$),
  (Q$_{\nu+1}$) --- available because $\nu + 1 \le \nu_0$ --- and let $C_r$ be as in
  \cref{lem_residual_constant}, so that $1 \le C_r \le 1 + C_{\varphi}$.  Put
  \[
    c \;:=\; \max\bigl\{\,(1+C_r)^{\nu/(\nu+1)}C_r^{1/(\nu+1)},\;
      \sqrt{C_{\varphi}C_F}\,(C_{\nu+1} + C_r)\,\bigr\} ,
  \]
  which depends only on the filter and on $\nu$.  Now fix $\rho,\delta > 0$, $y$,
  $y^{\delta}$, $\alpha^{*}$ and $\delta^{*}$ as in the statement.

  \emph{Preliminaries.}  Since $y \in \Ran(K)$ we have $P_Ry = y$, so
  $Kx\dg = KK\dg y = y$ by \cref{lem_mp_equations}(iv).  By
  \cref{lem_residual_constant}, for every $z \in Y$ and $\beta > 0$
  \begin{equation}\label{eq_hr_resid}
    \|z - KR_{\beta}z\|_Y \;\le\; C_r\,\|z\|_Y ;
  \end{equation}
  note that no density assumption is made or needed here.  Split
  $\|x^{\delta}_{\alpha^{*}} - x\dg\|_X \le \|x_{\alpha^{*}} - x\dg\|_X +
  \|x^{\delta}_{\alpha^{*}} - x_{\alpha^{*}}\|_X$.

  \emph{The approximation error.}  By \cref{def_source_set} write $x\dg = |K|^{\nu}w$
  with $\|w\|_X \le \rho$ and, as in \cref{thm_discrepancy_order_optimal}, with
  $w \in \Ker(K)^{\perp}$.  With $\xi := r_{\alpha^{*}}(K^{*}K)w$,
  \cref{lem_approx_error_spectral} and \cref{lem_fun_calc_hom} give $x\dg -
  x_{\alpha^{*}} = |K|^{\nu}\xi$, and \cref{def_fun_calc} together with
  \cref{lem_parseval_series} and \cref{thm_bessel} give $\|\xi\|_X \le C_r\rho$.
  Moreover, by \cref{lem_fun_calc_hom}, \cref{lem_abs_norm} and
  \eqref{eq_hr_resid} applied to $z := y - y^{\delta}$,
  \[
    \||K|^{\nu+1}\xi\|_X = \|K(x\dg - x_{\alpha^{*}})\|_Y = \|y - Kx_{\alpha^{*}}\|_Y
      \le \|y^{\delta} - Kx^{\delta}_{\alpha^{*}}\|_Y + C_r\delta
      = \delta^{*} + C_r\delta .
  \]
  Since $\delta^{*} + C_r\delta \le (1 + C_r)\max\{\delta,\delta^{*}\}$, the
  interpolation inequality \cref{thm_interpolation} with exponents $\nu$ and $\nu+1$ gives
  \[
    \|x_{\alpha^{*}} - x\dg\|_X \le \||K|^{\nu+1}\xi\|_X^{\nu/(\nu+1)}
      \|\xi\|_X^{1/(\nu+1)}
      \le c\,\max\{\delta,\delta^{*}\}^{\nu/(\nu+1)}\rho^{1/(\nu+1)} .
  \]

  \emph{The data error.}  By \cref{lem_filter_norm_bound} and (F),
  $\|R_{\alpha^{*}}\|_{\Ll(Y,X)} \le \sqrt{C_{\varphi}C_F}(\alpha^{*})^{-1/2}$, so, using
  $\delta^{*} > 0$ to insert the factor $\delta^{*}/\delta^{*}$ and the definition of
  $\Psi$ in \cref{def_hanke_raus},
  \[
    \|x^{\delta}_{\alpha^{*}} - x_{\alpha^{*}}\|_X \le \|R_{\alpha^{*}}\|_{\Ll(Y,X)}\,
      \delta \le \sqrt{C_{\varphi}C_F}\,\frac{\delta}{\delta^{*}}\,
      \frac{\|Kx^{\delta}_{\alpha^{*}} - y^{\delta}\|_Y}{\sqrt{\alpha^{*}}}
      = \sqrt{C_{\varphi}C_F}\,\frac{\delta}{\delta^{*}}\,\Psi(\alpha^{*}) .
  \]

  \emph{Bounding $\Psi(\alpha^{*})$.}  For any $\alpha \in (0,\kappa]$, the triangle
  inequality, \eqref{eq_hr_resid} and \cref{lem_omega_nu} at exponent $\nu+1$ together
  with (Q$_{\nu+1}$) give
  \[
    \|Kx^{\delta}_{\alpha} - y^{\delta}\|_Y \;\le\; \|Kx_{\alpha} - y\|_Y + C_r\delta
      \;\le\; C_{\nu+1}\alpha^{(\nu+1)/2}\rho + C_r\delta ,
  \]
  so that, dividing by $\sqrt{\alpha}$,
  \[
    \Psi(\alpha) \;\le\; C_{\nu+1}\alpha^{\nu/2}\rho + C_r\delta\,\alpha^{-1/2} .
  \]
  Take $\alpha_{\delta} := (\delta/\rho)^{2/(\nu+1)}$.  This is an admissible test
  parameter: by \cref{lem_fun_calc_bounded}, $\||K|^{\nu}\|_{\Ll(X,X)} \le \kappa^{\nu/2}$,
  so the hypothesis $\|y\|_Y > \delta$ gives
  \[
    \begin{aligned}
      \delta \;<\; \|y\|_Y \;=\; \|Kx\dg\|_Y
        &\;\le\; \|K\|_{\Ll(X,Y)}\,\bigl\||K|^{\nu}w\bigr\|_X \\
        &\;\le\; \sqrt{\kappa}\,\kappa^{\nu/2}\rho \;=\; \kappa^{(\nu+1)/2}\rho ,
    \end{aligned}
  \]
  whence $\alpha_{\delta} < \kappa$.   Since $\alpha_{\delta}^{\nu/2}\rho =
  \delta^{\nu/(\nu+1)}\rho^{1/(\nu+1)} = \delta\alpha_{\delta}^{-1/2}$, minimality of
  $\alpha^{*}$ (\cref{def_hanke_raus}) gives
  \[
    \Psi(\alpha^{*}) \;\le\; \Psi(\alpha_{\delta}) \;\le\;
      (C_{\nu+1} + C_r)\,\delta^{\nu/(\nu+1)}\rho^{1/(\nu+1)} .
  \]

  \emph{Conclusion.}  Combining, and using $\delta \le \max\{\delta,\delta^{*}\}$,
  \[
    \|x^{\delta}_{\alpha^{*}} - x\dg\|_X \le c\Bigl(1 + \frac{\delta}{\delta^{*}}\Bigr)
      \max\{\delta,\delta^{*}\}^{\nu/(\nu+1)}\rho^{1/(\nu+1)} . \qedhere
  \]
\end{proof}


\begin{theorem}[Convergence of Hanke--Raus on data bounded away from the range]
  \label{thm_hanke_raus_convergence}
  Fix $\varepsilon > 0$ and put
  $N_{\varepsilon} := \{\, y + \eta \in Y : \|(\Id - P_R)\eta\|_Y \ge \varepsilon
  \|\eta\|_Y \,\}$, where $P_R$ is the orthogonal projection onto
  $\overline{\Ran(K)}$.  Assume the hypotheses of \cref{thm_hanke_raus_conditional} and,
  in addition, that $\alpha \mapsto \varphi_{\alpha}(\lambda)$ is continuous on
  $(0,\infty)$ for each $\lambda \in (0,\kappa]$.  Then for every $x\dg \in X_{\nu,\rho}$
  with $\nu \in (0,\nu_0-1]$ and $\rho > 0$, and with $A(y^{\delta}) :=
  \arg\min_{\alpha \in (0,\kappa]}\Psi(\alpha)$ as in \cref{def_hanke_raus},
  \[
    \lim_{\delta \to 0}\ \sup\bigl\{\,\|R_{\alpha^{*}}y^{\delta} - x\dg\|_X :
      y^{\delta} \in N_{\varepsilon},\ \|y^{\delta} - Kx\dg\|_Y = \delta,\
      \alpha^{*} \in A(y^{\delta})\,\bigr\} = 0 ,
  \]
  and $A(y^{\delta}) \ne \emptyset$ for every such $y^{\delta}$, so the supremum is over
  a nonempty set whenever $\delta > 0$ admits admissible data.
\end{theorem}

\begin{proof}
  Write $y := Kx\dg$.  By \cref{lem_source_in_kerperp}, $x\dg \in \Ker(K)^{\perp}$, so
  $y \in \Ran(K)$ and $K\dg y = x\dg$ by \cref{lem_mp_equations}(iii); thus $x\dg$ is the
  minimum norm solution belonging to the exact datum $y$, as the statement intends.
  Let $C_F$, $C_{\nu+1}$, $C_r$ be as in \cref{thm_hanke_raus_conditional}, and let
  $\delta > 0$ and $y^{\delta} \in N_{\varepsilon}$ with $\|y^{\delta} - y\|_Y = \delta$
  be given; put $\eta := y^{\delta} - y$, so $\|\eta\|_Y = \delta$ and
  $\|(\Id - P_R)\eta\|_Y \ge \varepsilon\delta$.  Since $y \in \Ran(K) \subseteq
  \overline{\Ran(K)}$ we have $(\Id - P_R)y = 0$, hence
  \begin{equation}\label{eq_hr_noise}
    \|(\Id - P_R)y^{\delta}\|_Y \;=\; \|(\Id - P_R)\eta\|_Y \;\ge\; \varepsilon\delta .
  \end{equation}

  \emph{The residual is bounded below.}  For every $\alpha > 0$ we have
  $Kx^{\delta}_{\alpha} \in \Ran(K)$, so $(\Id - P_R)Kx^{\delta}_{\alpha} = 0$, and since
  $\Id - P_R$ is an orthogonal projection its norm is at most $1$
  (\cref{prop_proj_props}); therefore, by \eqref{eq_hr_noise},
  \begin{equation}\label{eq_hr_lower}
    \|Kx^{\delta}_{\alpha} - y^{\delta}\|_Y \;\ge\;
      \|(\Id - P_R)(Kx^{\delta}_{\alpha} - y^{\delta})\|_Y
      \;=\; \|(\Id - P_R)y^{\delta}\|_Y \;\ge\; \varepsilon\delta \;>\; 0 .
  \end{equation}
  In particular $\Psi(\alpha) \ge \varepsilon\delta\,\alpha^{-1/2}$ on $(0,\kappa]$, and
  $\delta^{*} := \|Kx^{\delta}_{\alpha^{*}} - y^{\delta}\|_Y \ge \varepsilon\delta > 0$
  for any minimiser $\alpha^{*}$, so $\delta/\delta^{*} \le \varepsilon^{-1}$.  This is
  precisely the hypothesis \eqref{eq_hr_lower} that
  \cref{thm_hanke_raus_conditional} needs and that
  \cref{def_hanke_raus} does not supply by itself; note also that
  \cref{thm_hanke_raus_conditional} was proved without any density assumption on
  $\Ran(K)$, so $N_{\varepsilon}$ is not forced to degenerate to $\{y\}$.

  \emph{Existence of the minimiser: $A(y^{\delta}) \ne \emptyset$.}  By
  \cref{lem_spectral_residual_form}(iii), $\alpha \mapsto \|Kx^{\delta}_{\alpha} -
  y^{\delta}\|_Y^{2}$ is the constant $\|(\Id-P_R)y^{\delta}\|_Y^{2}$ plus the series
  $\sum_n |r_{\alpha}(\sigma_n^{2})|^{2}|\langle u_n,y^{\delta}\rangle_Y|^{2}$.  Each
  summand is continuous in $\alpha$ by the assumed continuity of $\alpha \mapsto
  \varphi_{\alpha}(\lambda)$, and the tail beyond a finite $J \subseteq I$ is at most
  $C_r^{2}\sum_{n \in I \setminus J}|\langle u_n,y^{\delta}\rangle_Y|^{2}$ uniformly in
  $\alpha$ (\cref{lem_residual_constant}), which is arbitrarily small by
  \cref{thm_bessel}; exactly as in Step 2 of \cref{thm_discrepancy_order_optimal} this
  makes the sum, and hence $\Psi$, continuous on $(0,\kappa]$.  Putting $\beta :=
  \min\{\kappa, (\varepsilon\delta/(\Psi(\kappa)+1))^{2}\} > 0$, the bound
  $\Psi(\alpha) \ge \varepsilon\delta\alpha^{-1/2} > \Psi(\kappa)$ for $\alpha < \beta$
  shows that the infimum over $(0,\kappa]$ equals the infimum over the compact interval
  $[\beta,\kappa]$, where a continuous function attains it.  So $A(y^{\delta}) \ne
  \emptyset$, and every $\alpha^{*} \in A(y^{\delta})$ satisfies the estimates above.

  \emph{The degenerate case $y = 0$.}  Then $Kx\dg = 0$ and $x\dg \in \Ker(K) \cap
  \Ker(K)^{\perp}$, so $x\dg = 0$, and $\|y^{\delta}\|_Y = \delta$.  By
  \cref{lem_residual_constant}, $\|y^{\delta} - KR_{\kappa}y^{\delta}\|_Y \le
  C_r\|y^{\delta}\|_Y = C_r\delta$, so $\Psi(\kappa) \le C_r\delta\kappa^{-1/2}$.
  Minimality and \eqref{eq_hr_lower} give $\varepsilon\delta(\alpha^{*})^{-1/2} \le
  \Psi(\alpha^{*}) \le \Psi(\kappa) \le C_r\delta\kappa^{-1/2}$, i.e.\
  $(\alpha^{*})^{-1/2} \le C_r\kappa^{-1/2}\varepsilon^{-1}$.  Hence, by
  \cref{lem_filter_norm_bound} and (F),
  \[
    \|x^{\delta}_{\alpha^{*}} - x\dg\|_X = \|R_{\alpha^{*}}y^{\delta}\|_X
      \le \sqrt{C_{\varphi}C_F}\,(\alpha^{*})^{-1/2}\delta
      \le \frac{\sqrt{C_{\varphi}C_F}\,C_r}{\varepsilon\sqrt{\kappa}}\,\delta ,
  \]
  a bound independent of $y^{\delta}$ which tends to $0$ with $\delta$.

  \emph{The main case $y \ne 0$.}  Let $0 < \delta < \|y\|_Y$.  Then all hypotheses of
  \cref{thm_hanke_raus_conditional} hold, and $\delta/\delta^{*} \le \varepsilon^{-1}$,
  so
  \begin{equation}\label{eq_hr_conv_1}
    \|x^{\delta}_{\alpha^{*}} - x\dg\|_X \;\le\; c\,(1 + \varepsilon^{-1})\,
      \max\{\delta,\delta^{*}\}^{\nu/(\nu+1)}\rho^{1/(\nu+1)} .
  \end{equation}
  It remains to bound $\delta^{*}$ in terms of $\delta$ alone.  Exactly as in
  \cref{thm_hanke_raus_conditional}, $\alpha_{\delta} := (\delta/\rho)^{2/(\nu+1)}$ lies
  in $(0,\kappa]$ because $\delta < \|y\|_Y \le \kappa^{(\nu+1)/2}\rho$, and
  \cref{lem_omega_nu} at exponent $\nu+1$, with (Q$_{\nu+1}$) of
  \cref{def_qualification}, gives
  \[
    \Psi(\alpha_{\delta}) \;\le\;
      (C_{\nu+1} + C_r)\,\delta^{\nu/(\nu+1)}\rho^{1/(\nu+1)} .
  \]
  Since $\alpha^{*} \le \kappa$ and $\Psi(\alpha^{*}) \le \Psi(\alpha_{\delta})$,
  \[
    \delta^{*} \;=\; \sqrt{\alpha^{*}}\,\Psi(\alpha^{*}) \;\le\;
      \sqrt{\kappa}\,(C_{\nu+1} + C_r)\,\delta^{\nu/(\nu+1)}\rho^{1/(\nu+1)} .
  \]
  Thus $\max\{\delta,\delta^{*}\} \le h(\delta) := \max\{\delta,
  \sqrt{\kappa}(C_{\nu+1}+C_r)\delta^{\nu/(\nu+1)}\rho^{1/(\nu+1)}\}$, and $h(\delta) \to
  0$ as $\delta \to 0$ because $\nu/(\nu+1) > 0$.  The right-hand side of
  \eqref{eq_hr_conv_1} is therefore bounded by $c(1+\varepsilon^{-1})
  h(\delta)^{\nu/(\nu+1)}\rho^{1/(\nu+1)}$, a quantity independent of both $y^{\delta}$
  and $\alpha^{*}$.  It tends to $0$ as $\delta \to 0$, so taking the supremum over the
  admissible pairs $(y^{\delta},\alpha^{*})$ and letting $\delta \to 0$ gives the claim.
\end{proof}


\chapter{Tikhonov regularization}
\label{ch_16_tikhonov}

\noindent\emph{Background.} The variational functional
$J_{\alpha}(x) = \tfrac12\|Kx-y\|_Y^{2} + \tfrac{\alpha}{2}\|x\|_X^{2}$, whose unique
minimiser recovers the spectral filter $\varphi_{\alpha}(\lambda) = (\lambda+\alpha)^{-1}$
of the previous chapter, is Tikhonov's own regularization method \citep{Tikhonov1963},
developed systematically together with Arsenin in \citet{TikhonovArsenin1977}. That the
method saturates at qualification $2$ --- no source condition, however smooth, buys a
rate better than $\delta^{2/3}$, and insisting on one forces $x\dg = 0$ --- is the
converse and saturation result of \citet{Neubauer1997}.

\begin{assumptions}
  \label{asm_tikhonov_regularization}
  As in the previous chapter.  Throughout,
$\{\varphi_{\alpha}\}$ is the Tikhonov filter $\varphi_{\alpha}(\lambda) =
(\lambda+\alpha)^{-1}$ of \cref{def_standard_filters}(ii), and
$R_{\alpha} = \varphi_{\alpha}(K^{*}K)K^{*}$.  From \cref{def_value_functions} onwards
the datum $y^{\delta} \in Y$ is fixed and arbitrary.
\end{assumptions}

\begin{lemma}[Tikhonov has qualification $2$]
  \label{lem_tikhonov_qualification}
  The Tikhonov filter satisfies (F) with constant $C_F = 1$, and (Q$_{\nu}$) for every
  $\nu \in (0,2]$ with $C_{\nu} = 1$.  Its qualification is $\nu_0 = 2$.
\end{lemma}

\begin{proof}
  Throughout, $\varphi_{\alpha}(\lambda) = (\lambda+\alpha)^{-1}$ as in
  \cref{def_standard_filters}(ii), so that
  \[
    r_{\alpha}(\lambda) \;=\; 1 - \lambda\varphi_{\alpha}(\lambda)
      \;=\; \frac{\alpha}{\lambda + \alpha}, \qquad \lambda \in [0,\kappa],\ \alpha > 0 .
  \]

  \emph{Condition (F).}  For $\lambda \in (0,\kappa]$ and $\alpha > 0$ we have
  $\lambda + \alpha > \alpha > 0$, so $|\varphi_{\alpha}(\lambda)| = (\lambda +
  \alpha)^{-1} \le \alpha^{-1}$.  Hence (F) of \cref{def_qualification} holds with
  constant $C_F = 1$.

  \emph{An elementary inequality.}  Let $t \in (0,1]$ and $a, b > 0$; we claim
  $a^{t}b^{1-t} \le ta + (1-t)b$.  Dividing by $b$ and setting $s := a/b > 0$, this reads
  $s^{t} \le ts + 1 - t$.  Put $h(s) := ts + 1 - t - s^{t}$ on $(0,\infty)$; then
  $h(1) = 0$ and $h'(s) = t(1 - s^{t-1})$.  For $s \le 1$ one has $s^{t-1} \ge 1$ (the
  exponent $t-1$ is $\le 0$), so $h' \le 0$; for $s \ge 1$ one has $s^{t-1} \le 1$, so
  $h' \ge 0$.  Thus $h$ is nonincreasing on $(0,1]$ and nondecreasing on $[1,\infty)$,
  whence $h(s) \ge h(1) = 0$ for all $s > 0$, which is the claim.

  \emph{Condition (Q$_{\nu}$) for $\nu \in (0,2]$.}  Let $\alpha > 0$ and $\lambda \in
  (0,\kappa]$, and put $t := \nu/2 \in (0,1]$.  The inequality just proved with
  $a := \lambda$, $b := \alpha$ gives
  \[
    \lambda^{\nu/2}\alpha^{1-\nu/2} \;=\; \lambda^{t}\alpha^{1-t} \;\le\;
      t\lambda + (1-t)\alpha \;\le\; \lambda + \alpha ,
  \]
  and dividing by $\lambda + \alpha > 0$ and multiplying by $\alpha^{\nu/2} > 0$,
  \[
    \lambda^{\nu/2}\,|r_{\alpha}(\lambda)| \;=\;
      \frac{\lambda^{\nu/2}\alpha}{\lambda+\alpha} \;\le\; \alpha^{\nu/2} .
  \]
  Taking the supremum over $\lambda \in (0,\kappa]$ gives $\omega_{\nu}(\alpha) \le
  \alpha^{\nu/2}$, i.e.\ (Q$_{\nu}$) with $C_{\nu} = 1$.

  \emph{The qualification is exactly $2$.}  By the previous paragraph the set of $\nu >
  0$ satisfying (Q$_{\nu}$) contains $(0,2]$.  Suppose (Q$_{\nu}$) held for some
  $\nu > 2$, with constant $C_{\nu}$.  Evaluating $\omega_{\nu}$ at the single point
  $\lambda = \kappa$ gives $\omega_{\nu}(\alpha) \ge \kappa^{\nu/2}\alpha/(\kappa +
  \alpha)$, so for all $\alpha > 0$
  \[
    \frac{\kappa^{\nu/2}\alpha}{\kappa + \alpha} \;\le\; C_{\nu}\alpha^{\nu/2} ,
    \qquad\text{i.e.}\qquad
    \alpha^{1 - \nu/2} \;\le\; C_{\nu}\,\kappa^{-\nu/2}(\kappa + \alpha) .
  \]
  As $\alpha \to 0^{+}$ the left-hand side tends to $+\infty$, because $1 - \nu/2 < 0$,
  while the right-hand side tends to $C_{\nu}\kappa^{1-\nu/2} < \infty$ --- a
  contradiction.  Hence the set of admissible $\nu$ is contained in $(0,2]$ and contains
  it, so its supremum is $\nu_0 = 2$ (\cref{def_qualification}).
\end{proof}


\begin{corollary}[A priori rule for Tikhonov]
  \label{cor_tikhonov_apriori}
  For every $\nu \in (0,2]$, Tikhonov regularization with an a priori rule satisfying
  $c(\delta/\rho)^{2/(\nu+1)} \le \alpha(\delta) \le C(\delta/\rho)^{2/(\nu+1)}$ is an
  order optimal regularization method for that $\nu$ and $\rho$.  In particular, for
  $\alpha \sim \delta^{2/3}$ there is $c > 0$ with
  $\|x^{\delta}_{\alpha(\delta)} - x\dg\|_X \le c\,\delta^{2/3}$ for all
  $x\dg \in X_{2,\rho} \subseteq \Ran(K^{*}K)$ and $y^{\delta} \in B_{\delta}(Kx\dg)$.
\end{corollary}

\begin{proof}
  By \cref{def_standard_filters}(ii) the Tikhonov filter is a regularizing filter with
  $C_{\varphi} = 1$, and by \cref{lem_tikhonov_qualification} it satisfies (F) with
  constant $C_F = 1$ and (Q$_{\nu}$) with $C_{\nu} = 1$ for every $\nu \in (0,2]$.  With
  the a priori rule's own constants $c, C > 0$, \cref{thm_apriori_order_optimal} applies
  for every $\nu \in (0,2]$ and every $\rho > 0$ and yields, for all $\delta > 0$, all
  $x\dg \in X_{\nu,\rho}$ and all $y^{\delta} \in B_{\delta}(Kx\dg)$,
  \[
    \|x^{\delta}_{\alpha(\delta)} - x\dg\|_X \;\le\;
      \bigl(c^{-1/2} + C^{\nu/2}\bigr)\,
      \delta^{\nu/(\nu+1)}\rho^{1/(\nu+1)} ,
  \]
  together with order optimality for this $\nu$ and this $\rho$
  (\cref{def_order_optimal}) and convergence of the method.

  For the final clause take $\nu = 2$, so that $2/(\nu+1) = 2/3$ --- this is the meaning
  of $\alpha \sim \delta^{2/3}$ --- and $\nu/(\nu+1) = 2/3$ as well.  With $\rho > 0$
  fixed, the displayed estimate becomes
  \[
    \|x^{\delta}_{\alpha(\delta)} - x\dg\|_X \;\le\;
      \bigl(c^{-1/2} + C\bigr)\rho^{1/3}\,\delta^{2/3}
      \;=:\; c'\,\delta^{2/3}
  \]
  for all $x\dg \in X_{2,\rho}$ and $y^{\delta} \in B_{\delta}(Kx\dg)$.  Finally
  $X_{2,\rho} = \{|K|^{2}w : \|w\|_X \le \rho\} \subseteq \Ran(K^{*}K)$ by
  \cref{def_source_set}, since $|K|^{2} = (K^{*}K)^{1} = K^{*}K$ by \cref{def_frac_power}
  and \cref{lem_fun_calc_basic}.
\end{proof}




\begin{corollary}[Discrepancy principle for Tikhonov]
  \label{cor_tikhonov_discrepancy}
  Assume $\Ran(K)$ is dense in $Y$ and $\dim\Ran(K) = \infty$, and let $\tau > 1$.  Then
  the discrepancy principle defines an a posteriori rule $\alpha_{\tau}$ for Tikhonov
  regularization; $(R_{\alpha},\alpha_{\tau})$ is a convergent regularization method; and
  for every $\nu \in (0,1]$ it satisfies the rate bound of
  \cref{thm_discrepancy_order_optimal}(i).  In particular there is $c' > 0$, depending only on $\tau$ and on $K$, with
  \[
    \|x^{\delta}_{\alpha_{\tau}(\delta,y^{\delta})} - x\dg\|_X
      \;\le\; c'\,\rho^{1/2}\,\delta^{1/2}
  \]
  for every $\rho > 0$, all $x\dg \in X_{1,\rho} \subseteq \Ran(K^{*})$, all $\delta > 0$
  and all
  $y^{\delta} \in B_{\delta}(Kx\dg)$ \emph{with} $\|y^{\delta}\|_Y > \tau\delta$.
\end{corollary}

\begin{proof}
  We check the hypotheses of \cref{thm_discrepancy_order_optimal} for the Tikhonov
  filter $\varphi_{\alpha}(\lambda) = (\lambda+\alpha)^{-1}$ of
  \cref{def_standard_filters}(ii).  Density of $\Ran(K)$ and $\dim\Ran(K) = \infty$ are
  assumed.  By
  \cref{lem_tikhonov_qualification} the qualification is $\nu_0 = 2 > 1$, and (F) as well
  as (Q$_{\nu}$) hold for every $\nu \in (0,\nu_0] = (0,2]$
  (\cref{def_qualification}).  For each fixed $\lambda \in [0,\kappa]$ the map
  $\alpha \mapsto \varphi_{\alpha}(\lambda) = (\lambda+\alpha)^{-1}$ is continuous on
  $(0,\infty)$.  Finally, $r_{\alpha}(\lambda) = \alpha/(\lambda+\alpha) \in (0,1)$ for
  $\lambda \in (0,\kappa]$, so $C_r \le 1$; and $C_r \ge 1$ by
  \cref{lem_residual_constant}, so $C_r = 1$ and the condition $\tau > C_r$ is exactly
  $\tau > 1$.  \cref{thm_discrepancy_order_optimal} therefore applies: it produces the
  rule $\alpha_{\tau}$ of \cref{def_discrepancy_principle}, its part (ii) gives that
  $(R_{\alpha},\alpha_{\tau})$ is a convergent regularization method, and its part (i)
  gives, for every $\nu \in (0,\nu_0-1] = (0,1]$, a constant $c = c(\nu,\tau)$
  independent of $\rho$ with the stated rate.

  For the final clause take $\nu = 1$, so $\nu/(\nu+1) = 1/2$, and put $c' := c(1,\tau)$.
  \cref{thm_discrepancy_order_optimal}(i) then gives
  $\|x^{\delta}_{\alpha_{\tau}(\delta,y^{\delta})} - x\dg\|_X \le
  c'\,\rho^{1/2}\,\delta^{1/2}$ for every $\rho > 0$, every $x\dg \in X_{1,\rho}$, every
  $\delta > 0$ and every $y^{\delta} \in B_{\delta}(Kx\dg)$ with
  $\|y^{\delta}\|_Y > \tau\delta$ --- exactly the data on which
  \cref{def_discrepancy_principle} governs the rule.

  It remains to note that $X_{1,\rho} \subseteq \Ran(K^{*})$
  (\cref{def_source_set}).  Let $x = |K|w$ with $\|w\|_X \le \rho$.  By
  \cref{def_fun_calc} and \cref{def_frac_power},
  \[
    |K|w \;=\; \sum_{n \in I} \sigma_n\,\langle v_n,w\rangle_X\, v_n .
  \]  The coefficients $\langle v_n,w\rangle_X$ are
  square summable by Bessel's inequality (\cref{thm_bessel}), so $u := \sum_{n \in I}
  \langle v_n,w\rangle_X u_n$ exists in $Y$ (\cref{lem_riesz_fischer}); applying the
  bounded operator $K^{*}$ term by term (\cref{lem_operator_through_series}) and using
  $K^{*}u_n = \sigma_n v_n$ (\cref{def_singular_system}) gives $K^{*}u = \sum_n
  \sigma_n\langle v_n,w\rangle_X v_n = x$.
\end{proof}




\begin{lemma}[Tikhonov solves a regularised normal equation]
  \label{lem_tikhonov_normal}
  Let $y \in Y$ and $\alpha > 0$.  For $x \in X$ one has $x = R_{\alpha}y$ if and only if
  $(K^{*}K + \alpha\,\Id)x = K^{*}y$.
\end{lemma}

\begin{proof}
  \emph{$K^{*}K + \alpha\Id$ is injective.}  For $x \in X$,
  \[
    \alpha\|x\|_X^{2} \;\le\; \|Kx\|_Y^{2} + \alpha\|x\|_X^{2}
      \;=\; \langle x, (K^{*}K + \alpha\Id)x\rangle_X
      \;\le\; \|x\|_X\,\|(K^{*}K+\alpha\Id)x\|_X ,
  \]
  using the definition of the adjoint (\cref{def_adjoint}) for the equality and
  Cauchy--Schwarz (\cref{thm_cauchy_schwarz}) for the second inequality.  Hence
  $(K^{*}K+\alpha\Id)x = 0$ forces $\|x\|_X = 0$, i.e.\ $x = 0$.  Consequently the
  equation $(K^{*}K + \alpha\Id)x = K^{*}y$ has at most one solution.

  \emph{$R_{\alpha}y$ is a solution.}  By \cref{def_spectral_reg} together with
  \cref{def_standard_filters}(ii) we have $R_{\alpha}y = \varphi_{\alpha}(K^{*}K)K^{*}y$,
  where $\varphi_{\alpha}(\lambda) = (\lambda+\alpha)^{-1}$.  Let $\iota(\lambda) := \lambda$
  and $\mathbf{1}(\lambda) := 1$ on $[0,\kappa]$, so that $\iota(K^{*}K) = K^{*}K$ and
  $\mathbf{1}(K^{*}K) = \Id$ by \cref{lem_fun_calc_basic}.  By linearity and
  multiplicativity of the functional calculus (\cref{lem_fun_calc_hom}),
  \[
    (K^{*}K + \alpha\Id)\,\varphi_{\alpha}(K^{*}K)
      \;=\; \bigl((\iota + \alpha\mathbf{1})\varphi_{\alpha}\bigr)(K^{*}K) ,
  \]
  and $(\lambda + \alpha)\varphi_{\alpha}(\lambda) = 1$ for \emph{every} $\lambda \in
  [0,\kappa]$, including $\lambda = 0$, where it reads $\alpha\cdot\alpha^{-1} = 1$.  So
  the multiplier is the constant function $\mathbf{1}$ and the operator on the right is
  $\Id$.  Applying both sides to $K^{*}y$ gives $(K^{*}K+\alpha\Id)R_{\alpha}y = K^{*}y$.

  Combining the two paragraphs: $x = R_{\alpha}y$ implies $(K^{*}K+\alpha\Id)x = K^{*}y$,
  and conversely any solution of that equation equals the solution $R_{\alpha}y$ by
  uniqueness.
\end{proof}


\begin{theorem}[Tikhonov saturates at $\delta^{2/3}$]
  \label{thm_tikhonov_saturation}
  Let $K \in \Kk(X,Y)$ have infinite-dimensional range and let $y \in \Ran(K)$.  If there
  is an a priori rule $\alpha$ with $\alpha(\delta) \to 0$ and
  \[
    \lim_{\delta \to 0}\ \sup_{y^{\delta} \in B_{\delta}(y)}
      \|x^{\delta}_{\alpha(\delta)} - x\dg\|_X\,\delta^{-2/3} = 0 ,
  \]
  then $x\dg = 0$.  In particular the qualification of the Tikhonov filter is exactly $2$,
  which is \cref{lem_tikhonov_qualification}.
\end{theorem}

\begin{proof}
  Throughout, $R_{\alpha}$ is the \emph{Tikhonov} family fixed at the head of the
  chapter, $R_{\alpha} = \varphi_{\alpha}(K^{*}K)K^{*}$ with $\varphi_{\alpha}(\lambda) =
  (\lambda+\alpha)^{-1}$; the binding is used twice below, through
  \cref{lem_tikhonov_normal} and through the explicit action of $R_{\alpha}$ on singular
  vectors, and the statement is false without it.

  Write $E(\delta) := \sup_{y^{\delta} \in B_{\delta}(y)}\|x^{\delta}_{\alpha(\delta)} -
  x\dg\|_X$, so the hypothesis is $E(\delta)\delta^{-2/3} \to 0$ as $\delta \to 0$; in
  particular $E(\delta) < \infty$ for all small $\delta$.  Suppose, for contradiction,
  that $x\dg \ne 0$.

  \emph{Step 1: $\alpha(\delta)\delta^{-2/3} \to 0$.}  Let $\delta > 0$ and abbreviate
  $\alpha := \alpha(\delta)$.  Since $y \in \Ran(K) \subseteq \Dom(K\dg)$ and $x\dg =
  K\dg y$ is the minimum norm solution (\cref{thm_mp_minnorm}(ii)), hence a least squares
  solution (\cref{def_lsq_minnorm}), \cref{cor_normal_equation} gives $K^{*}Kx\dg =
  K^{*}y$ and therefore $(K^{*}K + \alpha\Id)x\dg = K^{*}y + \alpha x\dg$.  Subtracting
  this from $(K^{*}K + \alpha\Id)x^{\delta}_{\alpha} = K^{*}y^{\delta}$
  (\cref{lem_tikhonov_normal}) yields
  \[
    (K^{*}K + \alpha\Id)\bigl(x^{\delta}_{\alpha} - x\dg\bigr)
      \;=\; K^{*}(y^{\delta} - y) - \alpha\,x\dg ,
  \]
  whence, by the triangle inequality, $\|K^{*}\|_{\Ll(Y,X)} = \|K\|_{\Ll(X,Y)} =
  \sqrt{\kappa}$ (\cref{prop_adjoint_props}) and $\|K^{*}K + \alpha\Id\|_{\Ll(X,X)} \le
  \kappa + \alpha$,
  \[
    \alpha\,\|x\dg\|_X \;\le\; \sqrt{\kappa}\,\|y^{\delta}-y\|_Y
      + (\kappa+\alpha)\,\|x^{\delta}_{\alpha} - x\dg\|_X
      \;\le\; \sqrt{\kappa}\,\delta + (\kappa+\alpha)E(\delta) ,
  \]
  taking any $y^{\delta} \in B_{\delta}(y)$, for instance $y^{\delta} := y$.  Dividing by
  $\delta^{2/3}\|x\dg\|_X > 0$,
  \[
    \alpha(\delta)\,\delta^{-2/3} \;\le\; \|x\dg\|_X^{-1}\Bigl(
      \sqrt{\kappa}\,\delta^{1/3} + (\kappa + \alpha(\delta))\,E(\delta)\delta^{-2/3}
      \Bigr) \;\longrightarrow\; 0
  \]
  as $\delta \to 0$, since $\alpha(\delta) \to 0$ keeps the factor $\kappa+\alpha(\delta)$
  bounded and $E(\delta)\delta^{-2/3} \to 0$ by hypothesis.

  \emph{Step 2: a contradiction.}  Since $\dim\Ran(K) = \infty$ the index set $I$ is all
  of $\mathbb{N}$ and $\{\sigma_n\}$ is a null sequence of strictly positive numbers
  (\cref{def_singular_system}).  Put $\delta_n := \sigma_n^{3}$ and $y_n := y +
  \delta_n u_n$, so $\|y_n - y\|_Y = \delta_n \to 0$ and $y_n \in B_{\delta_n}(y)$; write
  $\alpha_n := \alpha(\delta_n)$.  By \cref{def_spectral_reg} and
  \cref{def_standard_filters}(ii) together with orthonormality of $\{u_m\}$,
  \[
    R_{\alpha_n}(y_n - y) \;=\; \delta_n \sum_{m \in I}
      \frac{\sigma_m}{\sigma_m^{2} + \alpha_n}\,\langle u_m, u_n\rangle_Y\,v_m
      \;=\; \frac{\delta_n\sigma_n}{\sigma_n^{2} + \alpha_n}\,v_n ,
  \]
  a vector of norm $\delta_n\sigma_n/(\sigma_n^{2}+\alpha_n)$ since $\|v_n\|_X = 1$.  By
  linearity of $R_{\alpha_n}$,
  \[
    R_{\alpha_n}(y_n - y) \;=\; \bigl(x^{\delta_n}_{\alpha_n} - x\dg\bigr)
      - \bigl(x_{\alpha_n} - x\dg\bigr) ,
  \]
  and both $y_n$ and $y$ lie in $B_{\delta_n}(y)$, so each of the two brackets has norm
  at most $E(\delta_n)$.  Hence
  \[
    \frac{\delta_n\sigma_n}{\sigma_n^{2} + \alpha_n} \;\le\; 2E(\delta_n) ,
    \qquad\text{i.e.}\qquad
    \frac{\sigma_n\,\delta_n^{1/3}}{\sigma_n^{2} + \alpha_n}
      \;\le\; 2E(\delta_n)\delta_n^{-2/3} \;\longrightarrow\; 0 .
  \]
  But $\delta_n^{1/3} = \sigma_n$ and $\sigma_n^{2} = \delta_n^{2/3}$, so the left-hand
  side equals
  \[
    \frac{\sigma_n^{2}}{\sigma_n^{2}+\alpha_n}
      \;=\; \frac{1}{1 + \alpha_n\delta_n^{-2/3}} \;\longrightarrow\; 1
  \]
  by Step 1 --- a contradiction.  Therefore $x\dg = 0$.

  The final clause is \cref{lem_tikhonov_qualification}, where $\nu_0 = 2$ was
  established directly from the filter.
\end{proof}


\begin{theorem}[Variational characterisation]
  \label{thm_tikhonov_variational}
  Let $y \in Y$ and $\alpha > 0$.  Then $x_{\alpha} := R_{\alpha}y$ is the unique
  minimiser over $X$ of the \emph{Tikhonov functional}
  \[
    J_{\alpha}(x) := \tfrac12\|Kx - y\|_Y^{2} + \tfrac{\alpha}{2}\|x\|_X^{2} .
  \]
\end{theorem}

\begin{proof}
  Write $x_{\alpha} := R_{\alpha}y$ and let $x \in X$ be arbitrary.  Put $z := x -
  x_{\alpha}$.  Expanding both squared norms in $J_{\alpha}(x)$ around $x_{\alpha}$,
  \[
    \|Kx - y\|_Y^{2} = \|Kz\|_Y^{2} + 2\,\mathrm{Re}\,\langle Kz, Kx_{\alpha}-y\rangle_Y
      + \|Kx_{\alpha}-y\|_Y^{2},
  \]
  \[
    \|x\|_X^{2} = \|z\|_X^{2} + 2\,\mathrm{Re}\,\langle z, x_{\alpha}\rangle_X
      + \|x_{\alpha}\|_X^{2},
  \]
  where we used $\|a+b\|^{2} = \|a\|^{2} + 2\,\mathrm{Re}\langle a,b\rangle + \|b\|^{2}$,
  itself immediate from sesquilinearity and $\langle b,a\rangle = \overline{\langle a,b
  \rangle}$.  Subtracting, and using the definition of the adjoint (\cref{def_adjoint}) in
  the form $\langle Kz, Kx_{\alpha}-y\rangle_Y = \langle z, K^{*}(Kx_{\alpha}-y)
  \rangle_X$,
  \[
    J_{\alpha}(x) - J_{\alpha}(x_{\alpha}) \;=\; \tfrac12\|Kz\|_Y^{2}
      + \tfrac{\alpha}{2}\|z\|_X^{2}
      + \mathrm{Re}\,\bigl\langle z,\, K^{*}(Kx_{\alpha}-y) + \alpha x_{\alpha}
        \bigr\rangle_X .
  \]
  By \cref{lem_tikhonov_normal}, $(K^{*}K+\alpha\Id)x_{\alpha} = K^{*}y$, i.e.\
  $K^{*}(Kx_{\alpha}-y) + \alpha x_{\alpha} = 0$, so the last term vanishes and
  \[
    J_{\alpha}(x) - J_{\alpha}(x_{\alpha}) \;=\; \tfrac12\|Kz\|_Y^{2}
      + \tfrac{\alpha}{2}\|z\|_X^{2} \;\ge\; \tfrac{\alpha}{2}\|x - x_{\alpha}\|_X^{2}
      \;\ge\; 0 .
  \]
  Hence $x_{\alpha}$ minimises $J_{\alpha}$, and since $\alpha > 0$ the right-hand side
  is strictly positive whenever $x \ne x_{\alpha}$; so the minimiser is unique.
\end{proof}


\begin{definition}[Value functions]
  \label{def_value_functions}
  For fixed $y^{\delta} \in Y$ put, for $\alpha > 0$,
  \[
    f(\alpha) := \tfrac12\|Kx^{\delta}_{\alpha} - y^{\delta}\|_Y^{2},
    \qquad
    g(\alpha) := \tfrac12\|x^{\delta}_{\alpha}\|_X^{2},
  \]
  \[
    j(\alpha) := J_{\alpha}(x^{\delta}_{\alpha}) = f(\alpha) + \alpha g(\alpha) .
  \]
\end{definition}


\begin{lemma}[Monotonicity of the value functions]
  \label{lem_value_monotone}
  For all $\alpha_1, \alpha_2 > 0$,
  \[
    (f(\alpha_1) - f(\alpha_2))(\alpha_1 - \alpha_2) \ge 0,
    \qquad
    (g(\alpha_1) - g(\alpha_2))(\alpha_1 - \alpha_2) \le 0 .
  \]
  That is, the residual increases and the solution norm decreases with $\alpha$.
\end{lemma}

\begin{proof}
  By \cref{thm_tikhonov_variational}, $x^{\delta}_{\alpha}$ is the minimiser of
  $J_{\alpha}$ over $X$, and by \cref{def_value_functions} $J_{\alpha}(x) =
  \tfrac12\|Kx-y^{\delta}\|_Y^{2} + \tfrac{\alpha}{2}\|x\|_X^{2}$ satisfies
  $J_{\alpha}(x^{\delta}_{\beta}) = f(\beta) + \alpha g(\beta)$ for all $\alpha,\beta>0$.
  Testing the minimality of $x^{\delta}_{\alpha_1}$ for $J_{\alpha_1}$ against
  $x^{\delta}_{\alpha_2}$, and conversely,
  \begin{equation}\label{eq_val_min_1}
    f(\alpha_1) + \alpha_1 g(\alpha_1) \;\le\; f(\alpha_2) + \alpha_1 g(\alpha_2),
  \end{equation}
  \begin{equation}\label{eq_val_min_2}
    f(\alpha_2) + \alpha_2 g(\alpha_2) \;\le\; f(\alpha_1) + \alpha_2 g(\alpha_1).
  \end{equation}

  \emph{Monotonicity of $g$.}  Adding \eqref{eq_val_min_1} and \eqref{eq_val_min_2} and
  cancelling $f(\alpha_1) + f(\alpha_2)$ gives $\alpha_1 g(\alpha_1) + \alpha_2
  g(\alpha_2) \le \alpha_1 g(\alpha_2) + \alpha_2 g(\alpha_1)$, i.e.\
  \[
    (\alpha_1 - \alpha_2)\bigl(g(\alpha_1) - g(\alpha_2)\bigr) \;\le\; 0 .
  \]

  \emph{Monotonicity of $f$.}  Rearranging \eqref{eq_val_min_1} and
  \eqref{eq_val_min_2} separately gives the sandwich
  \[
    \alpha_2\bigl(g(\alpha_2) - g(\alpha_1)\bigr) \;\le\; f(\alpha_1) - f(\alpha_2)
      \;\le\; \alpha_1\bigl(g(\alpha_2) - g(\alpha_1)\bigr) .
  \]
  If $\alpha_1 > \alpha_2$ then $g(\alpha_2) - g(\alpha_1) \ge 0$ by the previous step,
  so the left inequality gives $f(\alpha_1) - f(\alpha_2) \ge \alpha_2 \cdot 0 = 0$ and
  hence $(f(\alpha_1)-f(\alpha_2))(\alpha_1-\alpha_2) \ge 0$.  If $\alpha_1 < \alpha_2$
  then $g(\alpha_2) - g(\alpha_1) \le 0$, so the right inequality gives $f(\alpha_1) -
  f(\alpha_2) \le 0$ and again the product is nonnegative.  For $\alpha_1 = \alpha_2$ the
  product is $0$.
\end{proof}


\begin{lemma}[The value function is concave]
  \label{lem_value_concave}
  $j$ is concave and nondecreasing on $(0,\infty)$.  Consequently the one-sided
  derivatives $D^{+}j(\alpha)$ and $D^{-}j(\alpha)$ exist at every $\alpha > 0$ and
  satisfy $D^{+}j(\alpha) \le D^{-}j(\alpha)$.
\end{lemma}

\begin{proof}
  By \cref{thm_tikhonov_variational} and \cref{def_value_functions},
  \begin{equation}\label{eq_j_as_inf}
    j(\alpha) \;=\; J_{\alpha}(x^{\delta}_{\alpha}) \;=\; \min_{x \in X}
      \Bigl[\tfrac12\|Kx - y^{\delta}\|_Y^{2} + \tfrac{\alpha}{2}\|x\|_X^{2}\Bigr]
      \qquad (\alpha > 0),
  \end{equation}
  so for each fixed $x$ the map $\alpha \mapsto J_{\alpha}(x)$ is affine.

  \emph{Concavity.}  Let $\alpha_1,\alpha_2 > 0$, $t \in [0,1]$ and $\alpha :=
  t\alpha_1 + (1-t)\alpha_2$.  For every $x \in X$, affineness gives $J_{\alpha}(x) =
  tJ_{\alpha_1}(x) + (1-t)J_{\alpha_2}(x) \ge tj(\alpha_1) + (1-t)j(\alpha_2)$ by
  \eqref{eq_j_as_inf}.  Taking $x := x^{\delta}_{\alpha}$ gives $j(\alpha) \ge
  tj(\alpha_1) + (1-t)j(\alpha_2)$.

  \emph{Monotonicity.}  If $\alpha_1 \le \alpha_2$ then $J_{\alpha_1}(x) \le
  J_{\alpha_2}(x)$ for every $x$, so $j(\alpha_1) \le J_{\alpha_1}(x^{\delta}_{\alpha_2})
  \le J_{\alpha_2}(x^{\delta}_{\alpha_2}) = j(\alpha_2)$.

  \emph{One-sided derivatives.}  Fix $\alpha > 0$ and set, for $t \ne 0$ with $\alpha + t
  > 0$, $q(t) := \bigl(j(\alpha+t) - j(\alpha)\bigr)/t$.  We claim $q$ is nonincreasing,
  in the sense that $q(t_1) \ge q(t_2)$ whenever $t_1 < t_2$ (both admissible, both
  nonzero).  Indeed, if $0 < t_1 < t_2$ then $\alpha + t_1 = (1 - t_1/t_2)\alpha +
  (t_1/t_2)(\alpha+t_2)$, so concavity gives $j(\alpha+t_1) \ge (1-t_1/t_2)j(\alpha) +
  (t_1/t_2)j(\alpha+t_2)$, i.e.\ $j(\alpha+t_1) - j(\alpha) \ge (t_1/t_2)
  (j(\alpha+t_2)-j(\alpha))$; dividing by $t_1 > 0$ gives $q(t_1) \ge q(t_2)$.  The case
  $t_1 < t_2 < 0$ is symmetric: $\alpha + t_2 = (1-t_2/t_1)\alpha + (t_2/t_1)
  (\alpha+t_1)$ with $t_2/t_1 \in (0,1)$, and dividing the resulting inequality by
  $t_2 < 0$ reverses it to $q(t_1) \ge q(t_2)$.  Finally, if $t_1 < 0 < t_2$ then
  $\alpha = \frac{t_2}{t_2-t_1}(\alpha+t_1) + \frac{-t_1}{t_2-t_1}(\alpha+t_2)$ is a
  convex combination, so
  \[
    j(\alpha) \;\ge\; \frac{t_2}{t_2-t_1}j(\alpha+t_1)
      + \frac{-t_1}{t_2-t_1}j(\alpha+t_2) ,
  \]
  and multiplying by $(t_2-t_1)>0$ and rearranging gives $t_2(j(\alpha+t_1)-j(\alpha))
  \le t_1(j(\alpha+t_2)-j(\alpha))$, i.e.\ (dividing by $t_1t_2 < 0$) $q(t_1) \ge
  q(t_2)$.

  Consequently $q$ is nonincreasing on $(0,\infty)$ and bounded \emph{above} there by
  $q(t_0)$ for any fixed $t_0 \in (-\alpha,0)$ --- since $t_0 < t$ for every $t > 0$ ---
  so $D^{+}j(\alpha) = \lim_{t \to 0^{+}}q(t) = \sup_{t>0}q(t)$ exists and is finite;
  likewise $q$ is nonincreasing on $(-\alpha,0)$ and bounded \emph{below} there by
  $q(t_1)$ for any fixed $t_1 > 0$, so $D^{-}j(\alpha) = \lim_{t \to
  0^{-}}q(t) = \inf_{-\alpha<t<0}q(t)$ exists and is finite.  Since $q(t) \ge q(s)$ for
  $t < 0 < s$, passing to the two limits gives $D^{-}j(\alpha) \ge D^{+}j(\alpha)$.
\end{proof}


\begin{lemma}[Bounds on the one-sided derivatives]
  \label{lem_value_derivative_bounds}
  For all $\alpha > 0$,
  \[
    D^{+}j(\alpha) \le g(\alpha) \le D^{-}j(\alpha),
    \qquad
    j(\alpha) - \alpha D^{-}j(\alpha) \le f(\alpha) \le j(\alpha) - \alpha D^{+}j(\alpha) .
  \]
\end{lemma}

\begin{proof}
  Fix $\alpha > 0$.  For any $\tilde{\alpha} > 0$, minimality of $x^{\delta}
  _{\tilde{\alpha}}$ for $J_{\tilde{\alpha}}$ (\cref{thm_tikhonov_variational}) tested
  against $x^{\delta}_{\alpha}$ gives, with \cref{def_value_functions},
  \[
    j(\tilde{\alpha}) \;=\; f(\tilde{\alpha}) + \tilde{\alpha}g(\tilde{\alpha})
      \;\le\; f(\alpha) + \tilde{\alpha}g(\alpha) ,
  \]
  whence, subtracting from $j(\alpha) = f(\alpha) + \alpha g(\alpha)$,
  \begin{equation}\label{eq_j_subgrad}
    j(\alpha) - j(\tilde{\alpha}) \;\ge\; (\alpha - \tilde{\alpha})\,g(\alpha) .
  \end{equation}
  Put $\tilde{\alpha} := \alpha + t$.  For $t > 0$, \eqref{eq_j_subgrad} reads
  $j(\alpha) - j(\alpha+t) \ge -t\,g(\alpha)$; dividing by $t > 0$ gives
  $\bigl(j(\alpha+t)-j(\alpha)\bigr)/t \le g(\alpha)$, and letting $t \to 0^{+}$ --- the
  limit exists by \cref{lem_value_concave} --- gives $D^{+}j(\alpha) \le g(\alpha)$.  For
  $-\alpha < t < 0$ the same inequality divided by $t < 0$ reverses to
  $\bigl(j(\alpha+t)-j(\alpha)\bigr)/t \ge g(\alpha)$, and letting $t \to 0^{-}$ gives
  $D^{-}j(\alpha) \ge g(\alpha)$.  This is the first pair.

  For the second pair, $f(\alpha) = j(\alpha) - \alpha g(\alpha)$ by
  \cref{def_value_functions}, and $\alpha > 0$, so multiplying the first pair by
  $-\alpha$ and adding $j(\alpha)$ gives
  \[
    j(\alpha) - \alpha D^{-}j(\alpha) \;\le\; j(\alpha) - \alpha g(\alpha) = f(\alpha)
      \;\le\; j(\alpha) - \alpha D^{+}j(\alpha) . \qedhere
  \]
\end{proof}


\begin{corollary}[Differentiability of the value function]
  \label{cor_value_differentiable}
  For almost every $\alpha > 0$, $j$ is differentiable with $j'(\alpha) = g(\alpha)$.
\end{corollary}

\begin{proof}
  By \cref{lem_value_concave} the one-sided derivatives $D^{\pm}j(\alpha)$ exist and are
  finite at every $\alpha > 0$, with $D^{+}j(\alpha) \le D^{-}j(\alpha)$.  Let
  \[
    N \;:=\; \{\,\alpha > 0 : D^{+}j(\alpha) < D^{-}j(\alpha)\,\} .
  \]
  For $\alpha \notin N$ the two one-sided limits of the difference quotient coincide, so
  $j$ is differentiable at $\alpha$; and then \cref{lem_value_derivative_bounds}, which
  gives $D^{+}j(\alpha) \le g(\alpha) \le D^{-}j(\alpha)$ with $g$ as in
  \cref{def_value_functions}, forces $j'(\alpha) = D^{+}j(\alpha) = g(\alpha)$.

  It therefore suffices to show that $N$ is countable, hence a Lebesgue null set.  Let
  $0 < \alpha < \beta$.  With $q_{\alpha}(t) := (j(\alpha+t)-j(\alpha))/t$ as in
  \cref{lem_value_concave}, that lemma gives $D^{+}j(\alpha) = \sup_{t>0}q_{\alpha}(t)
  \ge q_{\alpha}(\beta-\alpha)$ and $D^{-}j(\beta) = \inf_{-\beta<t<0}q_{\beta}(t) \le
  q_{\beta}(\alpha-\beta)$.  Since
  \[
    q_{\alpha}(\beta-\alpha) \;=\; \frac{j(\beta)-j(\alpha)}{\beta-\alpha}
      \;=\; \frac{j(\alpha)-j(\beta)}{\alpha-\beta} \;=\; q_{\beta}(\alpha-\beta),
  \]
  we conclude $D^{+}j(\alpha) \ge D^{-}j(\beta)$ whenever $\alpha < \beta$.  Hence the
  nonempty open intervals $I_{\alpha} := \bigl(D^{+}j(\alpha), D^{-}j(\alpha)\bigr)$,
  $\alpha \in N$, are pairwise disjoint: for $\alpha < \beta$ in $N$, every point of
  $I_{\beta}$ is $< D^{-}j(\beta) \le D^{+}j(\alpha)$, while every point of $I_{\alpha}$
  is $> D^{+}j(\alpha)$.  Choosing a rational number in each $I_{\alpha}$ defines an
  injection of $N$ into $\mathbb{Q}$, so $N$ is countable.  A countable subset of
  $(0,\infty)$ is a Lebesgue null set, so $j$ is differentiable with $j' = g$ at almost
  every $\alpha > 0$ --- indeed at every $\alpha$ outside a countable set.
\end{proof}


\chapter{Landweber regularization}
\label{ch_17_landweber}

\noindent\emph{Background.} Landweber iteration --- the fixed-point recursion
$x_m = x_{m-1} + \omega K^{*}(y - Kx_{m-1})$ for the normal equation $K^{*}Kx = K^{*}y$
--- is due to \citet{Landweber1951}, who introduced it for Fredholm integral equations of
the first kind; as a regularization method it plays the role of an iterative alternative
to Tikhonov regularization, with the stopping index $m$ acting as $1/\alpha$ and,
crucially, no saturation at qualification $2$.

\begin{assumptions}
  \label{asm_landweber_regularization}
  $X$, $Y$ are Hilbert spaces, $K \in \Kk(X,Y)$ is
compact and nonzero, $\kappa = \|K\|^{2}$, and the relaxation parameter satisfies
$\omega \in (0,\kappa^{-1})$.  The initial iterate is $x_0 = 0$ throughout.
The stopping index $m$ plays the role of the regularization parameter, with
$\alpha = 1/m$.  All three of these are standing conventions.
\end{assumptions}

\begin{definition}[Landweber iteration]
  \label{def_landweber}
  For $y \in Y$ define $x_0 := 0$ and
  \[
    x_m := x_{m-1} + \omega K^{*}(y - Kx_{m-1}), \qquad m \in \mathbb{N} .
  \]
\end{definition}


\begin{lemma}[Closed form of the iterates]
  \label{lem_landweber_closed_form}
  For all $m \in \mathbb{N}$,
  \[
    x_m = \omega \sum_{n=0}^{m-1} (\Id - \omega K^{*}K)^{n} K^{*} y .
  \]
\end{lemma}

\begin{proof}
  We induct on $m$, with the convention that an empty sum is $0$.

  \emph{Base case.}  For $m = 0$ the right-hand side is the empty sum, which is $0$, and
  $x_0 = 0$ by the convention of this chapter.

  \emph{Induction step.}  Assume the formula for $m$.  By \cref{def_landweber},
  \[
    x_{m+1} = x_m + \omega K^{*}(y - Kx_m)
      = x_m - \omega K^{*}Kx_m + \omega K^{*}y
      = (\Id - \omega K^{*}K)x_m + \omega K^{*}y .
  \]
  Insert the induction hypothesis and move the linear operator
  $\Id - \omega K^{*}K$ through the finite sum:
  \[
    x_{m+1} = \omega \sum_{n=0}^{m-1}(\Id - \omega K^{*}K)^{n+1}K^{*}y
      \;+\; \omega\,(\Id - \omega K^{*}K)^{0}K^{*}y .
  \]
  Reindexing the first sum by $j := n+1$ turns it into $\omega\sum_{j=1}^{m}(\Id -
  \omega K^{*}K)^{j}K^{*}y$, and the second term is the missing summand $j = 0$; hence
  \[
    x_{m+1} = \omega \sum_{n=0}^{m}(\Id - \omega K^{*}K)^{n}K^{*}y . \qedhere
  \]
\end{proof}


\begin{theorem}[Landweber is a spectral regularization]
  \label{thm_landweber_regularization}
  Let $\omega \in (0,\kappa^{-1}]$ and let $\varphi_m$ be the Landweber filter of
  \cref{def_standard_filters}(iii) at $\alpha = 1/m$, i.e.\
  $\varphi_m(\lambda) = \lambda^{-1}(1 - (1-\omega\lambda)^{m})$ with
  $\varphi_m(0) = \omega m$.  Then $x_m = \varphi_m(K^{*}K)K^{*}y$, and
  $\{R_m\}_{m \in \mathbb{N}}$ with $R_m := \varphi_m(K^{*}K)K^{*}$ satisfies
  $R_m y \to K\dg y$ for every $y \in \Dom(K\dg)$.
\end{theorem}

\begin{proof}
  \emph{The Landweber filter is a polynomial.}  Fix $m \in \mathbb{N}$ and put
  \[
    p_m(\lambda) := \omega \sum_{n=0}^{m-1}(1-\omega\lambda)^{n},
    \qquad \lambda \in [0,\kappa] ,
  \]
  a polynomial in $\lambda$.  For $\lambda \in (0,\kappa]$ the ratio $q := 1-\omega\lambda$
  satisfies $1 - q = \omega\lambda \ne 0$, so the finite geometric sum evaluates to
  \[
    p_m(\lambda) = \omega\,\frac{1 - (1-\omega\lambda)^{m}}{\omega\lambda}
      = \frac{1 - (1-\omega\lambda)^{m}}{\lambda} = \varphi_m(\lambda) ,
  \]
  while $p_m(0) = \omega m = \varphi_m(0)$.  Hence $p_m = \varphi_m$ as functions on
  $[0,\kappa]$; in particular $\varphi_m$ is continuous and bounded there.

  \emph{Identification of the iterate.}  The functional calculus is linear and
  multiplicative (\cref{lem_fun_calc_hom}) and sends the constant function $1$ to $\Id$ and
  the identity function to $K^{*}K$ (\cref{lem_fun_calc_basic}, which is where both of
  those two values are recorded); consequently it sends
  $\lambda \mapsto (1-\omega\lambda)^{n}$ to $(\Id - \omega K^{*}K)^{n}$ and the
  polynomial $p_m$ to the operator polynomial (\cref{lem_fun_calc_basic} and
  \cref{lem_fun_calc_hom}, applied termwise)
  \[
    p_m(K^{*}K) = \omega \sum_{n=0}^{m-1}(\Id - \omega K^{*}K)^{n} .
  \]
  Since $\varphi_m = p_m$ on $[0,\kappa]$, \cref{lem_landweber_closed_form} gives
  \[
    x_m = \omega \sum_{n=0}^{m-1}(\Id - \omega K^{*}K)^{n}K^{*}y
      = p_m(K^{*}K)K^{*}y = \varphi_m(K^{*}K)K^{*}y = R_m y ,
  \]
  which is the spectral regularization operator of \cref{def_spectral_reg} for the filter
  $\varphi_m$.

  \emph{Convergence.}  That $\{\varphi_{\alpha}\}_{\alpha>0}$ of
  \cref{def_standard_filters}(iii) is a regularizing filter with $C_{\varphi} = 1$ is part
  of that node; the substance is the following two observations, valid because
  $\omega \le \kappa^{-1}$ forces $0 < \omega\lambda \le \omega\kappa \le 1$, i.e.\
  $0 \le 1 - \omega\lambda < 1$, for every $\lambda \in (0,\kappa]$.  First,
  $(1-\omega\lambda)^{1/\alpha} \to 0$ as $\alpha \to 0$, since $1/\alpha \to \infty$ and
  $q^{t} \to 0$ as $t \to \infty$ for $q \in [0,1)$; hence
  $\varphi_{\alpha}(\lambda) \to \lambda^{-1}$, which is \cref{def_filter}(i).  Second,
  $\lambda|\varphi_{\alpha}(\lambda)| = 1 - (1-\omega\lambda)^{1/\alpha} \in [0,1]$, which
  is \cref{def_filter}(ii) with $C_{\varphi} = 1$.

  By \cref{thm_spectral_is_regularization} the family $R_{\alpha} =
  \varphi_{\alpha}(K^{*}K)K^{*}$ satisfies $R_{\alpha}y \to K\dg y$ as $\alpha \to 0$ for
  every $y \in \Dom(K\dg)$.  The numbers $\alpha_m := 1/m$ are positive and tend to $0$,
  so the sequence $R_m y = R_{\alpha_m}y$ converges to $K\dg y$ as well.
\end{proof}


\begin{theorem}[Landweber with the discrepancy principle]
  \label{thm_landweber_order_optimal}
  Assume $\Ran(K)$ is dense in $Y$ and $\dim\Ran(K) = \infty$.  The Landweber filter has
  infinite qualification: it satisfies (F) with $C = \omega(1-\omega\kappa)^{-1}$ and
  (Q$_{\nu}$) for every $\nu > 0$ with $C_{\nu} = (\nu/(2\omega))^{\nu/2}e^{-\nu/2}$; its
  residual constant is $C_r = 1$.  Consequently, for every $\tau > 1$ the discrepancy
  principle defines an a posteriori rule $\alpha_{\tau}$ for the Landweber filter;
  $(R_{\alpha},\alpha_{\tau})$ is a convergent regularization method; and for every
  $\nu > 0$ there is $c = c(\nu,\tau) \ge 1$, independent of $\rho$, with
  \[
    \bigl\|R_{\alpha_{\tau}(\delta,y^{\delta})}y^{\delta} - x\dg\bigr\|_X
      \;\le\; c\,\delta^{\nu/(\nu+1)}\rho^{1/(\nu+1)}
  \]
  for every $\rho > 0$, every $x\dg \in X_{\nu,\rho}$, every $\delta > 0$ and every
  $y^{\delta} \in B_{\delta}(Kx\dg)$ \emph{with} $\|y^{\delta}\|_Y > \tau\delta$.
\end{theorem}

\begin{proof}
  Throughout, $t := 1/\alpha$, and by \cref{def_standard_filters}(iii) and
  \cref{lem_approx_error_spectral}
  \[
    \varphi_{\alpha}(\lambda) = \frac{1 - (1-\omega\lambda)^{t}}{\lambda},
    \qquad
    r_{\alpha}(\lambda) = 1 - \lambda\varphi_{\alpha}(\lambda) = (1-\omega\lambda)^{t},
    \qquad \lambda \in (0,\kappa] .
  \]
  Since $\omega < \kappa^{-1}$ we have $0 < \omega\lambda \le \omega\kappa < 1$, so
  $c := 1 - \omega\lambda \in (0,1)$ and both expressions are well defined and
  nonnegative.

  \emph{An elementary inequality.}  For all $x \in \mathbb{R}$ one has $e^{x} \ge 1+x$:
  the function $g(x) := e^{x} - 1 - x$ has $g(0) = 0$ and $g'(x) = e^{x}-1$, which is
  $\le 0$ for $x \le 0$ and $\ge 0$ for $x \ge 0$, so $g$ attains its minimum $0$ at $0$.
  Two consequences will be used: $1 - e^{-u} \le u$ for $u \ge 0$ (take $x = -u$), and
  $\ln s \le s - 1$ for $s > 0$ (take $x = \ln s$).

  \emph{Condition (F).}  Let $\alpha > 0$ and $\lambda \in (0,\kappa]$.  Writing
  $c^{t} = e^{-t\ln(1/c)}$ with $\ln(1/c) \ge 0$ and using $1 - e^{-u} \le u$ and then
  $\ln(1/c) \le 1/c - 1 = (1-c)/c = \omega\lambda/c$,
  \[
    1 - c^{t} \;\le\; t\,\ln(1/c) \;\le\; \frac{t\,\omega\lambda}{c}
      \;\le\; \frac{t\,\omega\lambda}{1-\omega\kappa} ,
  \]
  so $|\varphi_{\alpha}(\lambda)| \le \omega(1-\omega\kappa)^{-1}\alpha^{-1}$.  Since
  $\alpha > 0$ was arbitrary, (F) of \cref{def_qualification} holds with the constant
  $C = \omega(1-\omega\kappa)^{-1}$, which is the constant of the statement.

  For $\alpha \in (0,1]$ only --- in particular on the parameter set $\alpha = 1/m$ of
  this chapter, where $t = m \ge 1$ --- the sharper bound $\omega\alpha^{-1}$ holds: for
  $t \ge 1$ and $x \in [0,1]$ the function $h(x) := (1-x)^{t} - 1 + tx$ has $h(0) = 0$
  and $h'(x) = t\bigl(1 - (1-x)^{t-1}\bigr) \ge 0$, because $0 \le 1-x \le 1$ and
  $t - 1 \ge 0$; hence $(1-x)^{t} \ge 1 - tx$, and with $x = \omega\lambda$,
  \[
    |\varphi_{\alpha}(\lambda)| = \frac{1 - (1-\omega\lambda)^{t}}{\lambda}
      \le \frac{t\,\omega\lambda}{\lambda} = \omega\,\alpha^{-1} ,
    \qquad \alpha \in (0,1] .
  \]
  This is \emph{not} (F), which quantifies over all $\alpha > 0$: for $t = 1/\alpha \in
  (0,1)$ the function $f(x) := 1 - (1-x)^{t} - tx$ has $f(0) = 0$ and
  $f'(x) = t\bigl((1-x)^{t-1} - 1\bigr) > 0$ on $(0,1)$, so
  $|\varphi_{\alpha}(\lambda)| > \omega\alpha^{-1}$ for every $\lambda \in (0,\kappa]$
  and every $\alpha > 1$.

  \emph{Condition (Q$_{\nu}$).}  Fix $\nu > 0$ and $\alpha > 0$.  From $e^{x} \ge 1+x$
  with $x = -\omega\lambda$ we get $0 \le 1 - \omega\lambda \le e^{-\omega\lambda}$, and
  $s \mapsto s^{t}$ is nondecreasing on $[0,\infty)$, so
  $r_{\alpha}(\lambda) \le e^{-\omega\lambda t}$.  Hence, with
  $h_{t}(\lambda) := \lambda^{\nu/2}e^{-\omega t\lambda}$ on $(0,\infty)$, the
  quantity $\omega_{\nu}$ of \cref{lem_omega_nu} obeys
  \[
    \omega_{\nu}(\alpha) = \sup_{\lambda \in (0,\kappa]}
      \lambda^{\nu/2}|r_{\alpha}(\lambda)|
      \;\le\; \sup_{\lambda \in (0,\infty)} h_{t}(\lambda) .
  \]
  The function $h_{t}$ is differentiable with
  \[
    h_{t}'(\lambda) = \lambda^{\nu/2-1}e^{-\omega t\lambda}
      \Bigl(\frac{\nu}{2} - \omega t\lambda\Bigr) ,
  \]
  which is $> 0$ for $\lambda < \lambda^{*} := \nu/(2\omega t)$ and $< 0$ for
  $\lambda > \lambda^{*}$; so $h_{t}$ is increasing on $(0,\lambda^{*}]$ and decreasing on
  $[\lambda^{*},\infty)$, and $h_{t}(\lambda) \le h_{t}(\lambda^{*})$ for all
  $\lambda > 0$.  With $t = 1/\alpha$,
  \[
    h_{t}(\lambda^{*}) = \Bigl(\frac{\nu\alpha}{2\omega}\Bigr)^{\nu/2}e^{-\nu/2}
      = \Bigl(\frac{\nu}{2\omega}\Bigr)^{\nu/2}e^{-\nu/2}\,\alpha^{\nu/2} ,
  \]
  so (Q$_{\nu}$) holds with $C_{\nu} = (\nu/(2\omega))^{\nu/2}e^{-\nu/2}$, as claimed.

  \emph{Qualification.}  The previous paragraph verifies (Q$_{\nu}$) for every
  $\nu > 0$, so the supremum defining the qualification in \cref{def_qualification} is
  $\nu_0 = \infty$; in particular $\nu_0 > 1$ and $(0,\nu_0 - 1] = (0,\infty)$.

  \emph{The remaining hypotheses of \cref{thm_discrepancy_order_optimal}.}  Density of
  $\Ran(K)$ in $Y$ and $\dim\Ran(K) = \infty$ are assumed.  For fixed
  $\lambda \in (0,\kappa]$ the map $\alpha \mapsto \varphi_{\alpha}(\lambda) =
  \lambda^{-1}(1 - e^{\ln(1-\omega\lambda)/\alpha})$ is continuous on $(0,\infty)$,
  because $1 - \omega\lambda > 0$ makes $\ln(1-\omega\lambda)$ a finite constant.
  Finally $0 \le r_{\alpha}(\lambda) = (1-\omega\lambda)^{t} \le 1$ for all $\alpha$ and
  $\lambda \in (0,\kappa]$, so the constant $C_r$ of \cref{lem_residual_constant}
  satisfies $C_r \le 1$; and $C_r \ge 1$ by that same lemma, so $C_r = 1$ and the
  hypothesis $\tau > C_r$ is exactly $\tau > 1$.

  \emph{Conclusion.}  All hypotheses of \cref{thm_discrepancy_order_optimal} are met.
  That theorem therefore produces the a posteriori rule $\alpha_{\tau}$ of
  \cref{def_discrepancy_principle}; its part (ii) gives that
  $(R_{\alpha},\alpha_{\tau})$ is a convergent regularization method; and its part (i)
  gives, for every $\nu \in (0,\nu_0-1] = (0,\infty)$, a constant $c = c(\nu,\tau) \ge 1$
  independent of $\rho$ with
  \[
    \bigl\|R_{\alpha_{\tau}(\delta,y^{\delta})}y^{\delta} - x\dg\bigr\|_X
      \;\le\; c\,\delta^{\nu/(\nu+1)}\rho^{1/(\nu+1)}
  \]
  for every $\rho > 0$, every $x\dg \in X_{\nu,\rho}$, every $\delta > 0$ and every
  $y^{\delta} \in B_{\delta}(Kx\dg)$ with $\|y^{\delta}\|_Y > \tau\delta$ --- exactly the
  data on which \cref{def_discrepancy_principle} governs the rule.  By
  \cref{thm_landweber_regularization} the operator $R_{\alpha}$ at $\alpha = 1/m$ is
  exactly the $m$-th Landweber iterate; the rule $\alpha_{\tau}$ is, however, a parameter
  in the continuum $(0,\infty)$ and is not asserted here to be of the form $1/m$.
\end{proof}


\begin{theorem}[The residual decreases strictly]
  \label{thm_landweber_residual_monotone}
  Let $m \in \mathbb{N}$ and let $x^{\delta}_m$ be the Landweber iterates for the datum
  $y^{\delta}$.  If $K^{*}(Kx^{\delta}_m - y^{\delta}) \ne 0$, then
  \[
    \|Kx^{\delta}_{m+1} - y^{\delta}\|_Y \;<\; \|Kx^{\delta}_m - y^{\delta}\|_Y .
  \]
  With the convention $x_0 = 0$ the hypothesis is equivalent to
  $y^{\delta} \notin \Ker(K^{*})$; if moreover $\overline{\Ran(K)} = Y$ --- which is
  \emph{not} a convention of this chapter and is assumed only where said --- it reduces
  to $y^{\delta} \ne 0$.
\end{theorem}

\begin{proof}
  Write $r_m := Kx^{\delta}_m - y^{\delta}$ and let $P$ be the orthogonal projection onto
  $\overline{\Ran(K)}$.  Since $K$ is compact it has a singular system
  $\{(\sigma_n,u_n,v_n)\}_{n \in I}$ (\cref{thm_svd}); by \cref{def_singular_system},
  $\{u_n\}_{n \in I}$ is an orthonormal basis of $\overline{\Ran(K)}$, the $\sigma_n$ are
  positive and nonincreasing, and $Kv_n = \sigma_n u_n$, $K^{*}u_n = \sigma_n v_n$.
  \cref{def_singular_system} orders the $\sigma_n$ and carries their positivity;
  $\sigma_n^{2} \le \kappa = \|K\|^{2}$ is \cref{lem_sigma_bound}, so
  $\sigma_n^{2} \in (0,\kappa]$ for every $n \in I$.  Furthermore
  $\Ker(K^{*}) = \overline{\Ran(K)}^{\perp}$
  (\cref{thm_four_subspaces}).  In particular, for $r \in Y$,
  \begin{equation}\label{eq_lw_ker_adj}
    K^{*}r = 0
    \iff r \in \overline{\Ran(K)}^{\perp}
    \iff Pr = 0
    \iff \langle u_n, r\rangle = 0 \ \text{ for all } n \in I ,
  \end{equation}
  the last equivalence because $\{u_n\}$ is a basis of the range of $P$
  (\cref{thm_proj_expansion}).

  \emph{The residual is propagated by $\Id - \omega KK^{*}$.}  \cref{def_landweber} gives
  $x^{\delta}_{m+1} = x^{\delta}_m + \omega K^{*}(y^{\delta} - Kx^{\delta}_m) =
  x^{\delta}_m - \omega K^{*}r_m$, hence
  \begin{equation}\label{eq_lw_res_step}
    r_{m+1} = Kx^{\delta}_{m+1} - y^{\delta} = r_m - \omega KK^{*}r_m
      = (\Id - \omega KK^{*})r_m .
  \end{equation}

  \emph{A Parseval identity.}  For any $r \in Y$ the vectors $Pr$ and $(\Id - P)r$ are
  orthogonal (\cref{lem_proj_characterisation}), so Pythagoras (\cref{lem_pythagoras}) gives
  $\|r\|_Y^{2} = \|Pr\|_Y^{2} + \|(\Id - P)r\|_Y^{2}$.  The expansion
  $Pr = \sum_{n \in I}\langle u_n, r\rangle_Y u_n$ in the orthonormal basis $\{u_n\}$ of
  $\overline{\Ran(K)}$ is \cref{thm_proj_expansion}, and its coefficients are square
  summable by Bessel's inequality (\cref{thm_bessel}); the accompanying \emph{norm}
  identity $\|Pr\|_Y^{2} = \sum_{n \in I}|\langle u_n,r\rangle_Y|^{2}$ is not part of
  \cref{thm_proj_expansion} and is supplied by \cref{lem_parseval_series}.  Hence
  \begin{equation}\label{eq_lw_parseval}
    \|r\|_Y^{2} = \|Pr\|_Y^{2} + \|(\Id - P)r\|_Y^{2}
      = \sum_{n \in I}|\langle u_n, r\rangle_Y|^{2} + \|(\Id - P)r\|_Y^{2} ,
  \end{equation}
  where $\langle u_n, Pr\rangle = \langle Pu_n, r\rangle = \langle u_n, r\rangle$ because
  $P$ is self-adjoint (\cref{prop_proj_props}) and $Pu_n = u_n$ --- the latter because
  $u_n \in \overline{\Ran(K)}$ and $u_n - u_n = 0 \perp \overline{\Ran(K)}$, which is the
  criterion of \cref{lem_proj_characterisation}.

  \emph{The coefficients of the propagated residual.}  Since $KK^{*}u_n = \sigma_n^{2}u_n$
  (\cref{lem_svd_companion}) and $\Id - \omega KK^{*}$ is self-adjoint with real
  coefficients,
  \eqref{eq_lw_res_step} gives
  \[
    \langle u_n, r_{m+1}\rangle = \langle (\Id - \omega KK^{*})u_n, r_m\rangle
      = (1 - \omega\sigma_n^{2})\,\langle u_n, r_m\rangle ,
    \qquad n \in I .
  \]
  Moreover $KK^{*}r_m \in \Ran(K) \subseteq \overline{\Ran(K)}$, so
  $(\Id - P)KK^{*}r_m = 0$ and therefore $(\Id - P)r_{m+1} = (\Id - P)r_m$ by
  \eqref{eq_lw_res_step}.  Substituting both into \eqref{eq_lw_parseval} and subtracting
  the two (convergent, by \cref{thm_bessel}) series,
  \begin{equation}\label{eq_lw_res_gain}
    \|r_m\|_Y^{2} - \|r_{m+1}\|_Y^{2}
      = \sum_{n \in I}\bigl(1 - (1-\omega\sigma_n^{2})^{2}\bigr)\,
        |\langle u_n, r_m\rangle_Y|^{2} .
  \end{equation}
  Every $n \in I$ satisfies $0 < \sigma_n^{2} \le \kappa$ by the first
  paragraph (\cref{def_singular_system} and \cref{lem_sigma_bound}), so
  $0 < \omega\sigma_n^{2} \le \omega\kappa \le 1$ and hence
  $0 \le 1 - \omega\sigma_n^{2} < 1$; therefore $(1-\omega\sigma_n^{2})^{2} < 1$ and every
  coefficient in \eqref{eq_lw_res_gain} is strictly positive.

  \emph{Conclusion.}  If $K^{*}r_m \ne 0$ then by \eqref{eq_lw_ker_adj} some
  $\langle u_n, r_m\rangle \ne 0$, so the sum in \eqref{eq_lw_res_gain} is strictly
  positive and $\|r_{m+1}\|_Y < \|r_m\|_Y$, which is the claim.

  \emph{The hypothesis for $x_0 = 0$.}  Here $r_0 = -y^{\delta}$, so iterating
  \eqref{eq_lw_res_step} gives $r_m = -(\Id - \omega KK^{*})^{m}y^{\delta}$ and, by the
  coefficient formula above applied $m$ times,
  $\langle u_n, r_m\rangle = -(1-\omega\sigma_n^{2})^{m}\langle u_n, y^{\delta}\rangle$.
  Under the standing assumption $\omega < \kappa^{-1}$ we have
  $\omega\sigma_n^{2} \le \omega\kappa < 1$, so $(1-\omega\sigma_n^{2})^{m} \ne 0$ and
  $\langle u_n, r_m\rangle = 0$ if and only if $\langle u_n, y^{\delta}\rangle = 0$.  By
  \eqref{eq_lw_ker_adj}, applied once to $r_m$ and once to $y^{\delta}$, the hypothesis
  $K^{*}r_m \ne 0$ is therefore equivalent to $y^{\delta} \notin \Ker(K^{*})$.  If
  moreover $\overline{\Ran(K)} = Y$, then $\Ker(K^{*}) = Y^{\perp} = \{0\}$ and the
  condition reads $y^{\delta} \ne 0$.
\end{proof}


\begin{theorem}[The error decreases while the residual exceeds twice the noise]
  \label{thm_landweber_error_monotone}
  Let $K \in \Ll(X,Y)$ be bounded (compactness is not needed), let $x\dg \in X$ satisfy
  $Kx\dg = y$, let $\|y - y^{\delta}\|_Y \le \delta$, and let $m \in \mathbb{N}$.  If
  $\|Kx^{\delta}_m - y^{\delta}\|_Y > 2\delta$, then
  \[
    \|x^{\delta}_{m+1} - x\dg\|_X \;<\; \|x^{\delta}_m - x\dg\|_X .
  \]
\end{theorem}

\begin{proof}
  Put $\rho_m := y^{\delta} - Kx^{\delta}_m$, so that
  $x^{\delta}_{m+1} = x^{\delta}_m + \omega K^{*}\rho_m$ by \cref{def_landweber}.
  Expanding the square,
  \begin{equation}\label{eq_lw_err_expand}
    \|x^{\delta}_{m+1} - x\dg\|_X^{2}
      = \|x^{\delta}_m - x\dg\|_X^{2}
        + 2\omega\,\mathrm{Re}\,\langle x^{\delta}_m - x\dg, K^{*}\rho_m\rangle_X
        + \omega^{2}\|K^{*}\rho_m\|_X^{2} .
  \end{equation}

  \emph{The cross term.}  By the definition of the adjoint and $Kx\dg = y$,
  \[
    \langle x^{\delta}_m - x\dg, K^{*}\rho_m\rangle_X
      = \langle K x^{\delta}_m - Kx\dg, \rho_m\rangle_Y
      = \langle Kx^{\delta}_m - y^{\delta}, \rho_m\rangle_Y
        + \langle y^{\delta} - y, \rho_m\rangle_Y ,
  \]
  where the second line splits $Kx^{\delta}_m - y = (Kx^{\delta}_m - y^{\delta}) +
  (y^{\delta} - y)$.  The first summand is $-\|\rho_m\|_Y^{2}$, and the second is bounded
  by $\delta\|\rho_m\|_Y$ in modulus by Cauchy--Schwarz (\cref{thm_cauchy_schwarz}) and
  $\|y - y^{\delta}\|_Y \le \delta$.  Hence
  \[
    \mathrm{Re}\,\langle x^{\delta}_m - x\dg, K^{*}\rho_m\rangle_X
      \;\le\; -\|\rho_m\|_Y^{2} + \delta\|\rho_m\|_Y .
  \]

  \emph{The quadratic term.}  Since $\|K^{*}\|_{\Ll(Y,X)} = \|K\|_{\Ll(X,Y)}$
  (\cref{prop_adjoint_props}), $\|K^{*}\rho_m\|_X^{2} \le \kappa\|\rho_m\|_Y^{2}$, so
  $\omega^{2}\|K^{*}\rho_m\|_X^{2} \le \omega^{2}\kappa\|\rho_m\|_Y^{2} \le
  \omega\|\rho_m\|_Y^{2}$, using $\omega\kappa < 1$.

  \emph{Conclusion.}  Substituting both estimates into \eqref{eq_lw_err_expand},
  \[
    \begin{aligned}
      \|x^{\delta}_{m+1} - x\dg\|_X^{2}
        &\le \|x^{\delta}_m - x\dg\|_X^{2}
          + 2\omega\delta\|\rho_m\|_Y - 2\omega\|\rho_m\|_Y^{2}
          + \omega\|\rho_m\|_Y^{2} \\
        &= \|x^{\delta}_m - x\dg\|_X^{2}
          + \omega\|\rho_m\|_Y\bigl(2\delta - \|\rho_m\|_Y\bigr) .
    \end{aligned}
  \]
  By hypothesis $\|\rho_m\|_Y = \|Kx^{\delta}_m - y^{\delta}\|_Y > 2\delta \ge 0$, so
  $\|\rho_m\|_Y > 0$ and $2\delta - \|\rho_m\|_Y < 0$; the last term is therefore strictly
  negative and
  $\|x^{\delta}_{m+1} - x\dg\|_X^{2} < \|x^{\delta}_m - x\dg\|_X^{2}$.  Taking square
  roots of nonnegative numbers preserves the strict inequality.
\end{proof}


\begin{theorem}[The discrepancy principle stops in $O(\delta^{-2})$ steps]
  \label{thm_landweber_stopping_bound}
  Let $x\dg \in X$, $\tau > 1$, $\delta > 0$ and $y^{\delta} \in B_{\delta}(Kx\dg)$, and
  let $x^{\delta}_m$ be the Landweber iterates for $y^{\delta}$.  Then the first index
  $m(\delta,y^{\delta})$ with $\|Kx^{\delta}_m - y^{\delta}\|_Y \le \tau\delta$ exists and
  satisfies
  \[
    m(\delta,y^{\delta}) \;\le\;
      \Bigl\lceil \omega^{-1}(\tau-1)^{-2}\,\|x\dg\|_X^{2}\ \delta^{-2} \Bigr\rceil .
  \]
\end{theorem}

\begin{proof}
  Write $y := Kx\dg$; only $Kx\dg = y$ is used below, not that $x\dg$ has minimal norm.
  Let $\{x_n\}$ be the Landweber iterates for the exact datum $y$ and
  $\rho_n := y - Kx_n$, and let $\rho^{\delta}_n := y^{\delta} - Kx^{\delta}_n$ be the
  residuals for $y^{\delta}$.  As in \eqref{eq_lw_res_step} of
  \cref{thm_landweber_residual_monotone}, but read for the exact and the noisy datum
  separately, \cref{def_landweber} gives
  \begin{equation}\label{eq_lw_stop_step}
    \rho_{n+1} = (\Id - \omega KK^{*})\rho_n ,
    \qquad
    \rho^{\delta}_{n+1} = (\Id - \omega KK^{*})\rho^{\delta}_n ,
  \end{equation}
  and since $x_0 = x^{\delta}_0 = 0$, iterating gives $\rho_n = (\Id-\omega KK^{*})^{n}y$
  and $\rho^{\delta}_n = (\Id - \omega KK^{*})^{n}y^{\delta}$.

  \emph{$\Id - \omega KK^{*}$ is a contraction.}  For $r \in Y$,
  \[
    \|(\Id - \omega KK^{*})r\|_Y^{2}
      = \|r\|_Y^{2} - 2\omega\langle r, KK^{*}r\rangle_Y + \omega^{2}\|KK^{*}r\|_Y^{2} .
  \]
  Here $\langle r, KK^{*}r\rangle_Y = \|K^{*}r\|_X^{2} \ge 0$ and
  $\|KK^{*}r\|_Y^{2} \le \kappa\|K^{*}r\|_X^{2}$, so
  \[
    \|(\Id - \omega KK^{*})r\|_Y^{2}
      \le \|r\|_Y^{2} - \omega(2 - \omega\kappa)\|K^{*}r\|_X^{2} \le \|r\|_Y^{2}
  \]
  because $\omega\kappa \le 1 < 2$.  Thus $\|(\Id - \omega KK^{*})r\|_Y \le \|r\|_Y$ for
  every $r$; with \eqref{eq_lw_stop_step} this gives $\|\rho_{n+1}\|_Y \le \|\rho_n\|_Y$
  for all $n$, hence $\|\rho_n\|_Y \ge \|\rho_m\|_Y$ whenever $n \le m$.

  \emph{A telescoping estimate for the exact iterates.}  As in the proof of
  \cref{thm_landweber_error_monotone}, with $\delta = 0$ and $y^{\delta}$ replaced by $y$,
  \[
    \|x_{n+1} - x\dg\|_X^{2}
      = \|x_n - x\dg\|_X^{2}
        + 2\omega\,\mathrm{Re}\,\langle Kx_n - y, \rho_n\rangle_Y
        + \omega^{2}\|K^{*}\rho_n\|_X^{2} ,
  \]
  where $\langle Kx_n - y,\rho_n\rangle_Y = -\|\rho_n\|_Y^{2}$ and
  $\omega^{2}\|K^{*}\rho_n\|_X^{2} \le \omega^{2}\kappa\|\rho_n\|_Y^{2} \le
  \omega\|\rho_n\|_Y^{2}$ by $\omega\kappa \le 1$.  Hence
  \[
    \|x\dg - x_n\|_X^{2} - \|x\dg - x_{n+1}\|_X^{2} \;\ge\; \omega\|\rho_n\|_Y^{2} .
  \]
  Summing over $n = 0,\dots,m-1$, the left-hand side telescopes, and $x_0 = 0$ gives
  \[
    \|x\dg\|_X^{2} \;\ge\; \|x\dg - x_0\|_X^{2} - \|x\dg - x_m\|_X^{2}
      \;\ge\; \omega\sum_{n=0}^{m-1}\|\rho_n\|_Y^{2}
      \;\ge\; \omega\,m\,\|\rho_m\|_Y^{2} ,
  \]
  the last step by the monotonicity just proved.  Therefore
  \begin{equation}\label{eq_lw_exact_res}
    \|y - Kx_m\|_Y \;\le\; (\omega m)^{-1/2}\|x\dg\|_X , \qquad m \ge 1 .
  \end{equation}

  \emph{The noisy residual.}  Using the closed forms above and the triangle inequality,
  \[
    \begin{aligned}
      \|Kx^{\delta}_m - y^{\delta}\|_Y
        &= \|(\Id-\omega KK^{*})^{m}y^{\delta}\|_Y \\
        &\le \|(\Id-\omega KK^{*})^{m}y\|_Y
          + \|(\Id-\omega KK^{*})^{m}(y^{\delta}-y)\|_Y .
    \end{aligned}
  \]
  The first summand is $\|\rho_m\|_Y = \|y - Kx_m\|_Y$; the second is at most
  $\|y^{\delta}-y\|_Y \le \delta$ by applying the contraction property $m$ times.  With
  \eqref{eq_lw_exact_res},
  \begin{equation}\label{eq_lw_noisy_res}
    \|Kx^{\delta}_m - y^{\delta}\|_Y \;\le\; (\omega m)^{-1/2}\|x\dg\|_X + \delta ,
    \qquad m \ge 1 .
  \end{equation}

  \emph{Conclusion.}  Put
  $\bar{m} := \lceil \omega^{-1}(\tau-1)^{-2}\|x\dg\|_X^{2}\delta^{-2}\rceil$ and let $S$
  be the set of indices $m \ge 0$ with $\|Kx^{\delta}_m - y^{\delta}\|_Y \le \tau\delta$.
  If $\|x\dg\|_X = 0$ then $y = 0$, $\bar{m} = 0$ and
  $\|Kx^{\delta}_0 - y^{\delta}\|_Y = \|y^{\delta}\|_Y \le \delta \le \tau\delta$, so
  $\bar{m} \in S$.  If $\|x\dg\|_X > 0$ then $\bar{m} \ge \omega^{-1}(\tau-1)^{-2}
  \|x\dg\|_X^{2}\delta^{-2} > 0$, so $\omega\bar{m} \ge (\tau-1)^{-2}\|x\dg\|_X^{2}
  \delta^{-2}$ and therefore $(\omega\bar{m})^{-1/2}\|x\dg\|_X \le (\tau-1)\delta$;
  by \eqref{eq_lw_noisy_res},
  \[
    \|Kx^{\delta}_{\bar{m}} - y^{\delta}\|_Y \le (\tau-1)\delta + \delta = \tau\delta ,
  \]
  so again $\bar{m} \in S$.  In both cases $S \ne \emptyset$, so by well-ordering $S$ has
  a least element $m(\delta,y^{\delta})$ --- the first index with residual at most
  $\tau\delta$ --- and $m(\delta,y^{\delta}) \le \bar{m}$.
\end{proof}


\begin{corollary}[Faster stopping under a source condition]
  \label{cor_landweber_stopping_rate}
  Assume $\Ran(K)$ dense in $Y$ and $\dim\Ran(K) = \infty$.  Let $\nu > 0$, $\rho > 0$,
  $x\dg \in X_{\nu,\rho}$, $\tau > 1$, $\delta > 0$ and $y^{\delta} \in B_{\delta}(Kx\dg)$,
  and let $m(\delta,y^{\delta})$ be the stopping index of
  \cref{thm_landweber_stopping_bound}.  Then there is
  $C = C(\nu,\rho,\tau,\omega,\kappa) > 0$, independent of $\delta$ and of $y^{\delta}$,
  with $m(\delta,y^{\delta}) \le C\,\delta^{-2/(\nu+1)}$.
\end{corollary}

\begin{proof}


  Since $x\dg \in X_{\nu,\rho} \subseteq X$, \cref{thm_landweber_stopping_bound} applies
  and $m(\delta,y^{\delta})$ --- the least $m \ge 0$ with
  $\|Kx^{\delta}_m - y^{\delta}\|_Y \le \tau\delta$ --- exists.  Put
  \[
    C \;:=\; c_{\tau,\nu}^{-1}\rho^{2/(\nu+1)}
      \;+\; \Bigl(\frac{\kappa^{(\nu+1)/2}\rho}{\tau-1}\Bigr)^{2/(\nu+1)} \;>\; 0 ,
  \]
  with $c_{\tau,\nu}$ the constant produced in Step~3 below.

  \emph{Step 1: the data on which the principle says nothing.}  If
  $\|y^{\delta}\|_Y \le \tau\delta$ then, since $x^{\delta}_0 = 0$ by the convention of
  this chapter, the residual at $m = 0$ is $\|Kx^{\delta}_0 - y^{\delta}\|_Y =
  \|y^{\delta}\|_Y \le \tau\delta$, so $m(\delta,y^{\delta}) = 0$ and the asserted bound
  holds because $C\delta^{-2/(\nu+1)} > 0$.  Assume from now on
  $\|y^{\delta}\|_Y > \tau\delta$, which is the hypothesis under which
  \cref{def_discrepancy_principle} and \cref{lem_discrepancy_parameter_bound} operate.

  \emph{Step 2: the noise level is bounded above.}  Since $K$ is compact it has a
  singular system $\{(\sigma_n,u_n,v_n)\}_{n \in I}$ (\cref{thm_svd}), with $\{v_n\}$ an
  orthonormal basis of $\Ker(K)^{\perp}$, $\{u_n\}$ an orthonormal basis of
  $\overline{\Ran(K)}$ and $Kv_n = \sigma_nu_n$ (\cref{def_singular_system}), and with
  $\kappa = \|K\|^{2}$, hence $\sigma_n^{2} \in (0,\kappa]$ for every
  $n \in I$, by \cref{def_singular_system} and \cref{lem_sigma_bound} --- which is what
  the estimate $\sigma_n^{2(\nu+1)} \le \kappa^{\nu+1}$ below
  rests on.  This is the system used by
  \cref{lem_spectral_residual_form} below.  By \cref{def_source_set} write
  $x\dg = |K|^{\nu}w$ with $\|w\|_X \le \rho$.  By \cref{def_fun_calc} and
  \cref{def_frac_power}, $|K|^{\nu}w = \sum_{n \in I}\sigma_n^{\nu}\langle v_n,w\rangle_Xv_n$;
  applying the bounded operator $K$ term by term (\cref{lem_operator_through_series}) and
  using $Kv_n = \sigma_nu_n$ (\cref{def_singular_system}) gives
  $Kx\dg = \sum_{n \in I}\sigma_n^{\nu+1}\langle v_n,w\rangle_Xu_n$.  These coefficients
  are square summable, since $\sigma_n^{2(\nu+1)} \le \kappa^{\nu+1}$ and
  $\sum_{n}|\langle v_n,w\rangle_X|^{2} \le \|w\|_X^{2} \le \rho^{2}$ by Bessel's
  inequality (\cref{thm_bessel}); so \cref{lem_parseval_series} applies to the orthonormal
  system $\{u_n\}$ and gives
  \[
    \|Kx\dg\|_Y^{2} = \sum_{n \in I}\sigma_n^{2(\nu+1)}|\langle v_n,w\rangle_X|^{2}
      \;\le\; \kappa^{\nu+1}\sum_{n \in I}|\langle v_n,w\rangle_X|^{2}
      \;\le\; \kappa^{\nu+1}\rho^{2} .
  \]
  Since $\|y^{\delta} - Kx\dg\|_Y \le \delta$, Step~1 gives $\tau\delta <
  \|y^{\delta}\|_Y \le \|Kx\dg\|_Y + \delta$, so $(\tau-1)\delta < \|Kx\dg\|_Y$ and
  \begin{equation}\label{eq_lw_delta_small}
    \delta \;<\; D \;:=\; \frac{\kappa^{(\nu+1)/2}\rho}{\tau - 1} .
  \end{equation}

  \emph{Step 3: the parameter is bounded below.}  By
  \cref{thm_landweber_order_optimal} the Landweber filter satisfies (F) and
  (Q$_{\nu'}$) of \cref{def_qualification} for every $\nu' > 0$, has qualification
  $\nu_0 = \infty > 1$ and residual constant $C_r = 1 < \tau$, and its proof also
  verifies the continuity of $\alpha \mapsto \varphi_{\alpha}(\lambda)$; density of
  $\Ran(K)$ and $\dim\Ran(K) = \infty$ are assumed here.  So the hypotheses of
  \cref{thm_discrepancy_order_optimal} hold, and since $(0,\nu_0-1] = (0,\infty) \ni \nu$,
  \cref{lem_discrepancy_parameter_bound} --- whose hypothesis
  $\|y^{\delta}\|_Y > \tau\delta$ holds by Step~1 --- provides $c_{\tau,\nu} > 0$,
  depending only on $\tau$, $\nu$ and the filter, hence only on $\tau$, $\nu$, $\omega$
  and $\kappa$, with
  \begin{equation}\label{eq_lw_alpha_lower}
    \alpha \;:=\; \alpha_{\tau}(\delta,y^{\delta})
      \;\ge\; c_{\tau,\nu}\,(\delta/\rho)^{2/(\nu+1)} \;>\; 0 .
  \end{equation}

  \emph{Step 4: the residual is monotone in the parameter.}  By
  \cref{def_standard_filters}(iii) and \cref{lem_approx_error_spectral} the Landweber
  residual function is $r_{\beta}(\lambda) = (1-\omega\lambda)^{1/\beta}$, and
  $q := 1 - \omega\lambda \in (0,1)$ for $\lambda \in (0,\kappa]$ because
  $0 < \omega\lambda \le \omega\kappa < 1$.  As $s \mapsto q^{s}$ is decreasing for
  $q \in (0,1)$ and $\beta \mapsto 1/\beta$ is decreasing, $\beta \mapsto
  r_{\beta}(\lambda)$ is nondecreasing on $(0,\infty)$ for each such $\lambda$.  Writing
  $c_n := \langle u_n, y^{\delta}\rangle_Y$ and $P_R$ for the orthogonal projection onto
  $\overline{\Ran(K)}$, \cref{lem_spectral_residual_form}(iii) gives
  \[
    \|y^{\delta} - KR_{\beta}y^{\delta}\|_Y^{2}
      = \|(\Id - P_R)y^{\delta}\|_Y^{2}
        + \sum_{n \in I}|r_{\beta}(\sigma_n^{2})|^{2}|c_n|^{2} ,
  \]
  a sum of nonnegative terms each nondecreasing in $\beta$; so
  $\beta \mapsto \|y^{\delta} - KR_{\beta}y^{\delta}\|_Y$ is nondecreasing on
  $(0,\infty)$.

  \emph{Step 5: from the parameter to the stopping index.}  Put
  $M := \lceil 1/\alpha \rceil$, so $M \ge 1$ is an integer and $1/M \le \alpha$.  By
  Step~1 and \cref{def_discrepancy_principle}, $\|KR_{\alpha}y^{\delta} - y^{\delta}\|_Y
  \le \tau\delta$, so Step~4 applied to $1/M \le \alpha$ gives
  \[
    \|KR_{1/M}y^{\delta} - y^{\delta}\|_Y
      \;\le\; \|KR_{\alpha}y^{\delta} - y^{\delta}\|_Y \;\le\; \tau\delta .
  \]
  By \cref{thm_landweber_regularization}, $R_{1/M}y^{\delta} = x^{\delta}_M$; hence $M$
  is an index whose residual is at most $\tau\delta$, and minimality gives
  $m(\delta,y^{\delta}) \le M = \lceil 1/\alpha \rceil \le \alpha^{-1} + 1$.  Note that
  no claim is made that $\alpha$ itself is of the form $1/m$.

  \emph{Step 6: conclusion.}  By \eqref{eq_lw_alpha_lower}, $\alpha^{-1} \le
  c_{\tau,\nu}^{-1}\rho^{2/(\nu+1)}\delta^{-2/(\nu+1)}$; and by
  \eqref{eq_lw_delta_small}, $\delta^{2/(\nu+1)} < D^{2/(\nu+1)}$, i.e.\
  $1 < D^{2/(\nu+1)}\delta^{-2/(\nu+1)}$, the exponentiation being order preserving on
  $(0,\infty)$.  Adding the two,
  \[
    m(\delta,y^{\delta}) \;\le\; \alpha^{-1} + 1
      \;\le\; \bigl(c_{\tau,\nu}^{-1}\rho^{2/(\nu+1)} + D^{2/(\nu+1)}\bigr)
        \delta^{-2/(\nu+1)}
      \;=\; C\,\delta^{-2/(\nu+1)} . \qedhere
  \]
\end{proof}


\chapter{Discretization as regularization}
\label{ch_18_discretization}

\noindent\emph{Background.} That letting the discretization level itself act as the
regularization parameter --- solving finite-dimensional least-squares problems on an
increasing sequence of subspaces, with the dimension $n$ playing the role of
$1/\alpha$ --- already yields a convergent, order-optimal regularization method is due to
\citet{Natterer1977}. The dual variant of this chapter, projecting the data space $Y$
rather than the domain $X$, together with its parameter choice and rate theory, follows
\citet{PlatoVainikko1990}.

\begin{assumptions}
  \label{asm_discretization_as_regularization}
  $X$, $Y$ are Hilbert spaces and $T \in \Ll(X,Y)$;
no spectral theory is used in the first section, so $T$ need not be compact there.  The
two sections use \emph{different} finite-dimensional families, and the symbols $X_n$,
$P_n$, $T_n$ mean different things in each; they are bound separately in
\cref{def_lsq_projection} and \cref{def_dual_lsq_projection} and must not be conflated.
Both preambles require the union of the subspaces to have \emph{dense} span, never to
exhaust the space: a countable increasing union of finite-dimensional subspaces is never
an infinite-dimensional Banach space.
\end{assumptions}

\begin{definition}[Least-squares projection]
  \label{def_lsq_projection}
  Let $\dim X = \infty$ and let $X_1 \subset X_2 \subset \cdots \subset X$ be nested
  subspaces with $\dim X_n = n$ and $\overline{\bigcup_n X_n} = X$.  Let $P_n := P_{X_n}$, put
  $T_n := T P_n \in \Ll(X,Y)$, and for $y \in Y$ set $x_n := T_n\dg y$.
\end{definition}


\begin{lemma}[Discrete ranges are closed]
  \label{lem_lsq_range_closed}
  Let $T \in \Ll(X,Y)$ and let $X_1 \subset X_2 \subset \cdots \subset X$ be nested
  subspaces with $\dim X_n = n$, and write $P_n := P_{X_n}$.  Then $\Ran(TP_n)$ is closed
  for every $n$.
\end{lemma}

\begin{proof}
  $\Ran(TP_n) = T(X_n)$ is the image of a finite-dimensional space, hence
  finite-dimensional and therefore closed (\cref{lem_findim_closed}).
\end{proof}


\begin{lemma}[Comparison estimate]
  \label{lem_lsq_comparison}
  For $y \in \Dom(T\dg)$ and every $n$,
  \[
    \|Tx_n - Tx\dg\|_Y \;\le\; \|T\|_{\Ll(X,Y)}\,\|(\Id - P_n)x\dg\|_X ,
  \]
  and the right-hand side tends to $0$.
\end{lemma}

\begin{proof}
  Throughout, $P_n$, $T_n = TP_n$ and $x_n = T_n\dg y$ are those of
  \cref{def_lsq_projection}.

  \emph{Two standing facts about $T_n$.}  First, $\Ran(T_n)$ is closed
  (\cref{lem_lsq_range_closed}); by \cref{thm_mp_continuous_iff_closed_range} and
  \cref{lem_mp_domain_dense}, $T_n\dg \in \Ll(Y,X)$ with $\Dom(T_n\dg) = Y$, so $x_n$ is
  defined for every $y \in Y$.  Second, $T_n^{*} = (TP_n)^{*} = P_nT^{*}$ because $P_n$ is
  self-adjoint (\cref{prop_proj_props}, \cref{prop_adjoint_props}), so
  $\Ran(T_n^{*}) \subseteq X_n$; as $X_n$ is closed,
  \begin{equation}\label{eq_lsq_in_xn}
    x_n \in \Ran(T_n\dg) = \Ker(T_n)^{\perp} = \overline{\Ran(T_n^{*})} \subseteq X_n ,
    \qquad\text{hence } P_nx_n = x_n ,
  \end{equation}
  by \cref{lem_mp_equations} and \cref{thm_four_subspaces}.  In particular
  $T_nx_n = TP_nx_n = Tx_n$.

  \emph{The comparison.}  Write $P_R$ for the orthogonal projection onto
  $\overline{\Ran(T)}$.  For every $x \in X$ we have $T_nx \in \Ran(T) \subseteq
  \overline{\Ran(T)}$ and $y - P_Ry \perp \overline{\Ran(T)}$
  (\cref{lem_proj_characterisation}), so Pythagoras (\cref{lem_pythagoras}) gives
  \[
    \|T_nx - y\|_Y^{2} = \|T_nx - P_Ry\|_Y^{2} + \|(\Id - P_R)y\|_Y^{2} .
  \]
  Since $y \in \Dom(T\dg)$, \cref{lem_mp_equations}(iv) gives $P_Ry = TT\dg y = Tx\dg$.
  By \cref{thm_mp_minnorm}(ii) the element $x_n = T_n\dg y$ is in particular a least
  squares solution of $T_nx = y$, so $\|T_nx_n - y\|_Y \le \|T_nz - y\|_Y$ for every
  $z \in X$; taking $z := P_nx\dg$ and cancelling the common summand
  $\|(\Id-P_R)y\|_Y^{2}$ in the displayed identity yields
  \[
    \begin{aligned}
      \|Tx_n - Tx\dg\|_Y = \|T_nx_n - Tx\dg\|_Y
        &\le \|T_nP_nx\dg - Tx\dg\|_Y = \|T(P_n - \Id)x\dg\|_Y \\
        &\le \|T\|_{\Ll(X,Y)}\,\|(\Id - P_n)x\dg\|_X ,
    \end{aligned}
  \]
  using $T_nx_n = Tx_n$ from \eqref{eq_lsq_in_xn} and $T_nP_n = TP_nP_n = TP_n$ from
  $P_n^{2} = P_n$ (\cref{prop_proj_props}).

  \emph{The right-hand side is a null sequence.}  The $X_n$ are nested with
  $\overline{\bigcup_n X_n} = X$, so $P_nx \to x$ for every $x \in X$
  (\cref{thm_nested_proj}); applied to $x\dg$ this gives
  $\|(\Id - P_n)x\dg\|_X \to 0$.
\end{proof}


\begin{lemma}[Convergence criterion]
  \label{lem_lsq_convergence_criterion}
  Let $y \in \Dom(T\dg)$.  Then $x_n \to x\dg$ if and only if
  $\limsup_{n \to \infty}\|x_n\|_X \le \|x\dg\|_X$.
\end{lemma}

\begin{proof}
  The sequence $\{x_n\}$ is the one of \cref{def_lsq_projection}.

  \emph{($\Rightarrow$)}  If $x_n \to x\dg$ then, by the triangle inequality,
  \[
    \|x_n\|_X \le \|x_n - x\dg\|_X + \|x\dg\|_X \longrightarrow \|x\dg\|_X ,
  \]
  so $\limsup_n\|x_n\|_X \le \|x\dg\|_X$.

  \emph{($\Leftarrow$)}  Assume $\limsup_n\|x_n\|_X \le \|x\dg\|_X$.  The norms
  $\|x_n\|_X$ are then bounded: some tail is bounded by $\|x\dg\|_X + 1$, and the
  finitely many remaining terms are finite.  Suppose, for contradiction, that
  $x_n \not\to x\dg$.  Then there are $\varepsilon > 0$ and a subsequence with
  $\|x_{n_j} - x\dg\|_X \ge \varepsilon$ for all $j$.  That subsequence is bounded, so
  by \cref{thm_weak_compactness} it has a further subsequence, written $\{x_k\}$, with
  $x_k \wc \bar{x}$ for some $\bar{x} \in X$.

  \emph{$\bar{x} - x\dg \in \Ker(T)$.}  Weak continuity of the bounded operator $T$
  (\cref{lem_weak_seq_cont}) gives $Tx_k \wc T\bar{x}$.  On the other hand
  \cref{lem_lsq_comparison} gives $\|Tx_k - Tx\dg\|_Y \to 0$, hence $Tx_k \wc Tx\dg$
  (\cref{lem_strong_to_weak}).  Weak limits are unique (\cref{lem_weak_limit_unique}), so
  $T\bar{x} = Tx\dg$.

  \emph{$\bar{x} = x\dg$.}  By \cref{lem_mp_equations}, $x\dg = T\dg y \in \Ran(T\dg) =
  \Ker(T)^{\perp}$, so $\bar{x} - x\dg \perp x\dg$ and Pythagoras (\cref{lem_pythagoras})
  gives $\|\bar{x}\|_X^{2} = \|\bar{x} - x\dg\|_X^{2} + \|x\dg\|_X^{2}$.  Weak lower
  semicontinuity of the norm (\cref{lem_weak_lsc}) and the hypothesis give
  \[
    \|\bar{x}\|_X \le \liminf_{k}\|x_k\|_X \le \limsup_{k}\|x_k\|_X
      \le \limsup_{n}\|x_n\|_X \le \|x\dg\|_X ,
  \]
  whence $\|\bar{x} - x\dg\|_X^{2} = \|\bar{x}\|_X^{2} - \|x\dg\|_X^{2} \le 0$ and
  $\bar{x} = x\dg$.

  \emph{Strong convergence of the subsequence.}  The display above now reads
  $\|x\dg\|_X \le \liminf_k\|x_k\|_X \le \limsup_k\|x_k\|_X \le \|x\dg\|_X$, so
  $\|x_k\|_X \to \|x\dg\|_X = \|\bar{x}\|_X$.  Together with $x_k \wc \bar{x}$, the
  Radon--Riesz property (\cref{thm_radon_riesz}) gives $x_k \to \bar{x} = x\dg$ strongly,
  contradicting $\|x_k - x\dg\|_X \ge \varepsilon$.  Hence $x_n \to x\dg$.
\end{proof}


\begin{theorem}[Convergence of the least-squares projection]
  \label{thm_lsq_convergence}
  Let $y \in \Dom(T\dg)$ and suppose
  \[
    \limsup_{n \to \infty} \bigl\|(T_n\dg)^{*} x_n\bigr\|_Y < \infty .
  \]
  Then $x_n \to x\dg$.
\end{theorem}

\begin{proof}
  \emph{A representation of $x_n$.}  With $T_n$ and $x_n$ as in
  \cref{def_lsq_projection}, put $w_n := (T_n\dg)^{*}x_n \in Y$.  By
  \cref{lem_lsq_range_closed}, $\Ran(T_n)$ is closed, so \cref{lem_mp_adjoint}
  gives $(T_n\dg)^{*} = (T_n^{*})\dg$.  The operator $T_n^{*} = P_nT^{*}$ has range
  $\Ran(T_n^{*}) = \Ker(T_n)^{\perp}$, which is contained in the finite-dimensional ---
  hence closed --- space $X_n$; so $\Ran(T_n^{*})$ is closed, $\Dom((T_n^{*})\dg) = X$,
  and \cref{lem_mp_equations}(iv) applied to $T_n^{*}$ gives
  $T_n^{*}(T_n^{*})\dg = P_{\overline{\Ran(T_n^{*})}} = P_{\Ker(T_n)^{\perp}}$.  Since
  $x_n \in \Ker(T_n)^{\perp}$ by \eqref{eq_lsq_in_xn},
  \begin{equation}\label{eq_lsq_repr}
    T_n^{*}w_n = T_n^{*}(T_n^{*})\dg x_n = P_{\Ker(T_n)^{\perp}}x_n = x_n .
  \end{equation}

  \emph{The key inner product tends to $0$.}  By \eqref{eq_lsq_repr} and the definition of
  the adjoint,
  \[
    \langle x_n - x\dg, x_n\rangle_X = \langle x_n - x\dg, T_n^{*}w_n\rangle_X
      = \langle T_n(x_n - x\dg), w_n\rangle_Y ,
  \]
  and $T_n(x_n - x\dg) = Tx_n - TP_nx\dg = (Tx_n - Tx\dg) + T(\Id - P_n)x\dg$, using
  $T_nx_n = Tx_n$ from \eqref{eq_lsq_in_xn}.  Hence, by Cauchy--Schwarz
  (\cref{thm_cauchy_schwarz}) and \cref{lem_lsq_comparison},
  \begin{equation}\label{eq_lsq_cross}
    \begin{aligned}
      \bigl|\langle x_n - x\dg, x_n\rangle_X\bigr|
        &\le \bigl(\|Tx_n - Tx\dg\|_Y + \|T(\Id - P_n)x\dg\|_Y\bigr)\|w_n\|_Y \\
        &\le 2\|T\|_{\Ll(X,Y)}\,\|(\Id - P_n)x\dg\|_X\,\|w_n\|_Y .
    \end{aligned}
  \end{equation}
  By hypothesis $\limsup_n\|w_n\|_Y < \infty$, so $M := \sup_n\|w_n\|_Y < \infty$ (a tail
  is bounded and the finitely many remaining terms are finite), while
  $\|(\Id - P_n)x\dg\|_X \to 0$ by \cref{thm_nested_proj}.  Thus
  $\varepsilon_n := |\langle x_n - x\dg, x_n\rangle_X| \to 0$.

  \emph{Conclusion.}  Writing $a_n := \|x_n\|_X$ and using Cauchy--Schwarz once more,
  \[
    a_n^{2} = \langle x_n, x_n\rangle_X
      = \langle x_n - x\dg, x_n\rangle_X + \langle x\dg, x_n\rangle_X
      \le \varepsilon_n + \|x\dg\|_X\,a_n .
  \]
  The nonnegative number $a_n$ therefore satisfies
  $a_n^{2} - \|x\dg\|_Xa_n - \varepsilon_n \le 0$, so it lies below the larger root of the
  corresponding quadratic:
  \[
    a_n \;\le\; \tfrac12\Bigl(\|x\dg\|_X
      + \sqrt{\|x\dg\|_X^{2} + 4\varepsilon_n}\Bigr) .
  \]
  As $\varepsilon_n \to 0$, the right-hand side tends to $\|x\dg\|_X$, so
  $\limsup_n\|x_n\|_X \le \|x\dg\|_X$, and \cref{lem_lsq_convergence_criterion} gives
  $x_n \to x\dg$.
\end{proof}


\begin{theorem}[The condition forces a source condition]
  \label{thm_lsq_source_necessary}
  Let $K \in \Kk(X,Y)$ be compact, $y \in \Dom(K\dg)$, and suppose the hypothesis of
  \cref{thm_lsq_convergence} holds.  Then $x\dg \in \Ran(K^{*})$.
\end{theorem}

\begin{proof}
  Here $T = K$ in \cref{def_lsq_projection}, so $K_n = KP_n$ and $x_n = K_n\dg y$.  Put
  $w_n := (K_n\dg)^{*}x_n$, so that $K_n^{*}w_n = x_n$ by \eqref{eq_lsq_repr} in the
  proof of \cref{thm_lsq_convergence}, i.e.
  \begin{equation}\label{eq_lsq_split}
    P_nK^{*}w_n = x_n ,
  \end{equation}
  since $K_n^{*} = P_nK^{*}$.  The hypothesis makes $\{w_n\}$ bounded, so
  \cref{thm_weak_compactness} provides a subsequence $\{w_k\}$ and $\bar{w} \in Y$ with
  $w_k \wc \bar{w}$.

  \emph{$K^{*}w_k \to K^{*}\bar{w}$ strongly.}  $K$ is compact, hence so is $K^{*}$
  (\cref{thm_schauder}), and a compact operator is completely continuous
  (\cref{thm_complete_continuity}): it maps weakly convergent sequences to strongly
  convergent ones.

  \emph{Splitting off the projection.}  By \eqref{eq_lsq_split},
  \[
    K^{*}w_k = P_kK^{*}w_k + (\Id - P_k)K^{*}w_k
      = x_k + (\Id - P_k)(K^{*}w_k - K^{*}\bar{w}) + (\Id - P_k)K^{*}\bar{w} .
  \]
  For any $z \in X$ the decomposition $z = P_kz + (\Id - P_k)z$ is orthogonal
  (\cref{lem_proj_characterisation}), so $\|(\Id - P_k)z\|_X \le \|z\|_X$ by Pythagoras
  (\cref{lem_pythagoras}).  Hence the second term is bounded in norm by
  $\|K^{*}w_k - K^{*}\bar{w}\|_X$, which tends to $0$ by the previous paragraph, while the
  third term tends to $0$ because $P_k\bar{z} \to \bar{z}$ for the \emph{fixed} vector
  $\bar{z} := K^{*}\bar{w}$ (\cref{thm_nested_proj}).  Finally $x_k \to x\dg$ by
  \cref{thm_lsq_convergence}, whose hypothesis is assumed.

  Passing to the limit in the display therefore gives $K^{*}\bar{w} = x\dg$, i.e.\
  $x\dg \in \Ran(K^{*})$.
\end{proof}


\begin{theorem}[Rate for the least-squares projection]
  \label{thm_lsq_rate}
  Let $K \in \Kk(X,Y)$, $y \in \Dom(K\dg)$, and suppose the hypothesis of
  \cref{thm_lsq_convergence} holds.  Then there is
  $C = C(x\dg, y, \{X_n\}) > 0$ with
  \[
    \|x_n - x\dg\|_X \;\le\; C\,\bigl\|(\Id - P_n)K^{*}\bigr\|_{\Ll(Y,X)}
    \qquad \text{for all } n .
  \]
\end{theorem}

\begin{proof}
  Here $T = K$ in \cref{def_lsq_projection}.  By \cref{thm_lsq_source_necessary} there
  is $w \in Y$ with $x\dg = K^{*}w$, and as in the proof of \cref{thm_lsq_convergence},
  read with $T = K$, the vectors
  $w_n := (K_n\dg)^{*}x_n$ satisfy $M := \sup_n\|w_n\|_Y < \infty$ and the first
  inequality of \eqref{eq_lsq_cross}.

  \emph{A sharper form of \cref{lem_lsq_comparison}.}  The chain in the proof of
  \cref{lem_lsq_comparison}, stopped one step before the final crude estimate, reads
  \begin{equation}\label{eq_lsq_sharp}
    \|Kx_n - Kx\dg\|_Y \le \|K(\Id - P_n)x\dg\|_Y .
  \end{equation}

  \emph{The two inner products.}  Combining \eqref{eq_lsq_cross} with
  \eqref{eq_lsq_sharp} in place of the crude bound,
  \[
    \begin{aligned}
      \bigl|\langle x_n - x\dg, x_n\rangle_X\bigr|
        &\le \bigl(\|Kx_n - Kx\dg\|_Y + \|K(\Id-P_n)x\dg\|_Y\bigr)\|w_n\|_Y \\
        &\le 2M\,\|K(\Id - P_n)x\dg\|_Y ,
    \end{aligned}
  \]
  while the source representation and \eqref{eq_lsq_sharp} give
  \[
    \begin{aligned}
      \bigl|\langle x_n - x\dg, x\dg\rangle_X\bigr|
        &= \bigl|\langle K(x_n - x\dg), w\rangle_Y\bigr|
          \le \|Kx_n - Kx\dg\|_Y\,\|w\|_Y \\
        &\le \|K(\Id-P_n)x\dg\|_Y\,\|w\|_Y ,
    \end{aligned}
  \]
  by Cauchy--Schwarz (\cref{thm_cauchy_schwarz}).  Subtracting,
  \begin{equation}\label{eq_lsq_rate_sq}
    \|x_n - x\dg\|_X^{2}
      = \langle x_n - x\dg, x_n\rangle_X - \langle x_n - x\dg, x\dg\rangle_X
      \le \bigl(2M + \|w\|_Y\bigr)\,\|K(\Id - P_n)x\dg\|_Y .
  \end{equation}

  \emph{The projection defect.}  Since $P_n^{2} = P_n$ (\cref{lem_proj_nested_comp}) we have
  $(\Id-P_n)^{2} = \Id - P_n$, so $(\Id-P_n)x\dg = (\Id-P_n)(\Id-P_n)K^{*}w$ and
  \[
    \begin{aligned}
      \|K(\Id-P_n)x\dg\|_Y
        &\le \|K(\Id-P_n)\|_{\Ll(X,Y)}\,\|(\Id-P_n)K^{*}w\|_X \\
        &\le \|(\Id-P_n)K^{*}\|_{\Ll(Y,X)}^{2}\,\|w\|_Y ,
    \end{aligned}
  \]
  where $\|K(\Id-P_n)\|_{\Ll(X,Y)} = \|(K(\Id-P_n))^{*}\|_{\Ll(Y,X)} =
  \|(\Id-P_n)K^{*}\|_{\Ll(Y,X)}$ because $\|A\| = \|A^{*}\|$ and $(\Id-P_n)^{*} =
  \Id-P_n$ (\cref{prop_adjoint_props}, \cref{prop_proj_props}).

  \emph{Conclusion.}  Substituting into \eqref{eq_lsq_rate_sq} and taking square roots,
  \[
    \|x_n - x\dg\|_X \;\le\; C\,\|(\Id-P_n)K^{*}\|_{\Ll(Y,X)},
    \qquad C := 1 + \bigl((2M + \|w\|_Y)\|w\|_Y\bigr)^{1/2} > 0 ,
  \]
  for every $n$; the constant depends on $x\dg$, on $y$ and on the family $\{X_n\}$
  through $M$ and $w$, as the statement allows.
\end{proof}


\begin{definition}[Dual least-squares projection]
  \label{def_dual_lsq_projection}
  Let $\dim \overline{\Ran(T)} = \infty$ and let
  $Y_1 \subset Y_2 \subset \cdots \subset \Ker(T^{*})^{\perp} = \overline{\Ran(T)}$ be
  nested subspaces with $\dim Y_n = n$ and $\overline{\bigcup_n Y_n} =
  \Ker(T^{*})^{\perp}$.  Let $Q_n := P_{Y_n}$, put $T_n := Q_n T \in \Ll(X,Y)$, and for
  $y \in Y$ set $x_n := T_n\dg Q_n y$.  Put $X_n := T^{*}Y_n$ and $P_n := P_{X_n}$.
\end{definition}


\begin{lemma}[Discrete ranges are closed, dual method]
  \label{lem_dual_lsq_range_closed}
  Let $T \in \Ll(X,Y)$ and let $Y_1 \subset Y_2 \subset \cdots \subset Y$ be nested
  subspaces with $\dim Y_n = n$, and write $Q_n := P_{Y_n}$.  Then $\Ran(Q_nT)$ is closed
  for every $n$.
\end{lemma}

\begin{proof}
  $\Ran(Q_nT) \subseteq Y_n$, which is finite-dimensional and therefore closed
  (\cref{lem_findim_closed}); a subspace of a finite-dimensional space is
  finite-dimensional, so $\Ran(Q_nT)$ is closed.
\end{proof}


\begin{lemma}[The dual method is a projection]
  \label{lem_dual_lsq_formula}
  For $y \in \Dom(T\dg)$ one has $x_n = P_n x\dg$.
\end{lemma}

\begin{proof}
  All symbols are those of \cref{def_dual_lsq_projection}.

  \emph{Standing facts about $T_n = Q_nT$.}  Its range satisfies $\Ran(T_n) \subseteq Y_n$
  with $\dim Y_n = n$, so $\Ran(T_n)$ is finite-dimensional, hence closed
  (\cref{lem_findim_closed}); by \cref{thm_mp_continuous_iff_closed_range} and
  \cref{lem_mp_domain_dense}, $T_n\dg \in \Ll(Y,X)$ with $\Dom(T_n\dg) = Y$.  Next,
  $T_n^{*} = (Q_nT)^{*} = T^{*}Q_n$ (\cref{prop_adjoint_props}, \cref{prop_proj_props}), so
  $\Ran(T_n^{*}) = T^{*}Q_n(Y) = T^{*}Y_n = X_n$, a finite-dimensional and hence closed
  subspace.  Therefore, by \cref{thm_four_subspaces} and \cref{lem_mp_equations},
  \begin{equation}\label{eq_dual_xn}
    \Ker(T_n)^{\perp} = \overline{\Ran(T_n^{*})} = X_n ,
    \qquad \Ran(T_n\dg) = \Ker(T_n)^{\perp} = X_n .
  \end{equation}

  \emph{The datum is exact for $T_n$.}  Let $P_R$ be the orthogonal projection onto
  $\overline{\Ran(T)} = \Ker(T^{*})^{\perp}$ (\cref{thm_four_subspaces}).  For $z \in Y$ the
  vector $(\Id - P_R)z$ lies in $\overline{\Ran(T)}^{\perp} = \Ker(T^{*})$
  (\cref{lem_proj_characterisation}), hence is orthogonal to
  $Y_n \subseteq \Ker(T^{*})^{\perp}$, so $Q_n(\Id - P_R)z = 0$
  (\cref{lem_proj_characterisation}).  Thus $Q_n = Q_nP_R$, and \cref{lem_mp_equations}(iv)
  with $y \in \Dom(T\dg)$ gives
  \begin{equation}\label{eq_dual_datum}
    Q_ny = Q_nP_Ry = Q_nTT\dg y = Q_nTx\dg = T_nx\dg .
  \end{equation}

  \emph{$P_nx\dg$ is a least squares solution.}  By \eqref{eq_dual_xn},
  $\Ker(T_n) = X_n^{\perp}$, so $T_n(\Id - P_n)x = 0$ for every $x \in X$, i.e.\
  $T_nP_n = T_n$.  With \eqref{eq_dual_datum} this gives
  $T_nP_nx\dg = T_nx\dg = Q_ny$, so $P_nx\dg$ solves $T_nx = Q_ny$ exactly and is in
  particular a least squares solution of that equation.

  \emph{Identification.}  By \cref{thm_mp_minnorm}(iii), applied to $T_n$ and the datum
  $Q_ny \in Y = \Dom(T_n\dg)$, the set of least squares solutions is
  $x_n + \Ker(T_n)$ with $x_n = T_n\dg Q_ny$.  Hence $P_nx\dg = x_n + k$ for some
  $k \in \Ker(T_n) = X_n^{\perp}$.  But $P_nx\dg \in X_n$ and $x_n \in X_n$ by
  \eqref{eq_dual_xn}, so $k \in X_n \cap X_n^{\perp} = \{0\}$
  (\cref{thm_orth_decomp}).  Therefore $x_n = P_nx\dg$.
\end{proof}


\begin{theorem}[Convergence of the dual least-squares projection]
  \label{thm_dual_lsq_convergence}
  For every $y \in \Dom(T\dg)$ one has $x_n \to x\dg$.  No condition analogous to that of
  \cref{thm_lsq_convergence} is needed.
\end{theorem}

\begin{proof}
  All symbols are those of \cref{def_dual_lsq_projection}.

  \emph{The subspaces $X_n$ are nested.}  $Y_n \subseteq Y_{n+1}$ gives
  $X_n = T^{*}Y_n \subseteq T^{*}Y_{n+1} = X_{n+1}$, and each $X_n$ is
  finite-dimensional, hence closed (\cref{lem_findim_closed}).

  \emph{Their union is dense in $\Ker(T)^{\perp}$.}  Put $U := \overline{\bigcup_nX_n}$, a
  closed subspace.  Each $X_n = T^{*}Y_n \subseteq \Ran(T^{*}) \subseteq
  \overline{\Ran(T^{*})} = \Ker(T)^{\perp}$ (\cref{thm_four_subspaces}), and the latter is
  closed, so $U \subseteq \Ker(T)^{\perp}$.  Conversely, let $z \in Y$ and split
  $z = P_{\Ker(T^{*})^{\perp}}z + P_{\Ker(T^{*})}z$ (\cref{thm_orth_decomp}); the second
  summand is annihilated by $T^{*}$, so $T^{*}z = T^{*}P_{\Ker(T^{*})^{\perp}}z$.  Since
  $\overline{\bigcup_nY_n} = \Ker(T^{*})^{\perp}$ there are
  $z_j \in \bigcup_nY_n$ with $z_j \to P_{\Ker(T^{*})^{\perp}}z$, and then
  $T^{*}z_j \in \bigcup_nX_n$ with $T^{*}z_j \to T^{*}z$ by continuity of $T^{*}$.  Hence
  $\Ran(T^{*}) \subseteq U$, and as $U$ is closed,
  $\Ker(T)^{\perp} = \overline{\Ran(T^{*})} \subseteq U$.  Altogether
  $U = \Ker(T)^{\perp}$.

  \emph{Conclusion.}  The $X_n$ are nested closed subspaces whose union is dense in the
  closed subspace $U = \Ker(T)^{\perp}$, so $P_nz \to P_Uz$ for every $z \in X$
  (\cref{thm_nested_proj_proper}); note that $U$ is in general a \emph{proper} subspace of
  $X$, which is why the version of the theorem for a proper limit subspace is the one
  needed.  Now $x\dg = T\dg y \in \Ran(T\dg) = \Ker(T)^{\perp} = U$ by
  \cref{lem_mp_equations}, so $P_Ux\dg = x\dg$, and \cref{lem_dual_lsq_formula} gives
  \[
    x_n = P_nx\dg \longrightarrow P_Ux\dg = x\dg . \qedhere
  \]
\end{proof}


\begin{theorem}[Rate for the dual least-squares projection]
  \label{thm_dual_lsq_rate}
  Let $y \in \Dom(T\dg)$ and suppose $x\dg \in \Ran(T^{*})$, say $x\dg = T^{*}w$.  Then
  \[
    \|x_n - x\dg\|_X \;\le\; \|w\|_Y\,\bigl\|(\Id - P_n)T^{*}\bigr\|_{\Ll(Y,X)}
    \qquad \text{for all } n .
  \]
\end{theorem}

\begin{proof}
  With $P_n$ and $x_n$ as in \cref{def_dual_lsq_projection},
  \cref{lem_dual_lsq_formula} and $x\dg = T^{*}w$ give
  \[
    x_n - x\dg = P_nx\dg - x\dg = -(\Id - P_n)x\dg = -(\Id - P_n)T^{*}w ,
  \]
  so taking norms and using the definition of the operator norm,
  \[
    \|x_n - x\dg\|_X = \|(\Id - P_n)T^{*}w\|_X
      \le \bigl\|(\Id - P_n)T^{*}\bigr\|_{\Ll(Y,X)}\,\|w\|_Y . \qedhere
  \]
\end{proof}


\begin{lemma}[Norm of the discrete pseudo-inverse]
  \label{lem_discrete_opnorm}
  Let $\mu_1 \ge \cdots \ge \mu_n > 0$ be the singular values of the rank-$n$ operator
  $T_n = Q_n T$.  Then $\|T_n\dg\|_{\Ll(Y,X)} = \mu_n^{-1}$.
\end{lemma}

\begin{proof}
  All symbols are those of \cref{def_dual_lsq_projection}.

  \emph{$T_n$ has rank $n$.}  As in the proof of \cref{lem_dual_lsq_formula},
  $\Ker(T_n)^{\perp} = \Ran(T_n^{*}) = T^{*}Y_n = X_n$, $T_nP_n = T_n$, and
  $T_n\dg \in \Ll(Y,X)$ with $\Dom(T_n\dg) = Y$.  The map $T^{*}$ is injective on $Y_n$:
  if $z \in Y_n$ and $T^{*}z = 0$ then $z \in \Ker(T^{*}) \cap \Ker(T^{*})^{\perp} = \{0\}$
  (\cref{thm_orth_decomp}), using $Y_n \subseteq \Ker(T^{*})^{\perp}$.  Hence
  $\dim X_n = \dim Y_n = n$.  Moreover $T_n$ is injective on $\Ker(T_n)^{\perp} = X_n$
  (its kernel there is $\Ker(T_n) \cap \Ker(T_n)^{\perp} = \{0\}$) and
  $\Ran(T_n) = T_n(X_n)$ because $T_nP_n = T_n$, so
  $\dim\Ran(T_n) = \dim X_n = n$; in particular $T_n$ has finite rank $n$.

  \emph{The singular system.}  Being of finite rank, $T_n$ is compact
  (\cref{lem_finrank_compact}) and has a singular system
  $\{(\mu_k,\tilde{u}_k,\tilde{v}_k)\}_{k \in J}$ (\cref{thm_svd}), where
  $\{\tilde{v}_k\}_{k \in J}$ is an orthonormal basis of
  $\overline{\Ran(T_n^{*})} = X_n$; since $\dim X_n = n$ this forces
  $J = \{1,\dots,n\}$, and all singular values are positive, so
  $\mu_1 \ge \cdots \ge \mu_n > 0$.

  \emph{The upper bound.}  For $z \in Y = \Dom(T_n\dg)$, \cref{cor_picard_domain} gives
  \[
    T_n\dg z = \sum_{k=1}^{n}\mu_k^{-1}\langle \tilde{u}_k, z\rangle_Y\,\tilde{v}_k ,
    \qquad\text{so}\qquad
    \|T_n\dg z\|_X^{2} = \sum_{k=1}^{n}\mu_k^{-2}|\langle\tilde{u}_k,z\rangle_Y|^{2}
  \]
  by orthonormality of $\{\tilde{v}_k\}$ and Pythagoras (\cref{lem_pythagoras}).  Since
  $\mu_k \ge \mu_n$ for $k \le n$, i.e.\ $\mu_k^{-2} \le \mu_n^{-2}$, Bessel's inequality
  (\cref{thm_bessel}) for the orthonormal system $\{\tilde{u}_k\}$ gives
  \[
    \|T_n\dg z\|_X^{2} \le \mu_n^{-2}\sum_{k=1}^{n}|\langle\tilde{u}_k,z\rangle_Y|^{2}
      \le \mu_n^{-2}\|z\|_Y^{2} ,
  \]
  so $\|T_n\dg\|_{\Ll(Y,X)} \le \mu_n^{-1}$.

  \emph{The bound is attained.}  Taking $z := \tilde{u}_n$, a unit vector, the same
  formula gives $T_n\dg\tilde{u}_n = \mu_n^{-1}\tilde{v}_n$, of norm $\mu_n^{-1}$; hence
  $\|T_n\dg\|_{\Ll(Y,X)} \ge \mu_n^{-1}$ and equality holds.
\end{proof}


\begin{theorem}[Noisy data and an a priori stopping rule]
  \label{thm_dual_lsq_noisy}
  Let $y \in \Dom(T\dg)$, $y^{\delta} \in B_{\delta}(y)$, and set
  $x^{\delta}_n := T_n\dg Q_n y^{\delta}$.  If $n(\delta) \to \infty$ and
  $\delta/\mu_{n(\delta)} \to 0$ as $\delta \to 0$, then
  $x^{\delta}_{n(\delta)} \to x\dg$.
\end{theorem}

\begin{proof}
  With $T_n = Q_nT$ and $Q_n$ as in \cref{def_dual_lsq_projection}, put
  $R_n := T_n\dg Q_n \in \Ll(Y,X)$, so that $x_n = R_ny$ and
  $x^{\delta}_n = R_ny^{\delta}$.

  \emph{Error decomposition.}  By linearity of $R_n$ and the triangle inequality --- this
  is \cref{lem_error_decomposition} for the family $\{R_n\}$ ---
  \[
    \|x^{\delta}_n - x\dg\|_X \le \|R_n(y^{\delta} - y)\|_X + \|R_ny - x\dg\|_X
      \le \|R_n\|_{\Ll(Y,X)}\,\delta + \|x_n - x\dg\|_X .
  \]

  \emph{The data error.}  For any $z \in Y$ the decomposition $z = Q_nz + (\Id - Q_n)z$ is
  orthogonal (\cref{lem_proj_characterisation}), so $\|Q_nz\|_Y \le \|z\|_Y$ by Pythagoras
  (\cref{lem_pythagoras}), i.e.\ $\|Q_n\|_{\Ll(Y,Y)} \le 1$.  With
  \cref{lem_discrete_opnorm},
  \[
    \|R_n\|_{\Ll(Y,X)} \le \|T_n\dg\|_{\Ll(Y,X)}\|Q_n\|_{\Ll(Y,Y)} \le \mu_n^{-1} .
  \]

  \emph{Passing to the limit.}  Combining the two displays at $n = n(\delta)$,
  \[
    \|x^{\delta}_{n(\delta)} - x\dg\|_X
      \;\le\; \frac{\delta}{\mu_{n(\delta)}} + \|x_{n(\delta)} - x\dg\|_X .
  \]
  The first term tends to $0$ as $\delta \to 0$ by hypothesis.  For the second, let
  $\varepsilon > 0$; by \cref{thm_dual_lsq_convergence} there is $N$ with
  $\|x_n - x\dg\|_X < \varepsilon$ for all $n \ge N$, and since $n(\delta) \to \infty$
  there is $\delta_0 > 0$ with $n(\delta) \ge N$ for all $\delta \in (0,\delta_0)$; hence
  $\|x_{n(\delta)} - x\dg\|_X < \varepsilon$ there.  Both terms therefore tend to $0$, and
  $x^{\delta}_{n(\delta)} \to x\dg$ as $\delta \to 0$.
\end{proof}


\begin{theorem}[Discrete singular values undercut the true ones]
  \label{thm_mu_le_sigma}
  Let $K \in \Kk(X,Y)$ have a singular system $\{(\sigma_n,u_n,v_n)\}$, let $n \ge 1$, and
  let $Y_n \subseteq \Ker(K^{*})^{\perp}$ with $\dim Y_n = n$.  Then $\mu_n \le \sigma_n$,
  with equality in the case $Y_n = \mathrm{span}\{u_1,\dots,u_n\}$.
\end{theorem}
\begin{proof}
  Write $K_n := Q_nK$ and $\mu_n$ for the smallest singular value of $K_n$, as in
  \cref{def_dual_lsq_projection}, and let $P$ denote the orthogonal projection onto
  $\overline{\Ran(K)} = \Ker(K^{*})^{\perp}$ (\cref{thm_four_subspaces}), of which
  $\{u_k\}_{k \in I}$ is an orthonormal basis.  Since $Y_n \subseteq \Ker(K^{*})^{\perp}$
  has dimension $n$, that basis has at least $n$ elements, so $\sigma_n$ is defined.

  \emph{$\{\tilde{u}_k\}_{k=1}^{n}$ is an orthonormal basis of $Y_n$.}  As in the proof of
  \cref{lem_discrete_opnorm} --- whose hypothesis $Y_n \subseteq \Ker(K^{*})^{\perp}$ is
  assumed here --- $K_n$ has rank $n$ and a singular system
  $\{(\mu_k,\tilde{u}_k,\tilde{v}_k)\}_{k=1}^{n}$ with
  $\mu_1 \ge \cdots \ge \mu_n > 0$, where $\{\tilde{u}_k\}_{k=1}^{n}$ is an orthonormal
  basis of $\overline{\Ran(K_n)} = \Ran(K_n)$.  As $\Ran(K_n) \subseteq Y_n$ and both have
  dimension $n$, they coincide; so $\{\tilde{u}_k\}_{k=1}^{n}$ is an orthonormal basis of
  $Y_n$.

  \emph{A lower bound on $Y_n$.}  For $z \in Y_n$ we have $Q_nz = z$, so
  $K_n^{*}z = K^{*}Q_nz = K^{*}z$.  Expanding $z = \sum_{k\le n}\langle\tilde{u}_k,
  z\rangle\tilde{u}_k$ (\cref{thm_proj_expansion}, with $Q_nz = z$) and using
  $K_n^{*}\tilde{u}_k = \mu_k\tilde{v}_k$ with $\{\tilde{v}_k\}$ orthonormal,
  \begin{equation}\label{eq_mu_lower}
    \|K^{*}z\|_X^{2} = \sum_{k=1}^{n}\mu_k^{2}|\langle\tilde{u}_k,z\rangle_Y|^{2}
      \ge \mu_n^{2}\sum_{k=1}^{n}|\langle\tilde{u}_k,z\rangle_Y|^{2}
      = \mu_n^{2}\|z\|_Y^{2} ,
  \end{equation}
  by $\mu_k \ge \mu_n$ and, for the first and the last equality, the Parseval norm
  identity of \cref{lem_parseval_series} for the orthonormal systems $\{\tilde{v}_k\}$
  and $\{\tilde{u}_k\}$; \cref{thm_proj_expansion} supplies the expansion of $z$, not that
  identity.

  \emph{An upper bound off the first $n-1$ singular vectors.}  Put
  $U_{n-1} := \mathrm{span}\{u_1,\dots,u_{n-1}\}$, of dimension exactly $n-1$: the
  $u_1,\dots,u_{n-1}$ exist because $I$ is an initial segment with $|I| \ge n$
  (\cref{def_singular_system}), and they are independent, being orthonormal.  With
  $\dim U_{n-1} = n-1 < n = \dim Y_n$, which is what \cref{lem_dim_count} asks for,
  that node gives
  $0 \ne \bar{z} \in Y_n \cap U_{n-1}^{\perp}$, which we normalise to
  $\|\bar{z}\|_Y = 1$.
  Then $\bar{z} \in Y_n \subseteq \overline{\Ran(K)}$ has
  $\langle u_k,\bar{z}\rangle = 0$ for $k < n$.  The expansion of $K^{*}$ in the singular
  system,
  \[
    K^{*}\bar{z} = \sum_{k \in I}\sigma_k\langle u_k,\bar{z}\rangle_Y\,v_k
      = \sum_{k \ge n}\sigma_k\langle u_k,\bar{z}\rangle_Y\,v_k ,
  \]
  is \cref{lem_svd_companion}; with $\{v_k\}$ orthonormal, \cref{lem_parseval_series} and
  $\sigma_k \le \sigma_n$ for $k \ge n$ give
  \begin{equation}\label{eq_mu_upper}
    \|K^{*}\bar{z}\|_X^{2} = \sum_{k \ge n}\sigma_k^{2}|\langle u_k,\bar{z}\rangle_Y|^{2}
      \le \sigma_n^{2}\sum_{k \ge n}|\langle u_k,\bar{z}\rangle_Y|^{2}
      = \sigma_n^{2}\|\bar{z}\|_Y^{2} = \sigma_n^{2} ,
  \end{equation}
  the last equality by \cref{thm_proj_expansion} for $\bar{z} \in \overline{\Ran(K)}$ and
  \cref{lem_parseval_series} for $\{u_k\}$.

  \emph{The inequality.}  Applying \eqref{eq_mu_lower} to $z := \bar{z}$ and combining
  with \eqref{eq_mu_upper} gives $\mu_n^{2} \le \|K^{*}\bar{z}\|_X^{2} \le \sigma_n^{2}$,
  hence $\mu_n \le \sigma_n$.

  \emph{The equality case.}  Let $Y_n = U_n := \mathrm{span}\{u_1,\dots,u_n\}$, which is
  contained in $\overline{\Ran(K)} = \Ker(K^{*})^{\perp}$ and has dimension $n$.  For a
  unit vector $z \in U_n$ the computation \eqref{eq_mu_upper}, now with the expansion
  supported in $\{1,\dots,n\}$ and $\sigma_k \ge \sigma_n$ there, gives
  \[
    \|K^{*}z\|_X^{2} = \sum_{k=1}^{n}\sigma_k^{2}|\langle u_k,z\rangle_Y|^{2}
      \ge \sigma_n^{2}\|z\|_Y^{2} = \sigma_n^{2} .
  \]
  Taking $z := \tilde{u}_n \in Y_n$, which is a unit vector with
  $\|K^{*}\tilde{u}_n\|_X = \|K_n^{*}\tilde{u}_n\|_X = \mu_n\|\tilde{v}_n\|_X = \mu_n$,
  yields $\mu_n \ge \sigma_n$; with $\mu_n \le \sigma_n$ this gives $\mu_n = \sigma_n$.
\end{proof}


\begin{theorem}[Projection defect of the dual subspaces]
  \label{thm_projection_defect}
  Let $K \in \Kk(X,Y)$ have singular system $\{(\sigma_n,u_n,v_n)\}$ with
  $\mathrm{rank}\,K \ge n+1$, let $Y_n \subseteq Y$ be finite dimensional with
  $\dim Y_n = n$, and put
  $X_n := K^{*}Y_n$, $P_n := P_{X_n}$.  Then
  \[
    \bigl\|(\Id - P_n)K^{*}\bigr\|_{\Ll(Y,X)} \;\ge\; \sigma_{n+1} ,
  \]
  with equality when $Y_n = \mathrm{span}\{u_1,\dots,u_n\}$.
\end{theorem}
\begin{proof}
  The symbols $X_n = K^{*}Y_n$ and $P_n = P_{X_n}$ are those of
  \cref{def_dual_lsq_projection}.  Since $K$ restricted to $\Ker(K)^{\perp}$ is
  injective with range $\Ran(K)$, and
  $\{v_k\}_{k \in I}$ is an orthonormal basis of $\Ker(K)^{\perp}$, the hypothesis
  $\mathrm{rank}\,K \ge n+1$ gives $|I| \ge n+1$, so $\sigma_1,\dots,\sigma_{n+1}$ and
  $v_1,\dots,v_{n+1}$ exist.  The space $X_n = K^{*}Y_n$ is the image of an
  $n$-dimensional space, hence $\dim X_n \le n$ and $X_n$ is closed
  (\cref{lem_findim_closed}), so $P_n$ is defined.  Note also that $P_nx = 0$ if and only if
  $x \perp X_n$ (\cref{lem_proj_characterisation}), i.e.\ $\Ker(P_n) = X_n^{\perp}$.

  \emph{Reduction to a supremum.}  Since $(\Id - P_n)^{*} = \Id - P_n$
  (\cref{prop_proj_props}) we have $(K(\Id-P_n))^{*} = (\Id - P_n)K^{*}$, and
  $\|A\| = \|A^{*}\|$ (\cref{prop_adjoint_props}), so
  \begin{equation}\label{eq_pd_norms}
    \bigl\|(\Id - P_n)K^{*}\bigr\|_{\Ll(Y,X)} = \|K(\Id - P_n)\|_{\Ll(X,Y)}
      = \sup\bigl\{\|Kx\|_Y : x \in X_n^{\perp},\ \|x\|_X \le 1\bigr\} .
  \end{equation}
  For the last equality: if $\|x\|_X \le 1$ then $(\Id-P_n)x \in X_n^{\perp}$
  (\cref{lem_proj_characterisation}) with $\|(\Id-P_n)x\|_X \le \|x\|_X \le 1$ by Pythagoras
  (\cref{lem_pythagoras}), so the left supremum is at most the right one; conversely every
  $x \in X_n^{\perp}$ with $\|x\|_X \le 1$ satisfies $(\Id - P_n)x = x$ and so is
  admissible on the left.

  \emph{The lower bound.}  Let $V := \mathrm{span}\{v_1,\dots,v_{n+1}\}$, of dimension
  $n+1$.  Since $\dim X_n \le n$, choose any $n$-dimensional subspace $U \subseteq X$ with
  $X_n \subseteq U$ --- possible because $\dim X \ge \dim V = n+1$.  Now
  \cref{lem_dim_count}, which asks only for $\dim U = n < n+1 = \dim V$, gives
  $0 \ne \bar{x} \in V \cap U^{\perp} \subseteq V \cap X_n^{\perp} = V \cap \Ker(P_n)$.
  (The interposed $U$ is no longer forced: that node was once stated with
  $\dim U = \dim V - 1$ as an equality, which is why an $n$-dimensional $U \supseteq X_n$
  had to be produced rather than applying the node to $X_n$ directly, whose dimension may
  be strictly less than $n$.  Applying it to $X_n$ would now be the shorter route.)
  Rescale to $\|\bar{x}\|_X = 1$ and write
  $\bar{x} = \sum_{k=1}^{n+1}c_kv_k$, so that $\sum_k|c_k|^{2} = 1$.  Then
  $K\bar{x} = \sum_{k=1}^{n+1}c_k\sigma_ku_k$, so by orthonormality of $\{u_k\}$ and
  $\sigma_k \ge \sigma_{n+1}$ for $k \le n+1$,
  \[
    \|K\bar{x}\|_Y^{2} = \sum_{k=1}^{n+1}\sigma_k^{2}|c_k|^{2}
      \ge \sigma_{n+1}^{2}\sum_{k=1}^{n+1}|c_k|^{2} = \sigma_{n+1}^{2} .
  \]
  With \eqref{eq_pd_norms} this gives
  $\|(\Id - P_n)K^{*}\|_{\Ll(Y,X)} \ge \|K\bar{x}\|_Y \ge \sigma_{n+1}$.

  \emph{The equality case.}  Let $Y_n = \mathrm{span}\{u_1,\dots,u_n\}$.  Then
  \[
    X_n = K^{*}Y_n = \mathrm{span}\{\sigma_1v_1,\dots,\sigma_nv_n\}
      = \mathrm{span}\{v_1,\dots,v_n\} ,
  \]
  because $K^{*}u_k = \sigma_kv_k$ (\cref{def_singular_system}) and every $\sigma_k > 0$.
  Let $x \in X_n^{\perp}$ with $\|x\|_X \le 1$; then $\langle v_k,x\rangle = 0$ for
  $k \le n$, so the singular value expansion of $K$ (\cref{thm_svd}) reads
  \[
    Kx = \sum_{k \in I}\sigma_k\langle v_k,x\rangle_X\,u_k
      = \sum_{k \ge n+1}\sigma_k\langle v_k,x\rangle_X\,u_k ,
  \]
  whence, by \cref{lem_parseval_series} for the orthonormal system $\{u_k\}$,
  \[
    \|Kx\|_Y^{2} = \sum_{k \ge n+1}\sigma_k^{2}|\langle v_k,x\rangle_X|^{2}
      \le \sigma_{n+1}^{2}\sum_{k \ge n+1}|\langle v_k,x\rangle_X|^{2}
      \le \sigma_{n+1}^{2}\|x\|_X^{2} \le \sigma_{n+1}^{2} ,
  \]
  using $\sigma_k \le \sigma_{n+1}$ for $k \ge n+1$ and Bessel (\cref{thm_bessel}).  By
  \eqref{eq_pd_norms}, $\|(\Id-P_n)K^{*}\|_{\Ll(Y,X)} \le \sigma_{n+1}$, and with the
  lower bound already proved, equality holds.
\end{proof}


\begin{corollary}[Best rate for the dual least-squares projection]
  \label{cor_dual_lsq_best_rate}
  Let $K \in \Kk(X,Y)$, $y \in \Dom(K\dg)$, $y^{\delta} \in B_{\delta}(y)$, and
  $x\dg = K^{*}w \in \Ran(K^{*})$.  Then
  \[
    \|x^{\delta}_n - x\dg\|_X \;\le\;
      \|w\|_Y\,\bigl\|(\Id - P_n)K^{*}\bigr\|_{\Ll(Y,X)} \;+\; \frac{\delta}{\mu_n} .
  \]
  For the choice $Y_n = \mathrm{span}\{u_1,\dots,u_n\}$ this reads
  \[
    \|x^{\delta}_n - x\dg\|_X \;\le\; \|w\|_Y\,\sigma_{n+1} + \frac{\delta}{\sigma_n} ,
  \]
  and by \cref{thm_mu_le_sigma} and \cref{thm_projection_defect} no other admissible
  family $\{Y_n\}$ improves either term: $\mu_n \le \sigma_n$ and
  $\|(\Id-P_n)K^{*}\| \ge \sigma_{n+1}$ always.
\end{corollary}

\begin{proof}
  Throughout, $T = K$ in \cref{def_dual_lsq_projection}, so $\{Y_n\}$ is nested in
  $\Ker(K^{*})^{\perp} = \overline{\Ran(K)}$ with $\dim Y_n = n$ and dense union, and
  $\dim\overline{\Ran(K)} = \infty$; since $\{u_k\}_{k \in I}$ is an orthonormal basis of
  $\overline{\Ran(K)}$, this makes $I$ infinite, so every $\sigma_n$ is defined and
  $\mathrm{rank}\,K \ge n+1$ for every $n$.

  \emph{The estimate.}  By the triangle inequality,
  \[
    \|x^{\delta}_n - x\dg\|_X \le \|x_n - x\dg\|_X + \|x^{\delta}_n - x_n\|_X .
  \]
  The first term is at most $\|w\|_Y\|(\Id-P_n)K^{*}\|_{\Ll(Y,X)}$ by
  \cref{thm_dual_lsq_rate}, whose hypothesis $x\dg = K^{*}w$ is assumed.  For the second,
  $x^{\delta}_n - x_n = K_n\dg Q_n(y^{\delta} - y)$, and as in the proof of
  \cref{thm_dual_lsq_noisy} we have $\|Q_n\|_{\Ll(Y,Y)} \le 1$ and
  $\|K_n\dg\|_{\Ll(Y,X)} = \mu_n^{-1}$ by \cref{lem_discrete_opnorm}, so that term is at
  most $\delta/\mu_n$.  This is the first display.

  \emph{The truncated singular value decomposition.}  For
  $Y_n = \mathrm{span}\{u_1,\dots,u_n\}$ we have $Y_n \subseteq \overline{\Ran(K)} =
  \Ker(K^{*})^{\perp}$ and $\dim Y_n = n$, so \cref{thm_mu_le_sigma} applies in its
  equality case and gives $\mu_n = \sigma_n$; and \cref{thm_projection_defect} applies in
  its equality case --- its hypothesis $\mathrm{rank}\,K \ge n+1$ holds by the first
  paragraph --- and gives $\|(\Id - P_n)K^{*}\|_{\Ll(Y,X)} = \sigma_{n+1}$.  Substituting
  both into the first display yields the second.

  \emph{Optimality of the two terms.}  Let $\{Y_n\}$ be any admissible family.  Its
  members satisfy $Y_n \subseteq \Ker(K^{*})^{\perp}$ with $\dim Y_n = n$, so
  \cref{thm_mu_le_sigma} gives $\mu_n \le \sigma_n$ and hence
  $\delta/\mu_n \ge \delta/\sigma_n$, while \cref{thm_projection_defect} gives
  $\|(\Id-P_n)K^{*}\|_{\Ll(Y,X)} \ge \sigma_{n+1}$ and hence
  $\|w\|_Y\|(\Id-P_n)K^{*}\|_{\Ll(Y,X)} \ge \|w\|_Y\sigma_{n+1}$.  So neither term of the
  first display is ever smaller than the corresponding term of the second.  This is a
  comparison of the two \emph{bounds}: because \cref{thm_projection_defect} bounds
  $\|(\Id-P_n)K^{*}\|$ from below, $\sigma_{n+1}$ may not be substituted for it in the
  upper bound of the first display, so the rate $\|w\|_Y\sigma_{n+1} + \delta/\sigma_n$ is
  asserted only for the truncated singular value decomposition.
\end{proof}



\chapter{General source conditions}
\label{ch_19_general_source_conditions}

\noindent\emph{Background.} Replacing the Hölder scale $\lambda \mapsto \lambda^{\nu/2}$
by an arbitrary index function $\varphi$ generalizes the power-type source sets of
\cref{ch_12_picard} to conditions adapted to the actual decay of the singular values.
Logarithmic source conditions, needed for severely (exponentially) ill-posed problems
where every Hölder rate is unattainably slow, are due to \citet{Hohage2000}. That
essentially every $x\dg$ satisfies some source condition, for a suitable index function,
is the general converse result of \citet{MatheHofmann2008}.

\begin{assumptions}
  \label{asm_general_source_conditions}
  $K \in \Kk(X,Y)$ is compact with singular system
$\{(\sigma_n,v_n,u_n)\}_{n \in I}$ and $\kappa := \|K\|^{2}$, as in the chapters above.
$\lambda$ ranges over $(0,\kappa]$ and is the spectral variable of $K^{*}K$, so that
$\varphi(K^{*}K)x = \sum_{n \in I}\varphi(\sigma_n^{2})\langle v_n,x\rangle_X v_n$ by the
functional calculus \cref{def_fun_calc}.  Where a statement of Chapters 2--5 is recovered as
an instance, that is said explicitly rather than left to be noticed.
\end{assumptions}

\section{Index functions and the source sets they define}

\begin{definition}[Index function]
  \label{def_index_function}
  A function $\varphi \colon [0,\kappa] \to [0,\infty)$ is an \emph{index function} if it
  is continuous and strictly increasing with $\varphi(0) = 0$.
\end{definition}


\begin{definition}[General source set]
  \label{def_general_source}
  For an index function $\varphi$ and $\rho > 0$ put
  \[
    X_{\varphi,\rho} \;:=\; \{\, \varphi(K^{*}K)\,w \;:\; w \in X,\ \|w\|_X \le \rho \,\} .
  \]
\end{definition}


\begin{lemma}[The H\"older scale is an instance]
  \label{lem_holder_is_index}
  For $\nu > 0$ the function $\varphi_{\nu}(\lambda) := \lambda^{\nu/2}$ is an index
  function, and $X_{\varphi_{\nu},\rho} = X_{\nu,\rho}$ in the sense of
  \cref{def_source_set}.
\end{lemma}
\begin{proof}
  Continuity and strict monotonicity of $\lambda \mapsto \lambda^{\nu/2}$ on $[0,\kappa]$
  hold for $\nu > 0$, and $0^{\nu/2} = 0$, so $\varphi_{\nu}$ is an index function by
  \cref{def_index_function}.  For the second claim, $|K| = (K^{*}K)^{1/2}$ acts on the
  singular basis by $|K|v_n = \sigma_n v_n$, so $|K|^{\nu}$ and $\varphi_{\nu}(K^{*}K)$
  both send $v_n$ to $\sigma_n^{\nu}v_n$ and both annihilate $\Ker(K)$.  Two bounded
  operators agreeing on an orthonormal basis of $\Ker(K)^{\perp}$ and on $\Ker(K)$ agree
  on $X$, by \cref{thm_onb} and \cref{thm_orth_decomp}.  Hence the two sets
  $\{|K|^{\nu}w : \|w\| \le \rho\}$ and $\{\varphi_{\nu}(K^{*}K)w : \|w\| \le \rho\}$
  coincide.
\end{proof}


\begin{definition}[Logarithmic index function]
  \label{def_log_index}
  For $p > 0$ put
  \[
    \varphi^{\log}_{p}(\lambda) \;:=\;
      \Bigl(\log\frac{e\kappa}{\lambda}\Bigr)^{-p}
      \quad \text{for } \lambda \in (0,\kappa],
    \qquad \varphi^{\log}_{p}(0) := 0 .
  \]
\end{definition}


\begin{lemma}[The logarithmic scale is an index function]
  \label{lem_log_is_index}
  $\varphi^{\log}_{p}$ is an index function for every $p > 0$.
\end{lemma}
\begin{proof}
  Write $g(\lambda) := \log(e\kappa/\lambda) = 1 + \log(\kappa/\lambda)$ for
  $\lambda \in (0,\kappa]$.  Then $g$ is continuous and strictly decreasing on
  $(0,\kappa]$, with $g(\kappa) = 1$ and $g(\lambda) \to \infty$ as
  $\lambda \downarrow 0$; in particular $g \ge 1$ throughout, so $g^{-p}$ is defined and
  takes values in $(0,1]$.  Since $t \mapsto t^{-p}$ is strictly decreasing on
  $[1,\infty)$ for $p > 0$, the composite $\varphi^{\log}_p = g^{-p}$ is strictly
  increasing on $(0,\kappa]$, and it is continuous there as a composite of continuous
  maps.  As $\lambda \downarrow 0$ we have $g(\lambda) \to \infty$ and hence
  $\varphi^{\log}_p(\lambda) \to 0 = \varphi^{\log}_p(0)$, so the extension by $0$ is
  continuous at $0$ and the function is strictly increasing on all of $[0,\kappa]$.
\end{proof}


\begin{lemma}[Monotonicity of the source sets]
  \label{lem_source_monotone}
  Let $\varphi,\psi$ be index functions and $c > 0$ with
  $\varphi(\lambda) \le c\,\psi(\lambda)$ for every $\lambda \in (0,\kappa]$.  Then
  $X_{\varphi,\rho} \subseteq X_{\psi,c\rho}$ for every $\rho > 0$.
\end{lemma}
\begin{proof}
  Let $x \in X_{\varphi,\rho}$, say $x = \varphi(K^{*}K)w$ with $\|w\| \le \rho$.  Define
  $h \colon (0,\kappa] \to \mathbb{R}$ by $h(\lambda) := \varphi(\lambda)/\psi(\lambda)$,
  which is well defined because $\psi$ is strictly increasing with $\psi(0)=0$ and hence
  strictly positive on $(0,\kappa]$, and which satisfies $0 \le h \le c$ by hypothesis.
  Put $w' := h(K^{*}K)w$, using the functional calculus \cref{def_fun_calc}; by
  \cref{lem_fun_calc_bounded}, $\|w'\| \le \bigl(\sup_{\lambda}|h(\lambda)|\bigr)\|w\|
  \le c\rho$.

  It remains to check $\psi(K^{*}K)w' = x$.  Both sides are computed on the singular
  basis: for $n \in I$, $\langle v_n, \psi(K^{*}K)w'\rangle = \psi(\sigma_n^{2})
  h(\sigma_n^{2})\langle v_n,w\rangle = \varphi(\sigma_n^{2})\langle v_n,w\rangle
  = \langle v_n, x\rangle$, using $\sigma_n^{2} > 0$ so that $h$ is defined there.  On
  $\Ker(K)$ both $\psi(K^{*}K)w'$ and $x = \varphi(K^{*}K)w$ vanish, since
  $\psi(0) = \varphi(0) = 0$.  By \cref{thm_onb} and \cref{thm_orth_decomp} the two vectors
  agree, so $x = \psi(K^{*}K)w' \in X_{\psi,c\rho}$.
\end{proof}


\begin{corollary}[Every H\"older source set sits inside every logarithmic one]
  \label{cor_holder_in_log}
  For all $\nu, p, \rho > 0$ there is $c = c(\nu,p,\kappa) > 0$ with
  $X_{\nu,\rho} \subseteq X_{\varphi^{\log}_{p},\,c\rho}$.
\end{corollary}
\begin{proof}
  By \cref{lem_holder_is_index} and \cref{lem_log_is_index} both $\varphi_{\nu}$ and
  $\varphi^{\log}_{p}$ are index functions, so \cref{lem_source_monotone} applies once a
  constant $c$ with $\varphi_{\nu} \le c\,\varphi^{\log}_{p}$ on $(0,\kappa]$ is produced.
  Equivalently, put
  \[
    F(\lambda) \;:=\; \frac{\varphi_{\nu}(\lambda)}{\varphi^{\log}_{p}(\lambda)}
      \;=\; \lambda^{\nu/2}\Bigl(\log\frac{e\kappa}{\lambda}\Bigr)^{p} ,
      \qquad \lambda \in (0,\kappa] ,
  \]
  and show $F$ is bounded.  $F$ is continuous on $(0,\kappa]$ with
  $F(\kappa) = \kappa^{\nu/2}$.  Substituting $\lambda = e\kappa\,e^{-s}$ with
  $s \ge 1$ gives $F = (e\kappa)^{\nu/2}e^{-s\nu/2}s^{p}$, and $s \mapsto e^{-s\nu/2}s^{p}$
  is continuous on $[1,\infty)$ and tends to $0$ as $s \to \infty$, hence is bounded there
  by some $B$.  So $F \le (e\kappa)^{\nu/2}B =: c$ on $(0,\kappa]$, and
  \cref{lem_source_monotone} gives $X_{\nu,\rho} = X_{\varphi_{\nu},\rho} \subseteq
  X_{\varphi^{\log}_p, c\rho}$, the first equality by \cref{lem_holder_is_index}.
\end{proof}


\section{Qualification and rates at a general index function}

\begin{definition}[Qualification at an index function]
  \label{def_general_qualification}
  For a regularizing filter $\{\varphi_{\alpha}\}$ with residual functions
  $r_{\alpha}(\lambda) := 1 - \lambda\varphi_{\alpha}(\lambda)$ and an index function
  $\varphi$, put
  \[
    \omega_{\varphi}(\alpha) \;:=\;
      \sup_{\lambda \in (0,\kappa]} \varphi(\lambda)\,|r_{\alpha}(\lambda)| .
  \]
  The filter has \emph{qualification $\varphi$ with constant $C_{\varphi}^{\mathrm{q}}$},
  written (Q$_{\varphi}$), if
  $\omega_{\varphi}(\alpha) \le C_{\varphi}^{\mathrm{q}}\,\varphi(\alpha)$ for all
  $\alpha \in (0,\kappa]$.
\end{definition}


\begin{lemma}[Approximation error at a general index function]
  \label{lem_general_omega}
  Let $y \in \Dom(K\dg)$ with $x\dg = K\dg y \in X_{\varphi,\rho}$ for an index function
  $\varphi$ and $\rho > 0$.  Then for every $\alpha > 0$,
  \[
    \|x_{\alpha} - x\dg\|_X \;\le\; \omega_{\varphi}(\alpha)\,\rho ,
  \]
  where $x_{\alpha} := R_{\alpha}y$ is the spectral regularization of
  \cref{def_spectral_reg}.
\end{lemma}
\begin{proof}
  Write $x\dg = \varphi(K^{*}K)w$ with $\|w\| \le \rho$.  Since $y \in \Dom(K\dg)$ we have
  $x\dg \in \Ker(K)^{\perp}$, and $\{v_n\}_{n \in I}$ is an orthonormal basis of that
  space by the singular system.  By \cref{lem_spectral_residual_form} the error is
  diagonal:
  \[
    x_{\alpha} - x\dg \;=\; -\sum_{n \in I} r_{\alpha}(\sigma_n^{2})
      \langle v_n, x\dg\rangle_X\, v_n ,
  \]
  the series converging in $X$.  Now $\langle v_n, x\dg\rangle =
  \varphi(\sigma_n^{2})\langle v_n, w\rangle$ by the definition of the functional calculus
  on the singular basis, so by \cref{lem_parseval_series},
  \[
    \|x_{\alpha}-x\dg\|_X^{2}
      \;=\; \sum_{n \in I} |r_{\alpha}(\sigma_n^{2})|^{2}\varphi(\sigma_n^{2})^{2}
        |\langle v_n,w\rangle|^{2}
      \;\le\; \omega_{\varphi}(\alpha)^{2}\sum_{n \in I}|\langle v_n,w\rangle|^{2} ,
  \]
  using $\varphi(\sigma_n^{2})|r_{\alpha}(\sigma_n^{2})| \le \omega_{\varphi}(\alpha)$
  from \cref{def_general_qualification}, valid because $\sigma_n^{2} \in (0,\kappa]$.  The
  remaining sum is at most $\|w\|^{2} \le \rho^{2}$ by Bessel's inequality
  \cref{thm_bessel}.  Taking square roots gives the claim.
\end{proof}


\begin{definition}[The balancing function]
  \label{def_theta_function}
  For an index function $\varphi$ put
  $\Theta_{\varphi}(\alpha) := \sqrt{\alpha}\,\varphi(\alpha)$ for
  $\alpha \in [0,\kappa]$.  It is again an index function, so it has a continuous strictly
  increasing inverse $\Theta_{\varphi}^{-1}$ on $[0,\Theta_{\varphi}(\kappa)]$.
\end{definition}


\begin{theorem}[Rate for an a priori rule at a general index function]
  \label{thm_general_apriori_rate}
  Let $\varphi$ be an index function, let the filter satisfy (F) with constant $C_F$ and
  (Q$_{\varphi}$) with constant $C^{\mathrm{q}}_{\varphi}$, and let $y \in \Dom(K\dg)$
  with $x\dg = K\dg y \in X_{\varphi,\rho}$.  Suppose $\delta/\rho \le
  \Theta_{\varphi}(\kappa)$ and let the a priori rule be
  $\alpha(\delta) := \Theta_{\varphi}^{-1}(\delta/\rho)$.  Then for every
  $y^{\delta} \in B_{\delta}(y)$,
  \[
    \|x^{\delta}_{\alpha(\delta)} - x\dg\|_X \;\le\;
      \bigl(\sqrt{C_{\varphi}C_F} + C^{\mathrm{q}}_{\varphi}\bigr)\,
      \rho\,\varphi\bigl(\Theta_{\varphi}^{-1}(\delta/\rho)\bigr) .
  \]
\end{theorem}
\begin{proof}
  Write $\alpha := \alpha(\delta) = \Theta_{\varphi}^{-1}(\delta/\rho)$, which lies in
  $(0,\kappa]$ by the hypothesis $\delta/\rho \le \Theta_{\varphi}(\kappa)$ and
  \cref{def_theta_function}.  By \cref{lem_error_decomposition},
  \[
    \|x^{\delta}_{\alpha} - x\dg\|_X \;\le\;
      \delta\,\|R_{\alpha}\|_{\Ll(Y,X)} \;+\; \|x_{\alpha}-x\dg\|_X .
  \]

  \emph{The data term.}  Condition (F) of \cref{def_qualification} gives
  $\sup_{\lambda \in (0,\kappa]}|\varphi_{\alpha}(\lambda)| \le C_F\alpha^{-1}$, so
  \cref{lem_filter_norm_bound} yields
  $\|R_{\alpha}\| \le \sqrt{C_{\varphi}}\,(C_F\alpha^{-1})^{1/2}
  = \sqrt{C_{\varphi}C_F}\,\alpha^{-1/2}$.  By the definition of $\alpha$,
  $\sqrt{\alpha}\,\varphi(\alpha) = \delta/\rho$, hence
  $\delta\alpha^{-1/2} = \rho\,\varphi(\alpha)$, and the data term is at most
  $\sqrt{C_{\varphi}C_F}\,\rho\,\varphi(\alpha)$.

  \emph{The approximation term.}  By \cref{lem_general_omega} it is at most
  $\omega_{\varphi}(\alpha)\rho$, and by (Q$_{\varphi}$) of
  \cref{def_general_qualification} that is at most
  $C^{\mathrm{q}}_{\varphi}\varphi(\alpha)\rho$.

  Adding the two bounds and substituting
  $\alpha = \Theta_{\varphi}^{-1}(\delta/\rho)$ gives the display.
\end{proof}


\begin{corollary}[The H\"older rate is the instance $\varphi = \varphi_{\nu}$]
  \label{cor_general_holder_rate}
  For $\varphi = \varphi_{\nu}$ the rule of \cref{thm_general_apriori_rate} is
  $\alpha(\delta) = (\delta/\rho)^{2/(\nu+1)}$ and the bound reads
  \[
    \|x^{\delta}_{\alpha(\delta)} - x\dg\|_X \;\le\;
      \bigl(\sqrt{C_{\varphi}C_F} + C_{\nu}\bigr)\,
      \delta^{\nu/(\nu+1)}\rho^{1/(\nu+1)} ,
  \]
  which is the conclusion of \cref{thm_apriori_order_optimal} with $c = C = 1$.
\end{corollary}
\begin{proof}
  With $\varphi = \varphi_{\nu}$, an index function by \cref{lem_holder_is_index},
  \cref{def_theta_function} gives
  $\Theta_{\varphi_{\nu}}(\alpha) = \alpha^{1/2}\alpha^{\nu/2} = \alpha^{(\nu+1)/2}$,
  whose inverse is $t \mapsto t^{2/(\nu+1)}$.  So
  $\alpha(\delta) = (\delta/\rho)^{2/(\nu+1)}$ and
  \[
    \varphi_{\nu}\bigl(\alpha(\delta)\bigr)
      = \bigl((\delta/\rho)^{2/(\nu+1)}\bigr)^{\nu/2}
      = (\delta/\rho)^{\nu/(\nu+1)} .
  \]
  Multiplying by $\rho$ gives
  $\rho^{1-\nu/(\nu+1)}\delta^{\nu/(\nu+1)} = \rho^{1/(\nu+1)}\delta^{\nu/(\nu+1)}$, which
  with \cref{thm_general_apriori_rate} is the display.  The constant
  $C^{\mathrm{q}}_{\varphi_{\nu}}$ is the constant $C_{\nu}$ of
  \cref{def_qualification}, since (Q$_{\varphi_{\nu}}$) and (Q$_{\nu}$) are the same
  condition.
\end{proof}


\begin{corollary}[Logarithmic rate]
  \label{cor_log_rate}
  Let the filter satisfy (F) and (Q$_{\varphi^{\log}_p}$), and let
  $x\dg \in X_{\varphi^{\log}_{p},\rho}$.  Then the a priori rule of
  \cref{thm_general_apriori_rate} gives, for all sufficiently small $\delta > 0$,
  \[
    \|x^{\delta}_{\alpha(\delta)} - x\dg\|_X \;\le\;
      C\,\rho\,\Bigl(\log\frac{\rho}{\delta}\Bigr)^{-p}
  \]
  with $C$ depending only on $p$, $\kappa$, $C_{\varphi}$, $C_F$ and
  $C^{\mathrm{q}}_{\varphi^{\log}_p}$.
\end{corollary}
\begin{proof}
  Write $\varphi := \varphi^{\log}_p$, an index function by \cref{lem_log_is_index}, put
  $t := \delta/\rho$ and $\alpha := \Theta_{\varphi}^{-1}(t)$, so that
  $\sqrt{\alpha}\,\varphi(\alpha) = t$ by \cref{def_theta_function}.  By
  \cref{thm_general_apriori_rate} the error is at most
  $(\sqrt{C_{\varphi}C_F}+C^{\mathrm{q}}_{\varphi})\rho\,\varphi(\alpha)$, so it suffices
  to bound $\varphi(\alpha)$ by $(\log(1/t))^{-p}$ for small $t$.

  \emph{A lower bound for $\alpha$.}  By \cref{def_log_index},
  $\log(e\kappa/\lambda) \ge 1$ on $(0,\kappa]$, so $\varphi \le 1$ there.  Hence
  $\sqrt{\alpha} = t/\varphi(\alpha) \ge t$, that is $\alpha \ge t^{2}$.

  \emph{An upper bound for $\alpha$.}  Since $\varphi$ is increasing and
  $\alpha \ge t^{2}$, we have $\varphi(\alpha) \ge \varphi(t^{2})$, and substituting into
  $\sqrt{\alpha} = t/\varphi(\alpha)$ gives
  \[
    \alpha \;\le\; \frac{t^{2}}{\varphi(t^{2})^{2}}
      \;=\; t^{2}\Bigl(\log\frac{e\kappa}{t^{2}}\Bigr)^{2p} \;=:\; \beta(t) .
  \]

  \emph{Conclusion.}  The function $t \mapsto t(\log(e\kappa/t^{2}))^{2p}$ is continuous
  on $(0,1)$ and tends to $0$ as $t \downarrow 0$, because a power of a logarithm grows
  more slowly than $1/t$.  Choose $t_0 \in (0,1)$, depending only on $p$ and $\kappa$, so
  small that $t(\log(e\kappa/t^{2}))^{2p} \le \kappa$ and
  $t(\log(e\kappa/t^{2}))^{2p} \le e\kappa\,t$ fail nowhere on $(0,t_0]$; the first gives
  $\beta(t) \le \kappa t \le \kappa$, so $\beta(t)$ lies in the domain of $\varphi$, and
  the second gives $e\kappa/\beta(t) \ge 1/t$.  For $t \le t_0$, monotonicity of $\varphi$
  and $\alpha \le \beta(t)$ then give
  \[
    \varphi(\alpha) \;\le\; \varphi\bigl(\beta(t)\bigr)
      \;=\; \Bigl(\log\frac{e\kappa}{\beta(t)}\Bigr)^{-p}
      \;\le\; \Bigl(\log\frac{1}{t}\Bigr)^{-p} ,
  \]
  the last step because $\log$ is increasing.  Substituting $t = \delta/\rho$ and
  multiplying by $\rho$ and by the constant of \cref{thm_general_apriori_rate} gives the
  display.
\end{proof}


\chapter{Iterated Tikhonov and semi-iterative methods}
\label{ch_20_iterated_tikhonov}

\noindent\emph{Background.} Iterating the Tikhonov step $m$ times raises the achievable
qualification from $2$ to $2m$, postponing --- though never eliminating --- the
saturation of \cref{ch_16_tikhonov}; this goes back to
\citet{KingChillingworth1979}. The residual-polynomial framework of the second section,
of which Landweber iteration is the simplest instance, is the abstract setting for the
semi-iterative (Krylov-type) methods studied by \citet{Brakhage1987}, whose
conjugate-gradient method for ill-posed problems attains every finite qualification.

\begin{assumptions}
  \label{asm_iterated_tikhonov_semi}
  Those of the spectral chapter, unchanged: $K$, the
singular system, $\kappa = \|K\|^{2}$, the filter $\{\varphi_{\alpha}\}$ and its residual
functions $r_{\alpha}(\lambda) = 1 - \lambda\varphi_{\alpha}(\lambda)$.  $m \ge 1$ is a
fixed integer, the number of Tikhonov steps.
\end{assumptions}

\section{Iterated Tikhonov}

\begin{definition}[Iterated Tikhonov filter]
  \label{def_iterated_tikhonov}
  For $m \ge 1$ and $\alpha > 0$ put
  \[
    \varphi^{(m)}_{\alpha}(\lambda) \;:=\;
      \frac{1}{\lambda}\Bigl(1 - \Bigl(\frac{\alpha}{\alpha+\lambda}\Bigr)^{m}\Bigr)
      \quad (\lambda > 0),
    \qquad \varphi^{(m)}_{\alpha}(0) := \frac{m}{\alpha} ,
  \]
  so that $r^{(m)}_{\alpha}(\lambda) = (\alpha/(\alpha+\lambda))^{m}$.  For $m = 1$ this
  is the Tikhonov filter of \cref{def_standard_filters}.
\end{definition}


\begin{lemma}[The filter is $m$ Tikhonov steps]
  \label{lem_iterated_tikhonov_recursion}
  Fix $\alpha > 0$ and $y \in Y$, and define $z_0 := 0$ and, for $k = 1,\dots,m$,
  \[
    z_k \;:=\; z_{k-1} + (K^{*}K+\alpha)^{-1}K^{*}(y - Kz_{k-1}) .
  \]
  Then $z_m = \varphi^{(m)}_{\alpha}(K^{*}K)K^{*}y$.
\end{lemma}
\begin{proof}
  Both sides are diagonal in the singular basis, so it suffices to prove, for each
  $n \in I$ and $k = 0,\dots,m$, that
  \[
    \langle v_n, z_k\rangle_X \;=\;
      \frac{1}{\sigma_n}\Bigl(1-\Bigl(\frac{\alpha}{\alpha+\sigma_n^{2}}\Bigr)^{k}\Bigr)
      \langle u_n, y\rangle_Y ,
  \]
  and that $z_k \in \Ker(K)^{\perp}$ for every $k$.  The membership is immediate by
  induction, since $z_0 = 0$ and each increment lies in $\Ran(K^{*}) \subseteq
  \Ker(K)^{\perp}$.

  For $k = 0$ the display reads $0 = 0$.  Assume it for $k-1$.  By
  \cref{lem_tikhonov_normal} the operator $(K^{*}K+\alpha)^{-1}K^{*}$ acts on the singular
  basis by sending $u_n$ to $\sigma_n(\sigma_n^{2}+\alpha)^{-1}v_n$, and $K$ sends $v_n$
  to $\sigma_n u_n$.  Writing $c_n := \langle u_n,y\rangle$ and
  $a_n := \langle v_n, z_{k-1}\rangle$, the increment contributes
  \[
    \frac{\sigma_n}{\sigma_n^{2}+\alpha}\bigl(c_n - \sigma_n a_n\bigr) ,
  \]
  so $\langle v_n,z_k\rangle = a_n + \sigma_n(\sigma_n^{2}+\alpha)^{-1}(c_n-\sigma_na_n)$.
  Substituting the inductive value of $a_n$ and simplifying with
  $1 - \sigma_n^{2}(\sigma_n^{2}+\alpha)^{-1} = \alpha(\sigma_n^{2}+\alpha)^{-1}$ gives
  \[
    \langle v_n, z_k\rangle
      \;=\; \frac{c_n}{\sigma_n}\Bigl(1 -
        \Bigl(\frac{\alpha}{\alpha+\sigma_n^{2}}\Bigr)^{k}\Bigr) ,
  \]
  which is the display for $k$.  At $k = m$ the right-hand side is
  $\varphi^{(m)}_{\alpha}(\sigma_n^{2})\sigma_nc_n$ by \cref{def_iterated_tikhonov}, and
  \cref{def_spectral_reg} together with \cref{lem_spectral_residual_form} identifies the
  vector with those coefficients as $\varphi^{(m)}_{\alpha}(K^{*}K)K^{*}y$.
\end{proof}


\begin{lemma}[Iterated Tikhonov is a regularizing filter]
  \label{lem_iterated_tikhonov_filter}
  $\{\varphi^{(m)}_{\alpha}\}$ is a regularizing filter in the sense of
  \cref{def_filter}, with $C_{\varphi} = 1$, and it satisfies condition (F) with
  $C_F = m$.
\end{lemma}
\begin{proof}
  Write $s := \lambda/(\alpha+\lambda) \in [0,1)$ for $\lambda \in (0,\kappa]$, so that
  $r^{(m)}_{\alpha}(\lambda) = (1-s)^{m}$ and
  $\lambda\varphi^{(m)}_{\alpha}(\lambda) = 1-(1-s)^{m}$.

  \emph{Pointwise convergence.}  For fixed $\lambda > 0$, $s \to 1$ as $\alpha \to 0$, so
  $\lambda\varphi^{(m)}_{\alpha}(\lambda) \to 1$, i.e.\
  $\varphi^{(m)}_{\alpha}(\lambda) \to \lambda^{-1}$, which is the first clause of
  \cref{def_filter}.

  \emph{The bound $C_{\varphi} = 1$.}  Since $s \in [0,1)$ we have
  $(1-s)^{m} \in (0,1]$, so
  $\lambda|\varphi^{(m)}_{\alpha}(\lambda)| = 1-(1-s)^{m} \in [0,1)$, giving the second
  clause with $C_{\varphi}=1$.

  \emph{Condition (F).}  Bernoulli's inequality gives $(1-s)^{m} \ge 1-ms$ for
  $s \in [0,1]$ and $m \ge 1$, hence $1-(1-s)^{m} \le ms$ and
  \[
    |\varphi^{(m)}_{\alpha}(\lambda)| \;=\; \frac{1-(1-s)^{m}}{\lambda}
      \;\le\; \frac{ms}{\lambda} \;=\; \frac{m}{\alpha+\lambda} \;\le\; \frac{m}{\alpha} .
  \]
  Taking the supremum over $\lambda \in (0,\kappa]$ gives (F) with $C_F = m$ in the sense
  of \cref{def_qualification}.  Boundedness and piecewise continuity on $[0,\kappa]$ are
  visible from the closed form, the value at $0$ being the limit.
\end{proof}


\begin{theorem}[Qualification $2m$]
  \label{thm_iterated_tikhonov_qualification}
  For $0 < \nu \le 2m$ the filter $\{\varphi^{(m)}_{\alpha}\}$ satisfies (Q$_{\nu}$) with
  a constant $C_{\nu} \le 1$; that is,
  $\omega_{\nu}(\alpha) \le \alpha^{\nu/2}$ for all $\alpha > 0$.
\end{theorem}
\begin{proof}
  Let $\lambda \in (0,\kappa]$ and $\alpha > 0$.  With $t := \lambda/(\alpha+\lambda)$ and
  $u := \alpha/(\alpha+\lambda) = 1-t$,
  \[
    \lambda^{\nu/2}\,r^{(m)}_{\alpha}(\lambda)
      \;=\; \bigl(\lambda\bigr)^{\nu/2}u^{m}
      \;=\; \alpha^{\nu/2}\Bigl(\frac{\lambda}{\alpha}\Bigr)^{\nu/2}u^{m}
      \;=\; \alpha^{\nu/2}\,\Bigl(\frac{t}{u}\Bigr)^{\nu/2}u^{m}
      \;=\; \alpha^{\nu/2}\,t^{\nu/2}u^{\,m-\nu/2} ,
  \]
  using $\lambda/\alpha = t/u$.  Since $0 < \nu \le 2m$ the exponent $m-\nu/2$ is
  non-negative, and $t,u \in (0,1)$, so $t^{\nu/2}u^{m-\nu/2} \le 1$.  Hence
  $\lambda^{\nu/2}|r^{(m)}_{\alpha}(\lambda)| \le \alpha^{\nu/2}$ for every admissible
  $\lambda$, and taking the supremum gives
  $\omega_{\nu}(\alpha) \le \alpha^{\nu/2}$, which is (Q$_{\nu}$) with $C_{\nu} = 1$.
\end{proof}


\begin{lemma}[The qualification is not larger]
  \label{lem_iterated_tikhonov_qualification_sharp}
  Suppose $K$ has infinite-dimensional range and let $\nu > 2m$.  Then there is no
  constant $C$ with $\omega_{\nu}(\alpha) \le C\alpha^{\nu/2}$ for all
  $\alpha \in (0,\kappa]$.  The qualification of $\{\varphi^{(m)}_{\alpha}\}$ in the sense
  of \cref{def_qualification} is therefore exactly $2m$.
\end{lemma}
\begin{proof}
  Fix $\alpha \in (0,\kappa]$ and bound the supremum defining $\omega_{\nu}$ from below by
  its value at the single point $\lambda = \kappa$, which is admissible:
  \[
    \omega_{\nu}(\alpha) \;\ge\; \kappa^{\nu/2}
      \Bigl(\frac{\alpha}{\alpha+\kappa}\Bigr)^{m}
      \;\ge\; \kappa^{\nu/2}\Bigl(\frac{\alpha}{2\kappa}\Bigr)^{m}
      \;=\; \frac{\kappa^{\nu/2}}{(2\kappa)^{m}}\,\alpha^{m}
  \]
  for $\alpha \le \kappa$.  If (Q$_{\nu}$) held with a constant $C$ we would obtain
  $\kappa^{\nu/2}(2\kappa)^{-m}\alpha^{m} \le C\alpha^{\nu/2}$, that is
  $\alpha^{m-\nu/2} \le C(2\kappa)^{m}\kappa^{-\nu/2}$, for all $\alpha \in (0,\kappa]$.
  Since $\nu > 2m$ the exponent $m-\nu/2$ is strictly negative, so the left-hand side
  tends to $\infty$ as $\alpha \downarrow 0$, contradicting the fixed bound on the right.

  Hence (Q$_{\nu}$) fails for every $\nu > 2m$, while
  \cref{thm_iterated_tikhonov_qualification} gives it for every $\nu \le 2m$; the
  supremum in \cref{def_qualification} is therefore exactly $2m$.  The hypothesis that
  $\Ran(K)$ is infinite-dimensional is used only to know that $\kappa > 0$, i.e.\ that
  $K \ne 0$ and the supremum is over a non-degenerate interval.
\end{proof}


\begin{corollary}[Iterated Tikhonov is order optimal up to $\nu = 2m$]
  \label{cor_iterated_tikhonov_order_optimal}
  Let $0 < \nu \le 2m$, let $y \in \Dom(K\dg)$ with $x\dg = K\dg y \in X_{\nu,\rho}$, and
  let $\alpha(\delta) := (\delta/\rho)^{2/(\nu+1)}$.  Then
  \[
    \|x^{\delta}_{\alpha(\delta)} - x\dg\|_X \;\le\;
      \bigl(\sqrt{m}+1\bigr)\,\delta^{\nu/(\nu+1)}\rho^{1/(\nu+1)} ,
  \]
  so the method is order optimal for every $\nu$ up to $2m$.
\end{corollary}
\begin{proof}
  By \cref{lem_iterated_tikhonov_filter} the family is a regularizing filter with
  $C_{\varphi} = 1$ satisfying (F) with $C_F = m$, and by
  \cref{thm_iterated_tikhonov_qualification} it satisfies (Q$_{\nu}$) with $C_{\nu} = 1$
  for the given $\nu$.  Apply \cref{thm_apriori_order_optimal} with $c = C = 1$: its
  constant is $\sqrt{C_{\varphi}C_F}\,c^{-1/2} + C_{\nu}C^{\nu/2}
  = \sqrt{1\cdot m} + 1$.
\end{proof}


\section{Semi-iterative methods}

\begin{definition}[Residual polynomials and semi-iterative methods]
  \label{def_residual_polynomial}
  A family $\{q_k\}_{k \ge 0}$ of real polynomials is a \emph{residual family} if
  $q_k(0) = 1$ and $\deg q_k \le k$ for every $k$.  The associated \emph{semi-iterative
  filter} is
  \[
    \varphi_{k}(\lambda) \;:=\; \frac{1-q_k(\lambda)}{\lambda} \quad (\lambda > 0),
    \qquad \varphi_k(0) := -q_k'(0) ,
  \]
  which is a polynomial of degree at most $k-1$, and the method is
  $x_k := \varphi_k(K^{*}K)K^{*}y$, with residual function $r_k = q_k$.
\end{definition}


\begin{definition}[The Landweber residual family]
  \label{def_landweber_residual_poly}
  Fix $\omega \in (0,\kappa^{-1}]$.  The \emph{Landweber residual family} is the family
  of polynomials $q_k(\lambda) := (1-\omega\lambda)^{k}$, $k \ge 0$.
\end{definition}


\begin{lemma}[Landweber is a semi-iterative method]
  \label{lem_landweber_residual_polynomial}
  Fix $\omega \in (0,\kappa^{-1}]$ and put $q_k(\lambda) := (1-\omega\lambda)^{k}$.  Then
  $\{q_k\}$ is a residual family, and its semi-iterative filter is the Landweber filter of
  \cref{def_standard_filters} at $\alpha = 1/k$.
\end{lemma}
\begin{proof}
  $q_k(0) = 1^{k} = 1$ and $\deg q_k = k$, so $\{q_k\}$ is a residual family in the sense
  of \cref{def_residual_polynomial}.  Its filter is
  \[
    \varphi_k(\lambda) \;=\; \frac{1-(1-\omega\lambda)^{k}}{\lambda}
      \qquad (\lambda > 0) ,
  \]
  which is the Landweber filter of \cref{def_standard_filters} with $1/\alpha = k$.  At
  $\lambda = 0$ the definition of \cref{def_residual_polynomial} gives
  $-q_k'(0) = \omega k$, and \cref{def_standard_filters} gives $\omega/\alpha = \omega k$;
  the two agree, so the filters are equal as functions on $[0,\kappa]$.  The iterates then
  coincide with those of \cref{def_landweber} by \cref{lem_landweber_closed_form}.
\end{proof}


\begin{theorem}[When a semi-iterative method regularizes]
  \label{thm_semi_iterative_regularization}
  Let $\{q_k\}$ be a residual family and suppose
  \begin{enumerate}
    \item $\sup_{k}\sup_{\lambda \in [0,\kappa]}|q_k(\lambda)| \le C_q < \infty$, and
    \item $q_k(\lambda) \to 0$ as $k \to \infty$, for every $\lambda \in (0,\kappa]$.
  \end{enumerate}
  Then $\{\varphi_k\}$, reindexed by $\alpha := 1/k$, is a regularizing filter with
  $C_{\varphi} = 1 + C_q$, and the resulting family $\{R_{\alpha}\}$ is a regularization
  in the sense of \cref{def_regularization}.
\end{theorem}
\begin{proof}
  For the first clause of \cref{def_filter}, fix $\lambda \in (0,\kappa]$; then
  $\lambda\varphi_k(\lambda) = 1-q_k(\lambda) \to 1$ by hypothesis (2), so
  $\varphi_k(\lambda) \to \lambda^{-1}$, and in the reindexed parameter this is
  convergence as $\alpha \to 0$.

  For the second clause, hypothesis (1) gives
  \[
    \lambda\,|\varphi_k(\lambda)| \;=\; |1-q_k(\lambda)| \;\le\; 1 + C_q
  \]
  for every $\lambda \in (0,\kappa]$ and every $k$, which is the bound with
  $C_{\varphi} = 1+C_q$.  Each $\varphi_k$ is a polynomial, hence bounded and continuous
  on $[0,\kappa]$, so $\{\varphi_k\}$ is a regularizing filter.

  The final claim is \cref{thm_spectral_is_regularization} applied to this filter, which
  produces a regularization in the sense of \cref{def_regularization}.
\end{proof}


\chapter{Regularization with a perturbed operator}
\label{ch_21_perturbed_operator}

\noindent\emph{Background.} Splitting the total error into an operator error (from
knowing $K$ only up to an approximation $K_h$), a propagated data error, and the usual
approximation error is the standard treatment of regularization with an incompletely
known operator, systematized alongside the noisy-data theory in
\citet{EnglHankeNeubauer1996}.

\begin{assumptions}
  \label{asm_regularization_with_perturbed}
  $K \in \Kk(X,Y)$ is the exact operator and
$K_h \in \Kk(X,Y)$ an approximation with $\|K-K_h\|_{\Ll(X,Y)} \le h$; no relation between
their singular systems is assumed, but $K_h$ carries one of its own.  $y \in
\Ran(K)$, $x\dg := K\dg y$, and $y^{\delta} \in B_{\delta}(y)$.  For a filter
$\{\varphi_{\alpha}\}$ we write $R_{\alpha}[K]$ for the spectral regularization built from
$K$ and $R_{\alpha}[K_h]$ for the one built from $K_h$, and
$x^{h,\delta}_{\alpha} := R_{\alpha}[K_h]\,y^{\delta}$.
\end{assumptions}

\section{The three error sources}

\begin{definition}[Perturbed problem]
  \label{def_perturbed_problem}
  The pair $(K_h,y^{\delta})$ is \emph{$(h,\delta)$-admissible data} for $(K,y)$ if
  $\|K-K_h\| \le h$ and $\|y-y^{\delta}\| \le \delta$.
\end{definition}


\begin{lemma}[Three-term error decomposition]
  \label{lem_perturbed_decomposition}
  For $(h,\delta)$-admissible data and every $\alpha > 0$,
  \[
    \|x^{h,\delta}_{\alpha} - x\dg\|_X
      \;\le\; \underbrace{\bigl\|\bigl(R_{\alpha}[K_h]-R_{\alpha}[K]\bigr)y^{\delta}
        \bigr\|_X}_{\text{operator error}}
      \;+\; \underbrace{\delta\,\|R_{\alpha}[K]\|}_{\text{data error}}
      \;+\; \underbrace{\|R_{\alpha}[K]y - x\dg\|_X}_{\text{approximation error}} .
  \]
\end{lemma}
\begin{proof}
  Insert $R_{\alpha}[K]y^{\delta}$ and then $R_{\alpha}[K]y$:
  \[
    x^{h,\delta}_{\alpha} - x\dg
      \;=\; \bigl(R_{\alpha}[K_h]-R_{\alpha}[K]\bigr)y^{\delta}
        \;+\; R_{\alpha}[K]\bigl(y^{\delta}-y\bigr)
        \;+\; \bigl(R_{\alpha}[K]y - x\dg\bigr) .
  \]
  The triangle inequality bounds the norm by the sum of the three norms, and the middle
  one is at most $\|R_{\alpha}[K]\|\,\|y^{\delta}-y\| \le \delta\|R_{\alpha}[K]\|$ by
  \cref{def_perturbed_problem}.
\end{proof}


\section{Tikhonov with a perturbed operator}

\begin{lemma}[Perturbation bound for the Tikhonov operator]
  \label{lem_tikhonov_operator_perturbation}
  Let $\varphi_{\alpha}(\lambda) = (\lambda+\alpha)^{-1}$ be the Tikhonov filter, so that
  $R_{\alpha}[K] = (K^{*}K+\alpha)^{-1}K^{*}$.  Then for $\|K-K_h\| \le h$ and every
  $\alpha > 0$,
  \[
    \bigl\|R_{\alpha}[K_h] - R_{\alpha}[K]\bigr\|_{\Ll(Y,X)}
      \;\le\; \frac{h}{\alpha} \;+\; \frac{h\,(2\|K\|+h)}{2\,\alpha^{3/2}} .
  \]
\end{lemma}
\begin{proof}
  Write $A := K^{*}K+\alpha$ and $A_h := K_h^{*}K_h+\alpha$, both boundedly invertible
  with $\|A^{-1}\|,\|A_h^{-1}\| \le \alpha^{-1}$, since $K^{*}K$ and $K_h^{*}K_h$ are
  positive and self-adjoint.  Split
  \[
    R_{\alpha}[K_h]-R_{\alpha}[K]
      \;=\; A_h^{-1}K_h^{*} - A^{-1}K^{*}
      \;=\; A^{-1}\bigl(K_h^{*}-K^{*}\bigr)
        \;+\; \bigl(A_h^{-1}-A^{-1}\bigr)K_h^{*} .
  \]
  The first term has norm at most $\alpha^{-1}h$, since
  $\|K_h^{*}-K^{*}\| = \|K-K_h\| \le h$.

  For the second, the resolvent identity gives
  $A_h^{-1}-A^{-1} = A^{-1}(A-A_h)A_h^{-1}$, so
  \[
    \bigl\|\bigl(A_h^{-1}-A^{-1}\bigr)K_h^{*}\bigr\|
      \;\le\; \|A^{-1}\|\;\|A-A_h\|\;\bigl\|A_h^{-1}K_h^{*}\bigr\| .
  \]
  Here $A-A_h = K^{*}K-K_h^{*}K_h = K^{*}(K-K_h) + (K-K_h)^{*}K_h$, whence
  $\|A-A_h\| \le h\|K\| + h\|K_h\| \le h(2\|K\|+h)$ using
  $\|K_h\| \le \|K\|+h$.  Finally, $A_h^{-1}K_h^{*}$ is the Tikhonov operator for $K_h$,
  and by \cref{lem_tikhonov_normal} it acts on the singular basis of $K_h$ by the
  multipliers $s(s^{2}+\alpha)^{-1}$; since $s(s^{2}+\alpha)^{-1} \le
  (2\sqrt{\alpha})^{-1}$ for every $s \ge 0$, by the inequality
  $s^{2}+\alpha \ge 2s\sqrt{\alpha}$, we get
  $\|A_h^{-1}K_h^{*}\| \le (2\sqrt{\alpha})^{-1}$.

  Combining, the second term has norm at most
  $\alpha^{-1}h(2\|K\|+h)(2\sqrt{\alpha})^{-1}$, which is the second summand of the
  display.
\end{proof}


\begin{theorem}[Total error for perturbed Tikhonov]
  \label{thm_perturbed_total_error}
  Let $(K_h,y^{\delta})$ be $(h,\delta)$-admissible for $(K,y)$ with $y \in \Ran(K)$, let
  $x\dg = K\dg y \in X_{\nu,\rho}$ for $0 < \nu \le 2$, and let $x^{h,\delta}_{\alpha}$ be
  the Tikhonov reconstruction from the perturbed data.  Then for every
  $\alpha \in (0,1]$,
  \[
    \|x^{h,\delta}_{\alpha} - x\dg\|_X \;\le\;
      C_1\,\frac{h\,(\|y\|+\delta)}{\alpha^{3/2}}
      \;+\; \frac{\delta}{2\sqrt{\alpha}}
      \;+\; \rho\,\alpha^{\nu/2} ,
  \]
  with $C_1 := 1 + \tfrac{1}{2}(2\|K\|+h)$.
\end{theorem}
\begin{proof}
  Apply \cref{lem_perturbed_decomposition} and bound the three terms.

  \emph{Operator error.}  By \cref{lem_tikhonov_operator_perturbation} and
  $\|y^{\delta}\| \le \|y\|+\delta$,
  \[
    \bigl\|\bigl(R_{\alpha}[K_h]-R_{\alpha}[K]\bigr)y^{\delta}\bigr\|
      \;\le\; \Bigl(\frac{h}{\alpha}+\frac{h(2\|K\|+h)}{2\alpha^{3/2}}\Bigr)
        \bigl(\|y\|+\delta\bigr) .
  \]
  For $\alpha \le 1$ we have $\alpha^{-1} \le \alpha^{-3/2}$, so the bracket is at most
  $C_1 h\alpha^{-3/2}$ with $C_1 = 1+\tfrac{1}{2}(2\|K\|+h)$.

  \emph{Data error.}  The Tikhonov filter of \cref{def_standard_filters} satisfies
  $\lambda\varphi_{\alpha}(\lambda) \le 1$, so $C_{\varphi} = 1$, and
  $\sup_{\lambda}|\varphi_{\alpha}(\lambda)| \le \alpha^{-1}$.  By
  \cref{lem_filter_norm_bound}, $\|R_{\alpha}[K]\| \le \alpha^{-1/2}$; the sharper
  constant $\tfrac{1}{2}$ comes from the same computation as in
  \cref{lem_tikhonov_operator_perturbation}, namely
  $s(s^{2}+\alpha)^{-1} \le (2\sqrt{\alpha})^{-1}$.  Hence the data error is at most
  $\delta/(2\sqrt{\alpha})$.

  \emph{Approximation error.}  By \cref{lem_tikhonov_qualification} the Tikhonov filter
  satisfies (Q$_{\nu}$) with $C_{\nu} = 1$ for $\nu \le 2$, so
  $\omega_{\nu}(\alpha) \le \alpha^{\nu/2}$, and \cref{lem_omega_nu} gives
  $\|R_{\alpha}[K]y-x\dg\| \le \rho\,\alpha^{\nu/2}$.

  Summing the three bounds gives the display.
\end{proof}


\begin{corollary}[Rate when the operator error is small relative to the data error]
  \label{cor_perturbed_rate}
  In the setting of \cref{thm_perturbed_total_error}, put
  $\alpha(\delta) := (\delta/\rho)^{2/(\nu+1)}$ and suppose that
  $0 < \delta \le \rho$ and that the operator error satisfies
  \[
    h\,(\|y\|+\delta) \;\le\; \delta\,\alpha(\delta) .
  \]
  Then
  \[
    \|x^{h,\delta}_{\alpha(\delta)} - x\dg\|_X \;\le\;
      \Bigl(C_1 + \tfrac{3}{2}\Bigr)\,\delta^{\nu/(\nu+1)}\rho^{1/(\nu+1)} ,
  \]
  the same order as \cref{thm_apriori_order_optimal} gives for the exact operator.
\end{corollary}
\begin{proof}
  Write $\alpha := \alpha(\delta) = (\delta/\rho)^{2/(\nu+1)}$ and
  $D := \delta^{\nu/(\nu+1)}\rho^{1/(\nu+1)}$.  Since $\delta \le \rho$ we have
  $\alpha \in (0,1]$, so \cref{thm_perturbed_total_error} applies.

  Two identities do all the work.  First,
  $\alpha^{-1/2} = (\delta/\rho)^{-1/(\nu+1)}$, so
  $\delta\,\alpha^{-1/2} = \delta^{1-1/(\nu+1)}\rho^{1/(\nu+1)} = D$.  Second,
  $\rho\,\alpha^{\nu/2} = \rho(\delta/\rho)^{\nu/(\nu+1)} = D$.

  Now bound the three terms of \cref{thm_perturbed_total_error}.  The hypothesis
  $h(\|y\|+\delta) \le \delta\alpha$ turns the first into
  $C_1\,\delta\alpha\,\alpha^{-3/2} = C_1\,\delta\alpha^{-1/2} = C_1D$.  The second is
  $\delta/(2\sqrt{\alpha}) = D/2$.  The third is $\rho\alpha^{\nu/2} = D$.  Summing,
  \[
    \|x^{h,\delta}_{\alpha} - x\dg\|_X \;\le\;
      C_1D + \tfrac{1}{2}D + D \;=\; \Bigl(C_1+\tfrac{3}{2}\Bigr)D ,
  \]
  which is the display.  The comparison with \cref{thm_apriori_order_optimal} is that the
  exponent pair $(\nu/(\nu+1),\,1/(\nu+1))$ is the same; only the constant differs.
\end{proof}


\chapter{The balancing principle}
\label{ch_22_balancing}

\noindent\emph{Background.} The balancing principle --- choosing the regularization
parameter from a geometric grid by comparing reconstructions across scales, without
knowing the source condition in advance, and still attaining an order-optimal rate ---
adapts to ill-posed problems the adaptive estimation method of \citet{Lepskii1990} from
nonparametric statistics; its formulation here, through index functions and an explicit
oracle inequality, follows \citet{MathePereverzev2003}.

\begin{assumptions}
  \label{asm_balancing_principle}
  $\{\varphi_{\alpha}\}$ is a regularizing filter
satisfying (F) with constant $C_F$, so that by \cref{lem_filter_norm_bound}
$\|R_{\alpha}\| \le \sqrt{C_{\varphi}C_F}\,\alpha^{-1/2} =: C_R\alpha^{-1/2}$.  The
\emph{data majorant} is $\Delta(\alpha) := C_R\,\delta\,\alpha^{-1/2}$.  A \emph{geometric
grid} is a finite set $\alpha_0 > \alpha_1 > \cdots > \alpha_N > 0$ with
$\alpha_{i-1} = q\,\alpha_i$ for a fixed ratio $q > 1$; we write $\Delta_i :=
\Delta(\alpha_i)$, which is increasing in $i$, and $x_i := x^{\delta}_{\alpha_i} =
R_{\alpha_i}y^{\delta}$.
\end{assumptions}

\section{The rule and its oracle inequality}

\begin{definition}[Balancing principle]
  \label{def_balancing}
  On a geometric grid, the \emph{balancing index} is
  \[
    i_{+} \;:=\; \min\bigl\{\, i \in \{0,\dots,N\} \;:\;
      \|x_i - x_j\|_X \le 4\Delta_j \ \text{ for all } j \ge i \,\bigr\} ,
  \]
  and the balancing principle selects $\alpha_{i_{+}}$.
\end{definition}


\begin{lemma}[The rule is well defined]
  \label{lem_balancing_nonempty}
  The set in \cref{def_balancing} contains $i = N$, so $i_{+}$ exists and
  $i_{+} \le N$.
\end{lemma}
\begin{proof}
  Take $i = N$ in \cref{def_balancing}.  The condition ranges over $j \in \{0,\dots,N\}$
  with $j \ge N$, hence over $j = N$ only, and requires
  $\|x_N-x_N\|_X = 0 \le 4\Delta_N$, which holds since $\Delta_N \ge 0$.  So the set is a
  non-empty subset of the finite set $\{0,\dots,N\}$ and has a least element.
\end{proof}


\begin{lemma}[The balancing index does not overshoot]
  \label{lem_balancing_le_star}
  Let $\varphi$ be an index function, $x\dg \in X_{\varphi,\rho}$, and put
  $\omega_i := \omega_{\varphi}(\alpha_i)\rho$, so that
  $\|x_i - x\dg\| \le \omega_i + \Delta_i$ for every $i$.  Suppose $i_{*}$ is an index
  with $\omega_j \le \Delta_j$ for every $j \ge i_{*}$.  Then $i_{+} \le i_{*}$.
\end{lemma}
\begin{proof}
  By \cref{lem_error_decomposition} and \cref{lem_general_omega},
  $\|x_i-x\dg\| \le \Delta_i + \omega_i$ for every $i$, the first term bounding the
  propagated data error through $\|R_{\alpha_i}\| \le C_R\alpha_i^{-1/2}$.

  Let $j \ge i_{*}$.  Then also $\omega_{i_{*}} \le \Delta_{i_{*}}$ and
  $\omega_j \le \Delta_j$ by hypothesis, and $\Delta_{i_{*}} \le \Delta_j$ since $\Delta$
  is increasing in the index.  Hence
  \[
    \|x_{i_{*}} - x_j\|_X \;\le\; \|x_{i_{*}}-x\dg\| + \|x\dg-x_j\|
      \;\le\; (\omega_{i_{*}}+\Delta_{i_{*}}) + (\omega_j+\Delta_j)
      \;\le\; 2\Delta_{i_{*}} + 2\Delta_j \;\le\; 4\Delta_j .
  \]
  So $i_{*}$ satisfies the condition of \cref{def_balancing}, and $i_{+}$, being the least
  such index by \cref{lem_balancing_nonempty}, satisfies $i_{+} \le i_{*}$.
\end{proof}


\begin{theorem}[Oracle inequality for the balancing principle]
  \label{thm_balancing_oracle}
  In the setting of \cref{lem_balancing_le_star}, suppose in addition that $i_{*}$ is the
  \emph{least} index with $\omega_j \le \Delta_j$ for all $j \ge i_{*}$, and that the bias
  is non-increasing along the grid: $\omega_j \le \omega_i$ whenever $i \le j \le N$.  Then
  \[
    \|x_{i_{+}} - x\dg\|_X \;\le\; 6\sqrt{q}\ \min_{0 \le i \le N}
      \bigl(\omega_{\varphi}(\alpha_i)\rho + \Delta_i\bigr) .
  \]
\end{theorem}
\begin{proof}
  Write $e_i := \omega_i + \Delta_i$, so $\|x_i-x\dg\| \le e_i$ as in the proof of
  \cref{lem_balancing_le_star}.

  \emph{Step 1: the selected error is at most $6\Delta_{i_{*}}$.}  By
  \cref{lem_balancing_le_star}, $i_{+} \le i_{*}$, so $i_{*} \ge i_{+}$ is one of the
  indices $j$ over which the condition defining $i_{+}$ quantifies, and therefore
  $\|x_{i_{+}}-x_{i_{*}}\| \le 4\Delta_{i_{*}}$.  Hence
  \[
    \|x_{i_{+}}-x\dg\| \;\le\; \|x_{i_{+}}-x_{i_{*}}\| + \|x_{i_{*}}-x\dg\|
      \;\le\; 4\Delta_{i_{*}} + e_{i_{*}}
      \;\le\; 4\Delta_{i_{*}} + 2\Delta_{i_{*}} \;=\; 6\Delta_{i_{*}} ,
  \]
  using $\omega_{i_{*}} \le \Delta_{i_{*}}$.

  \emph{Step 2: $\Delta_{i_{*}} \le \sqrt{q}\,\min_i e_i$.}  Let $i$ be arbitrary.

  If $i \ge i_{*}$ then $e_i \ge \Delta_i \ge \Delta_{i_{*}}$, since $\Delta$ increases
  with the index.

  If $i < i_{*}$ then $i_{*} \ge 1$, and minimality of $i_{*}$ says that $i_{*}-1$ does
  not have the defining property, i.e.\ there is $j \ge i_{*}-1$ with
  $\omega_j > \Delta_j$.  Every $j \ge i_{*}$ has $\omega_j \le \Delta_j$, so that $j$ can
  only be $i_{*}-1$: thus $\omega_{i_{*}-1} > \Delta_{i_{*}-1}$.  Since $\omega$ decreases
  with the index and $i \le i_{*}-1$, we get $\omega_i \ge \omega_{i_{*}-1}
  > \Delta_{i_{*}-1}$.  On a geometric grid $\alpha_{i_{*}-1} = q\alpha_{i_{*}}$, so
  $\Delta_{i_{*}-1} = C_R\delta(q\alpha_{i_{*}})^{-1/2} = q^{-1/2}\Delta_{i_{*}}$, whence
  $e_i \ge \omega_i > q^{-1/2}\Delta_{i_{*}}$.

  In both cases $e_i \ge q^{-1/2}\Delta_{i_{*}}$, so
  $\Delta_{i_{*}} \le \sqrt{q}\,\min_i e_i$.

  Combining the two steps gives the display.
\end{proof}


\begin{corollary}[Adaptive order optimality]
  \label{cor_balancing_adaptive}
  Let $\nu \in (0,\nu_0]$ with $\nu_0$ the qualification of the filter and let
  $x\dg \in X_{\nu,\rho}$.  Suppose every grid point satisfies $\alpha_i \le \kappa$, and
  the grid contains an index $i$ with
  $\alpha_i \in [c(\delta/\rho)^{2/(\nu+1)}, C(\delta/\rho)^{2/(\nu+1)}]$ for constants
  $c, C > 0$.  Then
  \[
    \|x_{i_{+}} - x\dg\|_X \;\le\;
      6\sqrt{q}\bigl(C_R c^{-1/2}+C_{\nu}C^{\nu/2}\bigr)\,
      \delta^{\nu/(\nu+1)}\rho^{1/(\nu+1)} ,
  \]
  and the rule achieving this uses neither $\nu$ nor $\rho$.
\end{corollary}
\begin{proof}
  By \cref{lem_holder_is_index} the H\"older source set $X_{\nu,\rho}$ is the general
  source set at the index function $\varphi_{\nu}$, so \cref{thm_balancing_oracle}
  applies with $\omega_{\varphi_{\nu}} = \omega_{\nu}$.  Let $i$ be an index of the grid
  as in the hypothesis.  Then
  \[
    \Delta_i \;=\; C_R\delta\alpha_i^{-1/2}
      \;\le\; C_R\,c^{-1/2}\,\delta^{\nu/(\nu+1)}\rho^{1/(\nu+1)}
  \]
  and, by (Q$_{\nu}$) of \cref{def_qualification} with constant $C_{\nu}$, which applies
  since $\alpha_i \le \kappa$, together with \cref{lem_general_omega},
  \[
    \omega_{\nu}(\alpha_i)\rho \;\le\; C_{\nu}\alpha_i^{\nu/2}\rho
      \;\le\; C_{\nu}C^{\nu/2}\,\delta^{\nu/(\nu+1)}\rho^{1/(\nu+1)} ,
  \]
  both by the same exponent arithmetic as in \cref{thm_apriori_order_optimal}.  Hence
  $\min_j e_j \le e_i \le (C_Rc^{-1/2}+C_{\nu}C^{\nu/2})
  \delta^{\nu/(\nu+1)}\rho^{1/(\nu+1)}$, and \cref{thm_balancing_oracle} multiplies this
  by $6\sqrt{q}$.  The final clause is a reading of \cref{def_balancing}: the index
  $i_{+}$ is computed from $\delta$, the filter constants and the family $\{x_i\}$, none
  of which involves $\nu$ or $\rho$.
\end{proof}


 
\begin{thebibliography}{99}

\bibitem[Bakushinski{\u\i}, 1985]{Bakushinskii1985}
A.~B. Bakushinski{\u\i} (1985), Remarks on choosing a regularization parameter using the
quasioptimality and ratio criterion, \emph{USSR Comput. Math. Math. Phys.} 24(4),
181--182.

\bibitem[Brakhage, 1987]{Brakhage1987}
H. Brakhage (1987), On ill-posed problems and the method of conjugate gradients, in:
H.~W. Engl \& C.~W. Groetsch (eds.), \emph{Inverse and Ill-Posed Problems} (Sankt
Wolfgang, 1986), Notes Rep. Math. Sci. Engrg. 4, Academic Press, Boston, MA, 165--175.

\bibitem[Clason, 2021]{Clason2021}
C. Clason (2021), \emph{Regularization of Inverse Problems}, lecture notes,
arXiv:2001.00617v2 [math.FA].

\bibitem[Engl et al., 1996]{EnglHankeNeubauer1996}
H.~W. Engl, M. Hanke \& A. Neubauer (1996), \emph{Regularization of Inverse Problems},
Kluwer Academic Publishers, Dordrecht.

\bibitem[Hadamard, 1902]{Hadamard1902}
J. Hadamard (1902), Sur les probl\`emes aux d\'eriv\'ees partielles et leur signification
physique, \emph{Princeton University Bulletin} 13, 49--52.

\bibitem[Hanke \& Raus, 1996]{HankeRaus1996}
M. Hanke \& T. Raus (1996), A general heuristic for choosing the regularization
parameter in ill-posed problems, \emph{SIAM J. Sci. Comput.} 17(4), 956--972.

\bibitem[Hohage, 2000]{Hohage2000}
T. Hohage (2000), Regularization of exponentially ill-posed problems,
\emph{Numer. Funct. Anal. Optim.} 21(3--4), 439--464.

\bibitem[King \& Chillingworth, 1979]{KingChillingworth1979}
J.~T. King \& D. Chillingworth (1979), Approximation of generalized inverses by
iterated regularization, \emph{Numer. Funct. Anal. Optim.} 1(5), 499--513.

\bibitem[Landweber, 1951]{Landweber1951}
L. Landweber (1951), An iteration formula for Fredholm integral equations of the first
kind, \emph{Amer. J. Math.} 73(3), 615--624.

\bibitem[Lepski{\u\i}, 1990]{Lepskii1990}
O.~V. Lepski{\u\i} (1990), On a problem of adaptive estimation in Gaussian white noise,
\emph{Theory Probab. Appl.} 35(3), 454--466.

\bibitem[Math\'e \& Hofmann, 2008]{MatheHofmann2008}
P. Math\'e \& B. Hofmann (2008), How general are general source conditions?,
\emph{Inverse Problems} 24(1), 015009.

\bibitem[Math\'e \& Pereverzev, 2003]{MathePereverzev2003}
P. Math\'e \& S.~V. Pereverzev (2003), Geometry of linear ill-posed problems in variable
Hilbert scales, \emph{Inverse Problems} 19(3), 789--803.

\bibitem[Moore, 1920]{Moore1920}
E.~H. Moore (1920), On the reciprocal of the general algebraic matrix, \emph{Bull. Amer.
Math. Soc.} 26(9), 394--395.

\bibitem[Morozov, 1966]{Morozov1966}
V.~A. Morozov (1966), On the solution of functional equations by the method of
regularization, \emph{Soviet Math. Dokl.} 7, 414--417.

\bibitem[Nashed, 1987]{Nashed1987}
M.~Z. Nashed (1987), A new approach to classification and regularization of ill-posed
operator equations, in: H.~W. Engl \& C.~W. Groetsch (eds.), \emph{Inverse and Ill-Posed
Problems} (Sankt Wolfgang, 1986), Notes Rep. Math. Sci. Engrg. 4, Academic Press, Boston,
MA, 53--75.

\bibitem[Natterer, 1977]{Natterer1977}
F. Natterer (1977), Regularisierung schlecht gestellter Probleme durch
Projektionsverfahren, \emph{Numer. Math.} 28, 329--342.

\bibitem[Natterer, 1984]{Natterer1984}
F. Natterer (1984), Error bounds for Tikhonov regularization in Hilbert scales,
\emph{Applicable Anal.} 18(1--2), 29--37.

\bibitem[Neubauer, 1997]{Neubauer1997}
A. Neubauer (1997), On converse and saturation results for Tikhonov regularization of
linear ill-posed problems, \emph{SIAM J. Numer. Anal.} 34(2), 517--527.

\bibitem[Penrose, 1955]{Penrose1955}
R. Penrose (1955), A generalized inverse for matrices, \emph{Math. Proc. Cambridge
Philos. Soc.} 51(3), 406--413.

\bibitem[Picard, 1910]{Picard1910}
\'E. Picard (1910), Sur un th\'eor\`eme g\'en\'eral relatif aux \'equations int\'egrales
de premi\`ere esp\`ece et sur quelques probl\`emes de physique math\'ematique,
\emph{Rend. Circ. Mat. Palermo} 29, 79--97.

\bibitem[Plato \& Vainikko, 1990]{PlatoVainikko1990}
R. Plato \& G. Vainikko (1990), On the regularization of projection methods for solving
ill-posed problems, \emph{Numer. Math.} 57(1), 63--79.

\bibitem[Tikhonov, 1963]{Tikhonov1963}
A.~N. Tikhonov (1963), Solution of incorrectly formulated problems and the
regularization method, \emph{Soviet Math. Dokl.} 4, 1035--1038.

\bibitem[Tikhonov \& Arsenin, 1977]{TikhonovArsenin1977}
A.~N. Tikhonov \& V.~Y. Arsenin (1977), \emph{Solutions of Ill-Posed Problems},
V.~H. Winston \& Sons, Washington, D.C.

\end{thebibliography}

\end{document}
